

\documentclass{article}
\usepackage[utf8]{inputenc}

%% Packages
\usepackage[T1]{fontenc}
\usepackage{soul, comment} 
% Unused phfqit package omitted.
% Unused proba package omitted. 
\usepackage[sort,nocompress]{cite}
\usepackage{amssymb}
\usepackage{amsmath}
\usepackage{amsthm}
\usepackage{amsfonts}
\usepackage{mathtools}
\usepackage{xspace}
\usepackage{bm}
\usepackage{nicefrac}
\usepackage{commath}
\usepackage{bbm}
% Unused boxedminipage package omitted.
\usepackage{xparse}
\usepackage{xcolor}
\usepackage{float}
\usepackage{multirow}
\usepackage{makecell}
\usepackage{graphicx}
\usepackage{caption}
\usepackage{subcaption}
\usepackage{ifthen}
\usepackage{thm-restate}
%\usepackage{algorithmic}
\usepackage{algorithm}
\usepackage[noend]{algpseudocode}
\usepackage{colortbl} 
\usepackage{fancyhdr}   
  \pagestyle{fancy}
\usepackage{hyperref}
    \hypersetup{ colorlinks=true, linkcolor=blue, citecolor=magenta, urlcolor=blue,}
\pdfstringdefDisableCommands{%
  \def\Cref#1{<#1>}%
  \def\({}
  \def\){}
}
\usepackage{fullpage}
\usepackage[utf8]{inputenc}
\usepackage{environ}
\usepackage{mdframed}
\usepackage{cleveref}
% Unused optidef package omitted. 
\usepackage{enumitem}
\usepackage{tikz}
\usetikzlibrary{shapes}
\usetikzlibrary{positioning}
\usetikzlibrary{fit}
\usetikzlibrary{graphs,graphs.standard}
\usepackage{enumitem}
%% Marking
\newcommand\TODO[1]{{\color{red}{#1}}\xspace}
\newcommand\change[1]{{\color{red}\underline{#1}}\xspace}
\newcommand{\commentdg}[1]{{\color{blue}{#1}}\xspace}
\newcommand\commentop[1]{{\color{red}{#1}}\xspace}

%% Theorems and Formal Statements
\newtheorem{proposition}{Proposition}[section]
\newtheorem{lemma}{Lemma}[section]
\newtheorem{remark}{Remark}[section]
\newtheorem{theorem}{Theorem}[section]
\newtheorem{corollary}{Corollary}[section]
\newtheorem{thmcorollary}{Corollary}[section]
\newtheorem{definition}{Definition}[section]
\newtheorem{conjecture}{Conjecture}[section]
\newtheorem{observation}{Observation}[section]
\newtheorem{property}{Property}[section]
\newtheorem{claim}{Claim}[section]
\newtheorem{goal}{Goal}[section]
\newtheorem{invariant}{Invariant}
\newtheorem{question}{Question}
\newtheorem{case}{Case}
\usepackage{enumitem}

%\crefname{theorem}{Theorem}{Theorems}
%\crefname{enumi}{part}{parts}

%% English Abbreviations
\newcommand{\ie}{\text{i.e.}\xspace}
\newcommand{\etal}{\text{et al.}\xspace}
\newcommand{\naive}{\text{na\"ive}\xspace}
\newcommand{\wrt}{\text{w.r.t.}\xspace}
\newcommand{\whp}{\text{w.h.p.}\xspace}
\newcommand{\resp}{\text{resp.}\xspace}
\newcommand{\eg}{\text{e.g.}\xspace}
\newcommand{\cf}{\text{cf.}\xspace}
\newcommand{\visavis}{\text{vis-\`a-vis}\xspace}

%% Crypto
\def\alice{{{Alice}}\xspace}
\def\bob{{{Bob}}\xspace}
\newcommand{\polylog}{\ensuremath{\mathrm{polylog}\;}\xspace}
\newcommand{\negl}{\ensuremath{\mathsf{negl}}\xspace}
\newcommand{\secpar}{\ensuremath{{\lambda}}\xspace}
\newcommand{\zo}{\ensuremath{{\{0,1\}}}\xspace}
\newcommand{\attacker}{\ensuremath{\mathcal{A}}}

%% Math Abbreviations
\newcommand{\tuple}[1]{\ensuremath{\left(#1\right)}}
\newcommand{\eps}{\ensuremath{\varepsilon}}
\newcommand{\defeq}{\triangleq}
\newcommand{\ceil}[1]{\ensuremath{\left\lceil{#1}\right\rceil}\xspace}
\newcommand{\floor}[1]{\ensuremath{\left\lfloor{#1}\right\rfloor}\xspace}

%% Problem Box
\NewEnviron{problem}[1]{%
	\begin{center}\fbox{\parbox{6in}{%
				{\centering\scshape #1\par}%
				\parskip=1ex
				\everypar{\hangindent=1em}%
				\BODY
}}\end{center}}

\newcommand{\nnreal}{\ensuremath{\mathbb{R}_{\geq 0}}\xspace}

%% Random Variables
\newcommand{\rvA}{\rv{A}}
\newcommand{\rvB}{\rv{B}}
\newcommand{\rvC}{\rv{C}}
\newcommand{\rvD}{\rv{D}}
\newcommand{\rvG}{\rv{G}}
\newcommand{\rvL}{\rv{L}}
\newcommand{\rvP}{\rv{P}}
\newcommand{\rvQ}{\rv{Q}}
\newcommand{\rvR}{\rv{R}}
\newcommand{\rvS}{\rv{S}}
\newcommand{\rvT}{\rv{T}}
\newcommand{\rvU}{\rv{U}}
\newcommand{\rvV}{\rv{V}}
\newcommand{\rvW}{\rv{W}}
\newcommand{\rvX}{\rv{X}}
\newcommand{\rvY}{\rv{Y}}
\newcommand{\rvZ}{\rv{Z}}
\newcommand{\rvm}{\rv{m}}
\newcommand{\rvr}{\rv{r}}
\newcommand{\rvx}{\rv{x}}
\newcommand{\rvy}{\rv{y}}
\newcommand{\rvc}{\rv{c}}

%% Cal Alphabets
\newcommand{\cA}{\ensuremath{{\mathcal A}}\xspace}
\newcommand{\cB}{\ensuremath{{\mathcal B}}\xspace}
\newcommand{\cC}{\ensuremath{{\mathcal C}}\xspace}
\newcommand{\cD}{\ensuremath{{\mathcal D}}\xspace}
\newcommand{\cE}{\ensuremath{{\mathcal E}}\xspace}
\newcommand{\cF}{\ensuremath{{\mathcal F}}\xspace}
\newcommand{\cG}{\ensuremath{{\mathcal G}}\xspace}
\newcommand{\cH}{\ensuremath{{\mathcal H}}\xspace}
\newcommand{\cI}{\ensuremath{{\mathcal I}}\xspace}
\newcommand{\cK}{\ensuremath{{\mathcal K}}\xspace}
\newcommand{\cL}{\ensuremath{{\mathcal L}}\xspace}
\newcommand{\cM}{\ensuremath{{\mathcal M}}\xspace}
\newcommand{\cN}{\ensuremath{{\mathcal N}}\xspace}
\newcommand{\cO}{\ensuremath{{\mathcal O}}\xspace}
\newcommand{\cP}{\ensuremath{{\mathcal P}}\xspace}
\newcommand{\cQ}{\ensuremath{{\mathcal Q}}\xspace}
\newcommand{\cR}{\ensuremath{{\mathcal R}}\xspace}
\newcommand{\cS}{\ensuremath{{\mathcal S}}\xspace}
\newcommand{\cT}{\ensuremath{{\mathcal T}}\xspace}
\newcommand{\cU}{\ensuremath{{\mathcal U}}\xspace}
\newcommand{\cV}{\ensuremath{{\mathcal V}}\xspace}
\newcommand{\cW}{\ensuremath{{\mathcal W}}\xspace}
\newcommand{\cX}{\ensuremath{{\mathcal X}}\xspace}
\newcommand{\cY}{\ensuremath{{\mathcal Y}}\xspace}
\newcommand{\cZ}{\ensuremath{{\mathcal Z}}\xspace}
    
%% Bold Alphabets
\newcommand{\bA}{\ensuremath{{\mathbf A}}\xspace}
\newcommand{\bB}{\ensuremath{{\mathbf B}}\xspace}
\newcommand{\bC}{\ensuremath{{\mathbf C}}\xspace}
\newcommand{\bD}{\ensuremath{{\mathbf D}}\xspace}
\newcommand{\bE}{\ensuremath{{\mathbf E}}\xspace}
\newcommand{\bF}{\ensuremath{{\mathbf F}}\xspace}
\newcommand{\bG}{\ensuremath{{\mathbf G}}\xspace}
\newcommand{\bQ}{\ensuremath{{\mathbf Q}}\xspace}
\newcommand{\bU}{\ensuremath{{\mathbf U}}\xspace}
\newcommand{\bV}{\ensuremath{{\mathbf V}}\xspace}
\newcommand{\bX}{\ensuremath{{\mathbf X}}\xspace}
\newcommand{\bY}{\ensuremath{{\mathbf Y}}\xspace}
\newcommand{\br}{\ensuremath{{\mathbf r}}\xspace}

% Misc
\newcommand{\OPT}{\ensuremath{\mathtt{OPT}}\xspace}
\newcommand{\spanfunc}{\ensuremath{\mathsf{span}}\xspace}
%\newcommand{\R}{\mathbb{R}}
% \newcommand{\tO}{\tilde{O}}
% \newcommand{\eat}[1]{}
% \newcommand{\lca}{\operatorname{lca}}
\newcommand{\zeros}{\ensuremath{\mathcal{Z}}\xspace}
\newcommand{\demand}{\ensuremath{\mathsf{dem}}\xspace}
\newcommand{\Imerge}{\ensuremath{\mathcal{M}\xspace}}
\newcommand{\Iremove}{\ensuremath{\mathcal{R}\xspace}}

\newcommand{\indicator}{\ensuremath{\mathbbm{1}}\xspace}
\newcommand{\support}{\ensuremath{\mathtt{supp}\xspace}}
\newcommand{\contractionstep}{\textsc{Contraction Step}{}}
\newcommand{\createstep}{\textsc{Create Step}{}}
\newcommand{\algorithmname}{\textsc{Algorithm}{}}
\newcommand{\OracleMain}{\ensuremath{p\text{-max}\mathtt{-Oracle}}\xspace}
\newcommand{\OracleMainWithArg}[1]{\ensuremath{{#1}\text{-max}\mathtt{Oracle}}\xspace}
\newcommand{\OracleAlphaFour}{\ensuremath{\alpha^{(4)}\text{-}\mathtt{Oracle}}\xspace}
\newcommand{\OracleQpm}{\ensuremath{Q\text{-}\mathtt{Oracle}}\xspace}
\newcommand{\alphagood}{\ensuremath{\alpha\text{-good}}\xspace}


\newcommand{\functionMaximizationOracle}[1]{\ensuremath{{#1}\text{-}\mathtt{max}\text{-}\mathtt{wc}\text{-}\mathtt{Oracle}}\xspace}
\newcommand{\functionMaximizationOracleStrongCover}[1]{\ensuremath{{#1}\text{-}\mathtt{max}\text{-}\mathtt{sc}\text{-}\mathtt{Oracle}}\xspace}
\newcommand{\functionMaximizationEmptyOracle}[1]
{\ensuremath{{#1}\text{-}\mathtt{max}\text{-}\mathtt{Oracle}}\xspace}

%% Covering functions tight set families.
\newcommand{\maximalTightSetFamily}[1]{\ensuremath{\calT_{#1}}\xspace}
\newcommand{\maximalTightSetFamilyAfterCreate}[1]{\ensuremath{\calT_{#1}'}\xspace}
\newcommand{\tightSetFamily}[1]{\ensuremath{T_{#1}}\xspace}
\newcommand{\tightSetFamilyAfterCreate}[1]{\ensuremath{T_{#1}'}\xspace}
\newcommand{\projTightSetFamily}[2]{\ensuremath{T_{#1}|_{V_{#2}}}\xspace}
\newcommand{\projTightSetFamilyAfterCreate}[2]{\ensuremath{T_{#1}'|_{V_{#2}}}\xspace}
\newcommand{\projMaximalTightSetFamily}[2]{\ensuremath{\calT_{#1}|_{V_{#2}}}\xspace}
\newcommand{\projMaximalTightSetFamilyAfterCreate}[2]{\ensuremath{\calT_{#1}'|_{V_{#2}}}\xspace}
\newcommand{\cumulativeMaximalTightSetFamily}[1]{\ensuremath{\calT_{\leq #1}}}
\newcommand{\cumulativeMaximalTightSetFamilyAfterCreate}[1]{\ensuremath{\calT_{\leq #1}'}}
\newcommand{\projCumulativeMaximalTightSetFamily}[2]{\ensuremath{\calT_{\leq #1}|_{V_{#2}}}}
\newcommand{\projCumulativeMaximalTightSetFamilyAfterCreate}[2]{\ensuremath{\calT_{\leq #1}'|_{V_{#2}}}}

%% Weak to strong covering tight set families
\newcommand{\maximalTightSetFamilyAfterMerge}[1]{\ensuremath{\calT_{#1}^M}\xspace}
\newcommand{\maximalTightSetFamilyAfterReduce}[1]{\ensuremath{\calT_{#1}^R}\xspace}
\newcommand{\tightSetFamilyAfterMerge}[1]{\ensuremath{T_{#1}^M}\xspace}
\newcommand{\tightSetFamilyAfterReduce}[1]{\ensuremath{T_{#1}^R}\xspace}
\newcommand{\projTightSetFamilyAfterMerge}[2]{\ensuremath{T_{#1}^M|_{V_{#2}}}\xspace}
\newcommand{\projTightSetFamilyAfterReduce}[2]{\ensuremath{T_{#1}^R|_{V_{#2}}}\xspace}
\newcommand{\projMaximalTightSetFamilyAfterMerge}[2]{\ensuremath{\calT_{#1}^M|_{V_{#2}}}\xspace}
\newcommand{\projMaximalTightSetFamilyAfterReduce}[2]{\ensuremath{\calT_{#1}^R|_{V_{#2}}}\xspace}
\newcommand{\cumulativeMaximalTightSetFamilyAfterMerge}[1]{\ensuremath{\calT_{\leq #1}^M}}
\newcommand{\cumulativeMaximalTightSetFamilyAfterReduce}[1]{\ensuremath{\calT_{\leq #1}^R}}
\newcommand{\projCumulativeMaximalTightSetFamilyAfterMerge}[2]{\ensuremath{\calT_{\leq #1}^M|_{V_{#2}}}}
\newcommand{\projCumulativeMaximalTightSetFamilyAfterReduce}[2]{\ensuremath{\calT_{\leq #1}^R|_{V_{#2}}}}

\newcommand{\cumulativeZeros}[1]{\ensuremath{\calZ_{\leq #1}}\xspace}
\newcommand{\cumulativeZerosStrict}[1]{\ensuremath{\calZ_{< #1}}\xspace}

\newcommand{\cumulativeMinimalMaximizerFamily}[1]{\ensuremath{\calF_{\leq #1}}\xspace}
\newcommand{\minimalMaximizerFamily}[1]{\ensuremath{\calF_{#1}}\xspace}
\newcommand{\minimalMaximizerFamilyAfterCreate}[1]{\ensuremath{\calF_{#1}'}\xspace}
\newcommand{\cumulativeMinimalMaximizerFamilyAfterCreate}[1]{\ensuremath{\calF_{\leq #1}'}\xspace}
\newcommand{\maxFunctionValue}[1]{\ensuremath{K_{#1}}\xspace}
\newcommand{\maxFunctionValueAfterCreate}[1]{\ensuremath{K_{#1}'}\xspace}

\newcommand{\cumulativeFunctionMinimalMaximizerFamily}[2]{\ensuremath{\calF_{{#1}_{\leq #2}}\xspace}}
\newcommand{\cumulativeFunctionMinimalMaximizerFamilyAfterCreate}[2]{\ensuremath{\calF_{{#1}_{\leq #2}'}\xspace}}


\newcommand{\slack}{\ensuremath{\Gamma}\xspace}
\newcommand{\slackAfterCreate}{\ensuremath{\Gamma'}\xspace}
\newcommand{\slackAfterContract}{\ensuremath{\Gamma''}\xspace}
\newcommand{\pairwiseSlack}{\ensuremath{\gamma}\xspace}
\newcommand{\pairwiseSlackAfterCreate}{\ensuremath{\gamma'}\xspace}
\newcommand{\pairwiseSlackAfterContract}{\ensuremath{\gamma''}\xspace}

\newcommand{\shubhang}[1]{{\color{blue}[{\tiny Shubhang: \bf #1}]\marginpar{\color{blue}*}}}
\newcommand{\knote}[1]{{\color{red}[{\tiny Karthik: \bf #1}]\marginpar{\color{red}*}}}
\newcommand{\kristof}[1]{{\color{purple}[{\tiny Kristof: \bf #1}]\marginpar{\color{red}*}}}
\newcommand{\functionContract}[2]{\ensuremath{{#1}/_{#2}}\xspace}
\newcommand{\functionRestrict}[2]{\ensuremath{{#1}\backslash_{#2}}\xspace}

\newcommand{\maximalTightSetContainingHyperedge}[1]{\ensuremath{\tau_{#1}}\xspace}
\newcommand{\merge}{\ensuremath{\textsc{Merge}}\xspace}
\newcommand{\reduce}{\ensuremath{\textsc{Reduce}}\xspace}

% Applications of h-splitting-off
\newcommand{\head}{\ensuremath{\texttt{head}}\xspace}
\newcommand{\tails}{\ensuremath{\texttt{tails}}\xspace}
\newcommand{\indegree}{\ensuremath{d^{\text{in}}}\xspace}
\newcommand{\outdegree}{\ensuremath{d^{\text{out}}}\xspace}
\title{Splitting-off in Hypergraphs}
%\title{Hypergraph Splitting-off}
%\title{Hypergraph Splitting-off: Algorithm and Applications}

\author{
%\iffalse
Krist\'{o}f B\'{e}rczi\thanks{MTA-ELTE Matroid Optimization Research Group and HUN-REN-ELTE Egerv\'{a}ry Research Group, Department of Operations Research, E\"{o}tv\"{o}s Lor\'{a}nd
University, Budapest, Hungary, Email: \{kristof.berczi, tamas.kiraly\}@ttk.elte.hu.}
\and Karthekeyan Chandrasekaran\thanks{University of Illinois, Urbana-Champaign, USA, Email: \{karthe, smkulka2\}@illinois.edu. 
%Supported in part by NSF grants CCF-1814613 and CCF-1907937. Karthekeyan was additionally supported by the Distinguished Guest Scientist Fellowship of the Hungarian Academy of Sciences -- grant number VK-6/1/2022. 
Part of this work was done while Karthekeyan and Shubhang were visiting E\"{o}tv\"{o}s Lor\'{a}nd
University.}
\and Tam\'{a}s Kir\'{a}ly\footnotemark[1]
\and Shubhang Kulkarni\footnotemark[2]
%\fi
}
\date{}
\begin{document}
\maketitle
\begin{abstract}
The splitting-off operation in undirected graphs is a fundamental reduction operation that detaches all edges incident to a given vertex and adds new edges between the neighbors of that vertex while preserving their degrees. Lov\'{a}sz \cite{Lov74, Lovasz-problems-book} and Mader \cite{Mad78} showed the existence of this operation while preserving global and local connectivities respectively in graphs under certain conditions. These results have far-reaching applications in graph algorithms literature \cite{Lov76, Mad78, Frank93:Applications-of-submod-functions, FK02, Kiraly-Lau, Fra92, GB93, Fra94, JFJ95, Frank-book, nagamochi_ibaraki_2008_book, NNI97, HW96, Goe01, Jor03, Kri03, JMS03, CFLY11, BHTP03, Lau07, CS08, NZ20, BN23}. 
In this work, we introduce a splitting-off operation in hypergraphs. We show that there exists a local connectivity preserving complete splitting-off in hypergraphs and give a strongly polynomial-time algorithm to compute it in weighted hypergraphs. We illustrate the usefulness of our splitting-off operation in hypergraphs by showing two applications: (1) we give a constructive characterization of $k$-hyperedge-connected hypergraphs and (2) we give an alternate proof of an approximate min-max relation for max Steiner rooted-connected orientation of graphs and hypergraphs (due to Kir\'{a}ly and Lau \cite{Kiraly-Lau}). Our proof of the approximate min-max relation for graphs circumvents the Nash-Williams' strong orientation theorem and uses tools developed for hypergraphs.
\end{abstract}

\newpage
\tableofcontents
\newpage

\setcounter{page}{1}
\section{Introduction}\label{sec:introduction}
The splitting-off operation in undirected graphs is a simple yet powerful operation in graph theory.  
It is a reduction operation that detaches all edges incident to a given vertex and adds new edges between the neighbors of that vertex while preserving their degrees. 
%and edge-connectivity. 
Equivalently, it enables a vertex to exit the network by informing its neighbors  
how to reconfigure the lost links 
among themselves in order to preserve their degrees.
%edge-connectivity (as well as degrees). 
Lov\'{a}sz \cite{Lov74, Lovasz-problems-book} introduced the splitting-off operation and showed the existence of the operation to preserve global edge-connectivity under certain conditions. Mader \cite{Mad78} showed the existence of the splitting-off operation to preserve local edge-connectivity (i.e., all pairwise edge-connectivities) under certain conditions. 
Both Lov\'{a}sz's and Mader's results also admit strongly polynomial-time algorithms \cite{Frank-book, nagamochi_ibaraki_2008_book, Gab94}. 
%The inductive nature of the splitting-off operation leads to far ranging applications in graph theory, algorithms, and optimization. It enables fundamental graph theoretical results such as constructive characterization of $k$-edge-connected graphs \cite{??}. 
Owing to the inductive nature of the splitting-off operation, Lov\'{a}sz's and Mader's results have enabled fundamental results in graph theory as well as efficient algorithms and min-max relations for numerous graph optimization problems. 
In fact, Mader \cite{Mad78} illustrated the power of his local edge-connectivity preserving splitting-off result by deriving Nash-Williams' strong orientation theorem \cite{NW60} (also see Frank's exposition of this derivation \cite{Frank93:Applications-of-submod-functions}).  
%\shubhang{(side) note: \cite{Gab94} gives an efficient implementation of Frank's proof of the Strong Orientation Theorem.}). 
%In particular, it has been used 
Subsequently, the splitting-off operation 
has been used 
to give a constructive characterization of $k$-edge-connected graphs \cite{Frank-book} 
 and 
to address problems in edge-connectivity augmentation \cite{Fra92, Fra94, JFJ95, Frank-book, nagamochi_ibaraki_2008_book}, graph orientation \cite{FK02, Kiraly-Lau}, minimum cuts enumeration \cite{NNI97, HW96, Goe01}, network design \cite{Lov76, GB93, Jor03, CS08, CFLY11}, tree packing \cite{BHTP03, Kri03, JMS03, Lau07}, congruency-constrained cuts \cite{NZ20}, and approximation algorithms for TSP \cite{GB93, BN23}. 
Designing fast algorithms for global/local edge-connectivity preserving splitting-off remains an active area of research (e.g., see recent works \cite{LY13, CLP22-soda, CLP22-stoc, CHLP23}) due to these far-reaching applications. 
%We refer the reader to the textbooks by Frank \cite{Frank-book} and by Nagamochi and Ibaraki \cite{nagamochi_ibaraki_2008_book} for efficient algorithms to compute global/local edge-connectivity preserving splitting-off and some of their applications. 
In this work, we introduce a splitting-off operation in \emph{hypergraphs}, show the existence of local-connectivity preserving splitting-off operation, design a strongly polynomial-time algorithm to compute it in weighted hypergraphs, and illustrate its usefulness by showing two applications. 

\paragraph{Hypergraphs.}
Hypergraphs offer a richer and more accurate model than graphs for several applications. 
Consequently, hypergraphs have found applications in several modern areas (e.g., see \cite{VeldtBK22,Schlagetal23,LiM17, Ornes21}) and these applications have, in turn, driven exciting recent progress in algorithms for hypergraph optimization problems \cite{KK15, GKP17, CXY19-j, CX18, CC22, CKN20, SY19, BST19, KKTY21-FOCS, KKTY21-STOC, Lee22, JLS22, Quanrud22, GLSTW22, CC23, AGL23, FPZ23}. 
A hypergraph $G=(V, E)$ consists of a finite set $V$ of vertices and a set $E$ of hyperedges, where every hyperedge $e\in E$ is a subset of $V$. We will denote a hypergraph $G=(V, E)$ with hyperedge weights $w: E\rightarrow \Z_+$ by the tuple $(G, w)$ and use $G_w$ to denote the unweighted multi-hypergraph over vertex set $V$ containing $w(e)$ copies of every hyperedge $e\in E$.  
Throughout this work, we will be interested only in hypergraphs with positive integral weights and for algorithmic problems where the input/output is a hypergraph, we will require that the weights are represented in binary. 
If all hyperedges have size at most $2$, then the hyperedges are known as edges and we call such a hypergraph as a graph. 
We emphasize a subtle but important distinction between hypergraphs and graphs: the number of hyperedges in a hypergraph could be exponential in the number of vertices. This is in sharp contrast to graphs where the number of edges is at most the square of the number of vertices. Consequently, in hypergraph network design problems where the goal is to construct a hypergraph with certain properties,  
one needs to be mindful of the number of hyperedges in the hypergraph returned by the algorithm. %We mention that recent 
Recent works in hypergraph algorithms literature have 
%addressed this precise 
focused on this 
issue in the context of cut/spectral sparsification of hypergraphs \cite{KK15, CX18, SY19, BST19, CKN20, KKTY21-FOCS, KKTY21-STOC, Lee22, JLS22, Quanrud22}. 


\iffalse
A hypergraph $H=(V, E)$ consists of a finite set $V$ of vertices and a set $E$ of hyperedges, where every hyperedge $e\in E$ is a subset of $V$. We will denote a hypergraph $H=(V, E)$ with hyperedge weights $w: E\rightarrow \Z_+$ by the tuple $(H, w)$ and use $H_w$ to denote the unweighted multi-hypergraph over vertex set $V$ containing $w(e)$ copies of every hyperedge $e\in E$.  
Throughout this work, we will be interested only in hypergraphs with positive integral weights and for algorithmic problems where the input/output is a hypergraph, we will require that the weights are represented in binary. 
If all hyperedges have size at most $2$, then the hyperedges are known as edges and we call such a hypergraph as a graph. 
Hypergraphs offer a richer and more accurate model than graphs for several applications. 
Consequently, hypergraphs have found applications in several modern areas (e.g., see \cite{VeldtBK22,Schlagetal23,LiM17, Ornes21}) and these applications have, in turn, driven exciting recent progress in algorithms for hypergraph optimization problems \cite{KK15, GKP17, CXY19-j, CX18, CC22, CKN20, SY19, BST19, KKTY21-FOCS, KKTY21-STOC, Lee22, JLS22, Quanrud22}. 
We emphasize a subtle but important distinction between hypergraphs and graphs: the number of hyperedges in a hypergraph could be 
%We note that the number of hyperedges in a hypergraph could be 
exponential in the number of vertices. This is in sharp contrast to graphs where the number of edges is at most the square of the number of vertices. Consequently, in hypergraph network design problems where the goal is to construct a hypergraph with certain properties,  
one needs to be mindful of the number of hyperedges in the hypergraph returned by the algorithm. 
\fi

\iffalse
Recent works have explored this issue in the context of hypergraph sparsifiers \cite{KK15, CX18, SY19, BST19, CKN20, KKTY21-FOCS, KKTY21-STOC, Lee22, JLS22, Quanrud22}  
and shown that every hypergraph (possibly with exponential number of hyperedges) admits a cut/coverage/spectral sparsifier
in which the number of hyperedges is near-linear in the number of vertices and such a sparsifier can be constructed in time near-linear in the size of the input hypergraph. 
\knote{Perhaps, this discussion about exponential number of hyperedges is better done after defining splitting-off.} 
\fi

\paragraph{Notation.} Let $(G=(V, E), w: E\rightarrow \Z_+)$ be a hypergraph. 
For $X\subseteq V$, let $\delta_G(X)\coloneqq \{e\in E: e\cap X\neq \emptyset, e\setminus X\neq \emptyset\}$.  
We define the cut function $d_{(G, w)}: 2^V\rightarrow \Z_{\ge 0}$ by $d_{(G, w)}(X) \coloneqq \sum_{e\in \delta_G(X)}w(e)$ for every $X\subseteq V$. 
For a vertex $v\in V$, we use $d_{(G,w)}(v)$ to denote $d_{(G,w)}(\{v\})$. 
We define the degree of a vertex $v$ to be the sum of the weights of hyperedges containing $v$---we note that the degree of a vertex is not necessarily equal to $d_{(G, w)}(v)$ since we could have $\{v\}$ itself as a hyperedge (i.e., a singleton hyperedge that contains only the vertex $v$). 
For distinct vertices $u, v\in V$, let $\lambda_{(G, w)}(u, v)\coloneqq \min\{d_{(G, w)}(X): u\in X\subseteq V\setminus \{v\}\}$ -- i.e., $\lambda_{(G, w)}(u, v)$ is the value of a minimum $\{u, v\}$-cut in the hypergraph. If all hyperedge weights are unit, then we drop $w$ from the subscript and simply use $d_G$ and $\lambda_G$. 

\paragraph{Hypergraph Splitting-off.} 
%A notion of hypergraph splitting-off motivated by hypergraph connectivity augmentation applications has been studied in the literature before \cite{BJ99, Cosh-thesis, BK12}. These works have explored splitting-off a vertex $s$ from a hypergraph under the assumption that all hyperedges incident to the vertex $s$ are in fact edges. In contrast, we are interested in splitting-off at a vertex $s$ from a hypergraph without such an assumption (i.e., the vertex $s$ could have arbitrary-sized hyperedges incident to it) \knote{This may not be informative for a reader unfamiliar with the splitting-off operation.}. 
We now introduce our definition of splitting-off in hypergraphs. 
To compare and contrast our definition of splitting-off for hypergraphs with the classical definition of splitting-off for graphs, we include both our definition and the classical definition and distinguish them by identifying them as h-splitting-off and g-splitting-off. 
%(with the convention being that we add the prefix `h' for hypergraph related terminology and the prefix `g' for graph related terminology). 
In the definitions below, we encourage the reader to consider the input hypergraph $(G, w)$ to be a graph while considering g-splitting-off terminology and to be a hypergraph while considering h-splitting-off terminology. We encourage the reader to also assume unit weights during first read. See Figure \ref{figure:hypergraph-splitting-off-example-1} for an example. 

\begin{definition}\label{defn:splitting-off}
    Let $(G=(V, E), w: E\rightarrow \Z_+)$ be a hypergraph and $s\in V$. 
    \begin{enumerate}
            
        \item 
        In \emph{merge almost-disjoint hyperedges}, we pick a pair of hyperedges $e, f\in \delta_G(s)$ such that $e\cap f=\{s\}$, pick a positive integer $\alpha \in \Z_+$ such that $\alpha\le \min\{w(e), w(f)\}$, reduce the weights of hyperedges $e$ and $f$ by $\alpha$, and increase the weight of a hyperedge $g$ by $\alpha$. Here, 
        \begin{enumerate}
            \item if we choose $g\coloneqq e \cup f$, then the associated operation will be called as \emph{h-merge almost-disjoint hyperedges operation}.
            \item if we choose $g\coloneqq (e\cup f)\setminus \{s\}$, then the associated operation will be called as \emph{g-merge almost-disjoint hyperedges operation}.
        \end{enumerate}
        In the above, if $\alpha=w(e)$ (resp. if $\alpha = w(f)$), then we discard the hyperedge $e$ (resp. hyperedge $f$) from the hypergraph obtained after the operation; if the hyperedge $g\not\in E$, then we introduce $g$ as a new hyperedge with weight $w(g) := 0$ before performing the weight increase on the hyperedge $g$.
    
        \item 
        In \emph{trim hyperedge operation}, we pick a hyperedge $e\in \delta_G(s)$, pick a positive integer $\alpha \in \Z_+$,
        reduce the weight of the hyperedge $e$ and 
        increase the weight of the hyperedge $g\coloneqq e\setminus \{s\}$. Here, 
        \begin{enumerate}
            \item if we choose $\alpha \le w(e)$, reduce the weight of the hyperedge $e$ by $\alpha$,  and increase the weight of the hyperedge $g$ by $\alpha$, then the associated operation will be called as \emph{h-trim operation} (if $\alpha=w(e)$, then we discard $e$ from the hypergraph obtained after the operation; if $g\not\in E$, then we add $g$ as a new hyperedge with weight $w(g) := 0$ before performing the weight increase on the hyperedge $g$). 
            
            \item if we choose $\alpha \le w(e)/2$, reduce the weight of the hyperedge $e$ by $2\alpha$, and increase the weight of the hyperedge $g$ by $2\alpha$, then the associated operation will be called as \emph{g-trim operation} (if $\alpha=w(e)/2$, then we discard $e$ from the hypergraph obtained after the operation; if $g\not\in E$, then we add $g$ as a new hyperedge with weight $w(g) := 0$ before performing the weight increase on the hyperedge $g$).  
        \end{enumerate}

\item We say that a hypergraph $(H=(V, E_H), w_H: E_H\rightarrow \Z_+)$ is obtained by applying a 
        \begin{enumerate}
            \item \emph{h-splitting-off operation at $s$ from $(G, w)$} if $(H, w_H)$ is obtained from $(G, w)$ by either the h-merge almost-disjoint hyperedges operation or the h-trim hyperedge operation. 
            \item \emph{g-splitting-off operation at $s$ from $(G, w)$} if $(H, w_H)$ is obtained from $(G, w)$ by either the g-merge almost-disjoint hyperedges operation or the g-trim hyperedge operation. 
        \end{enumerate}

%\iffalse
\end{enumerate}
\end{definition}
Certain remarks regarding the definitions are in order. 
%h-splitting-off and g-splitting-off \emph{almost-disjoint} hyperedges operations ensure that the degree of every vertex $u\in V\setminus\{s\}$ is preserved. 
%However, if we allow the same operations for an arbitrary pair of hyperedges $e, f\in \delta_G(s)$, then degrees of vertices $u\in V\setminus \{s\}$ may not be preserved. 
%Similarly, both h-trimming and g-trimming operations preserve the degree of every vertex $u\in V\setminus \{s\}$. 
Firstly, the trim operation is valuable and unique to hypergraphs. It has been used in the hypergraph literature to obtain small-sized certificates for hypergraph connectivity \cite{CX18} and for certain notions of directed hypergraph connectivity \cite{Frank-book}. 
We note that the trim operation has limited value in graphs---trimming an edge leads to a singleton edge and singleton edges contribute only to the degree but not to the cut value of any set. 
%However, the trimming operation in graphs has limited value. 
Secondly, all operations mentioned above are degree preserving for vertices $u\in V\setminus \{s\}$: both h-trim and g-trim operations preserve degrees by definition; both h-merge and g-merge \emph{almost-disjoint} hyperedges operations preserve degrees due to the almost-disjoint property of the chosen hyperedges. 
Thirdly, all operations mentioned above 
%both h-merge operation and g-splitting-off operations 
\emph{do not} increase the %coverage/
cut values of subsets $X\subseteq V\setminus\{s\}$. Thus, the relevant goal with these operations is ensuring that the cut values do not decrease too much---i.e., preserving global/local connectivity. 
We will be interested in repeated application of h-splitting-off operations at a vertex from a given hypergraph to isolate that vertex while preserving global/local connectivity. We define these formally next. 
\begin{definition}
    Let $(G=(V, E), w: E\rightarrow \Z_+)$ be a hypergraph and $s\in V$. 
    \begin{enumerate}
    %\fi
        \item We say that a hypergraph $(G^*=(V, E^*), w^*: E^*\rightarrow \Z_+)$ is a 
        \begin{enumerate}
            \item \emph{complete h-splitting-off at $s$ from $(G, w)$} if $d_{(G^*, w^*)}(s) = 0$ and $(G^*, w^*)$ is obtained from $(G, w)$ by repeatedly applying h-splitting-off operations at $s$ from the current hypergraph. 
            \item \emph{complete g-splitting-off at $s$ from $(G, w)$} if $d_{(G^*, w^*)}(s) = 0$ and $(G^*, w^*)$ is obtained from $(G, w)$ by repeatedly applying g-splitting-off operations at $s$ from the current hypergraph. 
        \end{enumerate}

        \item Let $(G^*, w^*)$ be a complete h-splitting-off/g-splitting-off at $s$ from $(G, w)$. We say that $(G^*, w^*)$ 
        \begin{enumerate}
            \item 
            preserves local connectivity if $\lambda_{(G^*, w^*)}(u, v)=\lambda_{(G, w)}(u, v)$ for every distinct $u, v\in V\setminus\{s\}$ and 
            
            \item preserves global connectivity if \\
            $\min\{\lambda_{(G^*, w^*)}(u, v): u, v\in V\setminus \{s\}, u \neq v\}=\min\{\lambda_{(G, w)}(u, v): u, v\in V\setminus \{s\}, u \neq v\}$. 
            
        \end{enumerate}
    \end{enumerate}
\end{definition}

\iffalse
\begin{definition}\label{defn:splitting-off}
    Let $(G=(V, E), w: E\rightarrow \Z_+)$ be a hypergraph and $s\in V$. 
    \begin{enumerate}
        
        
        \item 
        In \emph{splitting-off almost-disjoint hyperedges} \knote{\emph{merge almost-disjoint hyperedges}}, we pick a pair of hyperedges $e, f\in \delta_G(s)$ such that $e\cap f=\{s\}$, pick a positive integer $\alpha \in \Z_+$ such that $\alpha\le \min\{w(e), w(f)\}$, reduce the weights of hyperedges $e$ and $f$ by $\alpha$, and increase the weight of a hyperedge $g$ by $\alpha$. Here, 
        \begin{enumerate}
            \item if we choose $g\coloneqq e \cup f$, then the associated operation will be called as \emph{h-splitting-off almost-disjoint hyperedges operation} \knote{\emph{h-merge almost-disjoint hyperedges operation}}.
            \item if we choose $g\coloneqq (e\cup f)\setminus \{s\}$, then the associated operation will be called as \emph{g-splitting-off almost-disjoint hyperedges operation} \knote{\emph{g-merge almost-disjoint hyperedges operation}}.
        \end{enumerate}
        In the above, if $\alpha=w(e)$ (resp. if $\alpha = w(f)$), then we discard the hyperedge $e$ (resp. hyperedge $f$) from the hypergraph obtained after the operation; if the hyperedge $g\not\in E$, then we introduce $g$ as a new hyperedge with weight $w(g) := 0$ before performing the weight increase on the hyperedge $g$.
    
        \item 
        In the \emph{trimming hyperedge operation}, we pick a hyperedge $e\in \delta_G(s)$, pick a positive integer $\alpha \in \Z_+$,
        reduce the weight of the hyperedge $e$ and 
        increase the weight of the hyperedge $g\coloneqq e\setminus \{s\}$. Here, 
        \begin{enumerate}
            \item if we choose $\alpha \le w(e)$, reduce the weight of the hyperedge $e$ by $\alpha$,  and increase the weight of the hyperedge $g$ by $\alpha$, then the associated operation will be called as \emph{h-trimming operation} (if $\alpha=w(e)$, then we discard $e$ from the hypergraph obtained after the operation; if $g\not\in E$, then we add $g$ as a new hyperedge with weight $w(g) := 0$ before performing the weight increase on the hyperedge $g$). 
            \item if we choose $\alpha \le w(e)/2$, reduce the weight of the hyperedge $e$ by $2\alpha$, and increase the weight of the hyperedge $g$ by $2\alpha$, then the associated operation will be called as \emph{g-trimming operation} (if $\alpha=w(e)/2$, then we discard $e$ from the hypergraph obtained after the operation; if $g\not\in E$, then we add $g$ as a new hyperedge with weight $w(g) := 0$ before performing the weight increase on the hyperedge $g$).  
        \end{enumerate}

        
       \item We say that a hypergraph $(H=(V, E_H), w_H: E_H\rightarrow \Z_+)$ is obtained by applying a 
        \begin{enumerate}
            \item \emph{h-splitting-off operation at $s$ from $(G, w)$} if $(H, w_H)$ is obtained from $(G, w)$ by either the h-splitting-off almost-disjoint hyperedges operation \knote{h-merge almost-disjoint hyperedges operation} or the h-trimming hyperedge operation. 
            \item \emph{g-splitting-off operation at $s$ from $(G, w)$} if $(H, w_H)$ is obtained from $(G, w)$ by either the g-splitting-off almost-disjoint hyperedges operation \knote{g-merge almost-disjoint hyperedges operation} or the g-trimming hyperedge operation. 
        \end{enumerate}
        
        \item We say that a hypergraph $(G^*=(V, E^*), w^*: E^*\rightarrow \Z_+)$ is a 
        \begin{enumerate}
            \item \emph{complete h-splitting-off at $s$ from $(G, w)$} if $d_{(G^*, w^*)}(s) = 0$ and $(G^*, w^*)$ is obtained from $(G, w)$ by repeatedly applying h-splitting-off operations at $s$ from the current hypergraph. 
            \item \emph{complete g-splitting-off at $s$ from $(G, w)$} if $d_{(G^*, w^*)}(s) = 0$ and $(G^*, w^*)$ is obtained from $(G, w)$ by repeatedly applying g-splitting-off operations at $s$ from the current hypergraph. 
        \end{enumerate}

        \item Let $(G^*, w^*)$ be a complete h-splitting-off/g-splitting-off at $s$ from $(G, w)$. We say that $(G^*, w^*)$ 
        \begin{enumerate}
            \item 
            preserves local connectivity if $\lambda_{(G^*, w^*)}(u, v)=\lambda_{(G, w)}(u, v)$ for every distinct $u, v\in V\setminus\{s\}$ and 
            
            \item preserves global connectivity if \\
            $\min\{\lambda_{(G^*, w^*)}(u, v): u, v\in V\setminus \{s\}, u \neq v\}=\min\{\lambda_{(G, w)}(u, v): u, v\in V\setminus \{s\}, u \neq v\}$. 
            
        \end{enumerate}
    \end{enumerate}
\end{definition}
\fi



Our first contribution in this work is the definition of h-splitting-off operations at a vertex from a hypergraph. 
%(i.e., Definition \ref{defn:splitting-off}). 
To the best of our knowledge, this definition 
%the definition of complete h-splitting-off at a vertex from a hypergraph 
has not appeared in the literature before. 
A notion of hypergraph splitting-off motivated by hypergraph connectivity augmentation applications has been studied in the literature before \cite{BJ99, Cosh-thesis, BK12}. 
These works have explored local connectivity preserving complete g-splitting-off at a vertex $s$ from a hypergraph under the assumption that \emph{all hyperedges incident to the vertex $s$ are edges (i.e., have size at most $2$)}. 
In contrast, our focus is on local connectivity preserving complete h-splitting-off at a vertex $s$ from a hypergraph without any assumption on the size of the hyperedges incident to the vertex $s$ (i.e., the vertex $s$ could have arbitrary-sized hyperedges incident to it). We refer the reader to Figure \ref{figure:hypergraph-splitting-off-example-1} for an example of complete h-splitting-off at a vertex from a hypergraph. 

\begin{figure}[thb]
\centering
\includegraphics[width=\textwidth]{hypergraph-splitting-off-example-1.jpg}
\caption{Example of (local connectivity preserving) complete h-splitting-off at a vertex $s$ from a hypergraph. Consider the leftmost hypergraph where all hyperedge weights are one and the vertex $s$ is as labeled. Operations (I) and (II) correspond to h-merge almost-disjoint hyperedges operations and Operation (III) corresponds to an h-trim hyperedge operation.}
\label{figure:hypergraph-splitting-off-example-1}
\end{figure}

We will primarily be concerned with complete \emph{h-splitting-off} at a vertex from a \emph{hypergraph} and complete \emph{g-splitting-off} at a vertex from a \emph{graph}.  
Complete g-splitting-off at a vertex from a \emph{graph} is equivalent to the classical and well-studied notion of complete splitting-off at a vertex from a graph (for the definition of the classical notion in graphs, see \cite{Frank-book, nagamochi_ibaraki_2008_book}). 
We cast the results of Lov\'{a}sz \cite{Lov74, Lovasz-problems-book} and Mader \cite{Mad78} in the framework of our definitions now. Let $(G, w)$ be a \emph{graph} and let $s$ be a vertex in $G$. 
Lov\'{a}sz \cite{Lov74, Lovasz-problems-book} showed that if $d_{(G,w)}(\{s\})$ is even and $\min\{\lambda_{(G, w)}(u,v): u, v\in V\setminus \{s\}\}\ge K$ for some $K\ge 2$, then there exists a \emph{global} connectivity preserving complete g-splitting-off at the vertex $s$ from $(G, w)$. 
Mader \cite{Mad78} showed that if $d_{(G,w)}(\{s\})$ is even, there is no cut-edge\footnote{Equivalently, for every edge $e\in \delta_G(s)$ with $w(e)=1$, the removal of that edge does not disconnect the graph.} incident to $s$,  
%(i.e., for every hyperedge $e\in \delta(s)$ with $w(e)=1$, the removal of that hyperedge does not disconnect the hypergraph), 
and $(G, w)$ is connected, then there exists a \emph{local} connectivity preserving complete g-splitting-off at the vertex $s$ from $(G, w)$. 


We compare and contrast complete h-splitting-off at a vertex from a hypergraph and complete g-splitting-off at a vertex from a graph.  
Complete h-splitting-off at a vertex $s$ from a hypergraph enables the vertex $s$ to exit the hypergraph by informing its incident hyperedges about how to merge and trim themselves in  order to preserve degrees. In this sense, the definition of complete h-splitting-off at a vertex from a hypergraph serves the same role as complete g-splitting-off at a vertex from a graph. On the other hand, there are important differences between the two notions. Firstly, complete h-splitting-off at a vertex from a \emph{graph} may not necessarily be a graph (owing to the creation of hyperedges of size at least $3$) while it is an easy exercise to show that complete g-splitting-off at a vertex from a \emph{graph} will necessarily be a graph. %\kristof{Is this footnote needed?}\shubhang{I think we can remove it since it doesn't add value.}. 
%\footnote{Complete g-splitting-off at a vertex from a graph will necessarily be a graph: we leave this as an exercise for the reader.} 
Secondly, local/global connectivity preserving complete g-splitting-off at a vertex from a graph may not exist---see Figure \ref{figure:hypergraph-splitting-off-example-2}. 
%\footnote{Consider the graph $(G=(V, E), c: E\rightarrow \{1\})$, where $G$ is the $4$-vertex star graph with the vertex $s$ being the center of the star. Local connectivity preserving complete g-splitting-off at the vertex $s$ from $(G, c)$ does not exist.}. 

\begin{figure}[thb]
\centering
\includegraphics[scale=0.5]{hypergraph-splitting-off-example-2.jpg}
\caption{An example showing that global connectivity preserving complete g-splitting-off at a vertex from a graph may not exist. All edge weights are one and the vertex $s$ is as labeled.}
\label{figure:hypergraph-splitting-off-example-2}
\end{figure}

As our second main contribution, we show that local connectivity preserving complete h-splitting-off at a vertex from a hypergraph always exists and can be computed in strongly polynomial time (the rightmost hypergraph in Figure \ref{figure:hypergraph-splitting-off-example-1} is a local connectivity preserving complete h-splitting-off at the vertex $s$ from the hypergraph in Figure \ref{figure:hypergraph-splitting-off-example-2}). 

\begin{theorem}\label{thm:complete-splitting-off}
    Given a hypergraph $(G=(V, E), w_G: E\rightarrow \Z_+)$ and a vertex $s\in V$, there exists a strongly polynomial-time algorithm to find a local connectivity preserving complete h-splitting-off at $s$ from $(G, w_G)$. 
\end{theorem}
%The run-time of our algorithm is a somewhat large polynomial in the size of the input hypergraph. A crude upper bound on the run-time of our algorithm is $O(n^8m)$, where $n$ is the number of vertices and $m$ is the number of hyperedges in the input hypergraph. 

A distinction between Theorem \ref{thm:complete-splitting-off} and the graph splitting-off results of Lov\'{a}sz and Mader is that Theorem \ref{thm:complete-splitting-off} shows the existence of a local connectivity preserving complete h-splitting-off at a vertex from a hypergraph \emph{without any assumptions on the hypergraph} whereas Lov\'{a}sz's and Mader's results hold only under certain technical assumptions on the graph. In several applications of their results, additional arguments are needed to address cases where those technical assumptions do not hold. For this reason, we believe that Theorem \ref{thm:complete-splitting-off} could be useful in simplifying the arguments involved in some of the applications of Lov\'{a}sz's and Mader's graph splitting-off results (e.g., we will later see that the edge version of Menger's theorem in undirected \emph{graphs} follows in a straightforward fashion from Theorem \ref{thm:complete-splitting-off}). 
%In various applications of their results, satisfying these mild conditions requires additional technical arguments whereas . 

A crude run-time of our algorithm that proves Theorem \ref{thm:complete-splitting-off} is $O(|V|^6(|V|+|E|)^3)$.  %$O(|V|^8|E| + |V|^6|E|^3)$. 
%We leave the problem of finding efficient implementations of our algorithm as an intriguing open question. 
%(and it can be improved slightly using the recently discovered near-linear time max-flow algorithm). 
We understand that this run-time is impractical for applications. Nevertheless, we mention it here explicitly for the sake of completeness and as a potential starting point for future work: it would be interesting to design a near-linear time algorithm to find a local connectivity preserving complete h-splitting-off at a vertex from a weighted hypergraph.


\iffalse
\knote{Would be good to state the run-time.}
\shubhang{Crude runtime bound of the algorithm can be broken down as follows: Let $(H = (V, E),w:E\rightarrow\Z_+)$ and $p:2^V\rightarrow\Z$ be inputs to the algorithm. At the topmost level of recursion, we have that:
\begin{enumerate}
    \item $\functionMaximizationOracleStrongCover{p}((H,w), S, T , y)$ can be implemented in $O(|V|^2)$ \emph{max-flow} calls on hypergraphs with sizes linear in that of $(H,w)$. (Note: This is because for every pair of vertices $u,v \in {V\choose 2}$, we can implement $\lambda_{((G, c), (G_0, c_0), y_0)}(u, v)$ in the proof of \Cref{lemma:helper-for-p-max-oracle} using a single max-flow computation on hypergraph with $|V|+1$ vertices and $|E| + |V|$ hyperedges. Then, follow the proof of \Cref{lemma:helper-for-p-max-oracle} to see that having this oracle for every pair of vertices suffices). We will crudely upper-bound the max-flow query time as $O((|V|+|E|)|V||E|)$.
    %$O((|V||E|)^{1+o(1)})$ (randomized) time assuming polynomial edge weights. 
    (Note: the $|V||E|$ term in the runtime is a crude upper bound on the area of the hypergraph.)
    \item $\functionMaximizationOracle{p}((H,w), S, T, y)$ can be implemented using $O(|V|)$ queries to $\functionMaximizationOracleStrongCover{p}$ and $O(|V| + |E|)$ time per query (to construct the input). 
    \end{enumerate}
    
    Note that by \Cref{claim:WeakToStrong:p_i-oracle-using-p-oracle}, \functionMaximizationOracle{p_i} for every input function $p_i$ encountered during the recursion can be implemented using $O(|V|)$ queries to $\functionMaximizationOracleStrongCover{p}$ (recall $p$ is the input function to the algorithm). Thus,  if we bound the time required for all steps of the algorithm (except recursion) at just the topmost recursive level in terms of maximization-weak-oracle calls, then multiplying the bound by the recursion depth gives us a crude runtime bound on the overall algorithm. To get the bound we observe that
    \begin{enumerate}
    \item The family $\maximalTightSetFamily{p, H, w}$ can be computed in $O(|V|^2)$ queries to $\functionMaximizationOracle{p}$ by \Cref{lem:FunctionMaximizetionOracles:computing-maximal-pHw-tight-sets-with-p-max-oracle}. Each query takes $O(|V| + |E|)$ time (to construct the query input).
     This is also the time needed for step (4) of the algorithm (note that during the first recursive call we will need to compute the entire family).
    \item $\beta^M$ can be computed in $O(|V|^2)$ queries to $\functionMaximizationOracle{p}$. Each query takes $O(|V| + |E|)$ time (to construct the query input).
    \item $\beta^R$ can be computed in $1$ query to $\functionMaximizationOracle{p}$ in $O(|V| + |E|)$ time (to construct the query input).
    \item All other steps of the algorithm can be computed in $O(|V| + |E|)$ time per step. Since these steps are inconsequential from runtime perspective, we will ignore them in our runtime bound below.
\end{enumerate}
Thus, time bound per recursive call is (ignoring big-Oh notation for convenience)
\begin{align*}
    &[\text{time per  \functionMaximizationOracle{p} call}] \cdot [\text{$\#$ of  \functionMaximizationOracle{p} calls during rec. call}] &\\
    & = (|V| \cdot (|V| + |E| + [\text{time per  \functionMaximizationOracleStrongCover{p} call}])) \cdot |V|^2&\\
    &= (|V| \cdot (|V| + |E| + |V|^2\cdot[\text{time per max-flow call}])) \cdot |V|^2&\\
    &= (|V| \cdot (|V| + |E| + |V|^2\cdot(|V|+|E|)|V||E|)) \cdot |V|^2&\\
    & = (|V| + |E|)\cdot|V|^{6}|E|&
\end{align*}
Finally, we note that the recursion depth of our algorithm is $|E| + O(|V|)$ by \Cref{lem:WeakToStrongCover:NumberOfHyperedges}. Thus overall runtime of the algorithm is $O((|V| + |E|)^2\cdot|V|^{6}|E|)$ as remarked above.
}
\fi

%\shubhang{Assuming polynomially bounded hyperedge weights, the implementation of \functionMaximizationOracleStrongCover{p} can be improved to $\Tilde{O}(|V|)$ max flow computations by computing $|V|$ Gomory-Hu trees, each using $\Tilde{O}(1)$ max-flow computations via Abboud et. al.. Perhaps a single Gomory-Hu tree of a modified hypergraph suffices? If so, then we can compute the required Gomory-Hu tree (without needing the polynomially bounded hyperedge-weights assumption) using the classical $O(|V|)$ max-flows algorithm.}

\begin{remark}\label{remark:exp-sized-conn-preserving-splitting-off}
%\paragraph{An exponential-sized connectivity preserving complete h-splitting-off.} 
We note that existence of a local/global connectivity preserving complete h-splitting-off at a vertex from a hypergraph does not necessarily imply a polynomial-time algorithm to find it. 
This is because, 
%In fact, it is possible that 
a local/global connectivity preserving complete h-splitting-off at a vertex from a hypergraph $(G, w_G)$ could contain exponential number of hyperedges although $G$ contains only polynomial number of hyperedges. 
%, it does not necessarily imply that the number of hyperedges in $G^*$ is polynomial in the number of hyperedges in $G$. 
We give an example to illustrate this issue. 
Consider the graph $(G=(V, E), w_G)$ where $G$ is the star graph on $n+1$ vertices with $s$ being the center of the star and all edge weights are $2^{n-1}-1$. 
Consider the hypergraph $(H=(V, E_H), w_H)$ such that 
%$E_H$ contains every non-empty subset of $V\setminus\{s\}$ of size at least $2$ 
$E_H:=\{e: e\subseteq V\setminus \{s\} \text{ and } |e|\ge 2\}$ 
with all hyperedge weights being one. 
The hypergraph $(H, w_H)$ is a local connectivity preserving complete h-splitting-off at $s$ from $(G, w_G)$, but $(H, w_H)$ has exponential number of hyperedges although $(G, w_G)$ has only $n$ edges. 
%\knote{Is this example also showing that the Bernath-Kiraly algorithm could give an exponential-sized hypergraph?}
In order to design a polynomial-time algorithm to find a local connectivity preserving complete h-splitting-off at a vertex from a hypergraph, a necessary step is to show the existence of a local connectivity preserving complete h-splitting-off at a vertex from a hypergraph that contains only polynomially many \emph{additional} hyperedges. 
For the star graph $(G,w_G)$ with edge weights $2^{n-1}-1$ mentioned above, the hypergraph $(H'=(V, E_{H'}), w_{H'})$ containing only one hyperedge, namely $E_{H'}:=\{V\setminus \{s\}\}$ with the weight of that hyperedge being $2^{n-1}-1$ is also a local connectivity preserving complete h-splitting-off at $s$ from $(G,w_G)$. 
One of the features of our algorithmic proof of Theorem \ref{thm:complete-splitting-off} is the existence of a local connectivity preserving complete h-splitting-off at a vertex from a hypergraph that contains only polynomially many additional hyperedges. 
%\knote{Say a bit about why existence does not necessarily imply an efficient algorithm to find it---expecially in the weighted case. Would be good to have an example of a weighted hypergraph with polynomial number of hyperedges, but a bad choice of sequence of splitting-off operations leads to an exponential number of hyperedges.}
\end{remark}

As our third main contribution, 
we present two applications of Theorem \ref{thm:complete-splitting-off}. 
%illustrate the usefulness of Theorem \ref{thm:complete-splitting-off} by presenting two applications. 

\paragraph{Application 1: Constructive characterization of $k$-hyperedge-connected hypergraphs.}
For the purposes of this application, 
graphs and hypergraphs will refer to multi-graphs and multi-hypergraphs, respectively.
%we will consider multi-hypergraphs and multi-graphs. 
Let $k$ be a positive integer. A graph $G=(V, E)$ is \emph{$k$-edge-connected} if $d_{G}(X)\ge k$ for every non-empty proper subset $X\subsetneq V$. 
Constructive characterization of $k$-edge-connected graphs is a central problem in graph theory. 
It is well-known that a graph is $1$-edge-connected if and only if it admits a spanning tree. 
Robbins' \cite{Rob39} showed that a graph is $2$-edge-connected if and only if it admits an ear decomposition (see \cite{Frank-book} for definition of ear decomposition). 
Generalizing Robbins' result, 
Lov\'{a}sz \cite{Lov74, Lovasz-problems-book} gave a constructive characterization of $k$-edge-connected graphs for \emph{even} $k$ using his result on global connectivity preserving complete g-splitting-off at a vertex from a graph. 
%An application of this characterization is an alternate proof of Nash-Williams' weak orientation theorem \cite{??} \knote{Do we have a citation for this?}. 
Mader \cite{Mad78} gave a constructive characterization of $k$-edge-connected graphs for \emph{odd} $k$ using his result on local connectivity preserving complete g-splitting-off at a vertex from a graph. 
Motivated by these results, we present a constructive characterization of $k$-hyperedge-connected hypergraphs using our splitting-off result in Theorem \ref{thm:complete-splitting-off}.
A hypergraph $G=(V, E)$ is defined to be \emph{$k$-hyperedge-connected} if $d_{G}(X)\ge k$ for every non-empty proper subset $X\subsetneq V$. 


Both Lov\'{a}sz's and Mader's constructive characterizations of $k$-edge-connected graphs are based on a pinching operation in graphs. Our constructive characterization of $k$-hyperedge-connected hypergraphs is also based on a pinching operation, but our pinching operation is defined for hypergraphs. We define this operation now (see \Cref{fig:hypergraph-pinching-example} for an example). %\knote{Include a figure to explain the pinching operation.}

\begin{definition}\label{def:pinching-operation}
    Let $G = (V, E)$ be a hypergraph and $p,k \in \Z_+$ be positive integers such that $p \leq k$. In \emph{$(k, p)$-pinching} hyperedges of $G$, we obtain a new hypergraph by performing the following sequence of operations:
    \begin{enumerate}
        \item pick $p$ distinct hyperedges $e_1, \ldots, e_p \in E$,
        \item pick $p$ positive integers $t_1, \ldots, t_p \in \Z_+$ such that $\sum_{i=1}^p t_i = k$,
        \item for each $i \in [p]$, choose a partition of the hyperedge $e_i$ into $t_i$ non-empty parts $e_i = \uplus_{j \in [t_i]} f_i^j$,
        \item remove the hyperedges $e_1, \ldots, e_p$ from the hypergraph $G$,
        \item add a new vertex $s$ and hyperedges $\{f_i^j \cup \{s\} : j \in [t_i], i \in [p]\}$ to the hypergraph $G$.
    \end{enumerate}
\end{definition}

%\shubhang{We note that the pinching operation is the \emph{inverse} of the h-splitting-off operation.}\knote{This might be too much detail for the reader to verify here. So, better not to state it here.} 
\begin{figure}[thb]
    \centering
    \includegraphics[width=0.7\textwidth]{hypergraph_pinching_example.png}
    \caption{An example of a $(4,2)$-pinching operation. Here, $t_1 = t_2 = 2$.}
    \label{fig:hypergraph-pinching-example}
\end{figure}
%\knote{Remove $H$ and the box from Figure \ref{fig:hypergraph-pinching-example}.}
With this definition of pinching, we show the following constructive characterization of $k$-hyperedge-connected hypergraphs. 
\begin{restatable}{theorem}{thmConstructiveCharacterization}\label{thm:k-EC-hypergraph-characterization}
    Let $k \in \Z_+$ be a positive integer. A hypergraph $G = (V, E)$ is $k$-hyperedge-connected if and only if $G$ can be obtained by starting from the single vertex hypergraph with no hyperedges and repeatedly applying one of the following two operations:
    \begin{enumerate}
        \item add a new hyperedge over a subset of vertices of the existing hypergraph, and
        \item $(k, p)$-pinching hyperedges of the existing hypergraph for some positive integer $p \leq k$.
    \end{enumerate}
\end{restatable}
Our proof of Theorem \ref{thm:k-EC-hypergraph-characterization} is constructive: i.e., given a $k$-hyperedge-connected hypergraph $G$, our proof gives a polynomial-time algorithm to construct a sequence of hypergraphs $G_0, G_1, G_2, \ldots, G_t$, where $G_0$ is the single vertex hypergraph with no hyperedges, $G_t=G$ and for each $i\in [t]$, the hypergraph $G_i$ is obtained from $G_{i-1}$ by either adding a new hyperedge over a subset of vertices in $G_{i-1}$ or by $(k,p)$-pinching hyperedges in $G_{i-1}$ for some positive integer $p\le k$. 

\begin{remark}
Robbins' constructive characterization of $2$-edge-connected graphs and Lov\'{a}sz's constructive characterization of $2k$-edge-connected graphs find applications in graph orientation problems---e.g., 
%Robbins' result implies that a graph has a strongly connected orientation if and only if it is $2$-edge-connected if and only if it has a strongly connected orientation and
%Lov\'{a}sz's result implies that a graph has a strongly $k$-arc-connected orientation if and only if it is $2k$-edge-connected. 
Robbins' result leads to an algorithm to find a strongly connected orientation of $2$-edge-connected graphs and 
Lov\'{a}sz's result leads to an algorithm to find a strongly $k$-arc-connected orientation of $2k$-edge-connected graphs. In fact, the latter leads to an alternative proof of Nash-Williams' \emph{weak orientation theorem} \cite{NW60}. 
%Alternative proof is here: https://math.mit.edu/~goemans/18438F09/lec17.pdf
%There is no known generalization of Nash-Williams' weak orientation theorem to hypergraphs. Our above characterization might be a helpful tool towards identifying possible generalizations to hypergraphs.
Along the same vein, we hope that our above characterization of $k$-edge-connected hypergraphs might find applications in hypergraph orientation problems. 
\end{remark}


\paragraph{Application 2.1: Steiner Rooted $k$-arc-connected Orientation of Graphs.}
Orienting a graph to achieve high connectivity is a fundamental area in graph theory, combinatorial optimization, and algorithms. Let $G=(V, E)$ be an undirected graph.  An \emph{orientation} $\overrightarrow{G}$ of $G$ is a directed graph obtained by assigning a direction to each edge of $G$. 
Let $G=(V, E)$ be an undirected graph, $T\subseteq V$ be a set of terminals, $r\in T$ be a root vertex, and $k$ be a positive integer. 
%In rooted Steiner $k$-arc-connected orientation of graphs, we are given an undirected graph $G=(V, E)$, a subset $T\subseteq V$ of terminals, a root vertex $r\in T$, and a positive integer $k$. 
An orientation $\overrightarrow{G}$ of $G$ is defined to be \emph{Steiner rooted $k$-arc-connected} if there exist $k$ arc-disjoint paths in $\overrightarrow{G}$ %from $r$ to $t$ 
from $t$ to $r$ 
for every terminal $t\in T\setminus \{r\}$. 
In \textsc{Max Steiner Rooted-Connected Orientation} problem, the goal is to find the maximum integer $k$ and an orientation $\overrightarrow{G}$ of $G$ such that $\overrightarrow{G}$ is Steiner rooted $k$-arc-connected. 
\textsc{Max Steiner Rooted-Connected Orientation} generalizes two classic problems in graph theory: The case of $|T|=2$ is the max edge-disjoint $\{r,v\}$-paths problem and is solved via Menger's theorem. The case of $T=V$ is the max edge-disjoint spanning trees problem and is solved via Tutte and Nash-Williams' theorem \cite{Tu61, NW61}. We mention that both these problems are also generalized by the Steiner Tree Packing problem and the associated Kriesell's conjecture \cite{Kri03, JMS03, Lau07}, but we will not focus on that generalization. 

Kir\'{a}ly and Lau \cite{Kiraly-Lau} introduced the \textsc{Max Steiner Rooted-Connected Orientation}, showed that it is NP-hard, and gave a $2$-approximation via an approximate min-max relation. We state their approximate min-max relation now. 
An undirected graph $G$ is \emph{Steiner $k$-edge-connected} if $\lambda_G(u,v)\ge k$ 
%there exist $k$ edge-disjoint paths between 
for every pair of distinct terminals $u, v\in T$. 
It is clear that if the graph $G$ has a Steiner rooted $k$-arc-connected orientation, then $G$ should be Steiner $k$-edge-connected. However, the converse is not necessarily true. Kir\'{a}ly and Lau observed that if the graph is Steiner $2k$-edge-connected, then it has a Steiner rooted $k$-arc-connected orientation. 

\begin{theorem}[Kir\'{a}ly and Lau \cite{Kiraly-Lau}]\label{thm:rooted-steiner-orientation-in-graphs}
Let $G=(V, E)$ be an undirected graph, $T\subseteq V$ be a subset of terminals, $r\in T$ be the root vertex, and $k$ be a positive integer. If $G$ is Steiner $2k$-edge-connected, then it has a Steiner rooted $k$-arc-connected orientation. 
\end{theorem}

Kir\'{a}ly and Lau observed that Theorem \ref{thm:rooted-steiner-orientation-in-graphs} follows immediately from Nash-Williams' \emph{strong orientation} theorem. 
Nash-Williams' \emph{strong orientation} theorem \cite{NW60} states that every undirected graph $G=(V, E)$ admits an orientation $\overrightarrow{G}$ such that $\lambda_{\overrightarrow{G}}(u,v)\ge \lfloor \lambda_G(u,v)/2\rfloor$ for every distinct $u, v\in V$, where $\lambda_{\overrightarrow{G}}(u,v)$ is the maximum number of arc-disjoint directed paths from $u$ to $v$ in $\overrightarrow{G}$. 
In this work, we give an alternative proof of Theorem \ref{thm:rooted-steiner-orientation-in-graphs} that does not rely on Nash-Williams' strong orientation theorem. Instead, we use Theorem \ref{thm:complete-splitting-off}. 
Our proof strategy is unique since it proves an orientation result for \emph{graphs} using tools developed for \emph{hypergraphs}. 

\begin{remark}
Nash-Williams' proof of the strong orientation theorem \cite{NW60} is a sophisticated inductive argument. Giving a simple and more insightful proof of the strong orientation theorem has been a central topic of interest in graph theory and combinatorial optimization (see \cite{Frank93:Applications-of-submod-functions}). Mader \cite{Mad78} gave a different proof of the strong orientation theorem using his local connectivity preserving splitting-off theorem, but his proof also involved sophisticated technical arguments. Frank \cite{Frank93:Applications-of-submod-functions} condensed the ideas of both Nash-Williams and Mader to present a proof of the strong orientation theorem using Mader's local connectivity preserving splitting-off, but it is still technically complicated. The technical complication in using Mader's local connectivity preserving splitting-off result arises from the assumptions that need to be satisfied by the vertex to be split-off. In contrast, our splitting-off result for hypergraphs (namely, Theorem \ref{thm:complete-splitting-off}) does not need any assumptions on the vertex to be split-off. In light of these considerations, our proof of Theorem \ref{thm:rooted-steiner-orientation-in-graphs} using Theorem \ref{thm:complete-splitting-off} provides hope that Theorem \ref{thm:complete-splitting-off} (or the ideas therein) could be used to give a conceptually simpler proof of Nash-Williams' strong orientation theorem. 
\end{remark}


\paragraph{Application 2.2: Steiner Rooted $k$-hyperarc-connected Orientation of Hypergraphs.}
Orienting \emph{hypergraphs} is also a fundamental area in graph theory and combinatorial optimization (see Frank's book \cite{Frank-book}) 
%. We note that fundamental results in hypergraph orientation have recently been used in coding theory \cite{AGL23}.
with far reaching implications. For example, Woodall's conjecture can be reformulated as a hypergraph orientation problem \cite{egres-open-problem-on-head-disjoint-sco}; moreover,  
hypergraph orientation results have recently been used in coding theory \cite{AGL23}.
%work in coding theory \cite{AGL23} that has exploited hypergraph orientation results). 
%the classic Woodall's conjecture can be phrased as a hypergraph orientation problem. 
Kir\'{a}ly and Lau \cite{Kiraly-Lau} showed that the approximate min-max relation in Theorem \ref{thm:rooted-steiner-orientation-in-graphs} also holds for hypergraphs. To state their result, we need some terminology in hypergraph orientations. 

Let $G=(V, E)$ be a hypergraph. An \emph{orientation} $\overrightarrow{G}=(V, E, \head:E\rightarrow V)$ of $G$ is a directed hypergraph obtained by assigning a unique head vertex $\head(e)\in e$ for each $e\in E$. 
%and setting $\tails(e):=e\setminus \head(e)$ for each $e\in E$. 
A pair $(e, \head(e))$ is denoted as a \emph{hyperarc} with the head of the hyperarc being $\head(e)$ and the tails of the hyperarc being $e\setminus \head(e)$. 
%A hyperarc is a hyperedge with a designated head vertex and other vertices as tail vertices. 
Let $G=(V, E)$ be a hypergraph, $T\subseteq V$ be a set of terminals, $r\in T$ be a root vertex, and $k$ be a positive integer. An orientation $\overrightarrow{G}$ of $G$ is defined to be \emph{Steiner rooted $k$-hyperarc-connected} if there exist $k$ hyperarc-disjoint paths in $\overrightarrow{G}$ 
%from $r$ to $t$ 
from $t$ to $r$ 
for every terminal $t\in T\setminus \{r\}$. Here, a path from $t$ to $r$ in a directed hypergraph is an alternating sequence of distinct vertices and hyperarcs $t=v_1, a_1, v_2, a_2, ..., a_{\ell-1}, v_{\ell}=r$ such that 
%$v_i$ is the head of $a_i$ and $v_{i+1}$ is a tail of $a_i$ \knote{Definition needs to be adjusted to go with convention.} 
$v_i$ is a tail of $a_i$ and $v_{i+1}$ is the head of $a_i$ 
for every $i\in [\ell-1]$. We say that a hypergraph $G$ is \emph{Steiner $k$-hyperedge-connected} if 
%there exist $k$ hyperedge-disjoint paths between 
$\lambda_G(u,v)\ge k$ for 
every pair of distinct terminals $u, v\in T$. 
%Here, a path in a hypergraph is an alternative sequence of distinct vertices and hyperedges $v_1, e_1, v_2, e_2, ..., e_{\ell-1}, v_{\ell}$ such that $v_i, v_{i+1}\in e_i$ for every $i\in [\ell-1]$. 
It is clear that if the hypergraph $G$ has a Steiner rooted $k$-hyperarc-connected orientation, then $G$ should be Steiner $k$-hyperedge-connected. However, the converse is not necessarily true. Kir\'{a}ly and Lau \cite{Kiraly-Lau} showed that if the hypergraph is Steiner $2k$-hyperedge-connected, then it has a Steiner rooted $k$-hyperarc-connected orientation. 

\begin{restatable}[Kir\'{a}ly and Lau \cite{Kiraly-Lau}]{theorem}{thmRootedSteinerArcConnectedOrientationsHypergraphs}\label{thm:rooted-steiner-orientation-in-hypergraphs}
Let $G=(V, E)$ be a hypergraph, $T\subseteq V$ be a subset of terminals, $r\in T$ be the root vertex, and $k$ be a positive integer. If $G$ is Steiner $2k$-hyperedge-connected, then it has a Steiner rooted $k$-hyperarc-connected orientation. 
\end{restatable}
Kir\'{a}ly and Lau's proof of Theorem \ref{thm:rooted-steiner-orientation-in-hypergraphs} was based on careful uncrossing and contractions. In this work, we give an alternative proof of Theorem \ref{thm:rooted-steiner-orientation-in-hypergraphs} using Theorem \ref{thm:complete-splitting-off}. 
%Our proof reveals the source of the $2$ factor in the min-max relation, namely, it is the gap between connectivity and weak-partition-connectivity of hypergraphs (see Definition \ref{defn:weak-partition-connectivity} for the definition of weak-partition-connectivity). 
Our proof of \Cref{thm:rooted-steiner-orientation-in-hypergraphs} reveals the source of the $2$-factor gap in the approximate min-max relation of Kir\'{a}ly and Lau for \textsc{Max Steiner Rooted-Connected Orientation Problem}: it arises from the $2$-factor gap between connectivity and weak-partition-connectivity of hypergraphs (see Definition \ref{defn:weak-partition-connectivity} for the definition of weak-partition-connectivity and Lemma \ref{lem:2k-conn-implies-k-weak-pc}). 

\iffalse
\knote{
\begin{remark}
As mentioned earlier, Theorem \ref{thm:rooted-steiner-orientation-in-graphs} follows immediately from Nash-Williams' strong orientation theorem for graphs. In a similar vein, one would hope that a generalization of Nash-Williams' strong orientation theorem for hypergraphs might help prove Theorem \ref{thm:rooted-steiner-orientation-in-hypergraphs}. However, there is no known generalization of Nash-Williams' strong orientation theorem for hypergraphs (also see Remark \ref{??}). We hope that 
\end{remark}
}
\fi

\begin{remark}
Our proof technique for Theorems \ref{thm:rooted-steiner-orientation-in-graphs} and \ref{thm:rooted-steiner-orientation-in-hypergraphs} using Theorem \ref{thm:complete-splitting-off}---i.e., via the local-connectivity preserving splitting-off operation in hypergraphs---also leads to an alternate proof of Menger's theorem in \emph{undirected} graphs and hypergraphs (edge-disjoint version). For details, we refer the reader to Section \ref{sec:rooted-steiner-connected-orientations} where we discuss a proof of Menger's theorem using \Cref{thm:complete-splitting-off} as a warm-up towards a proof of Theorems \ref{thm:rooted-steiner-orientation-in-graphs} and \ref{thm:rooted-steiner-orientation-in-hypergraphs}. 
\end{remark}

Both Theorems \ref{thm:rooted-steiner-orientation-in-graphs} and \ref{thm:rooted-steiner-orientation-in-hypergraphs} can be extended to weighted graphs/hypergraphs (by considering parallel copies of edges/hyperedges). The weighted version of Theorems \ref{thm:rooted-steiner-orientation-in-graphs} and \ref{thm:rooted-steiner-orientation-in-hypergraphs} can also be shown to admit strongly polynomial-time algorithms using our proof strategy as well as the proof strategy of Kir\'{a}ly and Lau. We avoid stating the weighted versions in the interests of brevity. 


%\knote{Might be good to discuss only the graph version of the result here instead of the hypergraph version. Graph orientations involve far less notation than hypergraph orientations. Readers might follow graph orientations better than hypergraph orientations and might appreciate the flavor of the proof. We can mention the more general result for hypergraph orientations at the end.}

%\subsection{Proof Techniques}\label{sec:proof-techniques}
%In this section, we outline our proof techniques for Theorems \ref{thm:complete-splitting-off}, \ref{thm:k-EC-hypergraph-characterization}, \ref{thm:rooted-steiner-orientation-in-graphs}, and \ref{thm:rooted-steiner-orientation-in-hypergraphs}.

\subsection{Proof Technique for Theorem \ref{thm:k-EC-hypergraph-characterization}}
We outline the proof technique for Theorem \ref{thm:k-EC-hypergraph-characterization}. The reverse direction follows by observing that if a hypergraph is $k$-hyperedge-connected, then both operations in the statement of the theorem preserve $k$-hyperedge-connectivity. We sketch a proof of the forward direction. 
The proof is by induction on the number of hyperedges plus the number of vertices. 
First, suppose that there exists a hyperedge $e \in E$ such that $G - e$ is still $k$‐hyperedge‐connected. We note that deleting the hyperedge $e$ is the inverse of operation (1). Consequently, the proof follows by deleting the hyperedge $e$, using the induction hypothesis on the resulting hypergraph $H$, and then noting that the hypergraph $G$ is obtained from $H$ by operation (1). 
Next, suppose that there does not exist a hyperedge $e \in E$ such that $G-e$ is $k$-hyperedge-connected. We call such a hypergraph to be minimally $k$-hyperedge-connected. In \Cref{lem:minimally-k-EC-hypergraph-has-degree-k-vertex}, we show that a minimally $k$-hyperedge-connected hypergraph contains a vertex $u$ with degree exactly $k$. By \Cref{thm:complete-splitting-off}, there exists a global-connectivity preserving complete h-splitting-off at the vertex $u$ from the hypergraph $G$. Let $H$ be a global-connectivity preserving complete h-splitting-off at the vertex $u$ from the hypergraph $G$. 
We note that complete h-splitting-off at $u$ followed by deletion of the vertex $u$ is the inverse of operation (2) at $u$. 
Consequently, the proof follows by using the induction hypothesis on the hypergraph $H-u$ and then noting that the hypergraph $G$ is obtained from $H-u$ by operation (2).

%\shubhang{sketch for forward direction: First, suppose that there exists a hyperedge $e \in E$ such that $H - e$ is still $k$‐hyperedge‐connected. We note that deleting the hyperedge $e$ is the inverse of operation (1). Consequently, the proof follows by induction. Alternatively, suppose that there does not exist a hyperedge $e \in E$ such that $H -e$ is $k$-hyperedge-connected, i.e. $H$ is minimally-k-hyperedge-connected. In \Cref{lem:minimally-k-EC-hypergraph-has-degree-k-vertex}, we show that there exists a vertex $u \in V$ with degree exactly $k$. By \Cref{thm:complete-splitting-off} there exists a global-connectivity preserving complete h-splitting-off of $u$ from the hypergraph $H$. We note that complete h-splitting-off at $u$ is the inverse of operation (2) at $u$. Consequently, the proof follows by induction.}


\subsection{Proof Technique for Theorems \ref{thm:rooted-steiner-orientation-in-graphs} and \ref{thm:rooted-steiner-orientation-in-hypergraphs}}
We outline the proof technique for Theorem \ref{thm:rooted-steiner-orientation-in-hypergraphs} and will remark after the proof about how it also implies a proof for Theorem \ref{thm:rooted-steiner-orientation-in-graphs}. 
Our proof of Theorem \ref{thm:rooted-steiner-orientation-in-hypergraphs} will be in three steps. 
Let us denote the set of non-terminals as \emph{Steiner vertices}. 
Our first step is to obtain a hypergraph $H = (T, E_H)$ by applying our local connectivity preserving complete h-splitting-off at each Steiner vertex of $G$ (sequentially, in arbitrary order of the Steiner vertices) and deleting the isolated  vertices. 
We note that deleting the isolated vertices ensures that the vertex set of $H$ is the set of terminals $T$. 
Moreover, our local connectivity preserving complete h-splitting-off ensures that the hypergraph $H$ is $2k$-hyperedge-connected since the hypergraph $G$ is Steiner $2k$-hyperedge-connected. 
Our second step is to show that this hypergraph $H=(T, E_H)$ admits a rooted $k$-hyperarc-connected orientation. 
A known characterization for the existence of a rooted $k$-hyperarc-connected orientation of a hypergraph is that the hypergraph is \emph{$k$-weak-partition-connected} (see Definition \ref{defn:weak-partition-connectivity} for the definition of weak-partition-connectivity and Theorem \ref{thm:frank:weak-pc-iff-rooted-hyparc-conn-orientation} for the characterization). We mention that the notion of weak-partition-connectivity in hypergraphs has been used recently in the  context of coding theory \cite{GLSTW22, AGL23}. 
In order to use the characterization for the existence of a rooted $k$-hyperarc-connected orientation of a hypergraph, we relate the connectivity of a hypergraph to its \emph{weak-partition-connectivity} and conclude that if $H$ is $2k$-hyperedge-connected, then it is $k$-weak-partition-connected (see Lemma \ref{lem:2k-conn-implies-k-weak-pc}). 
Consequently, the hypergraph $H$ admits a rooted $k$-hyperarc-connected orientation. 
We note that such an orientation of $H$ is equivalent to a Steiner rooted $k$-hyperarc-connected orientation of $H$ since the vertex set of $H$ is the set $T$ of terminals. 
Our third step is to use this Steiner rooted $k$-hyperarc-connected orientation of $H$ to obtain a Steiner rooted $k$-hyperarc-connected orientation of the hypergraph $G$: we will see that there is a natural way to extend the orientation of hyperedges while reversing the h-splitting-off operations to preserve Steiner rooted $k$-hyperarc-connected property (see Lemma \ref{lem:lifting-Steiner-orientation-through-splitting-off}). This would complete the proof of \Cref{thm:rooted-steiner-orientation-in-hypergraphs}. 

We note that if the hypergraph $G$ is a graph, then the same proof above obtains the required graph orientation, thus proving \Cref{thm:rooted-steiner-orientation-in-graphs}.
In particular, to prove Theorem \ref{thm:rooted-steiner-orientation-in-graphs}, we start from a graph $G=(V, E)$ that is Steiner $2k$-edge-connected, but our local connectivity preserving complete h-splitting-off operations at Steiner vertices results in a hypergraph $H=(T, E_H)$ that is $2k$-hyperedge-connected; by Lemma \ref{lem:2k-conn-implies-k-weak-pc}, the hypergraph $H$ is $k$-weak-partition-connected; now Theorem \ref{thm:frank:weak-pc-iff-rooted-hyparc-conn-orientation} gives a rooted $k$-hyperarc-connected orientation of the resulting hypergraph. Such an orientation is extended to a  Steiner rooted $k$-hyperarc-connected orientation of the \emph{graph} $G=(V, E)$ using Lemma \ref{lem:lifting-Steiner-orientation-through-splitting-off}. Essentially, the proof starts from the given graph, obtains a related hypergraph, orients that hypergraph, and extends that orientation of the hypergraph back into a desired orientation of the given graph. 

Our proof technique for Theorems \ref{thm:rooted-steiner-orientation-in-graphs} and \ref{thm:rooted-steiner-orientation-in-hypergraphs} also leads to an alternate proof of Menger's theorem in \emph{undirected} graphs and hypergraphs (edge-disjoint version)---see Section \ref{sec:rooted-steiner-connected-orientations}. 

\subsection{Proof Technique for Theorem \ref{thm:complete-splitting-off}}
We prove a more general statement that implies Theorem \ref{thm:complete-splitting-off}. We begin with the definitions needed for the more general statement. 
%Let $(H=(V, E), w:E\rightarrow \Z_+)$ be a hypergraph. For $X\subseteq V$, let $B_H(X):=\{e\in X: e\cap X\neq \emptyset\}$. We define the coverage function $b_{(H, w)}:2^V\rightarrow \Z_{\ge 0}$ by $b_{(H,w)}(X):=\sum_{e\in B_H(X)}w(e)$ for every $X\subseteq V$. For a vertex $v\in V$, we use $b_{(H, w)}(v)$ to denote $b_{(H, w)}(\{v\})$. We note that $b_{(H, w)}(v)$ denotes the degree of the vertex $v$. 

\begin{definition}
    Let $V$ be a finite set, $p: 2^V\rightarrow \Z$ be a set function,  and $(H=(V, E), w:E\rightarrow \Z_+)$ be a hypergraph. 
    %We will denote the maximum function value of $p$ by $K_p$, i.e., $K_p\coloneqq \max\{p(X): X\subseteq V\}$. 
    %\knote{Is it really worth defining $K_p$ here?}\kristof{We do not use it in this section, I think we could move its definition to the Preliminaries.}\shubhang{Moved to preliminaries.}
\begin{enumerate}
    \item The set function $p$
    \begin{enumerate}
    \item is \emph{symmetric} if $p(X)=p(V-X)$ for every $X\subseteq V$, and 
    \item is \emph{skew-supermodular} if for every $X, Y\subseteq V$, at least one of the following inequalities hold: 
    \begin{enumerate}
    \item $p(X) + p(Y) \leq p(X\cap Y) + p(X\cup Y)$. %If this inequality holds, then we say that $p$ is \emph{locally supermodular} at $X, Y$.
    \item $p(X) + p(Y) \leq p(X - Y) + p(Y - X)$. %If this inequality holds, then we say that $p$ is \emph{locally negamodular} at $X, Y$.
\end{enumerate}
\end{enumerate}
    \item The coverage function $b_{(H, w)}:2^V\rightarrow \Z_{\ge 0}$ is defined by $b_{(H,w)}(X):=\sum_{e\in B_H(X)}w(e)$ for every $X\subseteq V$, where $B_H(X):=\{e\in E: e\cap X\neq \emptyset\}$ for every $X\subseteq V$. 
    \item The hypergraph $(H, w)$ \emph{weakly covers} the function $p$ if $b_{(H, w)}(X)\ge p(X)$ for every $X\subseteq V$. 
    \item The hypergraph $(H, w)$ \emph{strongly covers} the function $p$ if $d_{(H, w)}(X)\ge p(X)$ for every $X\subseteq V$. 
\end{enumerate}
\end{definition}

If a hypergraph $(H, w)$ strongly covers a function $p:2^V\rightarrow \Z$, then it also weakly covers the function $p$. However, the converse is false -- i.e., a weak cover is not necessarily a strong cover\footnote{For example, consider the function $p: 2^V\rightarrow \Z$ defined by $p(X)\coloneqq 1$ for every non-empty proper subset $X\subsetneq V$ and $p(\emptyset)\coloneqq p(V)\coloneqq 0$, and the hypergraph $(H=(V, E\coloneqq \{\{u\}: u\in V\}), w: E\rightarrow \{1\})$.}.  
Bern\'{a}th and Kir\'{a}ly \cite{Bernath-Kiraly} showed 
that a weak cover of a \emph{symmetric skew-supermodular} function can be converted to a strong cover of the same function by repeated \emph{merging of disjoint hyperedges}. 
%a fundamental result that allows one to convert a weak cover to a strong cover via repeated \emph{merging operations}. %(which will be defined shortly).
We recall their definition of the merging operation, discuss their result, and its significance now. 
\begin{restatable}{definition}{defHyperedgeMerge}\label{def:WeakToStrongCover:hyperedge-merge}
%For a hyperedge-weighted hypergraph $(H=(V, E), w: E\rightarrow \Z_+)$, 
Let $(H=(V, E), w: E\rightarrow \Z_+)$ be a hypergraph. 
We use $H_w$ to denote the unweighted multi-hypergraph over vertex set $V$ containing $w(e)$ copies of every hyperedge $e\in E$.  
By \emph{merging two disjoint hyperedges} of $H_w$, we refer to the operation of replacing them by their union in $H_w$. 
We will say that a hypergraph $(G=(V, E_G), c: E_G\rightarrow \Z_+)$ is obtained from $(H, w)$ by \emph{merging hyperedges} if the multi-hypergraph $G_c$ is obtained from the multi-hypergraph $H_w$ by repeatedly merging two disjoint hyperedges in the current hypergraph. 
\end{restatable}
Bern\'{a}th and Kir\'{a}ly showed the following result: 
%Theorem \ref{thm:bernath-kiraly} stated below. Th
\begin{theorem}[Bern\'{a}th and Kir\'{a}ly \cite{Bernath-Kiraly}]\label{thm:bernath-kiraly}
    Let $(H=(V, E), w: E\rightarrow \Z_+)$ be a hypergraph and $p:2^V\rightarrow \Z$ be a symmetric skew-supermodular function such that 
    %$(H,w)$ weakly covers $p$. Then, there exists a hypergraph $(G^*=(V, E^*),w^*: E^*\rightarrow \Z_+)$ obtained by merging hyperedges of the hypergraph $(H,w)$ such that $(G,c)$ strongly covers $p$. 
$b_{(H, w)}(X) \geq p(X)$ for every $X \subseteq V$. 
Then, there exists a hypergraph $\left(H^* = (V, E^*), w^*:E^* \rightarrow\Z_+\right)$ 
such that
\begin{enumerate}[label=$(\arabic*)$]
    \item $d_{(H^*, w^*)}(X) \geq p(X)$ for every $X\subseteq V$ and 
    \item the hypergraph $(H^*,w^*)$ is obtained by merging hyperedges of the hypergraph $(H, w)$.
\end{enumerate}
\end{theorem}

We observe that Theorem \ref{thm:bernath-kiraly} can be used to prove the existential version of Theorem \ref{thm:complete-splitting-off}: namely, for every hypergraph $(G=(V, E), w_G: E\rightarrow \Z_+)$ and a vertex $s\in V$, \emph{there exists} a local connectivity preserving complete h-splitting-off at $s$ from $(G,w_G)$. This can be shown by setting up the hypergraph $(H,w)$ and the function $p$ suitably based on $(G,w_G)$ and using Theorem \ref{thm:bernath-kiraly} (see the first two paragraphs of the proof of Theorem \ref{thm:hypergraphs-splitting-off-stronger-form} in Section \ref{sec:hypergraph-splitting-off-proof}). 
We emphasize that this conclusion regarding hypergraph splitting-off from Bern\'{a}th and Kir\'{a}ly's result was not known before in the literature and is one of our contributions. 
\begin{remark}
    We were also able to prove the existential version of Theorem \ref{thm:complete-splitting-off} using \emph{element-connectivity preserving reduction operations} \cite{BCCKM-private} (see \cite{CK14} for the definition of element-connectivity and the notion of element-connectivity preserving reduction operations) -- we omit the details of this alternate proof in the interests of brevity. The alternate proof does not seem to be helpful for the purposes of a strongly polynomial time algorithm. In fact, it remains open to design a strongly polynomial-time algorithm to perform \emph{complete} element-connectivity preserving reduction operations in the weighted setting \cite{CK14}. 
\end{remark}

We recall that existence of a local connectivity preserving complete h-splitting-off at a vertex from a hypergraph does not immediately imply a polynomial-time algorithm---see the example in Remark \ref{remark:exp-sized-conn-preserving-splitting-off}. 
%algorithm. In particular, there exists a local connectivity preserving complete h-splitting-off at a vertex $s$ from a hypergraph $(G,w_G)$ that contains $2^{O(n)}$ hyperedges although $G$ has only $O(n)$ hyperedges, where $n$ is the number of vertices in $G$ (see the discussion after Theorem \ref{thm:complete-splitting-off} for such an example). 
However, the above-mentioned proof of 
existence of a local-connectivity preserving complete splitting-off at an arbitrary vertex from a hypergraph (i.e., existential version of Theorem \ref{thm:complete-splitting-off}) 
via Theorem \ref{thm:bernath-kiraly} suggests a natural approach towards designing a strongly polynomial time algorithm to find a local-connectivity preserving complete splitting-off at a given vertex from a given hypergraph: it suffices to prove a constructive version of Theorem \ref{thm:bernath-kiraly} via a strongly polynomial-time algorithm. Towards this end, the example in Remark \ref{remark:exp-sized-conn-preserving-splitting-off} suggests a necessary structural step towards a strongly polynomial-time algorithmic version of Theorem \ref{thm:bernath-kiraly}: we need to show Theorem \ref{thm:bernath-kiraly} with the extra conclusion that the number of additional hyperedges in $H^*$ is polynomial in the number of hyperedges and vertices in $H$. 

\iffalse
We observe that Theorem \ref{thm:bernath-kiraly} can be used to prove the existential version of Theorem \ref{thm:complete-splitting-off}: namely, for every hypergraph $(G=(V, E), w_G: E\rightarrow \Z_+)$ and a vertex $s\in V$, \emph{there exists} a local connectivity preserving complete h-splitting-off at $s$ from $(G,w_G)$. This can be shown by setting up the hypergraph $(H,w)$ and the function $p$ suitably based on $(G,w_G)$ and using Theorem \ref{thm:bernath-kiraly} (see the first two paragraphs of the proof of Theorem \ref{thm:hypergraphs-splitting-off-stronger-form} in Section \ref{sec:hypergraph-splitting-off-proof}). 
We emphasize that this conclusion regarding hypergraph splitting-off from Bern\'{a}th and Kir\'{a}ly's result was not known before in the literature and is one of our contributions. 
\knote{The existential version of Theorem \ref{thm:complete-splitting-off} can also be derived from element-connectivity preserving reduction operation. However, neither of these derivations lead to a strongly polynomial-time algorithm.}
However, we recall that existence does not immediately imply a polynomial-time algorithm---see the example in Remark \ref{remark:exp-sized-conn-preserving-splitting-off}. 
%algorithm. In particular, there exists a local connectivity preserving complete h-splitting-off at a vertex $s$ from a hypergraph $(G,w_G)$ that contains $2^{O(n)}$ hyperedges although $G$ has only $O(n)$ hyperedges, where $n$ is the number of vertices in $G$ (see the discussion after Theorem \ref{thm:complete-splitting-off} for such an example). 
This example suggests a necessary structural step towards a strongly polynomial-time algorithmic version of Theorem \ref{thm:bernath-kiraly}: we need to show Theorem \ref{thm:bernath-kiraly} with the extra conclusion that the number of additional hyperedges in $H^*$ is polynomial in the number of hyperedges and vertices in $H$. 
\fi

Bern\'{a}th and Kir\'{a}ly proved Theorem \ref{thm:bernath-kiraly} in the context of a reduction between certain hypergraph connectivity augmentation problems. For that reduction, the existential version of Theorem \ref{thm:bernath-kiraly} is sufficient. However, for the purposes of our application to hypergraph splitting-off, we need an algorithmic version of Theorem \ref{thm:bernath-kiraly}. 
Bern\'{a}th and Kir\'{a}ly's 
proof of Theorem \ref{thm:bernath-kiraly} is in fact algorithmic, but the run-time of the associated algorithm is not necessarily polynomial.  
Their proof implies that the number of additional hyperedges in the hypergraph $(H^*, c^*)$ returned by their algorithm is at most  $\sum_{e\in E}w(e)$ (i.e., $|E^*|-|E|\le \sum_{e\in E}w(e)$) and the run-time of the algorithm is $O(\sum_{e\in E}(|e|+w(e))$. In particular, their run-time is polynomial only if the input weights are given in unary. Moreover, the exponential-sized hypergraph given in the example in Remark \ref{remark:exp-sized-conn-preserving-splitting-off} could indeed arise as a consequence of their algorithm. 
%\knote{Add a reference to the example that appears after Theorem 1.1.}

We address both the structural and the algorithmic issues mentioned above by proving a stronger algorithmic 
%In order to design a strongly polynomial-time algorithm for local connectivity preserving complete h-splitting-off at a vertex $s$ from $(G,w)$, we prove a stronger algorithmic 
version of Theorem \ref{thm:bernath-kiraly}. 
In order to phrase an algorithmic version of Theorem \ref{thm:bernath-kiraly}, we need suitable access to the function $p$. Bern\'{a}th and Kir\'{a}ly \cite{Bernath-Kiraly} suggested access to a certain function maximization oracle associated with the function $p$ that we describe below. 

\begin{definition}
    Let $p: 2^V\rightarrow \Z$ be a set function.    $\functionMaximizationOracleStrongCover{p}\left(\left(G_0, c_0\right), S_0, T_0\right)$ takes as input a hypergraph $(G_0=(V, E_0), c_0: E_0\rightarrow \mathbb{Z}_{+})$ and disjoint sets $S_0, T_0\subseteq V$, and returns a tuple $(Z, p(Z))$, where $Z$ is an optimum solution to the following problem:
        \begin{align*}
        \max & \left\{p(Z)-d_{(G_0, c_0)}(Z): S_0\subseteq Z\subseteq V-T_0\right\}. \tag{\functionMaximizationOracleStrongCover{p}}\label{tag:maximization-oracle-strong-cover}
        \end{align*}
\end{definition}
For the purposes of our application (namely local connectivity preserving complete h-splitting-off at a vertex from a hypergraph), the above-mentioned function maximization oracle can be implemented to run in strongly polynomial time (see Lemma \ref{lemma:helper-for-p-max-oracle}). Using the above mentioned oracle, we prove the following algorithmic version of Theorem \ref{thm:bernath-kiraly}.

\begin{restatable}{theorem}{thmWeakToStrongCover}\label{thm:WeakToStrongCover:main}
Let $(H=(V, E), w: E\rightarrow \Z_+)$ be a hypergraph and $p:2^V\rightarrow \Z$ be a symmetric skew-supermodular function such that 
%$(H,w)$ weakly covers $p$. 
%Then, there exists a hypergraph $(G^*=(V, E^*),w^*: E^*\rightarrow \Z_+)$ obtained by merging hyperedges of the hypergraph $(H,w)$ such that $(G,c)$ strongly covers $p$ and \[|E^*|-|E|=O(|V|).\]
%Let $p:2^V\rightarrow\Z$ be a symmetric skew-supermodular function. Let $\left(H = (V, E), w:E \rightarrow\Z_+\right)$ be a hypergraph such that 
$b_{(H, w)}(X) \geq p(X)$ for every $X \subseteq V$. 
Then, there exists a hypergraph \linebreak $\left(H^* = (V, E^*), w^*:E^* \rightarrow\Z_+\right)$ 
such that
\begin{enumerate}[label=$(\arabic*)$, ref=(\arabic*)]
    \item \label{thm:WeakToStrongCover:main:(1)} $d_{(H^*, w^*)}(X) \geq p(X)$ for every $X\subseteq V$,
    \item \label{thm:WeakToStrongCover:main:(2)} the hypergraph $(H^*,w^*)$ is obtained by merging hyperedges of the hypergraph $(H, w)$, and
    \item \label{thm:WeakToStrongCover:main:(3)} $|E^*| - |E| = O(|V|)$.
\end{enumerate}
%Furthermore, given a hypergraph $(H=(V, E), w: E\rightarrow \Z_+)$ and access to \functionMaximizationOracleStrongCover{p} of a symmetric skew-supermodular function $p: 2^V\rightarrow \Z$ where $b_{(H, w)}(X)\ge p(X)$ for every $X\subseteq V$, there exists an algorithm that runs in time $\poly(|V|, |E|)$ using $\poly(|V|)$ queries to \functionMaximizationOracleStrongCover{p} to return a hypergraph $\left(G^* = (V, E^*), w^*:E^* \rightarrow\Z_+\right)$ satisfying the above three properties. 
Furthermore, given a hypergraph $(H=(V, E), w: E\rightarrow \Z_+)$ and access to \functionMaximizationOracleStrongCover{p} of a symmetric skew-supermodular function $p: 2^V\rightarrow \Z$ where $b_{(H, w)}(X)\ge p(X)$ for every $X\subseteq V$, there exists an algorithm that runs in $O(|V|^3(|V| + |E|)^2)$ time using $O(|V|^3(|V|+|E|))$ queries to \functionMaximizationOracleStrongCover{p} and returns a hypergraph $\left(H^* = (V, E^*), w^*:E^* \rightarrow\Z_+\right)$ satisfying the above three properties. The run-time includes the time to construct the hypergraphs used as input to the queries to \functionMaximizationOracleStrongCover{p}.
Moreover, for each query to \functionMaximizationOracleStrongCover{p}, the hypergraph $(G_0, c_0)$ used as input to the query has $O(|V|)$ vertices and $O(|V|+|E|)$ hyperedges.
\end{restatable}

%Theorem \ref{thm:complete-splitting-off} follows immediately from Theorem \ref{thm:WeakToStrongCover:main} (see Section \ref{sec:hypergraph-splitting-off-proof}). 
%We emphasize that the run-time in Theorem \ref{thm:WeakToStrongCover:main} includes the time to construct the hypergraphs $(G_0, c_0)$ that are used as inputs to \functionMaximizationOracleStrongCover{p}. 
\Cref{thm:WeakToStrongCover:main} is a strengthening of Theorem \ref{thm:bernath-kiraly} in two ways. Firstly, our theorem shows the existence of a hypergraph that not only satisfies properties \ref{thm:WeakToStrongCover:main:(1)} and \ref{thm:WeakToStrongCover:main:(2)}, but also satisfies property \ref{thm:WeakToStrongCover:main:(3)} -- i.e., the number of additional hyperedges in the returned hypergraph is \emph{linear} in the size of the vertex set. Secondly, our \Cref{thm:WeakToStrongCover:main} shows the existence of a strongly polynomial-time algorithm that returns a hypergraph satisfying the three properties. 
%Although Bern\'{a}th and Kir\'{a}ly's result is existential, their proof is algorithmic. However, their proof is based on induction on $\sum_{e\in E}w(e)$ and consequently, the straightforward bound on the number of additional hyperedges as well as the run-time is $O(\sum_{e\in E}w(e))$. Thus, the best run-time bound known for their algorithm was only pseudo-polynomial. 
Our main contribution is modifying Bern\'{a}th and Kir\'{a}ly's algorithm and analyzing the modified algorithm to bound the number of additional hyperedges and the run-time. 
%Our main contribution is analyzing an \emph{adaptation} of Bern\'{a}th and Kir\'{a}ly's algorithm to bound the number of additional hyperedges and the run-time. 
We mention that property \ref{thm:WeakToStrongCover:main:(3)} cannot be tightened to guarantee that $|E^*-E|=O(|V|)$ -- we were able to construct an example where 
%the only strong cover $(H^*=(V, E^*), w^*: E^*\rightarrow \Z_+)$ of a function $p: 2^V\rightarrow \Z$ that is obtained by merging hyperedges of a weak cover hypergraph $(H=(V, E), w: E\rightarrow \Z_+)$ of the function $p$ is such that 
$|E^*-E|=\Omega(|V|^2)$ (see \Cref{sec:appendix:tight-example-weak-to-strong}). 

%Our \Cref{thm:WeakToStrongCover:main} shows that there exists a hypergraph $(G, c)$ obtained from $(H, w)$ by merging hyperedges such that $(G, c)$ strongly covers $p$ and the number of additional hyperedges in $G$ compared to $H$ is $O(|V|)$ (i.e., $|E_G|-|E_H|=O(|V|)$); moreover, our proof is constructive and leads to a strongly polynomial-time algorithm to find such a strong cover $(G, c)$ of $p$ assuming access to the function $p$ via \functionMaximizationOracleStrongCover{p}. This result leads to a strongly polynomial-time algorithm for hypergraph splitting-off (see \Cref{sec:splitting-off}). 

\Cref{thm:WeakToStrongCover:main} immediately leads to  a proof of \Cref{thm:complete-splitting-off} (see Theorem \ref{thm:hypergraphs-splitting-off-stronger-form} and its proof in Section \ref{sec:hypergraph-splitting-off-proof}). Instead of using \Cref{thm:WeakToStrongCover:main} as a black-box, if we delve into the proof of it in the context of the proof of \Cref{thm:complete-splitting-off}, we obtain the following theorem: 

%In the proof of \Cref{thm:complete-splitting-off}, we used \Cref{thm:WeakToStrongCover:main} as a black-box. Instead, delving into the proof of \Cref{thm:WeakToStrongCover:main} in the context of the proof of \Cref{thm:complete-splitting-off} reveals that the following theorem also holds. 
\begin{theorem}\label{thm:splitting-off-or-trim-to-preserve-connectivity}
    Let $(G=(V, E), w: E\rightarrow \Z_+)$ be a hypergraph and $s\in V$. Then, there exists a hypergraph $(G'=(V, E'), w': E'\rightarrow \Z_+)$ obtained 
    by applying a h-splitting-off operation at $s$ from $(G, w)$ 
    %from $(G, w)$ by a h-splitting-off operation 
    such that $\lambda_{(G', w')}(u, v)=\lambda_{(G, w)}(u, v)$ for every distinct $u, v\in V\setminus \{s\}$. 
\end{theorem}
We omit the proof of \Cref{thm:splitting-off-or-trim-to-preserve-connectivity} in the interests of brevity. 
\Cref{thm:splitting-off-or-trim-to-preserve-connectivity} closely resembles the \emph{existential edge splitting-off} results of Lov\'{a}sz \cite{Lov74, Lovasz-problems-book} and Mader \cite{Mad78} for graphs. Lov\'{a}sz's and Mader's  existential edge splitting-off results for graphs are important since they have been used to simplify the proofs of fundamental results in graph theory---e.g., Nash-Williams' Strong Orientation Theorem. 
On the other hand, \Cref{thm:splitting-off-or-trim-to-preserve-connectivity} does not immediately imply a strongly polynomial-time algorithm for finding a local connectivity preserving complete h-splitting off at a vertex from a given weighted hypergraph. So, \Cref{thm:complete-splitting-off} may be useful in algorithmic contexts while \Cref{thm:splitting-off-or-trim-to-preserve-connectivity} may be useful in graph-theoretical contexts. 

\subsection{Proof Technique for \Cref{thm:WeakToStrongCover:main}}
In this section, we describe our proof technique for the existential result in \Cref{thm:WeakToStrongCover:main}. 
The algorithmic results in that theorem follow from the existential result using standard algorithmic tools for submodular functions, so we focus only on outlining a proof of the existential result. 
Let $(H=(V, E), w: E\rightarrow \Z_+)$ be a hypergraph and $p: 2^V\rightarrow \Z_+$ be a symmetric skew-supermodular function such that $(H, w)$ weakly covers the function $p$. Our goal is to show that there exists a hypergraph $(H^*=(V, E^*), w^*: E\rightarrow \Z_+)$ such that 
\begin{enumerate}[label=(\arabic*)]
    \item \hypertarget{WeakToStrong(1)}{$(H^*, w^*)$ strongly covers the function $p$},
    \item \hypertarget{WeakToStrong(2)}{$(H^*, w^*)$ is obtained by merging hyperedges of the hypergraph $(H, w)$}, and
    \item \hypertarget{WeakToStrong(3)}{$|E^*|-|E|=O(|V|)$}. 
\end{enumerate}
% \hypertarget{WeakToStrong(1)}{(1) $(H^*, w^*)$ strongly covers the function $p$}, \hypertarget{WeakToStrong(2)}{(2) $(H^*, w^*)$ is obtained by merging hyperedges of the hypergraph $(H, w)$}, and \hypertarget{WeakToStrong(3)}{(3) $|E^*|-|E|=O(|V|)$}. 

\paragraph{Preliminaries.} We define a set $X\subseteq V$ to be $(p, H, w)$-tight if $b_{(H, w)}(X) = p(X)$. For a function $p$ and hypergraph $(H, w)$, let $T_{p, H, w}$ denote the family of $(p,H, w)$-tight sets and let $\mathcal{T}_{p, H, w}$ be the family of inclusionwise maximal sets in $T_{p, H, w}$. We will need the following two operations: 
\begin{enumerate}[label=(\roman*)]
    \item For hyperedges $e, f\in E$ and a positive integer  $\alpha \le \min\{w(e), w(f)\}$, the operation \merge$((H, w), e, f, \alpha))$ returns the hypergraph obtained from $(H, w)$ by decreasing the weight of hyperedges $e$ and $f$ by $\alpha$ and increasing the weight of the hyperedge $e\cup f$ by $\alpha$. All hyperedges with zero weight are discarded.
    \item For a hyperedge $e\in E$ and a positive integer $\alpha\le w(e)$, the operation \reduce$((H, w), e, \alpha)$ returns the hypergraph obtained by decreasing the weight of the hyperedge $e$ by $\alpha$. All hyperedges with zero weight are discarded.
\end{enumerate}
 
%The sets in $\mathcal{T}_{p, H, w}$ can be shown to be pairwise disjoint using the hypotheses that $p$ is skew-supermodular and $(H,w)$ weakly covers $p$ \shubhang{The family $\maximalTightSetFamily{p, H, w}$ being disjoint requires a stronger hypothesis (see \Cref{lem:weak-to-strong:uncrossing:maximal-tight-set-family-disjoint}). This is the main motivation for enforcing that every hyperedge is contained in a tight set and every vertex is contained in a hyperedge during each recursive call of the algorithm. Consequently, this is also why we immediately do reduction and contraction steps after the merge step. This insight could also count as one of our key technical contributions.}. 

%For two hypergraphs $(G=(V, E_G), c_G:E_G\rightarrow \Z_+)$ and $(H=(V, E_H), w_H: E_H\rightarrow \Z_+)$ on the same vertex set $V$, we define the hypergraph $(G+H=(V, E_{G+H}), c_G+w_H)$ as the hypergraph with vertex set $V$ and hyperedge set $E_{G+H}\coloneqq E_G\cup E_H$ with the weight of every hyperedge $e\in E_G\cap E_H$ being $c_G(e) + w_H(e)$, the weight of every hyperedge $e\in E_G\setminus E_H$ being $c_G(e)$,  and the weight of every hyperedge $e\in E_H\setminus E_G$ being $w_H(e)$. 


\paragraph{Algorithm of \cite{Bernath-Kiraly}.} Our proof of the existential result builds on the techniques of Bern\'{a}th and Kir\'{a}ly \cite{Bernath-Kiraly} who proved the existence of a hypergraph $(H^*=(V, E^*), w^*: E^*\rightarrow \Z_+)$ satisfying properties \hyperlink{WeakToStrong(1)}{(1)} and \hyperlink{WeakToStrong(2)}{(2)}, so we briefly recall their techniques. 
We present the algorithmic version of their proof since it will be useful for our purposes. 

The proof in \cite{Bernath-Kiraly} is inductive, and consequently, the algorithm implicit in their proof is recursive. The algorithm takes as input a hypergraph $((H=(V, E), w: E\rightarrow \Z_+)$ and a symmetric skew-supermodular function $p: 2^V\rightarrow \Z$ such that the hypergraph $(H,w)$ weakly covers the function $p$. 
%The algorithm returns a hypergraph $(H^*, w^*)$ that satisfies properties (1) and (2).
If $w(E) = 0$, then the algorithm is in its base case and returns the empty hypergraph. 
Otherwise, $w(E) > 0$; the algorithm chooses an arbitrary hyperedge $e \in E$ and defines hypergraphs $(H_0, w_0)$ and $(H', w')$ and the function $p'$ by considering two cases.
First, suppose that the hyperedge $e$ is not contained in any set of the family $\maximalTightSetFamily{p, H, w}$. 
In this case, the algorithm defines $(H_0, w_0)$ to be the hypergraph on vertex set $V$ consisting of a single hyperedge $e$ with $w_0(e) = 1$, constructs the hypergraph $(H', w'):=\reduce((H, w), e, 1)$, and defines the function $p' := p - d_{(H_0, w_0)}$. Second, suppose that the hyperedge $e$ is contained in some set $X \in \mathcal{T}_{p, H, w}$. It can be shown that there exists a hyperedge $f \in E$ such that $f \subseteq V - X$. In this case, the algorithm defines $(H_0, w_0)$ to be the empty hypergraph on vertex set $V$, constructs the hypergraph $(H', w'):=\merge((H, w), e, f, 1)$, and defines the function $p' := p$. In both cases, the algorithm recurses on the inputs $(H',w')$ and $p'$ to obtain a hypergraph $(H^*_0, w^*_0)$ and returns the hypergraph $(H^*, w^*) = (H^*_0 + H_0, w^*_0 + w_0)$. 
Here, the hyperedges of $H^*$ are the union of the hyperedges of $H^*_0$ and $H_0$ with the weight $w^*(e)$ of a hyperedge $e$ being the sum of the weights $w^*_0(e)+w_0(e)$ if the hyperedge $e$ is present in both $H^*_0$ and $H_0$, being $w^*_0(e)$ if the hyperedge $e$ is present only in $H^*_0$, and being $w_0(e)$ if the hyperedge $e$ is present only in $H_0$. 
\iffalse
\knote{Edit this to be consistent with BK's proof.}\shubhang{Rest of this paragraph can be removed.}
First, suppose that there exists a hyperedge $e\in E$ that is not contained in any set of the family $\maximalTightSetFamily{p, H, w}$. 
In this case, the algorithm defines $(H_0, w_0)$ to be the hypergraph on vertex set $V$ consisting of a single hyperedge $e$ with $w_0(e) = 1$,
constructs the hypergraph $(H', w'):=\reduce((H, w), e, 1)$, and defines the function $p' := p - d_{(H_0, w_0)}$.
Second, suppose that every hyperedge $e\in E$ is contained in some set of the family $\mathcal{T}_{p, H, w}$. It can be shown that there exists a pair of disjoint hyperedges $e, f\in E$ contained in distinct sets of the family $\maximalTightSetFamily{p, H, w}$.
In this case, the algorithm defines $(H_0, w_0)$ to be the empty hypergraph on vertex set $V$, constructs the hypergraph $(H', w'):=\merge((H, w), e, f, 1)$, and defines the function $p' := p$.
In both cases, the algorithm recurses on the inputs $(H',w')$ and $p'$ to obtain a hypergraph $(H^*_0, w^*_0)$ and returns the hypergraph $(H^*, w^*) = (H^*_0 + H_0, w^*_0 + w_0)$. 
Here, the hyperedges of $H^*$ are the union of the hyperedges of $H^*_0$ and $H_0$ with the weight $w^*$ of a hyperedge $e$ being the sum of the weights $w^*_0(e)+w_0(e)$ if the hyperedge $e$ is present in both $H^*_0$ and $H_0$, being $w^*_0(e)$ if the hyperedge $e$ is present only in $H^*_0$, and being $w_0(e)$ if the hyperedge $e$ is present only in $H_0$. 
\fi

 We note that $w'(E') = w(E) - 1$. Furthermore, it can be shown that the function $p'$ is symmetric skew-supermodular and the hypergraph $(H', w')$ weakly covers the function $p'$. Consequently, by induction on $w(E)$, the algorithm can be shown to terminate in $w(E)$ recursive calls and returns a hypergraph satisfying properties \hyperlink{WeakToStrong(1)}{(1)} and \hyperlink{WeakToStrong(2)}{(2)}. Moreover, the number of additional hyperedges in the returned hypergraph is at most the
number of recursive calls where the \merge operation is performed, which is also at most $w(E)$. 
%\knote{Add a reference to the example that appears after Theorem 1.1.}
Thus, in order
to reduce the number of additional hyperedges and to design a strongly polynomial-time algorithm, our goal is to
reduce the recursion depth of the algorithm. We emphasize that the recursion depth of Bern\'{a}th and Kir\'{a}ly's algorithm could indeed be exponential (the exponential sized example mentioned in Remark \ref{remark:exp-sized-conn-preserving-splitting-off} could arise in the execution of their algorithm), so we need to necessarily modify their algorithm. 

\paragraph{Preprocessing for Additional Structure.}
%Our first technical insight is that by performing a preprocessing step (see Section \ref{sec:WeakToStrongCover}), we can assume that the inputs to our algorithm satisfy additional conditions. \shubhang{Furthermore, appropriate modifications to Bern\'{a}th and Kir\'{a}ly's algorithm ensures that these conditions hold during every recursive call of our recursive algorithm described in the paragraph below.}
%We now describe our adaptation of the above-mentioned Bern\'{a}th and Kir\'{a}ly's algorithm to reduce the recursion depth. 
Similar to Bern\'{a}th and Kir\'{a}ly's algorithm, our algorithm also takes as input a hypergraph $(H=(V, E), w: E\rightarrow \Z_+)$ and a symmetric skew-supermodular function $p: 2^V\rightarrow \Z$ such that the hypergraph $(H,w)$ weakly covers the function $p$. 
However, unlike Bern\'{a}th and Kir\'{a}ly's algorithm, our algorithm performs a preprocessing step so that the inputs $(H,w)$ and $p$ satisfy the following two additional conditions: 
\begin{enumerate}[label=(\alph*)]
    \item \hypertarget{WeakToStrong(a)}{every hyperedge $e\in E$ is contained in some set of $\mathcal{T}_{p, H, w}$} and 
    \item \hypertarget{WeakToStrong(b)}{the degree of every vertex in $(H, w)$ is non-zero.}
\end{enumerate}
%These two additional conditions will be inductively maintained by our algorithm. 
As a consequence of these additional conditions, the family $\maximalTightSetFamily{p, H, w}$ will be a \emph{disjoint family},  
%during every recursive call of our algorithm, 
a property that we leverage heavily throughout our analysis.
% \hypertarget{WeakToStrong(a)}{(a) every hyperedge $e\in E$ is contained in some set of $\mathcal{T}_{p, H, w}$} and \hypertarget{WeakToStrong(b)}{(b) the degree of every vertex in $(H, w)$ is non-zero} -- we show that these conditions are without loss of generality since an arbitrary instance can be reduced to an instance satisfying these conditions (see Section \ref{sec:WeakToStrongCover}).
Furthermore, we modify Bern\'{a}th and Kir\'{a}ly's algorithm to ensure that these conditions hold during every recursive call.

\paragraph{Our Algorithm.}
We now describe our modification of the above-mentioned Bern\'{a}th and Kir\'{a}ly's algorithm to reduce the recursion depth. Our algorithm is also recursive and its 
%, and returns a hypergraph $(H^*, w^*)$ that satisfies properties (1), (2) and (3).
base case is the same as that of Bern\'{a}th and Kir\'{a}ly's algorithm (i.e., $w(E)=0$). During recursive cases (i.e., if $w(E) > 0$), instead of performing one of the two (i.e., \merge or  \reduce) operations, our algorithm performs both operations in a sequential fashion. In particular, we find a pair of disjoint hyperedges $e, f\in E$ contained in distinct sets of $\maximalTightSetFamily{p, H, w}$ (such a pair exists by condition \hyperlink{WeakToStrong(a)}{(a)} and the arguments of Bern\'{a}th and Kir\'{a}ly mentioned above). Next, 
instead of performing one \merge operation (as was done by Bern\'{a}th and Kir\'{a}ly's algorithm), we perform as many  \merge operations as possible %\shubhang{(denoted by $\alpha^M \in \Z_+$)} 
using the hyperedges $e$ and $f$. 
Formally, let $$\alpha^M := \max\left\{\alpha \in \Z_+ : \text{hypergraph returned by } \merge((H, w), e, f, \alpha) \text{ weakly covers } p\right\}.$$
We let $(H^M, w^M):=\textsc{Merge}((H, w), e, f, \alpha^M)$ and $p^M:=p$.
%We describe this formally now: we obtain $(H^M, w^M):=\textsc{Merge}((H, w), e, f, \alpha^M)$ and $p^M:=p$, where $\alpha^M\le \min\{w(e), w(f)\}$ is the largest integer such that $(H^M, w^M)$ weakly covers the function $p^M$. 
Next, instead of recursing on $((H^M, w^M), p^M)$ (as was done by Bern\'{a}th and Kir\'{a}ly's algorithm), we perform as many \reduce operations as possible 
%\shubhang{(denoted by $\alpha^R \in \Z_{\geq 0}$)} 
on the newly created hyperedge $e\cup f$.
Formally, let $$\alpha^R := \max\left\{\alpha \in \Z_{\ge 0}: \text{hypergraph returned by } \reduce((H^M, w^M), e \cup f, \alpha) \text{ weakly covers } (p^M - d_{(H_0^\alpha, w_0^\alpha)})\right\},$$
%where $(H_0^{\alpha}:=(V, E_0^{\alpha}:=\{e\cup f\}),w_0^{\alpha}: E_0^{\alpha}\rightarrow \{\alpha\})$ 
where $(H_0^{\alpha}, w_0^{\alpha})$ denotes
the hypergraph on vertex set $V$ consisting of a single hyperedge $e\cup f$ with $w_0^{\alpha}(e\cup f) = \alpha$. 
We construct the hypergraph 
$(H_0, w_0):=(H_0^{\alpha^R}, w_0^{\alpha^R})$, the hypergraph $(H^R, w^R):=\reduce((H^M, w^M), e\cup f, \alpha^R)$ and define the function $p^R:=p^M - d_{(H_0^{\alpha^R}, w_0^{\alpha^R})}$.
%We describe this formally now: 
%for $\alpha\le w^M(e\cup f)$, let $(H_0^{\alpha}:=(V, E_0^{\alpha}:=\{e\cup f\}),w_0^{\alpha}: E_0^{\alpha}\rightarrow \{\alpha\})$. 
%Let $(H^R, w^R):=\textsc{Reduce}((H^M, w^M), e\cup f, \alpha^R)$ and $p^R:=p^M-d_{(H_0^{\alpha^R},w_0^{\alpha^R})}$, where $\alpha^R\le w^M(e\cup f)$ is the largest integer such that $(H^R, w^R)$ weakly covers the function $p^R$. 
%and $(H^R, w^R):=(H^R_{\alpha_R}, w^R_{\alpha_R})$, 
%and define the function $p^R:=p^R_{\alpha_R}$. 
This immediate reduce step ensures that the hypergraph $(H^R, w^R)$ and the function $p^R$ satisfy condition \hyperlink{WeakToStrong(a)}{(a)} -- i.e., every hyperedge in $(H^R, w^R)$ is contained in some set of $\mathcal{T}_{p^R, H^R, w^R}$. 
Finally, we compute sets $\zeros:=\{u\in V: b_{(H^R, w^R)}(u)=0\}$ and $V':=V-\zeros$, hypergraph $(H':=(V', E':=E^R), w':=w^R)$, and define the function $p':2^{V'}\rightarrow \Z$ by $p'(X):=\max\{p(X\cup R): R\subseteq \zeros\}$ for every $X\subseteq V'$ -- this final step can be viewed as a clean up step since it gets rid of vertices that are not incident to any hyperedges (and revises the function $p$ appropriately). 
This clean up step ensures that the hypergraph $(H', w')$ satisfies condition \hyperlink{WeakToStrong(b)}{(b)} -- i.e., the degree of every vertex in $(H', w')$ is non-zero. 
It can be shown that the function $p'$ is symmetric skew-supermodular and the hypergraph $(H', w')$ weakly covers the function $p'$. Furthermore, the function $p'$ and hypergraph $(H', w')$ satisfy conditions \hyperlink{WeakToStrong(a)}{(a)} and \hyperlink{WeakToStrong(b)}{(b)}. 
We recursively call the algorithm on input $((H', w'), p')$ to obtain a hypergraph $(H_0^*, w_0^*)$. We obtain the hypergraph $(G, c)$ from $(H_0^*, w_0^*)$ by adding the vertices $\zeros$ and return the hypergraph $(G+H_0, c+w_0)$.  
By induction on the total weight of hyperedges in the input hypergraph, it can be shown that our algorithm returns a hypergraph satisfying properties \hyperlink{WeakToStrong(1)}{(1)} and \hyperlink{WeakToStrong(2)}{(2)} and also terminates within a finite number of recursive calls. 

\paragraph{Recursion Depth and Potential Functions.} We now sketch our proof to show that the recursion depth of our algorithm is $|E|+O(|V|)$. We note that this also bounds the number of additional hyperedges in the hypergraph returned by the algorithm, and consequently this hypergraph also satisfies property \hyperlink{WeakToStrong(3)}{(3)}. Let $\ell$ be the number of recursive calls made by the algorithm on the input instance $((H, w), p)$. 
We partition 
%\shubhang{The proof doesn't exactly use a partition (i.e. there can be an overlap in the two sets of recursive calls that are analyzed separately in the proof). This should be fine though.} 
the set $[\ell]$ of recursive calls into two parts: let $P_1\subseteq [\ell]$ be the set of recursive calls during which 
the merged hyperedge $e\cup f$ survives in the hypergraph $(H', w')$ that is input to the subsequent recursive call and 
$P_2\subseteq [\ell]$ be the recursive calls during which 
the merged hyperedge $e\cup f$ does not survive in the hypergraph $(H', w')$ that is input to the subsequent recursive call. 
We note that $[\ell]=P_1\uplus P_2$. We bound $|P_1|$ and $|P_2|$ separately using certain carefully designed potential functions.

First, we show that $|P_1|=O(|V|)$ as follows: for $i\in [\ell]$, consider the maximal tight set family $\mathcal{T}_i=\mathcal{T}_{p_i, H_i, w_i}$ where $((H_i, w_i), p_i)$ is the input to the $i^{th}$ recursive call. Also, let $\mathcal{T}_{\le 1}:=\mathcal{T}_1$ and $\mathcal{T}_{\le i}:=\mathcal{T}_{i}\cup \{X\cap V_i: X\in \mathcal{T}_{\le i-1}\}$ for integers $i$ where $2 \le i\le \ell$ and $V_i$ is the ground set of the input to the $i^{th}$ recursive call. Thus, $\calT_{\leq i}$ is the \emph{projection} of all the maximal tight sets encountered in the first $i$ recursive calls of the algorithm onto the ground set of the inputs at the $i^{th}$ recursive call. We show that $\mathcal{T}_{\le i}$ is laminar for every $i\in [\ell]$ (\Cref{lem:WeakToStrong:properties-of-cumulative-tight-set-families}). 
However, $|\mathcal{T}_{\le i}|$ is not necessarily non-decreasing with $i$ since projection of a set family to a subset could result in the loss of sets from the family. Consequently, $|\mathcal{T}_{\le i}|$ is not suitable as a potential function to measure progress. Instead, we use the potential function $\phi(i):=|\mathcal{T}_{\le i}|+3|\zeros_{\le i-1}|$, where $\zeros_{\le i}$ is the union of the sets $\zeros$ computed up to the $i^{th}$ recursive call. We show that $\phi(i)$ is non-decreasing and strictly increases if $i\in P_1$ (\Cref{claim:WeakToStrong:hyperedge-support-size:alpha^R=weight-of-added-hyperedge} in \Cref{lem:WeakToStrongCover:NumberOfHyperedges}). Consequently, $|P_1|=O(|V|)$.

Secondly, we bound $|P_2|$ as follows. We use a lookahead-potential function: let $\Phi_1(i)$ be the number of recursive calls between $i$ and $\ell$
during which the merged hyperedge $e\cup f$ survives in the hypergraph $(H', w')$ that is input to the subsequent recursive call
and let $\Phi(i):=|E_i|+\Phi_1(i)$, where $E_i$ is the set of hyperedges in the hypergraph $(H_i, w_i)$ input to the $i^{th}$ recursive call. We show that $\Phi(i)$ is non-increasing and strictly decreases if 
$i\in P_2$ 
(\Cref{claim:WeakToStrong:hyperedge-support-size:alpha^R<weight-of-added-hyperedge} in \Cref{lem:WeakToStrongCover:NumberOfHyperedges}). Hence, $|P_2|\le \Phi(1)-\Phi(\ell)\le \Phi(1) \le |E_1|+|P_1|=|E|+O(|V))$, where the last equality is because of the bound on $|P_1|$ from the previous paragraph.  

\begin{remark}
    Our key technical contributions are twofold. Our first key technical contribution is identifying conditions \hyperlink{WeakToStrong(a)}{(a)} and \hyperlink{WeakToStrong(b)}{(b)} under which $\maximalTightSetFamily{p, H, w}$ becomes a disjoint family. We ensure that conditions \hyperlink{WeakToStrong(a)}{(a)} and \hyperlink{WeakToStrong(b)}{(b)} hold in every recursive call by performing \emph{immediate} reduction and clean-up steps in the algorithm. 
    Our second key technical contribution is identifying appropriate potential functions to measure progress of the algorithm. The disjointness of $\maximalTightSetFamily{p, H, w}$ was crucial for identifying the laminar structure of the family of projected maximal tight sets across recursive calls, which was subsequently helpful in bounding the number of recursive calls corresponding to $P_1$. 
    %Our third key technical contribution is bounding the number of recursive calls corresponding to $P_2$ required a lookahead-potential function that relates $|P_2|$ to $|P_1|$. As discussed above, the additive $|E|$ in the recursion depth comes from the bound on $|P_2|$.
    Moreover, we bound the number of recursive calls corresponding to $P_2$ using a lookahead-potential function that relates $|P_2|$ to $|P_1|$. As discussed above, the additive $|E|$ in the recursion depth comes from the bound on $|P_2|$.
\end{remark}
\section{Preliminaries}\label{sec:prelims}
%\kristof{$-$ and $\setminus$ are both used throughout the paper. We should unify the notation.}

%\knote{Consistently use ``for every''. Avoid ``for all'', ``for each'', etc. I have already done this in the intro and results section. }

We denote the sets of 
%\emph{reals}, \emph{non-negative reals}, \emph{positive reals}, 
\emph{integers}, \emph{non-negative integers}, and \emph{positive integers} by 
%$\mathbb{R}$, $\mathbb{R}_{\geq 0}$, $\mathbb{R}_+$, 
$\mathbb{Z}$, $\mathbb{Z}_{\geq 0}$, and $\mathbb{Z}_+$ respectively. For a positive integer $n$, we use $[n]$ to denote the set $\{1, 2, \ldots, n\}$. %We denote the value of the \emph{indicator function} for a statement $S$ by $\indicator_S$, i.e. $\indicator_S=1$ of $S$ is true and $0$ otherwise. \knote{State only what is used.}
Let $V$ be a finite set. We use $2^{V}$ to denote all subsets of $V$. 
%The \emph{characteristic vector} of a set $X\subseteq V$ is denoted by $\chi_X \in \{0, 1\}^V$ where $\chi_X(u)=1$ if $u\in X$ and $0$ otherwise for every $u\in V$. 
Let $X, Y\subseteq V$. We will denote $X\setminus Y$ by $X-Y$ and $X\cup Y$ by $X+Y$. If $Y$ consists of a single element $y$, then $X-\{y\}$ and $X+\{y\}$ are abbreviated as $X-y$ and $X+y$, respectively. 
Let $\calF\subseteq 2^V$ be a family of subsets of $V$. 
%A set $T\subseteq V$ is a \emph{transversal} for the family $\calF$ if each set in $\calF$ has non-empty intersection with the set $T$. 
The family $\calF$ is 
a \emph{disjoint family} if $X\cap Y=\emptyset$ for every $X, Y\in \calF$, 
a \emph{chain family} if either $X\subseteq Y$ or $Y\subseteq X$ for every $X, Y\in \calF$, and a \emph{laminar family} if either $X\cap Y=\emptyset$ or $X\subseteq Y$ or $Y\subseteq X$ for every $X, Y\in \calF$. 
We recall that the size of a laminar family $\mathcal{F}\subseteq 2^V$ is at most $2|V|$. 
For a subset $U\subseteq V$, we define the projection of $\mathcal{F}$ to $U$ as   $\mathcal{F}|_U:=\{X\cap U: X\in \mathcal{F}\}-\{\emptyset\}$ (we will use the convention that $\mathcal{F}|_U$ is a set and not a multi-set). 
%We will be interested in the \emph{projection} of a laminar family
%, i.e., the restriction of a laminar family $\calL \subseteq 2^V$ to a ground set $V - \calZ$ for some $\calZ\subseteq V$. 
The next lemma shows that the projection of a laminar family is also laminar and the size of the projection is comparable to that of the original family. We give a proof in \Cref{appendix:sec:uncrossing-properties:projection-of-laminar-families}.  %\knote{State only what is used.}\shubhang{I think we use everything except the set addition notation --- whether this is used anywhere still has to be verified}
%$\calL$.

\begin{restatable}{lemma}{projectionLaminarFamily}\label{lem:CoveringAlgorithm:projection-laminar-family}
    Let $\calL \subseteq 2^V$ be a laminar family and $\calZ \subseteq V$ be a subset of elements. 
    Let $\calL':=\calL|_{V-\zeros}$. 
    %Let $\calL' \coloneqq  \{X - \calZ : X \in \calL\} - \{\emptyset\}$. 
    Then, the family $\calL'$ is a laminar family and $|\calL| \leq |\calL'| + 3|\calZ|$.
\end{restatable}

%We prove \Cref{lem:CoveringAlgorithm:projection-laminar-family} in \Cref{appendix:sec:uncrossing-properties:projection-of-laminar-families}.


%\paragraph{Set Functions.} 
Let $p: 2^{V} \rightarrow \mathbb{R}$ be a set function over a finite ground set $V$. 
%We will denote the maximum function value of $p$ by $K_p$, i.e., $K_p\coloneqq \max\{p(X): X\subseteq V\}$ \knote{Is $K_p$ used at all?}.
The function $p$ is \emph{monotone} if  $p(X)\le p(Y)$ for every $X\subseteq Y\subseteq V$. 
%If the function $p$ is locally supermodular 
If $p(X)+p(Y)\le p(X\cap Y)+p(X\cup Y)$
for every $X,Y\subseteq V$, then the function $p$ is said to be \emph{supermodular}. 
%If the reverse of the same inequality holds for every $X,Y\subseteq V$, then the function $p$ is said to be \emph{submodular}. 
If $-p$ is supermodular, then the function $p$ is said to be \emph{submodular}. 
We recall that for a hypergraph $(H=(V, E), w: E\rightarrow \Z_+)$, the cut function $d_{(H, w)}: 2^V\rightarrow \Z_{\ge 0}$ is symmetric submodular and the coverage function $b_{(H, w)}: 2^V\rightarrow \Z_{\ge 0}$ is monotone submodular \cite{Frank-book}. 
%A set $X\subseteq V$ is said to be a \emph{$p$-maximizer} if $p(X) = \max\{p(Z): Z\subseteq V\}$. We use the notation $Z \coloneqq  \arg\max\{p(Z): Z\subseteq V\}$ to denote an arbitrary set $Z\subseteq V$ such that $p(Z)\ge p(X)$ for every $X\subseteq V$ (similarly for minimizers of the function $p$).
For a subset $Z \subseteq V$, the \emph{contracted function} $\functionContract{p}{Z} : 2^{V - Z}\rightarrow\Z$ is defined as $\functionContract{p}{Z}(X) \coloneqq  \max\{p(X\cup R) : R \subseteq Z\}$. We note that if $p$ is a skew-supermodular function and $Z\subseteq V$, then $\functionContract{p}{Z}$ is skew-supermodular. Moreover, if $p$ is a symmetric function and $Z\subseteq V$, then $\functionContract{p}{Z}$ is symmetric. 
%For a function $m: V\rightarrow \R$ and a subset $Z\subseteq V$, we define the \emph{restricted function} $\functionRestrict{m}{Z} : V-Z \rightarrow \R$ as $\functionRestrict{m}{Z}(u) \coloneqq  m(u)$ for every $u\in V - Z$.  
%\knote{State only what is used.}

We recall that in our algorithmic problems, we have access to the input set function $p$ via \functionMaximizationOracleStrongCover{p}. 
%We describe an additional oracle that  will simplify our proofs. We will show that both oracles defined below can be implemented in strongly polynomial time using \functionMaximizationOracleStrongCover{p}. 
We describe an additional oracle that will simplify our proofs. We will show that the oracle defined below can be implemented in strongly polynomial time using \functionMaximizationOracleStrongCover{p}. 
\begin{definition}[$p$-maximization oracle]
    Let $p: 2^V\rightarrow \Z$ be a set function. 
    \functionMaximizationOracle{p}$((G_0, c_0), S_0, T_0)$ takes as input a hypergraph $(G_0=(V, E_0), c_0: E_0\rightarrow \mathbb{Z}_{+})$  
    and disjoint sets $S_0, T_0\subseteq V$, and returns a tuple $(Z, p(Z))$, where $Z$ is an optimum solution to the following problem:
    \begin{align*}
        \max & \left\{p(Z)-b_{(G_0, c_0)}(Z): S_0\subseteq Z\subseteq V-T_0\right\}. \tag{\functionMaximizationOracle{p}}\label{tag:oracle-main}
        \end{align*} 
\end{definition}
%\knote{If we want strongly polynomial time algorithms, then we can query the above oracles only on hypergraphs $(G_0, c_0)$, where $|E_0|=$ poly$(|V|)$ and $c_0$ is also bounded by ...}
The evaluation oracle for a function $p$ can be implemented using one query to  
%\functionMaximizationEmptyOracle{p};  
%\functionMaximizationEmptyOracle{p} can be implemented using one query to 
\functionMaximizationOracle{p} where the input hypergraph $(G_0, c_0)$ is the empty hypergraph and $S_0 := Z$, $T_0 := V - Z$. The next lemma shows that \functionMaximizationOracle{p} can be implemented using at most $|V|+1$ queries to \functionMaximizationOracleStrongCover{p} where the hypergraphs used as input to  \functionMaximizationOracleStrongCover{p} have size of the order of the size of the hypergraph input to \functionMaximizationOracle{p}. We give a proof in \Cref{sec:wc-oracle-from-sc-oracle}.  

\begin{restatable}{lemma}{lemmaWCoraclefromSCoracle}\label{lem:Preliminaries:wc-oracle-from-sc-oracle}
Let $p: 2^V\rightarrow \Z$ be a set function, $(G = (V, E), w:E\rightarrow\Z_+)$ be a hypergraph, and $S, T \subseteq V$ be disjoint sets. Then, $\functionMaximizationOracle{p}((G, w), S, T)$ can be implemented to run in $O(|V|(|V|+|E|))$ time using at most $|V|+1$ queries to $\functionMaximizationOracleStrongCover{p}$.
The run-time includes the time to construct the hypergraphs used as input to the queries to \functionMaximizationOracleStrongCover{p}. Moreover, each query to \functionMaximizationOracleStrongCover{p} is on an input hypergraph $(G_0, c_0)$ that has at most $|V|$ vertices and $|E|$ hyperedges.
\end{restatable}


For two hypergraphs $(G=(V, E_G), c_G:E_G\rightarrow \Z_+)$ and $(H=(V, E_H), w_H: E_H\rightarrow \Z_+)$ on the same vertex set $V$, we recall that the hypergraph $(G+H=(V, E_{G+H}), c_G+w_H)$ is the hypergraph with vertex set $V$ and hyperedge set $E_{G+H}\coloneqq E_G\cup E_H$ with the weight of every hyperedge $e\in E_G\cap E_H$ being $c_G(e) + w_H(e)$, the weight of every hyperedge $e\in E_G\setminus E_H$ being $c_G(e)$,  and the weight of every hyperedge $e\in E_H\setminus E_G$ being $w_H(e)$. 
\section{Local Connectivity Preserving Complete h-Splitting-Off}
\label{sec:hypergraph-splitting-off-proof}
We prove Theorem \ref{thm:complete-splitting-off} using Theorem \ref{thm:WeakToStrongCover:main} in this section. 
We need the following lemma showing that a certain set function $p: 2^V\rightarrow \Z$ is symmetric skew-supermodular and admits an efficient algorithm to implement \functionMaximizationOracleStrongCover{p}. This lemma has been observed in the literature before. We refer the reader to \Cref{sec:helper-for-p-max-oracle} for its proof. 

\begin{restatable}{lemma}{lempmaxoracle}\label{lemma:helper-for-p-max-oracle}
    Let $(G=(V, E), c: E\rightarrow \Z_+)$ be a hypergraph. Let $r:\binom{V}{2}\rightarrow \Z_{\ge 0}$ be a function on pairs of elements of $V$ and $R, p_{(G, c, r)}:2^V\rightarrow \Z$ be functions defined by $R(X)\coloneqq \max\{r(\{u, v\}): u\in X, v\in V\setminus X\}$ for every $X\subseteq V$, $R(\emptyset)\coloneqq 0$, $R(V)\coloneqq 0$, and $p_{(G, c, r)}(X)\coloneqq R(X)-d_{(G, c)}(X)$ for every $X\subseteq V$. 
    \begin{enumerate}
        \item The function $p_{(G, c, r)}$ is symmetric skew-supermodular. 
        \item 
        Given hypergraph $(G, c)$, function $r$, hypergraph $(G_0=(V, E_0), c_0: E_0\rightarrow \Z_{+})$, and disjoint sets $S_0, T_0\subseteq V$, the oracle   \functionMaximizationOracleStrongCover{p_{(G, c, r)}}$((G_0, c_0), S_0, T_0)$ can be computed in $O(|V|^3(|V|+|E_0|+|E|)(|E_0|+|E|))$ time.
        \iffalse
        Given $(G, c)$ and the function $r$, there exists a strongly polynomial-time algorithm to implement \linebreak \functionMaximizationOracleStrongCover{p}. 
        In particular, for a given hypergraph $(G_0=(V, E_0), c_0: E_0\rightarrow \Z_{+})$, and disjoint sets $S_0, T_0\subseteq V$, the oracle   \functionMaximizationOracleStrongCover{p}$((G_0, c_0), S_0, T_0)$ can be implemented to run in $O(|V|^3(|V|+|E_0|+|E|)(|E_0|+|E|))$ time.
        \fi
    \end{enumerate}
\end{restatable}


We now prove the following stronger form of Theorem \ref{thm:complete-splitting-off}.
\begin{theorem}\label{thm:hypergraphs-splitting-off-stronger-form}
    Let $(G=(V, E), w_G: E\rightarrow \Z_+)$ be a hypergraph and $s\in V$. Then, there exists a hypergraph $(G^*=(V, E^*), c^*)$ such that 
    \begin{enumerate}
        \item $(G^*, c^*)$ is a local connectivity preserving complete h-splitting-off at $s$ from $(G, w_G)$ and 
        \item $|E^*|-|E|=O(|V|)$. 
    \end{enumerate}
    Moreover, given the hypergraph $(G, w_G)$, a hypergraph $(G^*, c^*)$ satisfying the above two properties can be obtained in $O(|V|^6(|V|+|E|)^3)$ time. 
\end{theorem}
\begin{proof}
    We begin by proving the existence of a local connectivity preserving complete h-splitting-off at $s$ from $(G=(V, E), w_G: E\rightarrow \Z_+)$ such that the number of additional hyperedges is $O(|V|)$. 
    We will use \Cref{thm:WeakToStrongCover:main} with an appropriate setting of the symmetric skew-supermodular function $p$ and the hypergraph $(H, w)$. 
    Let $V'\coloneqq V\setminus \{s\}$. 
    Let the hypergraph $(H:=(V', E_H), w_H: E_H\rightarrow \Z_+)$ be defined by $E_H\coloneqq \{e-\{s\}: e\in \delta_G(s)\}$ with $w_H(e-\{s\})\coloneqq w_G(e)$ for every $e\in \delta_G(s)$. 
    In order to define the function $p: 2^{V'}\rightarrow \Z$, we consider another hypergraph $(J:=({V'}, E_J), w_J: E_J\rightarrow \Z_+)$ defined by $E_J\coloneqq E\setminus B_G(s)$ with $w_J(e)\coloneqq w_G(e)$ for every $e\in E\setminus B_G(s)$. 
    %We recall that the cut function $d_{(J, w_J)}: 2^{V'}\rightarrow \Z_{\ge 0}$ of the hypergraph $(J, w_J)$ is symmetric submodular. 
    We consider the function $R:2^{V'}\rightarrow \Z_{\ge 0}$ defined by $R(X)\coloneqq \max\{\lambda_{(G, w_G)}(u, v): u\in X, v\in V'\setminus X\}$ for every non-empty proper subset $X\subsetneq  V'$, $R(\emptyset)\coloneqq 0$, and $R(V')\coloneqq 0$. 
    %We observe that the function $R$ is symmetric skew-supermodular \knote{need to prove or cite a reference}. 
    Now, consider the function $p: 2^{V'}\rightarrow \Z$ defined by $p(X)\coloneqq R(X)-d_{(J, w_J)}(X)$ for every $X\subseteq {V'}$. 
    %Since $p$ is the difference of a symmetric skew-supermodular function and a symmetric submodular function, it follows that 
    The function $p$ is symmetric skew-supermodular by \Cref{lemma:helper-for-p-max-oracle}. 
    We also note that $b_{(H, w_H)}(X)\ge p(X)$ for every $X\subseteq {V'}$: for an arbitrary $X\subseteq {V'}$, we have that $b_{(H, w_H)}(X)+d_{(J, w_J)}(X)=d_{(G, w_G)}(X)\ge R(X)$ where the last inequality is by definition of the function $R$ and hence, $b_{(H, w_H)}(X)\ge R(X)-d_{(J, w_J)}(X)=p(X)$. 
    Hence, $(H,w_H)$ weakly covers the symmetric skew-supermodular function $p$. 
    Applying \Cref{thm:WeakToStrongCover:main} to the hypergraph $(H, w_H)$ and the function $p$, we obtain that there exists a hypergraph $(H^* = (V', E_{H^*}), w_{H^*}: E_{H^*}\rightarrow \Z_+)$ satisfying the following three properties: (i) $d_{(H^*, w_{H^*})}\ge p(X)$ for every $X\subseteq V'$,  (ii) the hypergraph $(H^*, w_{H^*})$ is obtained by merging hyperedges of the hypergraph $(H, w_H)$, and (iii) $|E_{H^*}|-|E_H|=O(|V|)$. 
    Consider the hypergraph $(G^*=(V, E^*), c^*)$ obtained from the hypergraph $(J+H^*, w_J+w_{H^*})$ by adding the vertex $s$. 
    We have that 
    \begin{align*}
    |E^*|
    &=|E_J\cup E_{H^*}| &\text{(by definition of $E^*$)}\\
    &\le |E_J|+|E_{H^*}| &\\
    &=|E_J|+|E_H|+O(|V|) & \text{(by property (iii))}\\
    &=|E_J\cup E_H|+O(|V|) & \text{(by definitions of $E_J$ and $E_H$)}\\
    &=|E|+O(|V|), & \text{(by definitions of $E_J$ and $E_H$)}
    \end{align*}
    and hence, $|E^*|-|E|=O(|V|)$. 
    
    We now show that $(G^*, c^*)$ is a complete h-splitting-off at $s$ from $(G, w_G)$ and $(G^*, c^*)$ preserves local connectivity. 
    We note that property (ii) implies that $(G^*, c^*)$ is a complete h-splitting-off at $s$ from $(G, c)$: this is because, 
    hyperedges of $E_{H^*}\setminus E_H$ that are obtained by merging hyperedges of $E_H$ are equivalent to being obtained by a sequence of h-merge almost-disjoint hyperedges operation starting from the hypergraph $(G, w_G)$ and a final h-trim operation; 
    moreover, hyperedges of $E_{H^*}\cap E_H$ are equivalent to being obtained by the h-trim operation from the hypergraph $(G, w_G)$;  
    finally, we have that $d_{(G^*, c^*)}(s)=0$ by definition of $(G^*, c^*)$. 
    Next, we note that property (i) implies that $(G^*, c^*)$ preserves local connectivity: consider an arbitrary pair of distinct vertices $u, v\in V'$. We recall that h-merge almost-disjoint hyperedges as well as h-trim operations do not increase the cut value of any subset $X\subseteq V'$ and hence, $\lambda_{(G^*, c^*)}(u, v) \le d_{(G^*, c^*)}(X)\le d_{(G, w_G)}(X)$ for every $X\subseteq V'$ such that $u\in X, v\in V'\setminus X$ and in particular, $\lambda_{(G^*, c^*)}(u, v) \le \lambda_{(G, w_G)}(u, v)$. Hence, it suffices to show that $d_{(G^*, c^*)}(X)\ge \lambda_{(G, w_G)}(u, v)$ for every $X\subseteq V'$ such that $u\in X, v\in V'\setminus X$. Let $X\subseteq V'$ such that $u\in X, v\in V'\setminus X$. We have that 
    \begin{align*}
    d_{(G^*, c^*)}(X)
    &=d_{(J, w_J)}(X) + d_{(H^*, w_{H^*})}(X) & \text{(by definition of $(G^*, c^*)$)}\\
    &\ge d_{(J, w_J)}(X) +p(X) & \text{(by property (i) of $(H^*, w_{H^*})$)}\\
    &=d_{(J, w_J)}(X) +R(X)-d_{(J, w_J)}(X) & \text{(by definition of the function $p$)}\\
    &=R(X) \\
    &\ge \lambda_{(G, w_G)}(u,v). & \text{(by definition of the function $R$)}
    \end{align*}
    %Finally, we note that $b_{(G^*, c^*)}(s)=0$ by definition. 
    This completes the proof that $(G^*, c^*)$ is a complete splitting off at $s$ from $(G, w_G)$ and preserves local connectivity. 
    
    %We now show a strongly polynomial-time algorithm to obtain such a hypergraph $(G^*, c^*)$. Given $(G, w_G)$ as input, we can construct $(H, w_H)$ and $(J, w_J)$ as above in strongly polynomial time. There exists a strongly polynomial-time algorithm to implement \functionMaximizationOracleStrongCover{p}, where $p$ is the function defined above (see \Cref{lemma:helper-for-p-max-oracle}). By the algorithmic conclusion of \Cref{thm:WeakToStrongCover:main}, it follows that the hypergraph $(H^*, w_{H^*})$ mentioned in the previous paragraph can be obtained in strongly polynomial time. Finally, we can construct the hypergraph $(G^*, c^*)$ using the hypergraph $(H^*, w_{H^*})$ and the hypergraph $(J, w_J)$ in strongly polynomial time. 
   We now analyze the run-time of the algorithm to obtain such a hypergraph $(G^*, c^*)$. Given $(G=(V, E), w_G: E\rightarrow \Z_+)$ as input, we can construct $(H, w_H)$ and $(J, w_G)$ as above in time $O(|V| + |E|)$. We note that both $(H, w_H)$ and $(J, w_G)$ have vertex set $V-\{s\}$ and at most $|E|$ hyperedges. For the function $p$ defined as above, we consider the time to implement \functionMaximizationOracleStrongCover{p}: by Lemma \ref{lemma:helper-for-p-max-oracle}, for a given hypergraph $(G_0=(V-\{s\}, E_0), c_0: E_0\rightarrow \Z_+)$, sets $S_0, T_0\subseteq V$, and a vector $y_0\in \R^V$, the oracle \functionMaximizationOracleStrongCover{p}$((G_0, c_0), S_0, T_0, y_0)$ can be implemented to run in $O(|V|^3(|V|+|E_0|+|E|)(|E_0|+|E|))$ time. By the algorithmic conclusion of \Cref{thm:WeakToStrongCover:main}, it follows that the hypergraph $(H^*, w_{H^*})$ mentioned in the previous paragraph can be obtained in $O(|V|^6(|V|+|E|)^3)$ time. Finally, we can construct the hypergraph $(G^*, c^*)$ using the hypergraph $(H^*, w_{H^*})$ and the hypergraph $(J, w_J)$ in $O(|V|+|E_{H^*}|+|E|)=O(|V|+|E|)$ time. Thus, the overall run-time is $O(|V|^6(|V|+|E|)^3)$.
\end{proof}
\section{Constructive Characterization of \(k\)-Hyperedge-Connected Hypergraphs}
\label{sec:characterizing-k-edge-connected-hypergraphs}
%The splitting-off operation leads to a constructive characterization of $k$-edge-connected graphs \cite{Frank-book}.
In this section, we prove Theorem \ref{thm:k-EC-hypergraph-characterization} using Theorem \ref{thm:complete-splitting-off}. 
%use \Cref{thm:complete-splitting-off} to show a constructive characterization of $k$-hyperedge-connected hypergraphs (i.e., \Cref{thm:k-EC-hypergraph-characterization}). 
Throughout this section, we use the term \emph{hypergraph} to refer to a \emph{multi-hypergraph}.
%\knote{Throughout this section, hypergraph refers to multi-hypergraphs. }
%We remark that our constructive characterization can be obtained in polynomial time owing to the strongly polynomial time algorithm for h-splitting-off. We will also need the pinching operation in hypergraphs (recall \Cref{def:pinching-operation}) which is the inverse operation to complete h-splitting-off. The following is the main result of this section.
\iffalse
We first define a \emph{pinching} operation in hypergraphs analogous to the classical pinching operation encountered in the constructive characterization of $k$-edge-connected graphs. We then use this operation to state the main result of the section below.
\begin{definition}
    Let $H = (V, E)$ be a hypergraph and $p,k \in \Z_+$ be positive integers such that $p \leq k$. By \emph{$(k, p)$-pinching} hyperedges of $H$, we mean performing the following operation:
    \begin{enumerate}
        \item pick $p$ distinct hyperedges $e_1, \ldots, e_p \in E$,
        \item pick $p$ positive integers $t_1, \ldots, t_p \in \Z_+$ such that $\sum_{i \in [p]} t_i = k$,
        \item for $i \in [p]$, partition the hyperedge $e_i$ into $t_i$ non-empty parts $e_i = \uplus_{j \in [t_i]} f_i^j$,
        \item remove the hyperedges $e_1, \ldots e_p$ from the hypergraph $H$,
        \item add a new vertex $s$ and hyperedges $\{f_i^j \cup \{s\} : j \in [t_i], i \in [p]\}$ to the hypergraph $H$.
    \end{enumerate}
\end{definition}

\begin{theorem}\label{thm:k-EC-hypergraph-characterization}
    Let $k \in \Z_+$ be a positive integer. A hypergraph $H = (V, E)$ is $k$-hyperedge-connected if and only if $H$ can be obtained by starting from the single vertex hypergraph with no hyperedges and repeatedly applying one of the following two operations:
    \begin{enumerate}
        \item add a new hyperedge over a subset of (existing) vertices, and
        \item $(k, p)$-pinching hyperedges in the existing hypergraph for some $p \leq k$.
    \end{enumerate}
\end{theorem}
\fi
We define a hypergraph $H = (V, E)$ to be \emph{minimally} $k$-hyperedge-connected if $H-e$ is not $k$-hyperedge-connected for every $e\in E$. 
%the removal of any hyperedge results in a hypergraph that is not $k$-hyperedge-connected. 
%To show the characterization of \Cref{thm:k-EC-hypergraph-characterization}, 
We need the following observation about minimally $k$-hyperedge-connected hypergraphs. %We formally state and prove our structural observation in the next lemma.
\begin{lemma}\label{lem:minimally-k-EC-hypergraph-has-degree-k-vertex}
    Let $k \in \Z_+$ be a positive integer and $G = (V, E)$ be a minimally $k$-hyperedge-connected hypergraph with $|V| \geq 2$. Then, there exists a vertex $u \in V$ such that $d_G(u) = k$.
\end{lemma}
\begin{proof}
    Since $G$ is minimally $k$-hyperedge-connected, there exists a set $X \subseteq V$ such that $d_G(X) = k$. Let $X\subseteq V$ be an inclusion-wise minimal set such that $d_G(X)=k$. We note that if $|X| = 1$, then the claim holds. Suppose that $|X| \geq 2$.

    We first show that there exists a hyperedge $e\in E$ contained in the set $X$. By way of contradiction, suppose this is false. Let $u \in X$ be an arbitrary vertex. We note $X - \{u\} \not = \emptyset$ since $|X| \geq 2$. Furthermore, every hyperedge incident to $u$ is not contained in the set $X$, i.e. $\delta_G(u) \subseteq \delta_G(X)$. Thus, we have that $k \leq d_G(u) \leq d_G(X) = k$, where the first inequality is because $G$ is $k$-hyperedge-connected. Thus, all inequalities are equations, and we have that $d_G(u) = k$, contradicting minimality of the set $X$.

    Let $e\in E$ be a hyperedge contained in the set $X$. 
    Next, we show that %for the hyperedge $e$ contained in $X$, 
    there exists a set $Y \subseteq V$ such that $d_G(Y) = k$ and $e \in \delta_G(Y)$. 
    %(we note that this holds for every hyperedge in $H$ but we will not need the stronger statement for the purposes of this proof). 
    By way of contradiction, suppose this is false. Then, for every $Y \subseteq V$ such that $e \in \delta_G(Y)$, we have that $d_G(Y) > k$. Thus, the hypergraph $G - e$ is $k$-hyperedge-connected, contradicting the minimally $k$-hyperedge-connected property of $G$.

    We now complete the proof of the lemma. By minimality of the set $X$, we have that $Y-X\neq \emptyset$ and $(V-Y)-X\neq \emptyset$. 
    %\knote{$Y-X\neq \emptyset$ and $(V-Y)-X\neq \emptyset$. }
    Furthermore, since $e \subseteq X$ but $e \in \delta_G(Y)$, we have that $X-Y \not = \emptyset$ 
    %\knote{$X-Y\neq \emptyset$}
    and $X \cap Y \not = \emptyset$. 
    Thus, all four sets $Y-X, X-Y, X\cap Y, V-(X\cup Y)$ are non-empty.
    %\knote{Thus, all four sets $Y-X, X-Y, X\cap Y, V-(X\cup Y)$ are non-empty.} 
    Hence, we have the following: $2k = d_{G}(X) + d_G(Y) \geq d_G(X\cap Y) + d_G(X\cup Y) \geq 2k$, where the first inequality is by submodularity of the cut function and the second inequality is because $G$ is $k$-hyperedge-connected. 
    %\knote{For the second inequality to hold, we need $X\cup Y\neq V$. Why does this hold? Holds because $V-(X\cup Y)$ is also non-empty}
    Thus, all inequalities are equations and we have that $d_{G}(X\cap Y) = k$, contradicting minimality of the set $X$.
\end{proof}

We now restate and prove \Cref{thm:k-EC-hypergraph-characterization}. 
%We remark that the constructive characterization of \Cref{thm:k-EC-hypergraph-characterization} will rely on complete h-splitting-off and can be obtained in polynomial time owing to the strongly poly-time algorithm for complete h-splitting-off (\Cref{thm:complete-splitting-off}).

\thmConstructiveCharacterization*
\begin{proof}%[Proof of \Cref{thm:k-EC-hypergraph-characterization}.]
    %\knote{The reverse direction has some issue. Say we start from a single vertex hypergraph with no hyperedges, add a singleton hyperedge. The resulting hypergraph is not $k$-hyperedge-connected. }
    If $H$ is a $k$-hyperedge-connected hypergraph, then performing both operations mentioned in the theorem gives a $k$-hyperedge-connected hypergraph. Consequently, the reverse direction follows by induction on $|V| + |E|$. We prove the forward direction.  By way of contradiction, suppose the forward direction of the theorem is false. Let $G = (V, E_G)$ be a counter-example that minimizes $|V| + |E_G|$, i.e. $G$ is $k$-hyperedge-connected but cannot be obtained by applying operations (1) and (2). We note that $G$ is minimally $k$-hyperedge-connected since otherwise, there exists a hyperedge $e \in E_G$ such that $G - e$ is $k$-hyperedge-connected, and consequently, the hypergraph $G - e$ is also a counter-example to the forward direction of the theorem contradicting minimality of the counter-example $G$. Thus, by \Cref{lem:minimally-k-EC-hypergraph-has-degree-k-vertex}, there exists a vertex $u \in V$ such that $d_G(u) = k$. Let $H = (V - u, E_H)$ be the hypergraph obtained as a local connectivity preserving complete h-splitting-off at the vertex $u$ from $G$ as guaranteed to exist by \Cref{thm:complete-splitting-off}. Then, the hypergraph $H$ is $k$-hyperedge-connected. Moreover, $|V-u| + |E_H| < |V| + |E_G|$ because complete h-splitting-off at a vertex from a multi-hypergraph does not increase the total number of hyperedges. 
    %\knote{Why does this hold? We have that $|V-u|<|V|$. Moreover, $|E_G|\le |E_H|$ since complete h-splitting-off at a vertex from a multi-hypergraph does not increase the total number of hyperedges.}. 
    By minimality of the counter-example $G$, the hypergraph $H$ can be obtained by starting from the single vertex hypergraph with no hyperedges and repeatedly applying operations (1) and (2). 
    %We note that $H$ can be obtained from $G$ by operation (2) since it is the reverse operation to complete h-splitting-off 
    We note that $G$ can be obtained from $H$ by applying operation (2). This is because $H$ was obtained by complete h-splitting-off at vertex $u$ from $G$ where $d_G(u)=k$ and operation (2) is the inverse of a complete h-splitting-off operation at a vertex with degree exactly $k$. 
    %\knote{The proof is not informative regarding why we split-off at a degree $k$ vertex and not at an arbitrary vertex.}\shubhang{this is because a $(k, \cdot)$-pinching operation is the inverse of a complete h-splitting-off operation at a vertex with degree exactly $k$}. 
    Thus, $G$ is not a counter-example to the forward direction of the theorem, a contradiction.
\end{proof}

We remark that our proof of Theorem \ref{thm:k-EC-hypergraph-characterization} is constructive since it is equivalent to a proof by induction on $|V|+|E|$. 
Using the polynomial-time algorithm in Theorem \ref{thm:complete-splitting-off}, we get the following conclusion: given a $k$-hyperedge-connected hypergraph $H$, our proof gives a polynomial-time algorithm to construct a sequence of hypergraphs $H_0, H_1, H_2, \ldots, H_t$, where $H_0$ is the single vertex hypergraph with no hyperedges, $H_t=H$ and for each $i\in [t]$, the hypergraph $H_i$ is obtained from $H_{i-1}$ by either adding a new hyperedge over a subset of vertices in $H_{i-1}$ or by $(k,p)$-pinching hyperedges in $H_{i-1}$ for some positive integer $p\le k$. 
\section{Steiner Rooted Connected Orientation of Graphs and Hypergraphs}
\label{sec:rooted-steiner-connected-orientations}
In this section, we use  \Cref{thm:complete-splitting-off} to
prove Theorems \ref{thm:rooted-steiner-orientation-in-graphs} and \ref{thm:rooted-steiner-orientation-in-hypergraphs}. 
We recall the \textsc{Max Steiner Rooted-Connected Orientation} problem: the input here is a hypergraph $G=(V, E)$, a subset $T\subseteq V$ of terminals, and a root vertex $r\in T$. The goal is to find a maximum integer $k$ and an orientation $\overrightarrow{G}$ of $G$ such that $\overrightarrow{G}$ is Steiner rooted $k$-hyperarc-connected. 
%As warm-up to Theorems \ref{thm:rooted-steiner-orientation-in-graphs} and \ref{thm:rooted-steiner-orientation-in-hypergraphs}, we consider the 
Although \textsc{Max Steiner Rooted-Connected Orientation} is NP-hard, two extreme cases of the problem are polynomial-time solvable: the number of terminals $|T|=2$ and the number of terminals $|T|=|V|$. We discuss these polynomial-time solvable cases as warm-ups towards Theorems \ref{thm:rooted-steiner-orientation-in-graphs} and \ref{thm:rooted-steiner-orientation-in-hypergraphs}. The ideas underlying these cases will be useful in proving Theorems \ref{thm:rooted-steiner-orientation-in-graphs} and \ref{thm:rooted-steiner-orientation-in-hypergraphs}. 
We recall that non-terminals are denoted as Steiner vertices. 

\paragraph{The case of $|T|=2$.} \textsc{Max Steiner Rooted-Connected Orientation} problem where we have exactly $2$ terminals is polynomial-time solvable owing to Menger's theorem in hypergraphs. 
A path between $u$ and $v$ in a hypergraph is an alternative sequence of distinct vertices and hyperedges $v_1=u, e_1, v_2, e_2, ..., e_{\ell-1}, v_{\ell}=v$ such that $v_i, v_{i+1}\in e_i$ for every $i\in [\ell-1]$. 
%and $v_i\neq v_j$ for every distinct $i,j\in [\ell]$. 
We recall Menger's theorem: for a hypergraph $G=(V, E)$ and a pair of distinct vertices $r, v\in V$, there exist $k$ hyperedge-disjoint paths between $r$ and $v$ in $G$ if and only if $\lambda_G(r,v)\ge k$. We present a proof of this fundamental graph-theoretical result using \Cref{thm:complete-splitting-off}. Our proof also holds for undirected graphs. 
Parts of our proof will be useful in the proof of Theorems \ref{thm:rooted-steiner-orientation-in-graphs} and \ref{thm:rooted-steiner-orientation-in-hypergraphs}. 

Let $G=(V, E)$ be a hypergraph, $r,v\in V$ be distinct vertices, and $k$ be a positive integer. It is easy to see that if there exist $k$ hyperedge-disjoint paths between $r$ and $v$ in $G$, then $\lambda_G(r,v)\ge k$. We prove the converse. Suppose that $\lambda_G(r, v)\ge k$. We would like to show that there exist $k$ hyperedge-disjoint paths between $r$ and $v$ in $G$. 
Consider performing local connectivity preserving complete h-splitting-off operations at Steiner vertices in the current hypergraph and then deleting them. This will result in a $2$-vertex hypergraph $H=(\{r,v\}, E_H)$ such that $\lambda_H(r,v)\ge k$. Since $H$ contains only $2$ vertices, it immediately follows that all hyperedges of $H$ are in fact edges, and they can be oriented to obtain $k$ hyperarc-disjoint paths 
%from $r$ to $v$. 
from $v$ to $r$. 
Let $\overrightarrow{H}$ be such an orientation. 
We now extend this orientation $\overrightarrow{H}$ to an orientation $\overrightarrow{G}$ of the original hypergraph $G$ that still contains $k$ hyperarc-disjoint paths 
%from $r$ to $v$. 
from $v$ to $r$. 
We use Lemma \ref{lem:lifting-Steiner-orientation-through-splitting-off} below for this extension. 
Lemma \ref{lem:lifting-Steiner-orientation-through-splitting-off} below proves a more general statement for arbitrary number of terminals and for a single h-splitting-off operation. We recall that a complete h-splitting-off at a vertex from a hypergraph is a sequence of h-splitting-off operations. 
Lemma \ref{lem:lifting-Steiner-orientation-through-splitting-off} can be applied inductively starting from $\overrightarrow{H}$ to arrive at an orientation $\overrightarrow{G}$ of the hypergraph $G$ such that there exist $k$ hyperarc-disjoint paths 
%from $r$ to $v$ 
from $v$ to $r$ 
in $\overrightarrow{G}$. Such an orientation immediately gives the $k$ hyperedge-disjoint paths between $r$ and $v$ in $G$. This concludes our proof of Menger's theorem in hypergraphs and graphs (undirected). 
%We note that the proof of \Cref{thm:complete-splitting-off} does not rely on Menger's theorem, so this is a self-complete proof. 


\begin{lemma}\label{lem:lifting-Steiner-orientation-through-splitting-off}
    Let $G=(V, E_G)$ be a hypergraph, $T\subseteq V$ be a set of terminals, $r\in T$ be a specified root vertex, and $s\in V-T$. 
    Let $H = (V, E_H)$ be a hypergraph obtained by applying a h-splitting-off operation at the vertex $s$ from $G$. 
    Suppose there exists an orientation $\overrightarrow{H}$  
    %= (V, E_G, \head_G:E\rightarrow V)$ 
    of $H$ such that $\overrightarrow{H}$ is Steiner rooted $k$-hyperarc-connected. 
    Then, there exists an orientation 
    $\overrightarrow{G}$ 
    %= (V, E_H, \head_H:E_H\rightarrow V)$ 
    of $G$ such that $\overrightarrow{G}$ is Steiner rooted $k$-hyperarc-connected. 
\end{lemma}
\begin{proof}
 Let $\overrightarrow{H}= (V, E_H, \head_H:E_H\rightarrow V)$ be an orientation of $H$ that is Steiner rooted $k$-hyperarc-connected.  
  We will use the orientation $\overrightarrow{H}$ to define an orientation $\overrightarrow{G}=(V, E_G, \head_G: E_G\rightarrow V)$ for the hypergraph $G$. 
  We recall that $H$ is obtained by applying a h-splitting-off operation at $s$ from $G$. We have two cases: 
  %In particular, we define $\head_{H}:E_H\rightarrow V$ in two cases as follows:
\begin{enumerate}
    \item Suppose that $H$ is obtained 
    from $G$ by a h-merge almost disjoint hyperedges operation. Let $e, f \in \delta_{G}(s)$ be the hyperedges with $e\cap f = \{s\}$ that are h-merged during this operation. For notational convenience, we denote $g:= e\cup f$ as the hyperedge formed after the h-merge operation. Then, we define the orientation $\head_{G}:E_G\rightarrow V$ as follows: for each $e' \in E_G$,
$$\head_{G}(e') := \begin{cases}
    s& \text{ if } e' \in \{e, f\} \text{ and } \head_{H}(g) \not \in e'\\
    
    \head_{H}(g) &\text{ if } e' \in \{e, f\} \text{ and } \head_{H}(g) \in e'\\
    
    \head_{H}(e')& \text{ otherwise (i.e., $e'\in E_G-\{e, f\}$).}
\end{cases}$$
    %\knote{Replace ``otherwise'' above by ``$e'\in E_H-\{e, f\}$''}

    \item Suppose that $H$ is obtained from $G$ 
    %by a h-splitting-off at $s$ operation is a 
    by a h-trim hyperedge operation. Let $e \in \delta_{G}(s)$ be the hyperedge that is h-trimmed during this operation. For notational convenience, we denote $g:= e - \{s\}$ as the hyperedge formed after the h-trim operation. We define an orientation $\head_{G}:E_G\rightarrow V$ as follows: for each $e' \in E_G$,
$$\head_{G}(e') := \begin{cases}
    \head_{H}(g)& \text{ if } e' = e\\
    \head_{H}(e')& \text{ otherwise (i.e., $e'\in E_G-\{e\}$).}
\end{cases}$$  
%\knote{Replace ``otherwise'' above by ``$e'\in E_H-\{e\}$''}
\end{enumerate}
        We now show that the orientation $\overrightarrow{G} = (V, E_G, \head_G)$ is Steiner rooted $k$-hyperarc-connected. We will prove this only for the first case above, i.e. $H$ is obtained from $G$ by a h-merge almost disjoint hyperedges operation, and remark that the proof for the second case is along similar lines. Let $v \in T - r$ be a terminal. Since $\overrightarrow{H}$ is a Steiner rooted $k$-hyperarc-connected orientation of $H$, there exist $k$ hyperarc-disjoint paths $\calP := \{P_1, \ldots, P_k\}$ 
        %from $r$ to $v$ 
        from $v$ to $r$ 
        in $\overrightarrow{H}$. We will use the set of paths $\calP$ to show that there are also $k$ hyperarc-disjoint paths 
        %from $r$ to $v$ 
        from $v$ to $r$ 
        in $\overrightarrow{G}$. We recall that $e, f \in \delta_G(s)$ are the hyperedges that are h-merged and $g = e \cup f $ is the hyperedge formed during the splitting-off almost disjoint hyperedges operation. First, suppose that the hyperarc $(g, \head_H(g))$ is not present in any of the paths of $\calP$. Then, $\calP$ is a set of $k$ hyperarc-disjoint paths 
        %from $r$ to $v$ 
        from $v$ to $r$ 
        in $\overrightarrow{G}$ and the claim holds. Next, suppose that the hyperarc $(g, \head_H(g))$ is present in the path $P_j \in \calP$ for some $j \in [k]$. We note that since the paths in $\calP$ are hyperarc-disjoint, $P_j$ is the unique path containing the hyperarc $(g, \head_H(g))$ and the paths in $\calP - \{P_j\}$ are hyperarc-disjoint in $\overrightarrow{G}$. We now define a new path $\Tilde{P_j}$ 
        %from $r$ to $v$ 
        from $v$ to $r$ 
        in $\overrightarrow{G}$ that is hyperarc-disjoint from the paths in $\calP - \{P_j\}$. Let $\ell \in \Z_+$ denote the length of the path $P_j$ and let the vertices and hyperarcs along the path $P_j$ be 
        $$P_j = \left(v, \left(e_0, u_1\right), u_1, \left(e_1, u_2\right), u_2,  \ldots, u_{i-1}, \left(e_{i-1}, u_{i}\right), u_i, \left(g, u_{i+1}\right), u_{i+1}, \left(e_{i+1}, u_{i+2}\right), u_{i+1}, \ldots, u_{\ell}, \left(e_\ell, r\right), r\right).$$ 
        If the vertices $u_i, u_{i+1} \in e$, then we define the path $\Tilde{P}_j$ by replacing the hyperarc $(g, \head_H(g))=(g, u_{i+1})$ with the hyperarc $(e, \head_G(e))=(e, u_{i+1})$ in $P_j$ (the case where $u_i, u_{i+1} \in f$ can be shown similarly). 
        Suppose that $u_i \in e-f$ and $u_{i+1} \in f-e$: then, we define the path $\Tilde{P}_j$ by replacing the hyperarc $(g, \head_H(g))=(g,u_{i+1})$ with 
        %the hyperarcs $(e, s), (f, u_{i+1})$ in $P_j$ 
        the sub-path $(e, s), s, (f, u_{i+1})$ in $P_j$; we note that $\head_G(e)=s$, $s$ is a tail of $f$, and $\head_G(f)=u_{i+1}$ (the case where $u_i \in f-e$ and $u_{i+1} \in e-f$ can be shown similarly). 
        In both cases, the path $\Tilde{P}_j$ is a $v$ to $r$ path in the directed hypergraph $\overrightarrow{H}$. Moreover, the $k$ paths in $\calP - \{P_j\} + \{\Tilde{P}_j\}$ are hyperarc-disjoint. Consequently, $\overrightarrow{G}$ is a Steiner rooted $k$-hyperarc-connected orientation of $G$.
\end{proof}

\paragraph{The case of $|T|=|V|$. } \textsc{Max Steiner Rooted-Connected Orientation} problem where all vertices are terminals is polynomial-time solvable owing to a characterization for the existence of rooted $k$-hyperarc-connected orientation via weak-partition-connectivity. We define weak-partition-connectivity now. 
\begin{definition}\label{defn:weak-partition-connectivity}
    Let $k \in \Z_{\geq 0}$ be a non-negative integer. A hypergraph $H = (V, E)$ is \emph{$k$-weakly-partition-connected} if for every partition $\calP$ of $V$, we have that $$\sum_{e \in E} \left(\calP\left(e\right) - 1\right) \geq k (|\calP| - 1),$$
    where for each $e \in E$, $\calP(e)$ denotes the number of parts of $\calP$ that have non-empty intersection with the hyperedge $e$. 
    %We denote $\alpha_H := \min\{k \in \Z_{\geq 0} : H \text{ is $k$-weakly partition-connected}\}$ as the \emph{weak-partition-connectivity} of $H$. \knote{Could remove $\alpha_H$.}
\end{definition}

For a hypergraph $H=(V, E)$ and a root vertex $r\in V$, we recall that an orientation $\overrightarrow{H}=(V, E, \head: E\rightarrow V)$ of $H$ is rooted $k$-hyperarc-connected if for each $u\in V$, there exist $k$ hyperarc-disjoint paths 
%from the root vertex $r$ to the vertex $u$. 
from the vertex $u$ to the root vertex $r$. 
The next theorem is a consequence of Theorem 9.4.13 in \cite{Frank-book} and shows that weak-partition-connectivity characterizes the existence of rooted hyperarc-connected orientations. 

\begin{theorem}[Theorem 9.4.13 in \cite{Frank-book}]\label{thm:frank:weak-pc-iff-rooted-hyparc-conn-orientation}
    Let $H = (V, E)$ be an undirected hypergraph and $r \in V$ be a root vertex. There exists a rooted $k$-hyperarc-connected orientation of $H$ if and only if $H$ is $k$-weakly-partition-connected.
\end{theorem}
Theorem \ref{thm:frank:weak-pc-iff-rooted-hyparc-conn-orientation} leads to a polynomial-time algorithm to solve \textsc{Max Steiner Rooted-Connected Orientation} problem where all vertices are terminals. 
%We will now use Lemma \ref{lem:lifting-Steiner-orientation-through-splitting-off} and Theorem \ref{thm:frank:weak-pc-iff-rooted-hyparc-conn-orientation} to prove Theorems \ref{thm:rooted-steiner-orientation-in-hypergraphs} and \ref{thm:rooted-steiner-orientation-in-graphs}. 

\paragraph{Proof of Theorems \ref{thm:rooted-steiner-orientation-in-hypergraphs} and \ref{thm:rooted-steiner-orientation-in-graphs}.}
We will now use Lemma \ref{lem:lifting-Steiner-orientation-through-splitting-off} and Theorems \ref{thm:frank:weak-pc-iff-rooted-hyparc-conn-orientation} and \ref{thm:complete-splitting-off} to prove Theorems \ref{thm:rooted-steiner-orientation-in-graphs} and \ref{thm:rooted-steiner-orientation-in-hypergraphs}. 
We will prove Theorem \ref{thm:rooted-steiner-orientation-in-hypergraphs} and remark after the proof about how our proof also implies Theorem \ref{thm:rooted-steiner-orientation-in-graphs}. 
To prove Theorem \ref{thm:rooted-steiner-orientation-in-hypergraphs}, we relate connectivity and weak-partition-connectivity. We recall that a hypergraph $H=(V, E)$ is defined to be $k$-hyperedge-connected if $d_H(X)\ge k$ for every non-empty proper subset $X\subseteq V$. 
%By Menger's theorem, we know that a hypergraph $H=(V, E)$ is $k$-hyperedge-connected if and only if there exist $k$ hyperedge-disjoint paths between every pair of distinct vertices $u, v\in V$. 
It is easy to see that if a hypergraph is $k$-weakly-partition-connected, then it is also $k$-hyperedge-connected, but the converse is not necessarily true (counterexamples exist even for graphs). We show an approximate converse: if a hypergraph is $2k$-hyperedge-connected, then it is $k$-weak-partition-connected.
\begin{lemma}\label{lem:2k-conn-implies-k-weak-pc}
    %Let $H = (V, E)$ be an undirected hypergraph such that $\lambda_H$ is even. Then, $\lambda_H \leq 2 \alpha_H$.
    Let $H = (V, E)$ be an undirected hypergraph that is $2k$-hyperedge-connected for some non-negative integer $k$. Then, $H$ is $k$-weak-partition-connected.
\end{lemma}
\begin{proof}
    %The inequality $\alpha_H\leq \lambda_H$ is straightforward since the partition $\calP = (P_1, P_2)$ of $V$ where $P_1$ and $P_2$ are minimum cuts of $H$ (i.e. $d_H(P_1) = d_{H}(P_2) = \lambda_H$) satisfy $\sum_{e\in E}(\calP(e) - 1) = \lambda_H (|\calP| - 1)$. We now show the second inequality of the lemma.
%    \knote{Pick an arbitrary partition of $V$, contract the parts, and argue that the relevant ratio is at least $k$. To argue this, use the fact that the cut value of each part is at least $2k$. Such a proof would avoid the floor issue and would be simpler to follow.}
   % \shubhang{why does the proof strategy you mentioned avoid the floor issue? Perhaps having a single sequence of inequalities will make the proof easier to follow (see modification below)?}
    Let $\calP = (P_1, P_2, \ldots, P_t)$ %denote the partition of $V$ witnessing $\alpha_H$, 
    be a partition of $V$ such that $t > 1$.
    %i.e. $$\calP := \arg\min\left\{\left\lfloor \frac{\sum_{e \in E}(\calP(e) - 1)}{|\calP| - 1}\right\rfloor : \text{$\calP$ is a partition of $V$ s.t. $|\calP| > 1$}\right\}.$$
    Let $H' = (V', E')$ be the hypergraph obtained after contracting the parts of $\calP$ and discarding the singleton hyperedges. We note that $|V'| = t$ and $|e| \geq 2$ for each $e \in E'$. 
    %We first show a simple upper bound on $\lambda_H$.
    The following claim gives a useful intermediate step.
    \begin{claim} \label{claim:weak-partition-conn-vs-min-cut}
        $\frac{\sum_{e \in E'}|e|}{t} \leq \frac{2\sum_{e \in E'}(|e| - 1)}{t - 1}.$
    \end{claim}
    \begin{proof}
    First, we observe that $2t|E'| \leq \sum_{e\in E'}|e|$ because $|e| \geq 2$ for each $e \in E'$.  
    Adding $(t-1)\sum_{e \in E'}|e| - 2t|E'|$ to the LHS of the previous equation, and $(2t-1)\sum_{e \in E'}|e| - 2t|E'|$ to the RHS of the previous equation, we obtain the following inequality: $(t-1)\sum_{e \in E'}|e| \leq 2t\sum_{e \in E'}|e| - 2t|E'|$.
    This inequality implies the claim since $t>1$. 
    \end{proof}
    With the above claim, we have the following:
$$2k \leq \min_{i \in [t]}\{d_{H}(P_i)\} 
            = \min_{u \in V'}\{d_{H'}(u)\}
            \leq \frac{\sum_{u \in V'}d_{H'}(u)}{t}
            = \frac{\sum_{e \in E'}|e|}{t}
            \leq \frac{2\sum_{e \in E'}(|e| - 1)}{t - 1}
            =  \frac{2\sum_{e \in E'}(\calP(e) - 1)}{|P| - 1}.$$
    Here, the first inequality is because $H$ is $2k$-hyperedge-connected and the third inequality is by \Cref{claim:weak-partition-conn-vs-min-cut}. Thus, $H$ is $k$-weak-partition-connected.
\end{proof}

We now restate and prove \Cref{thm:rooted-steiner-orientation-in-hypergraphs}.

\thmRootedSteinerArcConnectedOrientationsHypergraphs*
\begin{proof}
    Let $u_1, \ldots, u_p$ denote an arbitrary ordering of the Steiner vertices. By repeated application of \Cref{thm:complete-splitting-off}, there exists a hypergraph $H = (V, E_H)$ that can be obtained by complete h-splitting-off vertices $u_1, \ldots, u_p$ in order while preserving local connectivity. Let $H' = (T, E_H)$ be the hypergraph obtained from $H$ by deleting the isolated vertices. We note that $H$ is Steiner $2k$-hyperedge-connected, and consequently, the hypergraph $H' = (T, E_H)$ obtained from $H$ by deleting the isolated Steiner vertices is $2k$-hyperedge-connected. By \Cref{lem:2k-conn-implies-k-weak-pc}, the hypergraph $H'$ is $k$-weak-partition-connected. By \Cref{thm:frank:weak-pc-iff-rooted-hyparc-conn-orientation}, there exists a rooted $k$-hyperarc-connected orientation $\overrightarrow{H'} = (T, E_{H}, \head_{H'}:E_H \rightarrow T)$ of the hypergraph $H'$. Consequently, the orientation $\overrightarrow{H} = (V, E_H, \head_H:E_H \rightarrow V)$, where $\head_H(e) := \head_{H'}(e)$ for each $e \in E_H$, is a Steiner rooted $k$-hyperarc-connected orientation of $H$. By \Cref{lem:lifting-Steiner-orientation-through-splitting-off} and induction on the number of h-splitting-off operations needed to obtain $H$ from $G$, we have that $G$ has a Steiner rooted $k$-hyperarc-connected orientation.
\end{proof}

\begin{remark}
Certain remarks regarding our proof of Theorem \ref{thm:rooted-steiner-orientation-in-hypergraphs} are in order.
\begin{enumerate}
\item Our proof of \Cref{thm:rooted-steiner-orientation-in-hypergraphs} is constructive and leads to a polynomial-time algorithm to find a Steiner rooted $k$-hyperarc-connected orientation of a $2k$-hyperedge-connected-hypergraph. 
\item Our proof of \Cref{thm:rooted-steiner-orientation-in-hypergraphs} also implies a proof of Theorem \ref{thm:rooted-steiner-orientation-in-graphs}: to prove Theorem \ref{thm:rooted-steiner-orientation-in-graphs}, we start from a graph $G=(V, E)$, but our local connectivity preserving complete h-splitting-off operations at Steiner vertices results in a hypergraph $H=(T, E_H)$ that is $2k$-hyperedge-connected; by Lemma \ref{lem:2k-conn-implies-k-weak-pc}, the hypergraph $H$ is $k$-weak-partition-connected;   
now Theorem \ref{thm:frank:weak-pc-iff-rooted-hyparc-conn-orientation} gives a rooted $k$-hyperarc-connected orientation of the resulting hypergraph. Such an orientation is extended to a Steiner rooted $k$-arc-connected orientation of the graph $G=(V, E)$ using Lemma \ref{lem:lifting-Steiner-orientation-through-splitting-off}. 
\item Our proof of \Cref{thm:rooted-steiner-orientation-in-hypergraphs} reveals the source of the $2$-factor gap in the approximate min-max relation of Kir\'{a}ly and Lau for \textsc{Max Steiner Rooted-Connected Orientation Problem}: it arises from the $2$-factor gap between connectivity and weak-partition-connectivity of hypergraphs. This insight immediately leads to a minimal example that exhibits tightness of the $2$-factor in Theorem %\ref{thm:rooted-steiner-orientation-in-hypergraphs} and 
\ref{thm:rooted-steiner-orientation-in-graphs}: consider the cycle graph on $3$ vertices with all vertices being terminals and an arbitrary vertex chosen as the root. This graph is $2$-edge-connected, $1$-weak-partition-connected but not $2$-weak-partition-connected, and admits a Steiner rooted $2$-arc-connected orientation but does not admit a Steiner rooted $2$-arc-connected orientation. 
%the example in Figure \ref{??} which exhibits the tightness of the $2$-factor in Theorems \ref{thm:rooted-steiner-orientation-in-hypergraphs} and \ref{thm:rooted-steiner-orientation-in-graphs}.
\end{enumerate}
\end{remark}

\section{Weak to Strong Cover in Strongly Polynomial Time}\label{sec:WeakToStrongCover}
%\knote{Address the comments from SODA reviews.}\shubhang{done.}
In this section, we 
prove Theorem \ref{thm:WeakToStrongCover:main}. We restate it below. 
%show the following main result:
\thmWeakToStrongCover*

Our proof of \Cref{thm:WeakToStrongCover:main} will require the following definitions. 
\begin{definition}\label{def:(pHw)-tight-sets}
    Let $p:2^V\rightarrow\Z$ be a set function and $(H=(V,E), w)$ be a hypergraph. We will say that the hypergraph $H$ is \emph{non-empty} if $E\neq \emptyset$. 
    %A set $X \subseteq V$ is said to be \emph{tight} w.r.t. the function $p$ and hypergraph $\left(H ,w\right)$ 
    A set $X \subseteq V$ is said to be \emph{$(p, H, w)$-tight} if $b_{(H,w)}(X)=p(X)$. We let $\tightSetFamily{p,H,w}$ denote the family of $(p, H, w)$-tight sets and $\maximalTightSetFamily{p, H, w}$  denote the family of \emph{inclusion-wise maximal} sets of $\tightSetFamily{p,H,w}$. 
    For a hyperedge $e\in E$, we let $\maximalTightSetContainingHyperedge{p, H, w}(e)$ denote an arbitrary set in $\tightSetFamily{p, H, w}$ that contains $e$; if there is no such set, then we will use the convention that $\maximalTightSetContainingHyperedge{p, H, w}(e)$ is undefined.  
    Additionally, we define two operations as follows:
    \begin{enumerate}
        \item For hyperedges $e, f\in E$ and an integer  $\alpha\le \min\{w(e), w(f)\}$, the operation $\merge((H, w), e, f, \alpha)$ returns the hypergraph obtained from $(H, w)$ by decreasing the weight of hyperedges $e$ and $f$ by $\alpha$ and increasing the weight of the hyperedge $e\cup f$ by $\alpha$ (if $\alpha=w(e)$, then we discard $e$; if $\alpha=w(f)$, then we discard $f$; if $e\cup f\not\in E$, then we add the hyperedge $e\cup f$ and set its weight to be $\alpha$). 
        \item For a hyperedge $e\in E$ and an integer $\alpha \le w(e)$, the operation $\reduce((H, w), e, \alpha)$ returns the hypergraph obtained by decreasing the weight of the hyperedge $e$ by $\alpha$ (if $\alpha=w(e)$, then we discard $e$). 
    \end{enumerate}
    %\shubhang{We should either handle the discarding of the 0 weight edges here or in the algorithm. Currently we do it in both places. (we set $E' := \supp(w^R)$ in the algorithm). I am making the edit to handle it here.}
    %For a hyperedge $e \in E$, we let $\maximalTightSetContainingHyperedge{p, H, w}(e)$ denote the \emph{maximal} tight set w.r.t. $(H, w)$ and $p$ containing the hyperedge $e$. We will use the convention that $\maximalTightSetContainingHyperedge{p}(e) \coloneqq  \emptyset$ if there are no tight sets containing $e$.
\end{definition}

We first describe how to prove \Cref{thm:WeakToStrongCover:main} under the assumptions that the input function $p$ and hypergraph $\left(H = (V,E),w\right)$ satisfy the following two technical conditions:
\begin{enumerate}[label=(\alph*), ref=(\alph*)]
    \item \label{cond:WeakToStrong:(a)} %$\maximalTightSetContainingHyperedge{p, H, w}(e) \not = \emptyset$ for every $e \in E$, i.e. 
    every hyperedge $e \in E$ is contained in some $(p, H, w)$-tight set, and
    \item \label{cond:WeakToStrong:(b)} $b_{(H,w)}(u) \not = 0$ for every $u \in V$, i.e. every vertex is contained in some hyperedge.
\end{enumerate}  
Under these assumptions, we will prove \Cref{thm:WeakToStrongCover:main} using the algorithm that we describe in \Cref{sec:WeakToStrongCover:Algorithm} (see \Cref{alg:WeakToStrongCover}). In \Cref{sec:WeakToStrongCover:UncrossingProperties}, we show certain properties of set families that will be used in the description and analysis of our algorithm. In \Cref{sec:WeakToStrongCover:PartialCorrectness}, we show that our algorithm terminates, and the hypergraph returned by the algorithm satisfies properties \ref{thm:WeakToStrongCover:main:(1)} and \ref{thm:WeakToStrongCover:main:(2)} (see \Cref{lem:weak-to-strong:correctness:main:strong-cover} and \Cref{lem:weak-to-strong:correctness:main:hyperedge-merge}). In \Cref{sec:WeakToStrongCover:NumberOfHyperedges}, we show that the hypergraph returned by the algorithm additionally satisfies property \ref{thm:WeakToStrongCover:main:(3)} (see \Cref{lem:WeakToStrongCover:NumberOfHyperedges}). In \Cref{sec:WeakToStrongCover:strongly-poly-runtime:main}, we show that our algorithm can be implemented to run in the claimed strongly polynomial time bound given oracle access to \functionMaximizationOracleStrongCover{p} (see \Cref{lem:WeakToStrong:strongly-polytime:main}).
\Cref{alg:WeakToStrongCover}, \Cref{lem:weak-to-strong:correctness:main:strong-cover}, \Cref{lem:weak-to-strong:correctness:main:hyperedge-merge}, \Cref{lem:WeakToStrongCover:NumberOfHyperedges}, and \Cref{lem:WeakToStrong:strongly-polytime:main} together complete the proof of \Cref{thm:WeakToStrongCover:main} under the assumptions that conditions \ref{cond:WeakToStrong:(a)} and \ref{cond:WeakToStrong:(b)} hold.

We now describe how to circumvent assumptions \ref{cond:WeakToStrong:(a)} and \ref{cond:WeakToStrong:(b)} on the input function $p$ and hypergraph $(H,w)$. Let $e \in E$ be a hyperedge not contained in any $(p, H, w)$-tight set. Let $\alpha_e^R \coloneqq  \min\{w(e), \beta^R(e)\}$, where $$\beta^R(e) \coloneqq  \min\{b_{(H,w)}(X) - p(X) : e \subseteq X \subseteq V\}.$$
We construct the hypergraph $(H_0^e:=(V, E_0^e:=\{e\}), w_0^e:E_0\rightarrow \{\alpha^R_e\})$. Now, consider the function $p' \coloneqq  p - d_{(H_0^e, w_0^e)}$ and the  hypergraph $(H', w'):=\reduce((H, w), e, \alpha^R_e)$; we term this operation as the reduce operation with respect to the hyperedge $e$. 
%(jumping ahead, we note that this operation is the \emph{reduce} operation which we describe in the next \Cref{sec:WeakToStrongCover:Algorithm}).
Iteratively performing this reduce operation with respect to every hyperedge $e \in E$ results in a function $p_1$ and a hypergraph $(H_1 = (V, E_1), w_1)$ such that $b_{(H_1, w_1)}(X)\ge p_1(X)$ for every $X\subseteq V$ and also satisfies condition \ref{cond:WeakToStrong:(a)}. Furthermore, suppose that $\zeros \coloneqq  \{u \in V : b_{H_1, w_1}(u) = 0\}$. Then, the function $p_2 \coloneqq  \functionContract{p_1}{\zeros}$ and hypergraph $(H_2 = (V - \zeros, E_1), w_1)$ is an instance with $b_{(H_2, w_1)}(X)\ge p_2(X)$ for every $X\subseteq V$ and also satisfies conditions \ref{cond:WeakToStrong:(a)} and \ref{cond:WeakToStrong:(b)}. Finally, we observe that if $(H^*, w^*)$ is a hypergraph that strongly covers the function $p_2$, then the hypergraph $(H^*+\sum_{e \in E}H_0^e, w^*+w_0^e)$ strongly covers the original function $p$. Moreover,  if $(H^*, w^*)$ is obtained by merging hyperedges of $(H_2, w_1)$, then the  hypergraph $(H^*+\sum_{e \in E}H_0^e, w^*+w_0^e)$ is obtained by merging hyperedges of $(H, w)$. Thus, our final algorithm is to first obtain the instance $(H_2, w_1)$ and an oracle for \functionMaximizationOracleStrongCover{p_2} using \functionMaximizationOracleStrongCover{p} such that the instance $(p_2, H_2, w_1)$ is such that $(H_2, w_1)$ weakly covers $p_2$ and satisfies the technical conditions \ref{cond:WeakToStrong:(a)} and \ref{cond:WeakToStrong:(b)}, obtain a hypergraph $(H^*, w^*)$ by the procedure described in the previous paragraph, and finally return $(H^*+\sum_{e \in E}H_0^e, w^*+w_0^e)$. 
%\kristof{We should define this operation, that is, the `addition' of edge-weighted hypergraphs in the preliminaries.}\shubhang{todo.}

\subsection{Weak to Strong Cover Algorithm}\label{sec:WeakToStrongCover:Algorithm}
Our algorithm is recursive and takes two inputs: (1) a hypergraph $(H = (V, E), w)$, and (2) a symmetric skew-supermodular function $p:2^V \rightarrow\Z$. We note that in contrast to the statement of \Cref{thm:WeakToStrongCover:main}, the inputs $(H,w)$ and $p$ to the algorithm are required to satisfy the two technical conditions 
\ref{cond:WeakToStrong:(a)} and \ref{cond:WeakToStrong:(b)} specified in the previous section. The algorithm returns a hypergraph $\left(H^*= \left(V, E^*\right), w^*\right)$. 

We now give an informal description of the algorithm (refer to \Cref{alg:WeakToStrongCover} for a formal description). If $E=\emptyset$, then the algorithm returns the empty hypergraph. Otherwise, the algorithm chooses two hyperedges $e, f \in E$  such that $\maximalTightSetContainingHyperedge{p, H, w}(e)\neq \maximalTightSetContainingHyperedge{p, H, w}(f)$ ($\maximalTightSetContainingHyperedge{p, H, w}(e)$ is well-defined for every hyperedge $e\in E$ owing to condition \ref{cond:WeakToStrong:(a)}; in the next section, we show that $\maximalTightSetContainingHyperedge{p, H, w}(e)$ is unique for every hyperedge $e\in E$ under condition \ref{cond:WeakToStrong:(a)} and also that if $(H, w)$ is not a strong cover of $p$, then there exists a pair of hyperedges $e, f\in E$ such that $\maximalTightSetContainingHyperedge{p, H, w}(e)\neq \maximalTightSetContainingHyperedge{p, H, w}(f)$). 
Next, 
%the algorithm performs \merge and \reduce operations (see \Cref{def:(pHw)-tight-sets}). In particular, 
the algorithm computes the value 
 $$\beta^M \coloneqq  \min\left\{b_{(H,w)}(X) - p(X) : X \subseteq V \text{ and } e \cap X, f \cap X \not = \emptyset\right\},$$ sets $\alpha^M \coloneqq  \min\{\beta^M, w(e), w(f)\}$, and defines the hypergraph $(H^M, w^M) := \merge((H, w), e, f, \alpha^M)$.
%The algorithm then computes an intermediate weight function $w^M \coloneqq  w - \alpha^M\chi_e -\alpha^M\chi_f + \alpha^M\chi_{e\cup f}$, where $\chi_e, \chi_f, \chi_{e\cup f}\in \{0, 1\}^{E\cup\{e\cup f\}}$ are the indicator vectors of hyperedges $e, f, e\cup f$. We will refer to these two steps as \emph{merge} steps. 
Next, the algorithm computes the value $$\beta^R \coloneqq  \min\left\{b_{(H^M, w^M)}\left(X\right) - p\left(X\right) : e \cup f \subseteq X \subseteq V\right\},$$ 
    sets $\alpha^R \coloneqq  \min\{\beta^R, w^M(e\cup f)\}$, and defines the hypergraph $(H^R, w^R) := \reduce((H^M, w^M), e \cup f, \alpha^R)$.
%The algorithm then computes the new weight function $w^R \coloneqq  w^M - \alpha^R\chi_{e\cup f}$. We will refer to these two steps as \emph{reduce} steps.
It computes hypergraph $(H_0=(V, E_0:=\{e\cup f\}), w_0: E_0\rightarrow \{\alpha^R\})$ and defines $p^R:=p-d_{(H_0, w_0)}$. 
Next, the algorithm computes the set $\zeros \coloneqq  \{u \in V : b_{(H^R, w^R)}(u) = 0\}$, and recursively calls itself with the inputs 
    $(H' = (V':=V-\zeros, E':=E^R),w':=w^R) $ and $p' \coloneqq  \functionContract{p^R}{\zeros}$ to obtain a hypergraph $\left(H_0^*, w_0^*\right)$. Finally, the algorithm obtains the hypergraph $(G, c)$ from $(H_0^*, w_0^*)$ by adding the set $\zeros$ of vertices and returns the hypergraph $(G+H_0, c + w_0)$.
    %, adds the hyperedge $e\cup f$ with (additional) weight $\alpha^R$, and returns the resulting hypergraph $(G, c) := $. 
    %We formally describe our algorithm in \Cref{alg:WeakToStrongCover}.


\begin{algorithm}[!h]
\caption{Weak to Strong Cover}\label{alg:WeakToStrongCover}
$\mathbf{\textsc{Algorithm}}\left(H = (V, E), w, p\right):$
\begin{algorithmic}[1]
    \If{$E = \emptyset$} \textsc{Base Case:}
    \State{\Return{$(H, w)$.}}
    \Else:
    \State{Pick a pair of hyperedges $e, f \in E $ s.t. $\maximalTightSetContainingHyperedge{p, H, w}(e) \not = \maximalTightSetContainingHyperedge{p, H, w}(f)$}\label{alg:WeakToStrongCover:(4)}
    \State{$\alpha^M \coloneqq  \min\begin{cases}
    \beta^M \coloneqq  \min\left\{b_{(H,w)}(X) - p(X) : X \subseteq V \text{ and } e \cap X, f \cap X \not = \emptyset\right\} & \\
    \min\left\{w(e), w(f)\right\}&
    \end{cases}
    $}\label{alg:WeakToStrongCover:(5)}
    % \vspace*{0.5mm}
    %\State{$\alpha^M \coloneqq  \left\{w(e), w(f), \min\left\{b_{(H,w)}(X) - p(X) : e\cup f \subseteq X \subseteq V\right\}\right\}$}
    \State{$(H^M, w^M):=\merge\left(\left(H, w\right), e, f, \alpha^M\right)$}\label{alg:WeakToStrongCover:(6)}
    %\State{Let $(H^M, w^M)$ be the hypergraph obtained from $(H, w)$ by reducing the weight of hyperedges $e$ and $f$ by $\alpha^M$ and increasing the weight of the hyperedge $e\cup f$ by $\alpha^M$ (discard hyperedge $e$ and/or $f$ if the weight of that hyperedge has decreased to zero;  introduce a new hyperedge $e\cup f$ with weight $\alpha^M$ if $e\cup f\not\in E$)}\label{alg:WeakToStrongCover:(6)}
    %\State{$w^M \coloneqq  w - \alpha^M\chi_e - \alpha^M\chi_f + \alpha^M\chi_{e\cup f}$}\label{alg:WeakToStrongCover:(6)}
    %\vspace*{0.5mm}
    \State{$\alpha^R \coloneqq  \min \begin{cases}
        \beta^R \coloneqq  \min\left\{b_{(H, w^M)}\left(X\right) - p\left(X\right) : e \cup f \subseteq X \subseteq V\right\} &\\
        w^M(e \cup f) &
    \end{cases}$}\label{alg:WeakToStrongCover:(7)}
    %\State{$\alpha^R \coloneqq  \min\left\{w^M(e \cup f), \min\left\{b_{(H^M, w^M)}(X) - p(X) : e\cup f \subseteq X \subseteq V\right\}\right\}$}
    %\vspace*{0.5mm}
    \State{$(H^R, w^R) \coloneqq \reduce\left(\left(H^M, w^M\right), e\cup f, \alpha^R\right)$}\label{alg:WeakToStrongCover:(8)}
    %\vspace*{0.5mm}
    \State{Construct $\left(H_0:=\left(V, E_0:=\left\{e\cup f\right\}\right), w_0: E_0\rightarrow \left\{\alpha^R\right\}\right)$}\label{alg:WeakToStrongCover:(9)}
    \State{$p^R\coloneqq p - d_{(H_0, w_0)}$}\label{alg:WeakToStrongCover:(10)}
    %\State{$w^R \coloneqq  w^M - \alpha^R\chi_{e\cup f}$ and $p^R\coloneqq p - \alpha^R d_{(V, \{e\cup f\})}$}\label{alg:WeakToStrongCover:(8)}
    %\vspace*{0.5mm}
    
    \State{$\zeros \coloneqq  \left\{v \in V : b_{(H^R, w^R)}(v) = 0\right\}$}\label{alg:WeakToStrongCover:(11)}
    % \vspace*{1mm}
    
    \State{$V' \coloneqq  V - \zeros$}\label{alg:WeakToStrongCover:(12)}
    \State{$p' \coloneqq  \functionContract{p^R}{\zeros}$}\label{alg:WeakToStrongCover:(13)}
    
    \State{$(H_0^*, w_0^*):=\textsc{Algorithm}\left(H' = \left(V', E' := E^R\right), w' := w^R , p'\right)$}\label{alg:WeakToStrongCover:(14)}
    
    \State{Obtain $(G, c)$ from $(H_0^*, w_0^*)$ by adding vertices \zeros}\label{alg:WeakToStrongCover:(15)}
    
    \State{\Return{$(G+H_0, c+w_0 )$}}\label{alg:WeakToStrongCover:(16)}
    \EndIf
\end{algorithmic}
\end{algorithm}

\subsection{Properties of Tight Sets}\label{sec:WeakToStrongCover:UncrossingProperties}
%\knote{Can this section be folded into Section 4 and the entire section be moved before Section 3?}
In this section, we show properties of the set families $\tightSetFamily{p, H, w}$ and $\maximalTightSetFamily{p, H, w}$ which will be useful to show correctness and convergence of \Cref{alg:WeakToStrongCover} in subsequent sections.
%We recall that  a set $X\subseteq V$ is said to be \emph{tight} w.r.t. a function $p:2^V\rightarrow\Z$ and a hypergraph $(H = (V, E), w)$ if $p(X) = b_{(H,w)}(X)$. We let $\tightSetFamily{p,H,w}$ denote the family of tight sets and $\maximalTightSetFamily{p, H, w}$ denote the family of \emph{inclusion-wise maximal} tight sets w.r.t. the function $p$ and hypergraph $(H,w)$. We also recall that for a hyperedge $e \in E$, the set $\maximalTightSetContainingHyperedge{p, H, w}(e)$ denotes the \emph{maximal} set in $\tightSetFamily{p, H, w}$ containing $e$. 
%\knote{Include a claim that $\maximalTightSetContainingHyperedge{p, H, w}(e)$ is unique for every hyperedge $e$.} 
%We use the convention that $\maximalTightSetContainingHyperedge{p, H, w}(e) \coloneqq  \emptyset$ if no such set exists. 
In particular, \Cref{lem:weak-to-strong:uncrossing:uncrossing-lemma} below shows an important property of the family of tight sets and \Cref{cor:weak-to-strong:uncrossing:maximal-tight-set-containing-hyperedge-unique} shows that for every hyperedge $e\in E$, there exists a unique maximal $(p, H, w)$-tight set containing $e$ under condition \ref{cond:WeakToStrong:(a)}. Furthermore, \Cref{lem:weak-to-strong:uncrossing:maximal-tight-set-family-disjoint} shows that the family $\maximalTightSetFamily{p, H, w}$ is disjoint under condition \ref{cond:WeakToStrong:(b)}, and consequently, \Cref{cor:weak-to-strong:uncrossing:exist-two-hyperedges-contained-in-disjoint-maximal-tight-sets} implies  that \Cref{alg:WeakToStrongCover} \Cref{alg:WeakToStrongCover:(4)} is well-defined.

\begin{lemma}\label{lem:weak-to-strong:uncrossing:uncrossing-lemma}
Let $p:2^V\rightarrow\Z$ be a skew-supermodular function and let $\left(H = (V, E), w:E \rightarrow\Z_+\right)$ be a hypergraph such that $b_{(H,w)}(Z) \geq p(Z)$ for every $Z \subseteq V$. Let $X, Y \in \tightSetFamily{p, H, w}$ be sets such that $X\cap Y \not = \emptyset$.
If there exists a hyperedge $e \in E$ such that (1) $e$ intersects the set $X\cap Y$ and 
(2) $e$ intersects at most one of the sets $X - Y$ and $Y - X$, then $X\cup Y, X \cap Y \in \tightSetFamily{p, H, w}$.
\end{lemma}
\begin{proof}
We note that the claim holds if $X\subseteq Y$ or $Y \subseteq X$. Suppose that $X - Y, Y - X \not= \emptyset$. First, suppose that the function $p$ satisfies $p(X) + p(Y) \leq p(X - Y) + p(Y - X)$. Then, the following gives us a contradiction:
\begin{align*}
    b_{(H,w)}(X) + b_{(H,w)}(Y) 
    &= p(X) + p(Y) &\\
    & \leq p(X - Y) + p(Y - X) & \\
    & \leq b_{(H,w)}(X - Y) + b_{(H,w)}(Y - X) & \\
    & < b_{(H,w)}(X) + b_{(H,w)}(Y).
\end{align*}
Here, the first equality is because $X, Y \in \tightSetFamily{p, H, w}$. The second inequality is because $b_{(H,w)}(Z) \geq p(Z)$ for every $Z\subseteq V$. The final inequality is because the hyperedge $e$ intersects at most one of the sets $X - Y$ and $Y - X$.
Thus, the function $p$ satisfies $p(X) + p(Y) \leq p(X\cap Y) + p(X \cup Y)$. Consequently, we have that
\begin{align*}
    b_{(H,w)}(X) + b_{(H,w)}(Y) 
    &= p(X) + p(Y) &\\
    & \leq p(X \cup Y) + p(X \cap Y) &\\
    & \leq b_{(H,w)}(X\cup Y) + b_{(H,w)}(X\cap Y) &\\
    & \leq b_{(H,w)}(X) + b_{(H,w)}(Y).&
\end{align*}
Here, the first equality is because $X, Y \in \tightSetFamily{p, H, w}$. The second inequality is because $b_{(H,w)}(Z) \geq p(Z)$ for every $Z\subseteq V$. The final inequality is because the coverage function $b_{(H, w)}$ is submodular. Thus, all inequalities in the above sequence are in fact equalities, and consequently we have that $X\cup Y, X\cap Y \in \tightSetFamily{p, H, w}$.
\end{proof}

\begin{corollary}\label{cor:weak-to-strong:uncrossing:maximal-tight-set-containing-hyperedge-unique}
Let $p:2^V\rightarrow\Z$ be a skew-supermodular function and let $\left(H = (V, E), w:E \rightarrow\Z_+\right)$ be a hypergraph such that  
every hyperedge $e \in E$ is contained in some $(p, H, w)$-tight set
%condition \ref{cond:WeakToStrong:(a)} such that 
and $b_{(H,w)}(Z) \geq p(Z)$ for every $Z \subseteq V$. 
%and every hyperedge $e \in E$ is contained in some set from the family $\tightSetFamily{p, H, w}$. 
Then, for every hyperedge $e \in E$, there is a unique set from the family $\maximalTightSetFamily{p, H, w}$ containing the hyperedge $e$.
\end{corollary}
\begin{proof}
    Let $e \in E$ be a hyperedge. We note that there exists a $(p, H, w)$-tight set containing the hyperedge $e$. Consequently, there exists a set in the family \maximalTightSetFamily{p, H, w} containing the hyperedge $e$. By way of contradiction, suppose that $X, Y \in \maximalTightSetFamily{p, H, w}$ are distinct sets containing the hyperedge $e$. We note that $X - Y, Y - X \not = \emptyset$ by the maximality of the sets $X, Y \in \maximalTightSetFamily{p, H, w}$.
    Also, $X\cap Y\neq \emptyset$ since $e$ is contained in $X\cap Y$. Moreover, $e$ is disjoint from both $X-Y$ and $Y-X$. Hence, by \Cref{lem:weak-to-strong:uncrossing:uncrossing-lemma}, we have that $X\cup Y \in \tightSetFamily{p, H, w}$, contradicting maximality of the sets $X, Y \in \maximalTightSetFamily{p, H, w}$.
\end{proof}

\iffalse
By \Cref{cor:weak-to-strong:uncrossing:maximal-tight-set-containing-hyperedge-unique}, $\maximalTightSetContainingHyperedge{p, H, w}(e)$ denotes the  unique maximal $(p, H, w)$-tight set containing a hyperedge $e\in E$ when condition \ref{cond:WeakToStrong:(a)} holds.
\begin{definition}
Let $p:2^V\rightarrow\Z$ be a skew-supermodular function and let $\left(H = (V, E), w:E \rightarrow\Z_+\right)$ be a hypergraph such that $b_{(H,w)}(Z) \geq p(Z)$ for every $Z \subseteq V$ and every hyperedge $e \in E$ is contained in some $(p, H, w)$-tight set. For a hyperedge $e \in E$, we let $\maximalTightSetContainingHyperedge{p, H, w}(e) \in \maximalTightSetFamily{p, H, w}$ denote the \emph{unique maximal} $(p, H, w)$-tight containing the hyperedge $e$. 
    %We will use the convention that $\maximalTightSetContainingHyperedge{p, H, w}(e) \coloneqq  \emptyset$ if there are no tight sets containing $e$.
\end{definition}
\fi
The next lemma shows that for a \emph{skew-supermodular} function $p$ and hypergraph $(H, w)$ satisfying conditions  \ref{cond:WeakToStrong:(a)} and \ref{cond:WeakToStrong:(b)}, the family $\maximalTightSetFamily{p, H, w}$ is disjoint. The corollary following the lemma says that if the function $p$ is also \emph{symmetric} and the hypergraph $H$ is non-empty, then there exist two hyperedges that are each contained in distinct maximal $(p, H, w)$-tight sets w.r.t. $p$ and $(H, w)$.  We note that this makes \Cref{alg:WeakToStrongCover} \Cref{alg:WeakToStrongCover:(4)} well-defined.

\begin{lemma}\label{lem:weak-to-strong:uncrossing:maximal-tight-set-family-disjoint}
    Let $p:2^V\rightarrow\Z$ be a skew-supermodular function and let $\left(H = (V, E), w:E \rightarrow\Z_+\right)$ be a hypergraph 
    %satisfying conditions \ref{cond:WeakToStrong:(a)} and \ref{cond:WeakToStrong:(b)} 
    such that every hyperedge $e \in E$ is contained in some $(p, H, w)$-tight set, $b_{(H,w)}(u) \not = 0$ for every $u \in V$, and $b_{(H,w)}(X) \geq p(X)$ for every $X \subseteq V$. 
    %Furthermore, suppose that $\maximalTightSetContainingHyperedge{p, H, w}(e) \not = \emptyset$ for every $e \in E$ and $b_{(H,w)}(u) > 0$ for every $u \in V$. 
    Then, the family $\maximalTightSetFamily{p, H, w}$ is a disjoint family.
\end{lemma}
\begin{proof}
By way of contradiction, let $X, Y \in \maximalTightSetFamily{p, H, w}$ be distinct sets such that $X\cap Y \not = \emptyset$. By maximality of the sets $X, Y$, we also have that $X - Y, Y - X \not = \emptyset$. Let $u \in X\cap Y$. Since $b_{(H,w)}(u) > 0$, there exists a hyperedge $e \in E$ such that $u \in e$. We note that if either $e\cap (X - Y) = \emptyset$ or $e\cap (Y - X) = \emptyset$, then the set $X \cup Y \in \tightSetFamily{p, H, w}$ by \Cref{lem:weak-to-strong:uncrossing:uncrossing-lemma}, contradicting maximality of $X, Y \in \maximalTightSetFamily{p, H, w}$. Thus, we have that $e \cap (X - Y), e \cap (Y - X) \not= \emptyset$. Furthermore, since the hyperedge $e$ is contained in some $(p, H, w)$-tight set, we have that  $\maximalTightSetContainingHyperedge{p, H, w}(e) \not = \emptyset$. Consequently, we have that $\maximalTightSetContainingHyperedge{p, H, w}(e)\cap (X - Y), \maximalTightSetContainingHyperedge{p, H, w}(e) \cap (Y - X) \not = \emptyset$. However, $e \subseteq \maximalTightSetContainingHyperedge{p, H, w}(e)$ and thus, $e \cap (X - \maximalTightSetContainingHyperedge{p, H, w}(e)) = \emptyset$. By \Cref{lem:weak-to-strong:uncrossing:uncrossing-lemma}, we have that the set $X \cup \maximalTightSetContainingHyperedge{p, H, w}(e) \in \tightSetFamily{p, H, w}$, contradicting maximality of the set $X \in \maximalTightSetFamily{p, H, w}$.
\end{proof}

We recall that a hypergraph is non-empty if it contains at least one hyperedge. 
\begin{corollary}[\hspace{-1sp}\cite{Bernath-Kiraly}]
\label{cor:weak-to-strong:uncrossing:exist-two-hyperedges-contained-in-disjoint-maximal-tight-sets}
Let $p:2^V\rightarrow\Z$ be a symmetric skew-supermodular function and $\left(H = (V, E), w:E \rightarrow\Z_+\right)$ be a non-empty hypergraph 
%\knote{Non-emptiness is not needed since $b_{(H, w)}(u)\neq 0$ for every $u\in V$?}\shubhang{$b_{(H, w)}(u)\neq 0$ for every $u\in V$ still allows $V = \emptyset$.} \knote{Ensure that the use of the corollary guarantees this hypothesis.} 
%satisfying conditions \ref{cond:WeakToStrong:(a)} and \ref{cond:WeakToStrong:(b)} 
such that every hyperedge $e \in E$ is contained in some $(p, H, w)$-tight set, $b_{(H,w)}(u) \not = 0$ for every $u \in V$, and $b_{(H,w)}(X) \geq p(X)$ for every $X \subseteq V$. 
%Let $p:2^V\rightarrow\Z$ be a symmetric skew-supermodular function. Furthermore, let $\left(H = (V, E), w:E \rightarrow\Z\right)$ be a non-empty hypergraph such that $b_{(H,w)}(X) \geq p(X)$ for every $X \subseteq V$ such that $\maximalTightSetContainingHyperedge{p, H, w}(e) \not = \emptyset$ for every hyperedge $e \in E$ and $b_{(H,w)}(u) > 0$ for every vertex $u \in V$. 
Then, there exist distinct hyperedges $e,f \in E$ such that $\maximalTightSetContainingHyperedge{p, H, w}(e) \cap \maximalTightSetContainingHyperedge{p, H, w}(f) = \emptyset$.
\end{corollary}
\begin{proof}
Let $e \in E$ be a hyperedge.
We will show that that there exists a hyperedge $f \in E$ such that $f \subseteq V - \maximalTightSetContainingHyperedge{p, H, w}(e)$. The claim will then follow by \Cref{lem:weak-to-strong:uncrossing:maximal-tight-set-family-disjoint} since $\maximalTightSetContainingHyperedge{p, H, w}(f) \not = \emptyset$ (because the hyperedge $f$ is contained in some $(p, H, w)$-tight set) and $\maximalTightSetContainingHyperedge{p, H, w}(e), \maximalTightSetContainingHyperedge{p, H , w}(f) \in \maximalTightSetFamily{p, H, w}$.
%$, the claim follows from \Cref{lem:weak-to-strong:uncrossing:maximal-tight-set-family-disjoint}. 
By way of contradiction, suppose that there is no hyperedge contained in the set $V - \maximalTightSetContainingHyperedge{p, H, w}(e)$. Then, we have that $p(V - \maximalTightSetContainingHyperedge{p, H, w}(e)) = p(\maximalTightSetContainingHyperedge{p, H, w}(e)) = b_{(H,w)}(\maximalTightSetContainingHyperedge{p, H, w}(e)) > b_{(H,w)}(V - \maximalTightSetContainingHyperedge{p, H, w}(e))$, a contradiction. Here, the first equality is because the function $p$ is symmetric and the final inequality is because every hyperedge that intersects the set $V - \maximalTightSetContainingHyperedge{p, H, w}(e)$ also intersects the set $\maximalTightSetContainingHyperedge{p, H, w}(e)$ by our contradiction assumption.
\end{proof}


\subsection{Termination and Partial Correctness of the Algorithm}
%\Cref{alg:WeakToStrongCover}}
\label{sec:WeakToStrongCover:PartialCorrectness}
In this section, we show that every recursive call of \Cref{alg:WeakToStrongCover} is well-defined. 
%when the inputs to the base recursive call satisfy the aforementioned properties. 
Furthermore, we show that \Cref{alg:WeakToStrongCover} terminates and that the hypergraph returned by the algorithm satisfies properties \ref{thm:WeakToStrongCover:main:(1)} and \ref{thm:WeakToStrongCover:main:(2)} of \Cref{thm:WeakToStrongCover:main}. The main results of this section are \Cref{lem:weak-to-strong:correctness:main:strong-cover} and \Cref{lem:weak-to-strong:correctness:main:hyperedge-merge}. We prove certain additional lemmas in this section which will be useful in subsequent sections. 
%We will use the convention that $\Tilde{e}$ refers to an arbitrary hyperedge in $E$, where $E$ is the hyperedge set of the hypergraph $H$ input to the current recursive call and use $e, f$ to denote the hyperedges picked by \Cref{alg:WeakToStrongCover} \Cref{alg:WeakToStrongCover:(4)} in the current recursive call. \knote{The last sentence can be deleted.}
 
    Suppose that the input to \Cref{alg:WeakToStrongCover} is a non-empty hypergraph $(H = (V, E), w)$ and a symmetric skew-supermodular function $p:2^V\rightarrow \Z$ 
%satisfying conditions \ref{cond:WeakToStrong:(a)} and \ref{cond:WeakToStrong:(b)} 
such that every hyperedge $e \in E$ is contained in some $(p, H, w)$-tight set, $b_{(H,w)}(u) \not = 0$ for every $u \in V$, and $b_{(H,w)}(X) \geq p(X)$ for every $X \subseteq V$. 
  Then, by \Cref{cor:weak-to-strong:uncrossing:exist-two-hyperedges-contained-in-disjoint-maximal-tight-sets}, there exist hyperedges $e , f \in E$ such that $\maximalTightSetContainingHyperedge{p, H, w}(e) \not = \maximalTightSetContainingHyperedge{p, H, w}(f)$. In particular, \Cref{alg:WeakToStrongCover} \Cref{alg:WeakToStrongCover:(4)} is well-defined. 
  %In this section, we show that every recursive call of \Cref{alg:WeakToStrongCover} is well-defined when the inputs to the base recursive call satisfy the aforementioned properties. Furthermore, we show that \Cref{alg:WeakToStrongCover} terminates and that the hypergraph returned by the algorithm satisfies properties (1) and (2) of \Cref{thm:WeakToStrongCover:main}. The main results of this section are \Cref{lem:weak-to-strong:correctness:main:strong-cover} and \Cref{lem:weak-to-strong:correctness:main:hyperedge-merge}. 
  The next two lemmas establish several useful properties of the intermediate quantities used in the algorithm.

\begin{lemma}\label{lem:WeakToStrong:properties-of-alpha^M-and-alpha^R}
    %Consider a recursive call of \Cref{alg:WeakToStrongCover}. 
    Suppose that the input to \Cref{alg:WeakToStrongCover} is a non-empty hypergraph $(H = (V, E), w)$ and a symmetric skew-supermodular function $p:2^V\rightarrow \Z$ 
such that every hyperedge $\Tilde{e} \in E$ is contained in some $(p, H, w)$-tight set, $b_{(H,w)}(u) \not = 0$ for every $u \in V$, and $b_{(H,w)}(X) \geq p(X)$ for every $X \subseteq V$. 
    %Furthermore, suppose that $\maximalTightSetContainingHyperedge{p, H, w}(\Tilde{e}) \not = \emptyset$ for every hyperedge $\Tilde{e} \in E$ and $b_{(H,w)}(u) > 0$ for every vertex $u \in V$. 
    Let $e, f, \alpha^M, (H^M, w^M), \alpha^R$, $(H^R, w^R),(H_0, w_0), p^R$ and $(H' = (V', E'), w')$ be as defined by \Cref{alg:WeakToStrongCover}. 
    %Let $(H' = (V', E'), w^R)$ be the hypergraph and $p':2^{V'}\rightarrow\Z$ be the function computed by \Cref{alg:WeakToStrongCover}. 
    Then, we have the following:
    \begin{enumerate}[label=(\arabic*)]
        \item $\alpha^M \in \Z_+$,
        \item 
        $E^M=E-\{e\}\cdot \indicator_{\alpha^M= w(e)}- \{f\}\cdot \indicator_{\alpha^M = w(f)} + \{e\cup f\}$ with the weights of the hyperedges in $E^M$ being  
        \[
        w^M(\Tilde{e})=
        \begin{cases}
            w(\Tilde{e}) &\text { if } \Tilde{e}\in E-\{e, f, e\cup f\},\\
            w(e\cup f) + \alpha^M &\text { if } \Tilde{e}=e\cup f \text{ and } e\cup f \in E, \\
            \alpha^M &\text { if } \Tilde{e}=e\cup f \text{ and } e\cup f \not \in E, \\
            w(e) - \alpha^M &\text { if } \Tilde{e}=e \text{ and } \alpha^M<w(e), \\
            w(f) - \alpha^M &\text { if } \Tilde{e}=f \text{ and } \alpha^M<w(f).
        \end{cases}
        \]
        %$w^M = w - w_0^M + w_1^M$, where $w_0^M:\{e, f\}\rightarrow\{\alpha^M\}$ and $w_1^M:\{e\cup f\} \rightarrow\{\alpha^M\}$, 
        %\knote{I am not a fan of this notation. Instead, can we say: $E^M=E-\{e\}\cdot \indicator_{\alpha^M= w(e)}- \{f\}\cdot \indicator_{\alpha^M = w(f)} + \{e\cup f\}$ with the weights of the hyperedges in $E^M$ being  
        % \[
        % w^M(\Tilde{e})=
        % \begin{cases}
        %     w(\Tilde{e}) &\text { if } \Tilde{e}\in E-\{e, f, e\cup f\},\\
        %     w(e\cup f) + \alpha^M &\text { if } \Tilde{e}=e\cup f \text{ and } e\cup f \in E, \\
        %     \alpha^M &\text { if } \Tilde{e}=e\cup f \text{ and } e\cup f \not \in E, \\
        %     w(e) - \alpha^M &\text { if } \Tilde{e}=e \text{ and } \alpha^M<w(e), \\
        %     w(f) - \alpha^M &\text { if } \Tilde{e}=f \text{ and } \alpha^M<w(f).
        % \end{cases}
        % \]
        % }
        \item $b_{(H^M, w^M)}(X) \geq p(X)$ for every $X \subseteq V$,
        \item $\alpha^R \in \Z_{\geq 0}$, 
        \item 
        $E' = E^R = E^M - \{e\cup f\}\cdot \indicator_{\alpha^R=w^M(e\cup f)}$ with the weights of the hyperedges in $E' = E^R$ being 
        \[
        w'(\Tilde{e}) = w^R(\Tilde{e})=
        \begin{cases}
            w^M(\Tilde{e}) &\text { if } \Tilde{e}\in E^M-\{e\cup f\},\\
            w^M(e\cup f) - \alpha^R &\text { if } \Tilde{e}=e\cup f \text{ and } \alpha^R< w^M(e\cup f), \text{ and}  \\
        \end{cases}
        \]
        %$w^R = w^M - w_0$, \knote{Again, not a fan of the notation. Instead, can we say: 
        % $E^R = E^M - \{e\cup f\}\cdot \indicator_{\alpha^R<w^M(e\cup f)}$ with the weights of the hyperedges in $E^R$ being 
        % \[
        % w^R(\Tilde{e})=
        % \begin{cases}
        %     w^M(\Tilde{e}) &\text { if } \Tilde{e}\in E^M-\{e\cup f\},\\
        %     w^M(e\cup f) - \alpha^R &\text { if } \Tilde{e}=e\cup f \text{ and } \alpha^R< w^M(e\cup f), \\
        % \end{cases}
        % \]
        % }
        \item $b_{(H^R, w^R)}(X) \geq p^R(X)$ for every $X \subseteq V$. 
        %\item $\alpha^R = 0$ if and only if $\alpha^M = \beta^M$ ( i.e., $\alpha^R > 0$ if and only if $\alpha^M < \beta^M$) , and
        % \item $E' = E - \{e\}\cdot \indicator_{\alpha^M = w(e)} - \{f\}\cdot \indicator_{\alpha^M = w(f)} + \{e\cup f\}\cdot \indicator_{\alpha^R < w^M(e\cup f)}$.
        %- \{e, f\} \subseteq E' \subseteq E \cup \{e \cup f\}$,
        % \item if $\alpha^R = \beta^R < w^M(e\cup f)$, then $e \cup f \in E'$,
        % \item if $\alpha^R < \beta^R$, then $E' \subseteq E$.
    \end{enumerate}
\end{lemma}
\begin{proof} We prove each claim separately below.
\begin{enumerate}
    \item We recall that the weight function $w:E\rightarrow\Z_+$ is a positive integer-valued function, and the function $p$ is an integer-valued function. Thus, it suffices to show that $\beta^M > 0$. We note that $\beta^M \geq 0$ since $b_{(H,w)}(X) \ge p(X)$ for every $X \subseteq V$. By way of contradiction, suppose that $\beta^M = 0$. Let $X \subseteq V$ witness $\beta^M$, i.e. $e \cap X, f\cap X \not = \emptyset$ and $b_{(H,w)}(X) -p(X)= 0$.
    The latter property implies that $X \in \tightSetFamily{p, H, w}$ and thus, there exists $Y \in \maximalTightSetFamily{p, H, w}$ such that $X \subseteq Y$. 
    We recall that by the algorithm's choice of hyperedges $e, f \in E$, we have that $\maximalTightSetContainingHyperedge{p, H, w}(e) \not = \maximalTightSetContainingHyperedge{p, H, w}(f)$. Since, $\maximalTightSetContainingHyperedge{p, H, w}(e), \maximalTightSetContainingHyperedge{p, H, w}(f) \in \maximalTightSetFamily{p, H, w}$, we have that $\maximalTightSetContainingHyperedge{p, H, w}(e) \cap \maximalTightSetContainingHyperedge{p, H, w}(f) = \emptyset$ by \Cref{lem:weak-to-strong:uncrossing:maximal-tight-set-family-disjoint}. In particular, we have that $\maximalTightSetContainingHyperedge{p, H, w}(e) \not = Y$ since
    $f \subseteq V - \maximalTightSetContainingHyperedge{p, H, w}(e)$ but $f \cap Y \not = \emptyset$.
    Thus, $Y\cap \maximalTightSetContainingHyperedge{p, H, w}(e) = \emptyset$ by \Cref{lem:weak-to-strong:uncrossing:maximal-tight-set-family-disjoint}, contradicting $\emptyset\neq e\cap X\subseteq e \cap Y \subseteq \maximalTightSetContainingHyperedge{p, H, w}(e)\cap Y$. 

    \item Follows from the definition of \merge operation (see \Cref{def:(pHw)-tight-sets}).
    
    \item 
    %We recall that the weight function $w^M = w - \alpha^M\chi_e - \alpha^M \chi_f + \alpha^M\chi_{e\cup f}$. 
    Let $X \subseteq V$. If at most one of the sets $X\cap e$ or $X\cap f$ is non-empty, then by part (2) of the current lemma (shown above) we have that $b_{(H^M, w^M)}(X) = b_{(H,w)}(X) \geq p(X)$. Suppose that $e\cap X, f\cap X \not = \emptyset$. Then, we have that $\alpha^M \leq \beta^M \leq b_{(H,w)}(X) - p(X)$.
    %\min\left\{b_{(H,w)}(X) - p(X) : X \subseteq V \text{ and } e \cap X, f \cap X \not = \emptyset\right\}$. 
    Thus,  
    $b_{(H^M, w^M)}(X) = b_{(H,w)}(X) - \alpha^M \geq b_{(H,w)}(X) - \beta^M \geq p(X)$, where the equality is by part (2) of the current lemma (shown above).

    \item We note that $w^M(e\cup f) = \alpha^M \in \Z_+$ by parts (1) and (2) of the current lemma (shown above). Furthermore, $b_{H^M,w^M}(X) \geq p(X)$ for every $X\subseteq V$ by part (3) of the current lemma (shown above). Thus, $\beta^R \in \Z_{\geq 0}$ since $p$ is an integer-valued function.

    \item Follows from the definition of \reduce operation (see \Cref{def:(pHw)-tight-sets}).

    \item Let $X\subseteq V$. First, suppose that the hyperedge $e\cup f$ is not contained in the set $X$. 
    %Let $(H_0, w_0)$ be the hypergraph defined by \Cref{alg:WeakToStrongCover}. 
    Then, we have that $b_{(H^R, w^R)}(X) = b_{(H^M, w^M)}(X) - d_{(H_0, w_0)}(X) \geq p(X) - d_{(H_0, w_0)}(X) = p^R(X)$. Here, the first equality is by part (5) of the current lemma and the inequality is by part (3) of the current lemma (proofs of both parts shown above).  Next, suppose that $e\cup f \subseteq X$. Then, we have that $b_{(H^R, w^R)}(X) = b_{(H^M, w^M)}(X) - \alpha^R \geq b_{H^M,w^M}(X) - \beta^R \geq b_{(H^M, w^M)}(X) - (b_{(H^M, w^M)}(X) - p(X)) = p(X) = p^R(X)$, where the first equality is by part (5) of the current lemma (shown above) and because $e\cup f \subseteq X$, and the final equality is because $e\cup f \subseteq X$.
\iffalse
    \item We prove the claim by showing the implication in both directions.
    First, suppose that $\alpha^M = \beta^M$. Let $X \subseteq V$ witness $\beta^M$, i.e. $\beta^M = b_{(H,w)}(X) - p(X)$ and $e\cap X, f\cap X \not = \emptyset$. Then, we have that
    $b_{H^M,w^M}(X) = b_{(H,w)}(X) - \alpha^M = p(X)$, where the first equality is by part (2) of the claim (shown above). Thus, $\alpha^R \leq \beta^R \leq b_{H^M,w^M}(X) - p(X) = 0$ \shubhang{the second inequality is a bug? Why does $e\cup f \subseteq X$? -- this can be shown using lemma 8 and lemma 4.}\shubhang{todo: move this statement to the place that actually uses it (I think claim 14.2 is the only place -- need to verify.) and prove it there correctly as a claim} and consequently $\alpha^R = 0$ by part (4) of the claim (shown above). Next, suppose that $\alpha^R = 0$. Then, since $w^M(e\cup f) \geq \alpha^M > 0$ by part (1), we have that $\alpha^R = \beta^R = 0$. Thus, there exists a set $X \subseteq V$ such that $e\cup f \subseteq X$ and $b_{H,w^M}(X) = p(X)$. 
    %We note that by the algorithm's choice of hyperedges $e, f$, we have that $\maximalTightSetContainingHyperedge{p, H, w}(e) \not = \maximalTightSetContainingHyperedge{p, H, w}(f)$ and consequently $\maximalTightSetContainingHyperedge{p, H, w}(e)\cap \maximalTightSetContainingHyperedge{p, H, w}(f) = \emptyset$ by \Cref{lem:weak-to-strong:uncrossing:maximal-tight-set-family-disjoint}. Thus, there does not exist $Y \in \tightSetFamily{p, H, w}$ such that $e\cap Y, f\cap Y \not= \emptyset$. In particular, $X \not \in \tightSetFamily{p, H, w}$. 
     Then, we have the following: $$p(X) = b_{(H^M, w^M)}(X) = b_{(H,w)}(X) - \alpha^M \geq b_{(H,w)}(X) - \beta^M \geq b_{(H,w)}(X) - (b_{(H,w)}(X) - p(X)) = p(X).$$
    Thus, all inequalities are equations, and we have that $\alpha^M = \beta^M$.

    \item 
    %The claim follows from the construction of the hyperedge set $E'$.
    We recall that $E' = E^R = \supp({w^R})$ by the definition of the \reduce operation. Furthermore, $w^R = w - w_0^M + w_1^M - w_0$, where $w_0^M:\{e, f\}\rightarrow\{\alpha^M\}$, $w_1^M:\{e\cup f\} \rightarrow\{\alpha^M\}$, and $w_0:\{e\cup f\}\rightarrow\{\alpha^R\}$ as defined by \Cref{alg:WeakToStrongCover}. Thus, if $w(e) = \alpha^M$, then $w^R(e) = 0$ and consequently $e \not \in E'$. Using similar reasoning with hyperedges $f$ and $e\cup f$ gives the claim.
    \fi
\end{enumerate}
\end{proof}

With the previous properties established, we now show that the function $p'$ and hypergraph $(H', w')$ constructed by \Cref{alg:WeakToStrongCover}  as input to the subsequent recursive call of the algorithm satisfy the two hypothesis conditions of \Cref{thm:WeakToStrongCover:main}  --  the function $p'$ is symmetric skew-supermodular, and the hypergraph $(H', w')$ weakly covers the function $p'$. This is a necessary step in order to give an inductive proof of \Cref{thm:WeakToStrongCover:main}. 

\begin{lemma}\label{lem:weak-to-strong-cover:next-recursive-call-hypothesis-satification:symm-skew-supermod-strong-cover}
%Consider a recursive call of \Cref{alg:WeakToStrongCover}. 
Suppose that the input to \Cref{alg:WeakToStrongCover} is a non-empty hypergraph $(H = (V, E), w)$ and a symmetric skew-supermodular function $p:2^V\rightarrow \Z$ %satisfying conditions \ref{cond:WeakToStrong:(a)} and \ref{cond:WeakToStrong:(b)} 
such that every hyperedge $\Tilde{e} \in E$ is contained in some $(p, H, w)$-tight set, $b_{(H,w)}(u) \not = 0$ for every $u \in V$, and $b_{(H,w)}(X) \geq p(X)$ for every $X \subseteq V$. 
Let $p'$ and $(H' = (V', E'), w')$ be as defined by \Cref{alg:WeakToStrongCover}. Then, we have the following:
\begin{enumerate}
    \item the function $p'$ is symmetric skew-supermodular, and
    \item $b_{(H', w')}(X) \geq p'(X)$ for every $X \subseteq V'$.
\end{enumerate}
\end{lemma}
\begin{proof} We prove each claim separately below. Let $(H_0, w_0), p^R, \zeros,$ and $(H^R, w^R)$ be as defined by \Cref{alg:WeakToStrongCover}.
\begin{enumerate}
    \item We recall that $p^R =  p - d_{(H_0, w_0)}$. We first observe that the function $p^R$ is symmetric skew-supermodular because the function $p$ is symmetric skew-supermodular and the cut function $d_{(H_0, w_0)}$ is symmetric submodular. Next, we observe that the function $p' = \functionContract{p^R}{\zeros}$. The claim then follows because the contraction operation preserves skew-supermodularity and symmetry.
    \item Let $X \subseteq V'$. Let $R \subseteq \zeros$ be such that $p'(X) = \functionContract{p^R}{\zeros}(X) = p^R(X \cup R)$. Then, we have the following: $b_{(H', w')}(X) = b_{(H^R, w^R)}(X \cup R) \geq p^R(X \cup R) = p'(X)$. Here, the first equality is because $w' = w^R$ and $b_{(H^R, w^R)}(u) = 0$ for every $u \in R$ by the definition of $\zeros$. The inequality is by \Cref{lem:WeakToStrong:properties-of-alpha^M-and-alpha^R}(6).  
\end{enumerate}
\end{proof}

In order to show that the function $p'$ and $(H',w')$ constructed by \Cref{alg:WeakToStrongCover} satisfy conditions \ref{cond:WeakToStrong:(a)} and \ref{cond:WeakToStrong:(b)}, we examine how the families of $(p, H, w)$-tight sets change during a recursive call of the algorithm. 
%We recall that  a set $X\subseteq V$ is said to be \emph{tight} w.r.t. a function $p:2^V\rightarrow\Z$ and a hypergraph $(H = (V, E), w)$ if $p(X) = b_{(H,w)}(X)$. We also recall that $\tightSetFamily{p,H,w}$ denotes the family of tight sets w.r.t. the function $p$ and hypergraph $(H,w)$. 
The next lemma shows that $(p, H, w)$-tight sets persist through the merge and reduce steps of a recursive call. Furthermore, the lemma also shows that the projection of a $(p, H, w)$-tight set onto the ground set of the subsequent recursive call is  $(p', H', w')$-tight.

\begin{lemma}\label{lem:WeakToStrongCover:TightSetFamiliesDuringAlgorithm}
%Consider a recursive call of \Cref{alg:WeakToStrongCover}. Suppose that the input to the recursive call is a non-empty hypergraph $(H = (V, E), w)$ and a symmetric skew-supermodular function $p:2^V\rightarrow \Z$ such that $b_{(H,w)}(X) \geq p(X)$ for every $X \subseteq V$. Furthermore, suppose that $\maximalTightSetContainingHyperedge{p, H, w}(\Tilde{e}) \not = \emptyset$ for every hyperedge $\Tilde{e} \in E$ and $b_{(H,w)}(u) > 0$ for every vertex $u \in V$.
%Consider a recursive call of \Cref{alg:WeakToStrongCover}. 
Suppose that the input to \Cref{alg:WeakToStrongCover} is a non-empty hypergraph $(H = (V, E), w)$ and a symmetric skew-supermodular function $p:2^V\rightarrow \Z$ %satisfying conditions \ref{cond:WeakToStrong:(a)} and \ref{cond:WeakToStrong:(b)} 
such that every hyperedge $\Tilde{e} \in E$ is contained in some $(p, H, w)$-tight set, $b_{(H,w)}(u) \not = 0$ for every $u \in V$, and $b_{(H,w)}(X) \geq p(X)$ for every $X \subseteq V$. 
Let $(H^M, w^M), (H^R, w^R), p'$ 
and $(H' = (V', E'), w')$ be as defined by \Cref{alg:WeakToStrongCover}.
Then, we have the following:
     \begin{enumerate}
        \item $\tightSetFamily{p, H, w} \subseteq \tightSetFamily{p, H^M, w^M} \subseteq \tightSetFamily{p^R, H^R, w^R}$, and 
        \item $\tightSetFamily{p^R, H^R, w^R}|_{V'} \subseteq \tightSetFamily{p', H', w'}$.
    \end{enumerate}
\end{lemma}
\begin{proof} We prove the claims separately below. Let $e, f, \alpha^M, \zeros$ be as defined by \Cref{alg:WeakToStrongCover}.
\begin{enumerate}
    \item We will first show that $\tightSetFamily{p, H, w} \subseteq \tightSetFamily{p, H^M, w^M}$. Let $X \in \tightSetFamily{p, H, w}$. 
    %Here, we observe that $e\cap X, f\cap X \not = \emptyset$ as otherwise 
    Suppose $e\cap X\neq \emptyset$ and $f\cap X\neq \emptyset$, then 
        $b_{(H^M, w^M)}(X) = b_{(H,w)}(X) - \alpha^M = p(X) - \alpha^M < p(X)$, which is a contradiction. Here, the first equality is because $e\cap X\neq \emptyset$ and $f\cap X\neq \emptyset$ and the final inequality is by \Cref{lem:WeakToStrong:properties-of-alpha^M-and-alpha^R}(1). 
        So, at most one of $e\cap X$ and $f\cap X$ is non-empty and consequently, $b_{(H,w)}(X)=b_{(H^M, w^M)}(X)$. 
        Thus, we have that $p(X) = b_{(H,w)}(X) = b_{(H^M, w^M)}(X)$, i.e. $X \in \tightSetFamily{p, H^M, w^M}$. 

        We will now show the second inclusion in the claim, i.e. $\tightSetFamily{p, H^M, w^M} \subseteq \tightSetFamily{p^R, H^R, w^R}$. Let $X \in \tightSetFamily{p, H^M, w^M} $. First, suppose that $\alpha^R = 0$. Then, we have that $b_{(H^R, w^R)}(X) = b_{(H^M, w^M)}(X) = p(X) = p^R(X)$, where the first and final equalities are because $\alpha^R = 0$. Thus, $X \in \tightSetFamily{p^R, H^R, w^R}$. Next, suppose that $\alpha^R > 0$. We note that if $(e\cup f) \cap X = \emptyset$, then $b_{(H^R,w^R)}(X) = b_{(H^M,w^M)}(X) = p(X) = p^R(X)$ and so $X \in \tightSetFamily{p^R, H^R, w^R}$. Consider the case where $(e\cup f) \cap X \not = \emptyset$. Here, we may assume that 
        $e\cup f\in \delta(X)$ since otherwise, 
        we obtain the following contradiction  $b_{(H^R, w^R)}(X) = b_{(H^M, w^M)}(X) - \alpha^R = p(X) - \alpha^R < p(X) = p^R(X)$, where the final equality is because $d_{(H_0, w_0)}(X) = 0$ since $e\cup f \not \in \delta_H(X)$. 
        Consequently, we have that $b_{(H^R, w^R)}(X) = b_{(H^M, w^M)}(X) - \alpha^R = p(X) - \alpha^R = p^R(X)$, and the claim holds. 
        
        \item Let $X \in \tightSetFamily{p^R, H^R, w^R}$.  Let $\zeros$ be as defined by \Cref{alg:WeakToStrongCover}. Then, we have the following:
        $$p'(X - \zeros) = \functionContract{p^R}{\zeros}(X - \zeros) \geq p^R(X) = b_{(H^R, w^R)}(X) = b_{(H^R, w^R)}(X - \zeros) = b_{(H', w')}(X - \zeros),$$
        where the penultimate equality is because $\zeros = \{u \in V : b_{(H^R, w^R)}(u) = 0\}$. However, by \Cref{lem:weak-to-strong-cover:next-recursive-call-hypothesis-satification:symm-skew-supermod-strong-cover}, we have that $b_{H',w'}(X - \zeros) \geq p'(X - \zeros)$. Thus, the inequality in the above sequence is an equation and we have that $X - \zeros \in \tightSetFamily{p', H', w'}$.
\end{enumerate}
\end{proof}
%With the previous relationships between the families $\tightSetFamily{p, H, w}, \tightSetFamily{p, H, w^M}, \tightSetFamily{p^R, H, w^R}$ and $\tightSetFamily{p', H', w'}$ established, 
We now show another useful property of \Cref{alg:WeakToStrongCover} which we will need while reasoning about the convergence of \Cref{alg:WeakToStrongCover} in the next section.
\begin{lemma}\label{lem:WeakToStrong:alpha^R=0-iff-alpha^M=beta^M}
Suppose that the input to \Cref{alg:WeakToStrongCover} is a non-empty hypergraph $(H = (V, E), w)$ and a symmetric skew-supermodular function $p:2^V\rightarrow \Z$
such that every hyperedge $\Tilde{e} \in E$ is contained in some $(p, H, w)$-tight set, $b_{(H,w)}(u) \not = 0$ for every $u \in V$, and $b_{(H,w)}(X) \geq p(X)$ for every $X \subseteq V$. 
 Let $\alpha^M, \alpha^R$ and $\beta^M$ be as defined by \Cref{alg:WeakToStrongCover}.  
 Then, $\alpha^R = 0$ if and only if $\alpha^M = \beta^M$.
 %(i.e., $\alpha^R > 0$ if and only if $\alpha^M < \beta^M$) 
\end{lemma}
\begin{proof}
    We prove the claim by showing the implication in both directions. Let $e, f, (H^M, w^M), \beta^R$ be as defined by \Cref{alg:WeakToStrongCover}. 
    %\knote{Add all other quantities from the algorithm that are used in the proof to the above sentence.}
    First, suppose that $\alpha^M = \beta^M$. Let $X \subseteq V$ be such that $\beta^M = b_{(H,w)}(X) - p(X)$ and $e\cap X, f\cap X \not = \emptyset$ (such a set exists by definition of $\beta^M$). Then, we have that
    $b_{(H^M,w^M)}(X) = b_{(H,w)}(X) - \alpha^M = p(X)$, where the first equality is by \Cref{lem:WeakToStrong:properties-of-alpha^M-and-alpha^R}(2). Thus, $X \in \tightSetFamily{p, H^M, w^M}$. Let $Y \in \maximalTightSetFamily{p, H^M, w^M}$ such that $X \subseteq Y$. \Cref{claim:WeakToStrong:new-tight-set-containing-e-and-f} below says that the set $Y$ contains the hyperedge $e \cup f$. With this claim, we complete the proof of the forward direction as follows: By \Cref{claim:WeakToStrong:new-tight-set-containing-e-and-f}, we have that $\alpha^R \leq \beta^R \leq b_{H^M,w^M}(Y) - p(Y) = 0$, where the final equality is because $Y \in \tightSetFamily{p, H^M, w^M}$. Consequently, $\alpha^R = 0$ because $\alpha^R \geq 0$ by \Cref{lem:WeakToStrong:properties-of-alpha^M-and-alpha^R}(4). 
\begin{claim}\label{claim:WeakToStrong:new-tight-set-containing-e-and-f}
       $e \cup f \subseteq Y$.
    \end{claim}
    \begin{proof}
    By way of contradiction, say $(e\cup f) - Y \not = \emptyset$. Thus, $e - Y \not = \emptyset$ or $f - Y \not = \emptyset$ (or both). We consider the case where $e - Y \not = \emptyset$ and remark that the proof of the case where $f - Y \not = \emptyset$ is along the same lines.
        We recall that the hyperedge $e$ is contained in some $(p, H, w)$-tight set, and thus the set $\maximalTightSetContainingHyperedge{p, H, w}(e)$ is well-defined and non-empty. Furthermore, by \Cref{lem:WeakToStrongCover:TightSetFamiliesDuringAlgorithm}(1), we have that the set $\maximalTightSetContainingHyperedge{p, H, w}(e) \in \tightSetFamily{p, H^M, w^M}$. Consequently, there exists a set $Z \in \maximalTightSetFamily{p, H^M, w^M}$ such that $e \subseteq \maximalTightSetContainingHyperedge{p, H, w}(e) \subseteq Z$. Since $e - Y \not = \emptyset$ but $e \subseteq Z$, we have that $Y \not = Z$.  We recall that $e \cap X \not = \emptyset$ by our choice of the set $X$, and hence, $e \cap Y \not = \emptyset$. Furthermore, by \Cref{lem:WeakToStrong:properties-of-alpha^M-and-alpha^R}(3), we have that $b_{(H^M, w^M)}(A) \geq p(A)$ for every $A \subseteq V$. Thus, by \Cref{lem:weak-to-strong:uncrossing:uncrossing-lemma} applied to the function $p$, hypergraph $(H^M, w^M)$, sets $Y$ and $Z$, and the hyperedge $e$, we have that the set $Y \cup Z \in \maximalTightSetFamily{p, H^M, w^M}$, a contradiction to the maximality of the sets $Y, Z \in \maximalTightSetFamily{p, H^M, w^M}$.
    \end{proof}
    
    Next, we prove the reverse direction of the claim. Suppose that $\alpha^R = 0$. 
    %By way of contradiction, say $\alpha^M < \beta^M$. 
    We note that $w^M(e\cup f) \geq \alpha^M > 0$ by Lemma \ref{lem:WeakToStrong:properties-of-alpha^M-and-alpha^R}(2) and \ref{lem:WeakToStrong:properties-of-alpha^M-and-alpha^R}(1). Thus, we have that $\alpha^R < w^M(e\cup f)$ and so $\alpha^R = \beta^R$. It follows that $\beta^R=0$. Consequently, there exists a set $X \subseteq V$ such that $e\cup f \subseteq X$ and $b_{(H^M,w^M)}(X) = p(X)$ (by definition of $\beta^R$). 
     Then, we have the following: $$p(X) = b_{(H^M, w^M)}(X) = b_{(H,w)}(X) - \alpha^M \geq b_{(H,w)}(X) - \beta^M \geq b_{(H,w)}(X) - (b_{(H,w)}(X) - p(X)) = p(X).$$
    Thus, all inequalities are equations, and we have that $\alpha^M = \beta^M$. Here, the second equality is by \Cref{lem:WeakToStrong:properties-of-alpha^M-and-alpha^R}(2) and $e\cup f \subseteq X$. The second inequality is also because $e\cup f \subseteq X$.
\end{proof}

 We now show that if the inputs to \Cref{alg:WeakToStrongCover} satisfy conditions \ref{cond:WeakToStrong:(a)} and \ref{cond:WeakToStrong:(b)}, then the inputs to every subsequent recursive call of \Cref{alg:WeakToStrongCover} also satisfy the two conditions.

%\knote{$e$ is overloaded}\shubhang{changed to $\Tilde{e}$. Also added a comment at the start of the section about this.}
\begin{lemma}\label{lem:weak-to-strong-cover:next-recursive-call-hypothesis-satification:hyperedge-in-tight-set-positive-coverage}
%Consider a recursive call of \Cref{alg:WeakToStrongCover}. Suppose that the input to the recursive call is a non-empty hypergraph $(H = (V, E), w)$ and a symmetric skew-supermodular function $p:2^V\rightarrow \Z$ such that $b_{(H,w)}(X) \geq p(X)$ for every $X \subseteq V$. Furthermore, suppose that $\maximalTightSetContainingHyperedge{p, H, w}(\Tilde{e}) \not = \emptyset$ for every hyperedge $\Tilde{e} \in E$ and $b_{(H,w)}(u) > 0$ for every vertex $u \in V$. 
%Consider a recursive call of \Cref{alg:WeakToStrongCover}. 
Suppose that the input to \Cref{alg:WeakToStrongCover} is a non-empty hypergraph $(H = (V, E), w)$ and a symmetric skew-supermodular function $p:2^V\rightarrow \Z$ %satisfying conditions \ref{cond:WeakToStrong:(a)} and \ref{cond:WeakToStrong:(b)} 
such that every hyperedge $\Tilde{e} \in E$ is contained in some $(p, H, w)$-tight set, $b_{(H,w)}(u) \not = 0$ for every $u \in V$, and $b_{(H,w)}(X) \geq p(X)$ for every $X \subseteq V$. 
Let $(H' = (V', E'), w')$ and $p'$ be as defined by \Cref{alg:WeakToStrongCover}. Then, we have the following:
\begin{enumerate}
    \item every hyperedge $\Tilde{e} \in E'$ is contained in some $(p', H', w')$-tight set, and
%$\maximalTightSetContainingHyperedge{p'}(\tilde{e}) \not = \emptyset$ for every hyperedge $\tilde{e} \in E'$, and
    \item  $b_{H',w'}(u) > 0$ for every vertex $u \in V'$.
\end{enumerate}
%In particular, we have that the hypergraph $(H',w')$ and the function $p'$ satisfy conditions \ref{cond:WeakToStrong:(a)} and \ref{cond:WeakToStrong:(b)}.
\end{lemma}
\begin{proof} We prove the claims separately below. Let $e, f, \zeros, \alpha^R, \beta^R, p', (H^M, w^M), (H^R, w^R)$ be as defined by \Cref{alg:WeakToStrongCover}.
%\knote{Let ... be as defined by \Cref{alg:WeakToStrongCover}.}
\begin{enumerate}
    \item We consider two cases. First, suppose that $\Tilde{e} \in E \cap E'$. Since the hyperedge $\Tilde{e}$ is contained in some $(p, H, w)$-tight set, we have that $\maximalTightSetContainingHyperedge{p, H, w}(\Tilde{e})$ is well-defined and non-empty. Furthermore, by \Cref{lem:WeakToStrongCover:TightSetFamiliesDuringAlgorithm}(2), we have that $\maximalTightSetContainingHyperedge{p, H, w}(\Tilde{e})|_{\zeros} \in \tightSetFamily{p', H', w'}$. Consequently, $\Tilde{e} \subseteq \maximalTightSetContainingHyperedge{p, H, w}(\Tilde{e})|_{\zeros} \in \tightSetFamily{p', H', w'}$. Next, suppose that $\Tilde{e} \in E' - E$. In particular, this implies that $E' - E \not= \emptyset$. 
    We recall that $e, f \in E$ are the hyperedges chosen by  \Cref{alg:WeakToStrongCover}  \Cref{alg:WeakToStrongCover:(4)}. 
    By Lemma \ref{lem:WeakToStrong:properties-of-alpha^M-and-alpha^R}(5) and \ref{lem:WeakToStrong:properties-of-alpha^M-and-alpha^R}(2), we have that $\Tilde{e} = e\cup f$ and $\alpha^R = \beta^R < w^M(e \cup f)$. Consequently, $\zeros = \emptyset$ and $V' = V$. Let $Z \subseteq V$ be a set such that 
    %witnessing $\alpha^R = \beta^R$, i.e. 
    $\alpha^R = b_{(H^M, w^M)}(Z) - p(Z)$ and $e\cup f \subseteq Z$ (such a set exists by definition of $\beta^R$).
    %\knote{Avoid using this phrasing: Let .. be a set witnessing .... It is non-standard. Instead, just say, let ... be a set such that ``conditions'' hold (such a set exists by definition of $\beta^R$}\shubhang{changed. leaving the comment here as a reminder to change in future proofs.}. 
    Then, we have the following: 
    $$b_{H',w'}(Z) = b_{(H^R, w^R)}(Z) = b_{(H^M, w^M)}(Z) - \alpha^R = p(Z) = p^R(Z) = p'(Z),$$
    and consequently the set $Z \in \tightSetFamily{p', H', w'}$.
    Here, the first equality is by definitions of the weight function $w^R$ and the set $\zeros$. The second and the fourth equalities are because $e \cup f \subseteq Z$. The third equality is because the set $Z$ is a witness set for $\alpha^R = \beta^R$. 
    The fifth equality is because $\zeros = \emptyset$.  %Thus, there exists a set $Y \in \maximalTightSetFamily{p',H',w'}$ such that $e \cup f \subseteq Z \subseteq Y$ and consequently $\maximalTightSetContainingHyperedge{p'}(e \cup f) \not = \emptyset$.

    \item We recall that $V' = V - \zeros = V - \{u\in V : b_{(H^R, w^R)}(u) = 0\}$. The claim then follows because $w' = w^R$. 
    %\knote{$w'=w^R$.}.
    
\end{enumerate}
\end{proof}

We now show the first main result of the section which says that \Cref{alg:WeakToStrongCover} terminates, and the hypergraph returned by the algorithm indeed strongly covers the input symmetric skew-supermodular function, i.e. satisfies property \ref{thm:WeakToStrongCover:main:(1)} of \Cref{thm:WeakToStrongCover:main}.

\begin{lemma}\label{lem:weak-to-strong:correctness:main:strong-cover}
    %Consider a recursive call of \Cref{alg:WeakToStrongCover}. 
    Suppose that the input to \Cref{alg:WeakToStrongCover} is a  hypergraph $(H = (V, E), w)$ and a symmetric skew-supermodular function $p:2^V\rightarrow \Z$ %satisfying conditions \ref{cond:WeakToStrong:(a)} and \ref{cond:WeakToStrong:(b)} 
such that every hyperedge $\Tilde{e} \in E$ is contained in some $(p, H, w)$-tight set, $b_{(H,w)}(u) \not = 0$ for every $u \in V$, and $b_{(H,w)}(X) \geq p(X)$ for every $X \subseteq V$.  Then, \Cref{alg:WeakToStrongCover} terminates in finite number of recursive calls. Furthermore, let $(H^*, w^*)$ be the hypergraph returned by the algorithm. Then, $d_{(H^*, w^*)}(X) \geq p(X)$ for every $X\subseteq V$.
    % Then, we have the following:
    % \begin{enumerate}
    %     \item $d_{(H^*, w^*)}(X) \geq p(X)$ for every $X\subseteq V$, and
    %     \item the hypergraph $(H^*, w^*)$ is formed by merging hyperedges of the hypergraph $(H,w)$.
    % \end{enumerate}
\end{lemma}
\begin{proof}
%\Cref{lem:weak-to-strong-cover:next-recursive-call-hypothesis-satification:symm-skew-supermod-strong-cover} and \Cref{lem:weak-to-strong-cover:next-recursive-call-hypothesis-satification:hyperedge-in-tight-set-positive-coverage} 
By Lemmas \ref{lem:weak-to-strong-cover:next-recursive-call-hypothesis-satification:symm-skew-supermod-strong-cover} and \ref{lem:weak-to-strong-cover:next-recursive-call-hypothesis-satification:hyperedge-in-tight-set-positive-coverage}, the input to every recursive call (except the last one) of the execution satisfies the hypothesis of \Cref{cor:weak-to-strong:uncrossing:exist-two-hyperedges-contained-in-disjoint-maximal-tight-sets}, and consequently, during every recursive call of the execution (except the last one), there exist two hyperedges such that \Cref{alg:WeakToStrongCover} \Cref{alg:WeakToStrongCover:(4)} is well-defined.
We first show that 
%the execution is finite, i.e. the 
algorithm terminates in a finite number of recursive calls. 
Consider the potential function $\phi(H=(V, E), w) \coloneqq  \sum_{e \in E}w(e)$. We note that $\phi(H, w) \in \Z_{\geq 0}$ since the weight function $w:E\rightarrow\Z_+$ is a positive integral function. Furthermore, if $\phi(H, w) = 0$, then we have that $E = \emptyset$ and the algorithm terminates. We recall that $p'$ and $(H', w')$ denote the inputs to the subsequent recursive call of \Cref{alg:WeakToStrongCover}. 
By Lemmas \ref{lem:weak-to-strong-cover:next-recursive-call-hypothesis-satification:symm-skew-supermod-strong-cover} and \ref{lem:weak-to-strong-cover:next-recursive-call-hypothesis-satification:hyperedge-in-tight-set-positive-coverage}, we have that $p'$ is a symmetric skew-supermodular function and $(H' = (V', E'), w')$ is a non-empty hypergraph such that every hyperedge $\Tilde{e} \in E'$ is contained in some $(p', H', w')$-tight set, $b_{(H', w')}(u) \not = 0$ for every $u \in V'$, and $b_{(H', w')}(X) \geq p'(X)$ for every $X \subseteq V'$.
Then, by Lemma \ref{lem:WeakToStrong:properties-of-alpha^M-and-alpha^R}(1), \ref{lem:WeakToStrong:properties-of-alpha^M-and-alpha^R}(2) and \ref{lem:WeakToStrong:properties-of-alpha^M-and-alpha^R}(5), we have that $\phi(H, w) > \phi(H', w')$, i.e. the value of the potential function strictly decreases between subsequent recursive calls. Consequently, the algorithm terminates in a finite number of recursive calls.


We now show that the hypergraph $(H^*,w^*)$ returned by \Cref{alg:WeakToStrongCover} strongly covers the function $p$, i.e. $d_{(H^*, w^*)}(X) \geq p(X)$ for every $X\subseteq V$. We show this by induction on $w(E)$. For the base case, consider $w(E) = 0$. Since the weight function $w$ is positive, we have that $E = \emptyset$.
Let $X \subseteq V$. We have that $d_{(H, w)}(X) = 0 = b_{(H,w)}(X) \geq p(X)$ and so the claim holds. For the inductive case, consider $w(E) > 0$. Let $X\subseteq V$ be an arbitrary subset. 
Let $(H_0^* = (V', E_0^*), w_0^*)$ be the hypergraph returned by the subsequent recursive call on input hypergraph $(H' = (V', E'), w')$ and function $p'$. 
%Let $(\Tilde{H} = (V', \Tilde{E}), \Tilde{w})$ be the hypergraph returned by the subsequent recursive call on input hypergraph $(H' = (V', E'), w')$ and function $p'$. 
We recall that $w'(E') < w(E)$ and so by the inductive hypothesis, we have that $d_{(H_0^*, w_0^*)}(X) \geq p'(X)$ for every $X\subseteq V'$. We also recall that the hypergraph $(H^*, w^*) = (G+H_0, c + w_0)$, where the hypergraph $(G, c)$ is obtained from $(H_0^*, w_0^*)$ by adding the vertices $\zeros$. Then, we have the following:
\begin{align*}
    d_{(H^*, w^*)}(X)& = d_{G,c}(X) + d_{(H_0, w_0)}(X)&\\
    & = d_{H_0^*, w_0^*}(X) + \alpha^R\indicator_{e\cup f\in \delta(X)}&\\
    &\geq p'(X) + \alpha^R\indicator_{e\cup f\in \delta(X)}& \text{(by inductive hypothesis)}\\
    & = \functionContract{\left(p - \alpha^R\cdot \indicator_{e\cup f\in \delta(X)}\right)}{\zeros}(X) + \alpha^R\indicator_{e\cup f\in \delta(X)} &\\
    &\geq \left(p - \alpha^R\cdot  \indicator_{e\cup f\in \delta(X)}\right)(X) + \alpha^R\indicator_{e\cup f\in \delta(X)}&\\
    &= p(X).&
\end{align*}
\end{proof}

We now show that the hypergraph returned by \Cref{alg:WeakToStrongCover} is obtained by merging hyperedges of the input hypergraph (see \Cref{def:WeakToStrongCover:hyperedge-merge} for the definition of merging hyperedges). 
This will imply that the hypergraph returned by the algorithm satisfies property \ref{thm:WeakToStrongCover:main:(2)} of \Cref{thm:WeakToStrongCover:main}. To show this property, we will need the following lemma about the choice of hyperedges $e$ and $f$ being merged by the algorithm.

\begin{lemma}\label{lem:WeakToStrongCover:properties-of-e-and-f-chosen-during-alg}
    %Consider a recursive call of \Cref{alg:WeakToStrongCover}. 
    Suppose that the input \Cref{alg:WeakToStrongCover} call is a non-empty hypergraph $(H = (V, E), w)$ and a symmetric skew-supermodular function $p:2^V\rightarrow \Z$ %satisfying conditions \ref{cond:WeakToStrong:(a)} and \ref{cond:WeakToStrong:(b)} 
such that every hyperedge $\Tilde{e} \in E$ is contained in some $(p, H, w)$-tight set, $b_{(H,w)}(u) \not = 0$ for every $u \in V$, and $b_{(H,w)}(X) \geq p(X)$ for every $X \subseteq V$.  Let $e, f$ be as defined by \Cref{alg:WeakToStrongCover}. Then,
\begin{enumerate}
    \item $e \cap f = \emptyset$, and
    \item $e\cup f \not \in E$.
\end{enumerate}
\end{lemma}
\begin{proof} We prove the claims separately below.
    \begin{enumerate}
        \item By the choice of hyperedges $e$ and $f$, we have that $\maximalTightSetContainingHyperedge{p, H, w}(e) \not = \maximalTightSetContainingHyperedge{p, H, w}(f)$. We observe that by maximality of the sets $\maximalTightSetContainingHyperedge{p, H, w}(e)$ and $\maximalTightSetContainingHyperedge{p, H, w}(f)$, we have that $\maximalTightSetContainingHyperedge{p, H, w}(e), \maximalTightSetContainingHyperedge{p, H, w}(f) \in \maximalTightSetFamily{p, H, w}$. By \Cref{lem:weak-to-strong:uncrossing:maximal-tight-set-family-disjoint}, we have that $\maximalTightSetContainingHyperedge{p, H, w}(e) \cap  \maximalTightSetContainingHyperedge{p, H, w}(f) = \emptyset$. Consequently, $e \cap f = \emptyset$.
        \item By way of contradiction, suppose that $e\cup f \in E$. Thus, $\maximalTightSetContainingHyperedge{p, H, w}(e\cup f) \not = \emptyset$. Since $e, f \subseteq \maximalTightSetContainingHyperedge{p, H, w}(e\cup f)$, we have that $\maximalTightSetContainingHyperedge{p, H, w}(e) = \maximalTightSetContainingHyperedge{p, H, w}(f) = \maximalTightSetContainingHyperedge{p, H, w}(e\cup f)$ by the maximality of the three sets, contradicting the choice of hyperedges $e, f$ made by the algorithm.
    \end{enumerate}
\end{proof}

We now conclude the section by showing that the hypergraph returned by \Cref{alg:WeakToStrongCover} indeed satisfies property \ref{thm:WeakToStrongCover:main:(2)} of \Cref{thm:WeakToStrongCover:main}.

\begin{lemma}\label{lem:weak-to-strong:correctness:main:hyperedge-merge}
    %Consider a recursive call of \Cref{alg:WeakToStrongCover}. 
    Suppose that the input to \Cref{alg:WeakToStrongCover} is a hypergraph $(H = (V, E), w)$ and a symmetric skew-supermodular function $p:2^V\rightarrow \Z$ %satisfying conditions \ref{cond:WeakToStrong:(a)} and \ref{cond:WeakToStrong:(b)} 
such that every hyperedge $\Tilde{e} \in E$ is contained in some $(p, H, w)$-tight set, $b_{(H,w)}(u) \not = 0$ for every $u \in V$, and $b_{(H,w)}(X) \geq p(X)$ for every $X \subseteq V$. 
    Then, the hypergraph $(H^* = (V, E^*), w^*)$ returned by the algorithm is obtained by merging hyperedges of the hypergraph $(H, w)$.
\end{lemma}
\begin{proof}
     We show the claim by induction on $w(E)$. For the base case, consider $w(E) = 0$. Since the weight function $w$ is positive, we have that $E = \emptyset$ and hence $(H, w)$ is the empty hypergraph. In this case, the hypergraph $(H^*, w^*)$ returned by the algorithm is also the empty hypergraph and so the claim holds. 

    For the inductive case, suppose that $w(E) > 0$. Let $(H^M, w^M), \alpha^M$, $(H' = (V', E'), w'), p',(H^R, w^R)$, $(H_0 = (V, E_0), w_0), \zeros, (H^*_0, w^*_0), (G, c)$ be as defined by \Cref{alg:WeakToStrongCover}. By Lemmas \ref{lem:weak-to-strong-cover:next-recursive-call-hypothesis-satification:symm-skew-supermod-strong-cover} and \ref{lem:weak-to-strong-cover:next-recursive-call-hypothesis-satification:hyperedge-in-tight-set-positive-coverage}, the inputs $(H',w')$ and $p'$ satisfy the hypothesis of \Cref{lem:weak-to-strong:correctness:main:strong-cover}. 
    Thus, the execution of the algorithm from this recursive call terminates in a finite number of recursive calls and returns the hypergraph $(H_0^* = (V', E_0^*), w_0^*)$. 
    %The following two claims show that the hypergraphs maintained and returned by the algorithm can be obtained by merging hyperedges of appropriate hypergraphs.
    Then, by Claims \ref{claim:WeakToStrong:MergeProperty:H^*w^*-from-H^Mw^M} and \ref{claim:WeakToStrong:MergeProperty:H^Mw^M-from-Hw} below, we have that the hypergraph $(H^*, w^*)$ is obtained by merging hyperedges of $(H, w)$.
    \end{proof}
    \iffalse
\begin{claim}\label{claim:WeakToStrong:Merge-Property:Gc-from-H^Rw^R}
        $(G, c)$ is formed by merging hyperedges of $(H^R, w^R)$.
    \end{claim}
    \begin{proof}
        By Lemma \ref{lem:WeakToStrong:properties-of-alpha^M-and-alpha^R}(1), \ref{lem:WeakToStrong:properties-of-alpha^M-and-alpha^R}(2) and \ref{lem:WeakToStrong:properties-of-alpha^M-and-alpha^R}(5), we have that $w'(E') < w(E)$. Thus, by the inductive hypothesis, we have that the hypergraph $(H_0^*, w_0^*)$ is formed by merging hyperedges of $(H', w')$. The claim then follows because of the following two observations: $(H', w')$ is obtained from $(H^R, w^R)$ by removing the vertices $\zeros$ and $(G, c)$ is formed by adding vertices $\zeros$ to $(H_0^*, w_0^*)$.
    \end{proof}
\fi
    \begin{claim}\label{claim:WeakToStrong:MergeProperty:H^*w^*-from-H^Mw^M}
        $(H^*, w^*)$ is obtained by merging hyperedges of $(H^M, w^M)$.
    \end{claim}
    \begin{proof}
        We first show that the hypergraph $(G, c)$ is obtained by merging hyperedges of $(H^R, w^R)$.
         By Lemma \ref{lem:WeakToStrong:properties-of-alpha^M-and-alpha^R}(1), (2), and (5), we have that $w'(E') < w(E)$. Thus, by the inductive hypothesis, we have that the hypergraph $(H_0^*, w_0^*)$ is obtained by merging hyperedges of $(H', w')$. Consequently, $(G, c)$ is formed by merging hyperedges of $(H^R, w^R)$ because of the following two observations: $(H', w')$ is obtained from $(H^R, w^R)$ by removing the vertices $\zeros$ and $(G, c)$ is obtained by adding vertices $\zeros$ to $(H_0^*, w_0^*)$.
        
        Next, we recall that $(H^*, w^*) = (G + H_0, c + w_0)$. Thus, by the conclusion of the previous paragraph, we have that $(H^*, w^*)$ is obtained by merging hyperedges of $(H^R + H_0, w^R + w_0)$. The claim then follows because $(H^R + H_0, w^R + w_0) = (H^M, w^M)$ by \Cref{lem:WeakToStrong:properties-of-alpha^M-and-alpha^R}(5).
    \end{proof}

    \begin{claim}\label{claim:WeakToStrong:MergeProperty:H^Mw^M-from-Hw}
        $(H^M,w^M)$ is obtained by merging hyperedges of $(H, w)$.
    \end{claim}
    \begin{proof}
        We recall that $(H^M, w^M) = \merge((H,w), e, f, \alpha^M)$ (see \Cref{alg:WeakToStrongCover} \Cref{alg:WeakToStrongCover:(6)}).
        Thus, by definition of the \merge operation, the hypergraph $(H^M, w^M)$ is obtained by decreasing the weight of hyperedges $e$ and $f$ by $\alpha^M$ and increasing the weight of the hyperedge $e\cup f$ by $\alpha^M$. Consequently, the multi-hypergraph $H^M_{w^M}$ can be obtained from the multi-hypergraph $H_w$ by repeatedly merging hyperedges $e$ and $f$ exactly $\alpha^M$ times. We note that the hyperedges $e$ and $f$ are disjoint by \Cref{lem:WeakToStrongCover:properties-of-e-and-f-chosen-during-alg}(1). Thus, the hypergraph $(H^M,w^M)$ is obtained by merging hyperedges of $(H, w)$ by definition. 
        %We recall that $H_0 = (V, E_0 = \{e\cup f\})$. Let $m_0 : E_0 \rightarrow \{\alpha^M\}$. Furthermore,  let $H_1 := (V, E_1:=\{e, f\})$ and $m_1: E_1 \rightarrow \{\alpha^M\}$. By \Cref{lem:WeakToStrong:properties-of-alpha^M-and-alpha^R}(2), we have that $(H^M, w^M) = (H + H_0 - H_1, w + m_0 - m_1)$ \shubhang{Is the ``hypergraph subtraction'' notation ok to use?} \knote{No, don't use subtraction notation. It raises confusion about whether the hyperedge weights are negative.}. We note that the hypergraph $(H_0, m_0)$ can be formed by merging hyperedges of $(H_1, m_1)$  since the hyperedges $e$ and $f$ are disjoint by \Cref{lem:WeakToStrongCover:properties-of-e-and-f-chosen-during-alg}(1).
    \end{proof}

    %By Claims \ref{claim:WeakToStrong:MergeProperty:H^*w^*-from-H^Mw^M} and \ref{claim:WeakToStrong:MergeProperty:H^*w^*-from-H^Mw^M} above, we have that the hypergraph $(H^*, w^*)$ is formed from the hypergraph $(H, w)$ by merging hyperedges.

    \iffalse
    For notational convenience, we use  $(G^* = (V', F^*), c^*)$ to denote the  hypergraph $ (H_0^* = (V', E_0^*), w_0^*)$. By Lemma \ref{lem:WeakToStrong:properties-of-alpha^M-and-alpha^R}(1), \ref{lem:WeakToStrong:properties-of-alpha^M-and-alpha^R}(2) and \ref{lem:WeakToStrong:properties-of-alpha^M-and-alpha^R}(5), we have that $w'(E') < w(E)$. Thus, by the inductive hypothesis, we have that the hypergraph $(G^*, c^*)$ is formed by merging hyperedges of the hypergraph $(H', w')$. In particular, the multi-hypergraph $G^*_{c^*}$ is obtained from the multi-hypergraph $H'_{w'}$ by repeatedly merging two disjoint hyperedges. Consider a sequence of $k \in \Z_{\geq 0}$ operations of merging disjoint hyperedges starting from the multi-hypergraph $H'_{w'}$ and resulting in the multi-hypergraph $G^*_{c^*}$. Let $\sigma = (e_1, f_1), (e_2, f_2), \ldots, (e_k, f_k)$ denote the collection of pairs of (disjoint) hyperedges merged in order in this sequence. We now use $\sigma$ to construct a new collection $\sigma'$ such that starting from the multi-hypergraph $H_w$ and sequentially merging the pairs of hyperedges in $\sigma'$   results in the multi-hypergraph $H^*_{w^*}$. We define the collection $\sigma'$ in two steps as follows.
    \begin{enumerate}
        \item Consider the sequence $\sigma$. For every $i \in [k]$ such that $e\cup f \in \{e_i, f_i\}$, add the pair $(e, f)$ before the pair $(e_i, f_i)$. Furthermore, let $q := |\{i \in [k] : e\cup f \in \{e_i, f_i\}\}|$; append $q$ pairs of $(e, f)$ at the end of the previous sequence. Denote the resulting sequence as $\sigma_1'$.
        \item Let $p := \alpha^M - \alpha^R$. Then, we define the sequence $\sigma_1'$ with an additional $p$ pairs of $(e, f)$ appended at the end.
    \end{enumerate}
    
    We recall that by Lemma \ref{lem:WeakToStrong:properties-of-alpha^M-and-alpha^R}(2) and \ref{lem:WeakToStrong:properties-of-alpha^M-and-alpha^R}(5), we have that $w'(e\cup f) = \alpha^M - \alpha^R = p$. Consequently, there are $p$ copies of the hyperedge $e\cup f$ in the multi-hypergraph $H'_{w'}$. 
\shubhang{todo: new proof below}
    \begin{claim}
        The multi-hypergraph $G_c$ can be obtained from the multi-hypergraph $H'_{w'}$ by merging hyperedges.
    \end{claim}
    \begin{proof}
        By the inductive hypothesis, the hypergraph $(H^*_0, w^*_0)$ can be obtained from the hypergraph $(H', w')$ by merging hyperedges. Since the hypergraph $(G, c)$ is formed from $(H_0^*, w_0^*)$ by adding vertices $\zeros$, it follows that the hypergraph $(G, c)$ can be obtained from the hypergraph $(H', w')$ by merging hyperedges.            
    \end{proof}

    Let $(\Tilde{H}_0 := (V, \Tilde{E}_0 :=\{e, f\})$ and $\Tilde{w}_0 : \Tilde{E_0} \rightarrow\{ \alpha^R\}\}$ 
    \begin{claim}
        The multi-hypergraph $H_{0_{w_0}}$ can be obtained from the multi-hypergraph $\Tilde{H}_{0_{\Tilde{w}_0}}$.
    \end{claim}
    \begin{proof}
        \shubhang{todo}
    \end{proof}

    \begin{claim}
        $H_w = G_c + \Tilde{H}_{0_{\Tilde{w}_0}}$.
    \end{claim}
    \begin{proof}
        \shubhang{todo}.
    \end{proof}

\shubhang{end new proof}
    
\shubhang{old proof below}

    
    $(G^* = (V', F^*), c^*)$. Furthermore, the inputs 
$(H', w')$ and $p'$ also satisfy the hypothesis of the claim, and $w'(E') < w(E)$ by \Cref{lem:WeakToStrong:properties-of-alpha^M-and-alpha^R}(2), \Cref{lem:WeakToStrong:properties-of-alpha^M-and-alpha^R}(5) and \Cref{lem:WeakToStrong:properties-of-alpha^M-and-alpha^R}(7). Thus, by the inductive hypothesis, we have that the hypergraph $(G^*, c^*)$ is formed by merging hyperedges of the hypergraph $(H', w')$. In particular, the multi-hypergraph $G^*_{c^*}$ is obtained from the multi-hypergraph $H'_{w'}$ by repeatedly merging two disjoint hyperedges. Consider a sequence of operations of merging disjoint hyperedges starting from the hypergraph $H'_{w'}$ and resulting in the hypergraph $G^*_{c^*}$. We now use this sequence to construct a new sequence of merge operations such that starting from the multi-hypergraph $H_w$ and performing the operations along the sequence results in the multi-hypergraph $H^*_{w^*}$. For this, we start with the multihypergraph $H_w$ and consider the aforementioned sequence guaranteed by the inductive hypothesis  --  we note that this may not a valid sequence of merge operations for the mutihypergraph $H_w$ as the hyperedge $e\cup f \not \in E$ by \Cref{lem:WeakToStrongCover:properties-of-e-and-f-chosen-during-alg}(2), but the hyperedge $e\cup f$ may be used in a merge operation in the sequence. By \Cref{lem:WeakToStrong:properties-of-alpha^M-and-alpha^R}(2) and \Cref{lem:WeakToStrong:properties-of-alpha^M-and-alpha^R}(5), we have that $w' = w^R = w - w_0^M + w_1^M - w_0$, where $w_0^M:\{e, f\}\rightarrow\{\alpha^M\}$, $w_1^M:\{e\cup f\} \rightarrow\{\alpha^M\}$, and $w_0:\{e\cup f\}\rightarrow\{\alpha^R\}$ as defined by \Cref{alg:WeakToStrongCover}. Consequently there are $\alpha^M - \alpha^R$ copies of the hyperedge $e\cup f$ in the multihypergraph $H'_{w'}$. Suppose that $k$ of these copies are used during merge operations along the sequence. We preface each of these $k$ operations with the operation of merging hyperedges $e$ and $f$. Furthermore, at the end of the sequence, we include $\alpha^M - \alpha^R - k$ new operations of merging hyperedges $e$ and $f$. By \Cref{lem:WeakToStrongCover:properties-of-e-and-f-chosen-during-alg}(1), the hyperedges $e$ and $f$ are disjoint. Thus, the resulting sequence of merge operations becomes a valid sequence of merge operations starting from the multihypergraph $H_w$. Furthermore, this new sequence of merge operations results in the multihypergraph $G^*_{c^*}$ along with an additional $\alpha^R$ copies of the hyperedges $e$ and $f$. By \Cref{alg:WeakToStrongCover} \Cref{alg:WeakToStrongCover:(11)}, the hypergraph $(H^*, w^*) = (G^* + H_0, c^* + w_0)$. Thus, we append the sequence with $\alpha^R$ new operations of merging the hyperedges $e$ and $f$ to get the multihypergraph $H^*_{w^*}$ as desired.
\fi
%\end{proof}

%% proof using induction on total weight and old definition of merging
\iffalse
\begin{proof}
    We show the claim by induction on $w(E)$. For the base case, consider $w(E) = 0$. Since the weight function $w$ is positive, we have that $E = \emptyset$ and hence $(H, w)$ is the empty hypergraph. In this case, the hypergraph $(H^*, w^*)$ returned by the algorithm is also the empty hypergraph and so the claim holds. 

    For the inductive case, suppose that $w(E) > 0$. Let the hypergraph $(H' = (V', E'), w')$ and the function $p':2^{V'}\rightarrow\Z$ be the inputs to the subsequent recursive call of the algorithm. By \Cref{lem:weak-to-strong-cover:next-recursive-call-hypothesis-satification:symm-skew-supermod-strong-cover} and \Cref{lem:weak-to-strong-cover:next-recursive-call-hypothesis-satification:hyperedge-in-tight-set-positive-coverage}, the inputs $(H',w')$ and $p'$ satisfy the hypothesis of \Cref{lem:weak-to-strong:correctness:main:strong-cover}. Thus, the the execution of the algorithm from this recursive call terminates in finite time and returns a hypergraph $(G^* = (V', F^*), c^*:F^*\rightarrow\Z_+)$. Furthermore, the inputs 
$(H', w')$ and $p'$ also satisfy the hypothesis of the claim, and $w'(E') < w(E)$ by \Cref{lem:WeakToStrong:properties-of-alpha^M-and-alpha^R} part (6). By the inductive hypothesis, we have that the hypergraph $(G^*, c^*)$ is formed by merging hyperedges of the hypergraph $(H', w')$. We now make some assumptions that will simplify the exposition of the proof. We henceforth assume that  $e, f \in E'$. We can make this assumption because if the hyperedge $e \not \in E'$, then we can add the hyperedge into $E'$ and set $w'(e) \coloneqq  0$. We similarly handle the case when the hyperedge $f \not \in E'$. We also assume that the hyperedge $e \cup f \in F^*$ and $e \cup f \in E$. As in the explanation of the previous assumption, if $e\cup f \not \in F^*$, then we add the hyperedge $e \cup f$ to $F^*$ and set $c^*(e\cup f) \coloneqq  0$. We note that $e\cup f \not \in E$ and so we add the hyperedge to $E$ with $w(e\cup f) \coloneqq  0$. With these assumption and by \Cref{lem:WeakToStrong:properties-of-alpha^M-and-alpha^R}, we have that $E = E'$. Furthermore, by the algorithm description, we have that  $E^* = F^*$. \knote{These assumptions are confusing to track inside the claims. For example, claim 42.1 is already confusing. Can the proof be written without making these assumptions? Perhaps, use indicator functions if you need to condition? Or, it is okay if you want to extend the domain of $\tilde{c}$ to $E^*\times E$ by defining its values to be zero on the new elements of the domain. }

    %Let $e^* \in E^*$ be an arbitrary hyperedge.
    By the inductive hypothesis, there exists a function $\Tilde{c}:F^*\times E' \rightarrow \Z_{\geq 0}$ satisfying property (1) of \Cref{def:WeakToStrongCover:hyperedge-merge}, and for every hyperedge $e^*\in E^*$, there exist positive integers $k' \in \Z_+$ and $\gamma_1', \ldots, \gamma_{k'}' \in \Z_+$, and a collection of hyperedge subsets $S_1', \ldots S_{k'}' \subseteq E'$  satisfying properties (2a)-(2c) of \Cref{def:WeakToStrongCover:hyperedge-merge}.

    We now extend these  quantities to simplify the exposition of the proof. First, suppose that the hyperedge $e \not \in E'$. Then, we extend the function $\Tilde{c}$ by defining $\Tilde{c}(e^*, e) \coloneqq  0$ for every $e^* \in F^* \cup \{e\cup f\}$. We perform a similar extension in the case where the hyperedge $f \not \in E'$. Next, suppose that the hyperedge $e \cup f \not \in F^*$. We extend the function $\Tilde{c}$ by defining the function value $\Tilde{c}(e\cup f, \Tilde{e}) \coloneqq  0$ for every $\Tilde{e} \in E' \cup \{e, f\}$. Here, we also define $k' \coloneqq  0$. We emphasize that these extensions are merely for notational convenience.
   
    We now define the weight function $\Tilde{w}:E^*\times E \rightarrow \Z_{\geq 0}$ as follows. Let $\alpha^R$ be the quantity defined by the recursive call of \Cref{alg:WeakToStrongCover}.
    %Let $e^* \in E_{i+1}^*$, and let $k', \gamma_1', \ldots, \gamma_{k'}' \in \Z_+$, and  $S_1', \ldots S_{k'}' \subseteq E_{i+1}$ be  values for $e^*$ satisfying (\ref{tag:hyperedge-merge-properties}). 
    For every $e^*\in E^*$ and $\Tilde{e} \in E$, we let



$$\Tilde{w}(e^*, \Tilde{e}) \coloneqq  \begin{cases}
        \Tilde{c}(e^*, \Tilde{e})&\text{ if } \Tilde{e} \not\in \{e, f\}\\
        \Tilde{c}(e^*, \Tilde{e}) + \Tilde{c}(e^*, e \cup f)&\text{ if } \Tilde{e} \in \{e, f\} \text{ and } e^* \not = e\cup f\\
        \Tilde{c}(e^*, \Tilde{e}) + \alpha^R& \text{ otherwise (i.e., if $\Tilde{e} \in \{e, f\}$ and $e^* = e \cup f$).}
    \end{cases}$$

    \Cref{claim:weak-to-strong:correctness:main:hyperedge-merge:property1} below says that the weight function $\Tilde{w}$ satisfies property (1) of \Cref{def:WeakToStrongCover:hyperedge-merge}. For the hyperedge $e^*$, we now define positive integers $k \in \Z_+$ and $\gamma_1, \ldots, \gamma_{k} \in \Z_+$, and the collections of hyperedge subsets $S_1, \ldots S_{k} \subseteq E$ as follows:
    \begin{align*}
        k &: =\begin{cases}
            k' + 1& \text{ if } e^* = e\cup f\\
            k'&\text{ otherwise.}
        \end{cases}&\\
        \gamma_j&\coloneqq  \begin{cases}
            \alpha^R& \text{ if } e^* = e\cup f \text{ and } j = k\\
            \gamma_j'&\text{ otherwise}.
        \end{cases}&\\
        S_j& \coloneqq  \begin{cases}
            \{e, f\}&\text{ if } e^* = e\cup f \text{ and } j = k\\
            S_j' - \{e\cup f\} + \{e, f\}&\text{ if } e\cup f \in S_j'\\
            S_j'& \text{ otherwise.}
        \end{cases}&
    \end{align*}
    \knote{Should the conditions for $\gamma_j$ and $S_j$ be for $j=k+1$?} 
    \Cref{claim:weak-to-strong:correctness:main:hyperedge-merge:property2}, \Cref{claim:weak-to-strong:correctness:main:hyperedge-merge:property3} and \Cref{claim:weak-to-strong:correctness:main:hyperedge-merge:property4} below show that the function $\Tilde{w}$, and the quantities $k$, $\gamma_1, \ldots, \gamma_k$ and $S_1, \ldots, S_k$ satisfy properties (2a)-(2c) of \Cref{def:WeakToStrongCover:hyperedge-merge} for the hyperedge $e^*$ and thus complete the proof of the claim. \knote{Phrase the claims to resemble Definition 4. }
\end{proof}

    \begin{claim}\label{claim:weak-to-strong:correctness:main:hyperedge-merge:property1}
        For every $\Tilde{e} \in E$, we have that $\sum_{e^* \in E^*}\Tilde{w}(e^*, \Tilde{e}) = w(\Tilde{e})$.
    \end{claim}
    \begin{proof}
        We note that $e \cup f \not \in E$ \knote{$w(e\cup f)=0$?} by \Cref{lem:WeakToStrongCover:properties-of-e-and-f-chosen-during-alg}. Let $\Tilde{e} \in E$. First, suppose that $\Tilde{e} \not \in \{e, f\}$. Then, we have the following:
        $$\sum_{e^* \in E^*}\Tilde{w}(e^*, \Tilde{e}) = \sum_{e^* \in F^*}\Tilde{c}(e^*, \Tilde{e}) = w'(\Tilde{e}) = w(\Tilde{e}).$$
        Here, the first and final equalities are because $\Tilde{e} \not \in \{e, f, e\cup f\}$ and the second equality is by the inductive hypothesis. Next, suppose that $\Tilde{e} \in \{e, f\}$. Then, we have the following:
        \begin{align*}
            \sum_{e^* \in E^*}\Tilde{w}(e^*, \Tilde{e})&= \sum_{e^* \in F^*}\Tilde{c}(e^*, \Tilde{e}) + \sum_{e^* \in F^*}\Tilde{c}(e^*, e\cup f) + \alpha^R&\\
            &= w'(\Tilde{e}) + w'(e\cup f) + \alpha^R&\\
            & = w(\Tilde{e}) - \alpha^M + (\alpha^M - \alpha^R) + \alpha^R&\\
            &= w(\Tilde{e}).  &
        \end{align*}
        Here, the second equality is by the inductive hypothesis and the third equality is by \Cref{lem:WeakToStrong:properties-of-alpha^M-and-alpha^R}.
    \end{proof}

\begin{claim}\label{claim:weak-to-strong:correctness:main:hyperedge-merge:property2}
    For every $j \in [k]$, we have that $e^* = \biguplus_{\Tilde{e} \in S_j}\Tilde{e}$.
\end{claim}
\begin{proof}
    We consider three cases based on the definition of the collection $S_j$. First, suppose that $e^* = e \cup f$ and $j = k$. Then, we have that $S_j = \{e, f\}$. We recall that the hyperedges $e$ and $f$ are disjoint by \Cref{lem:WeakToStrongCover:properties-of-e-and-f-chosen-during-alg} and so the claim holds. Next, suppose that $e \cup f \in S_j'$. Then, we have that $\uplus_{\Tilde{e} \in S_j} \Tilde{e} = \uplus_{\Tilde{e} \in S_j'} \Tilde{e} = e^*$, where the second equality is  because the hyperedges $e$ and $f$ are disjoint by \Cref{lem:WeakToStrongCover:properties-of-e-and-f-chosen-during-alg}, and the final equality is by the inductive hypothesis. Finally, suppose that $e\cup f \not \in S_j'$. Then, we have that $S_j = S_j'$ and the claim follows by the inductive hypothesis.
\end{proof}

\begin{claim}\label{claim:weak-to-strong:correctness:main:hyperedge-merge:property3}
     $\sum_{j \in [k]}\gamma_j\cdot \chi_{S_j} = \Tilde{w}_i(e^*, \cdot)$.
\end{claim}
\begin{proof}
    We consider two cases. First, suppose that $e^* \not = e\cup f$. Then, we have the following:
    $$\sum_{j \in [k]}\gamma_j\cdot \chi_{S_j} = \sum_{j \in [k']}\gamma_j'\cdot \chi_{S_j} = \Tilde{c}(e^*, \cdot) = \Tilde{w}(e^*, \cdot).$$
    Here, the first equality is because $\chi_{S_j} = \chi_{S_j'}$ since the hyperedges $e$ and $f$ are disjoint by \Cref{lem:WeakToStrongCover:properties-of-e-and-f-chosen-during-alg}. The second equality is by the inductive hypothesis. The final equality is by definition of the function $\Tilde{w}$. Next, suppose that the hyperedge $e^* = e\cup f$. We have the following:
    \begin{align*}
        \sum_{j \in [k]}\gamma_j\cdot \chi_{S_j}& = \sum_{j \in [k']}\gamma_j'\cdot \chi_{S_j'} + \alpha^R\chi_{e\cup f} &\\
        &= \sum_{j \in [k']}\gamma_j'\cdot \chi_{S_j'} + \alpha^R\chi_{S_k}&\\
        &= \Tilde{c}(e\cup f, \cdot) + \alpha^R\chi_{S_k}&\\
        &= \Tilde{w}(e\cup f, \cdot).&
    \end{align*}
    Here, the second equality is by definition of the set $S_i$. The third equality is by the inductive hypothesis. The final equality is by definition of the function $\Tilde{w}$.
\end{proof}

\begin{claim}\label{claim:weak-to-strong:correctness:main:hyperedge-merge:property4}
    $\sum_{j \in [k]}\gamma_j = w^*(e^*)$.
\end{claim}
\begin{proof}
    We consider two cases. First, suppose that $e^* \not = e\cup f$. Then, we have the following:
    $$\sum_{j \in [k]}\gamma_j = \sum_{j \in [k']}\gamma_j = c^*(e^*) = w^*(e^*).$$
    Here, the second equality is by the inductive hypothesis and the final equality is by \Cref{lem:WeakToStrong:properties-of-alpha^M-and-alpha^R}. Next, suppose that the hyperedge $e^* = e \cup f$. Then, we have the following:
    $$\sum_{j \in [k]}\gamma_j = \sum_{j \in [k']}\gamma_j' = c^*(e\cup f) + \alpha^R= w^*(e\cup f),$$
    where the second equality is by the inductive hypothesis.
\end{proof}
\fi

%% proof using recursion depth
\iffalse
\begin{proof}
    By \Cref{lem:weak-to-strong:correctness:main:strong-cover} the recursion depth of the algorithm is finite. Let $\ell \in \Z_+$ denote the recursion depth of the algorithm and let $i \in [\ell]$ be a recursive call. Let $(H_i = (V_i, E_i), w_i)$ and $p_i:2^{V_i}\rightarrow\Z$ denote the input hypergraph and input function to the $i^{th}$ recursive call. By \Cref{lem:weak-to-strong-cover:next-recursive-call-hypothesis-satification:symm-skew-supermod-strong-cover} and \Cref{lem:weak-to-strong-cover:next-recursive-call-hypothesis-satification:hyperedge-in-tight-set-positive-coverage} and induction on $i$, we have that the function $p_i$ is symmetric skew-supermodular, $b_{H_i, w_i}(X) \geq p_i(X)$ for every $X \subseteq V_i$, $\tau_{p_i}(e)\neq \emptyset$ for every hyperedge $e\in E_i$, and $b_{H_i, w_i}(u) > 0$ for every $u \in V_i$. Let $e_i, f_i\in E_i$ be the hyperedges chosen by the algorithm in the $i$th recursive call. Let $(H_i^*=(V, E_i^*), w_i^*)$ denote the hypergraph returned by the $i^{th}$ recursive call. We now show the claim by (reverse) induction on $i \in [\ell]$. For the base case, consider $i = \ell$. Here, we have that $(H_\ell^*, w_\ell^*) = (H_\ell, w_\ell) = (\emptyset, \emptyset)$ and so the claim holds. \knote{It would be better to do induction on  $\sum_{e\in E}w(e)$ instead of $i$. Such a proof will resemble the rest of the proofs in this section and would be easier to follow than figuring out this new index.} 
    
    For the inductive case, suppose that $i < \ell$.
    %By the inductive hypothesis, there exists a function $\Tilde{w}_{i+1}:E^*_{i+1}\times E_{i+1} \rightarrow \Z_{\geq 0}$, and for every $e^* \in E_{i+1}^*$ there exist positive integers $k_{i+1} \in \Z_+$ and $\gamma_1, \ldots, \gamma_{k_{i+1}} \in \Z_+$, and a collection of hyperedge subsets $S_1, \ldots S_{k_{i+1}} \subseteq E_{i+1}$ satisfying properties (1)-(4) of (\ref{tag:hyperedge-merge-properties}).
    For simplicity, we henceforth assume that  $e_i, f_i \in E_{i+1}$. We can make the previous assumption because if the hyperedge $e_i \not \in E_{i+1}$, then we can add the hyperedge into $E_{i+1}$ and set $w_{i+1}(e_i) \coloneqq  0$ \knote{Shouldn't the weight function be strictly positive?}. We similarly handle the case when the hyperedge $f_i \not \in E_{i+1}$. We also assume that the hyperedge $e_{i} \cup f_i \in E_{i+1}^*$ and $e_{i} \cup f_i \in E_{i}$. As in the explanation of the previous assumption, if $e_i\cup f_i \not \in E_{i+1^*}$, then we add the hyperedge $e_i \cup f_i$ to $E^*_{i+1}$ and set $w^*_{i+1}(e_i\cup f_i) \coloneqq  0$. We note that $e_{i}\cup f_i \not \in E_i$ and so we add the hyperedge to $E_i$ with $w_i(e_i\cup f_i) \coloneqq  0$. With these assumption and by \Cref{lem:WeakToStrong:properties-of-alpha^M-and-alpha^R}, we have that $E_i = E_{i+1}$. Furthermore, by the algorithm description, we have that  $E_{i}^* = E_{i+1}^*$.

    Let $e^* \in E_i^*$ be an arbitrary hyperedge.
    By the inductive hypothesis, there exists a function $\Tilde{w}_{i+1}:E^*_{i+1}\times E_{i+1} \rightarrow \Z_{\geq 0}$, positive integers $k' \in \Z_+$ and $\gamma_1', \ldots, \gamma_{k'}' \in \Z_+$, and a collection of hyperedge subsets $S_1', \ldots S_{k'}' \subseteq E_{i+1}$ satisfying properties (1)-(4) of (\ref{tag:hyperedge-merge-properties}).
    We now define the weight function $\Tilde{w}_i:E^*_{i}\times E_{i} \rightarrow \Z_{\geq 0}$ as follows. 
    %Let $e^* \in E_{i+1}^*$, and let $k', \gamma_1', \ldots, \gamma_{k'}' \in \Z_+$, and  $S_1', \ldots S_{k'}' \subseteq E_{i+1}$ be  values for $e^*$ satisfying (\ref{tag:hyperedge-merge-properties}). 
    For every $e \in E_i$ we let



$$\Tilde{w}_i(e^*, e) \coloneqq  \begin{cases}
        \Tilde{w}_{i+1}(e^*, e)&\text{ if } e \not\in \{e_i, f_i\}\\
        \Tilde{w}_{i+1}(e^*, e) + \Tilde{w}_{i+1}(e^*, e_i \cup f_i)&\text{ if } e \in \{e_i, f_i\} \text{ and } e^* \not = e_i\cup f_i\\
        \Tilde{w}_{i+1}(e^*, e) + \alpha_i^R& \text{ otherwise (i.e., if $e \in \{e_i, f_i\}$ and $e^* = e_i \cup f_i$).}
    \end{cases}$$

    \Cref{claim:weak-to-strong:correctness:main:hyperedge-merge:property1} below says that the weight function $\Tilde{w}_i$ satisfies property (1) of \Cref{def:WeakToStrongCover:hyperedge-merge}. For the hyperedge $e^*$, we now define positive integers $k \in \Z_+$ and $\gamma_1, \ldots, \gamma_{k} \in \Z_+$, and the collections of hyperedge subsets $S_1, \ldots S_{k} \subseteq E_{i}$ as follows:
    \begin{align*}
        k &: =\begin{cases}
            k' + 1& \text{ if } e^* = e_i\cup f_i\\
            k'&\text{ otherwise.}
        \end{cases}&\\
        \gamma_j&\coloneqq  \begin{cases}
            \alpha_i^R& \text{ if } e^* = e_i\cup f_i \text{ and } j = k \text{ for the hyperedge $e_i\cup f_i$}\\
            \gamma_j'&\text{ otherwise}.
        \end{cases}&\\
        S_j& \coloneqq  \begin{cases}
            \{e_i, f_i\}&\text{ if } e^* = e_i\cup f_i \text{ and } j = k \text{ for the hyperedge $e_i\cup f_i$}\\
            S_j' - \{e_i\cup f_i\} + \{e_i, f_i\}&\text{ if } e_i\cup f_i \in S_j'\\
            S_j'& \text{ otherwise.}
        \end{cases}&
    \end{align*}
    \Cref{claim:weak-to-strong:correctness:main:hyperedge-merge:property2}, \Cref{claim:weak-to-strong:correctness:main:hyperedge-merge:property3} and \Cref{claim:weak-to-strong:correctness:main:hyperedge-merge:property4} below show that the function $\Tilde{w}_i$, and the quantities $k$, $\gamma_1, \ldots, \gamma_k$ and $S_1, \ldots, S_k$ satisfy properties (2a)-(2c) of \Cref{def:WeakToStrongCover:hyperedge-merge} for the hyperedge $e^*$ and thus complete the proof of the claim.
\end{proof}

    \begin{claim}\label{claim:weak-to-strong:correctness:main:hyperedge-merge:property1}
        For every $e \in E_i$, we have that $\sum_{e^* \in E_i^*}\Tilde{w}_{i}(e^*, e) = w_i(e)$.
    \end{claim}
    \begin{proof}
        We note that $e_i \cup f_i \not \in E_i$ by \Cref{lem:WeakToStrongCover:properties-of-e-and-f-chosen-during-alg}. Let $e \in E_i$. First, suppose that $e \not \in \{e_i, f_i\}$. Then, we have the following:
        $$\sum_{e^* \in E_i^*}\Tilde{w}_{i}(e^*, e) = \sum_{e^* \in E_{i+1}^*}\Tilde{w}_{i+1}(e^*, e) = w_{i+1}(e) = w_i(e).$$
        Here, the first and final equalities are because $e \not \in \{e_i, f_i, e_i\cup f_i\}$ and he second equality is by the inductive hypothesis. Next, suppose that $e \in \{e_i, f_i\}$. Then, we have the following:
        \begin{align*}
            \sum_{e^* \in E_i^*}\Tilde{w}_i(e^*, e)&= \sum_{e^* \in E_{i+1}^*}\Tilde{w}_{i+1}(e^*, e) + \sum_{e^* \in E_{i+1}^*}\Tilde{w}_{i+1}(e^*, e_i\cup f_i) + \alpha_i^R&\\
            &= w_{i+1}(e) + w_{i+1}(e_i\cup f_i) + \alpha_i^R&\\
            & = w_i(e) - \alpha_i^M + (\alpha_i^M - \alpha_i^R) + \alpha_i^R&\\
            &= w_i(e).  &
        \end{align*}
        Here, the second equality is by the inductive hypothesis and the third equality is by \Cref{lem:WeakToStrong:properties-of-alpha^M-and-alpha^R}.
    \end{proof}

\begin{claim}\label{claim:weak-to-strong:correctness:main:hyperedge-merge:property2}
    For every $j \in [k]$, we have that $e^* = \biguplus_{e \in S_j}e$.
\end{claim}
\begin{proof}
    We consider three cases based on the definition of the collection $S_i$. First, suppose that $e^* = e_i \cup f_i$ and $j = k$ (for the hyperedge $e_i\cup f_i$). Then, we have that $S_j = \{e_i, f_i\}$. We recall that $e_i$ and $f_i$ are disjoint hyperedges by \Cref{lem:WeakToStrongCover:properties-of-e-and-f-chosen-during-alg} and so the claim holds. Next, suppose that $e_i \cup f_i \in S_j'$. Then, we have that $\uplus_{e \in S_j} e = \uplus_{e \in S_j'} e = e^*$, where the second equality is  because $e_i$ and $f_i$ are disjoint by \Cref{lem:WeakToStrongCover:properties-of-e-and-f-chosen-during-alg} , and the final equality is by the inductive hypothesis. Finally, suppose that $e_i\cup f_i \not \in S_j'$. Then, we have that $S_j = S_j'$ and the claim follows by the inductive hypothesis.
\end{proof}

\begin{claim}\label{claim:weak-to-strong:correctness:main:hyperedge-merge:property3}
     $\sum_{j \in [k]}\gamma_j\cdot \chi_{S_j} = \Tilde{w}_i(e^*, \cdot)$.
\end{claim}
\begin{proof}
    We consider two cases. First, suppose that $e^* \not = e_i\cup f_i$. Then, we have the following:
    $$\sum_{j \in [k]}\gamma_j\cdot \chi_{S_j} = \sum_{j \in [k']}\gamma_j'\cdot \chi_{S_j} = \Tilde{w}_{i+1}(e^*, \cdot) = \Tilde{w}_i(e^*, \cdot).$$
    Here, the first equality is because $\chi_{S_j} = \chi_{S_j'}$ since the hyperedges $e_i$ and $f_i$ are disjoint by \Cref{lem:WeakToStrongCover:properties-of-e-and-f-chosen-during-alg}. The second equality is by the inductive hypothesis. The final equality is by definition of the function $\Tilde{w}_i$. Next, suppose that $e^* = e_i\cup f_i$. We have the following:
    \begin{align*}
        \sum_{j \in [k]}\gamma_j\cdot \chi_{S_j}& = \sum_{j \in [k']}\gamma_j'\cdot \chi_{S_j'} + \alpha_i^R\chi_{e_i\cup f_i} &\\
        &= \sum_{j \in [k']}\gamma_j'\cdot \chi_{S_j'} + \alpha_i^R\chi_{S_k}&\\
        &= \Tilde{w}_{i+1}(e_i\cup f_i, \cdot) + \alpha_i^R\chi_{S_k}&\\
        &= \Tilde{w}_i(e_i\cup f_i, \cdot).&
    \end{align*}
    Here, the second equality is by definition of the set $S_i$. The third equality is by the inductive hypothesis. The final equality is by definition of the function $\Tilde{w}_i$.
\end{proof}

\begin{claim}\label{claim:weak-to-strong:correctness:main:hyperedge-merge:property4}
    $\sum_{j \in [k]}\gamma_j = w_i^*(e^*)$.
\end{claim}
\begin{proof}
    We consider two cases. First, suppose that $e^* \not = e_i\cup f_i$. Then, we have the following:
    $$\sum_{j \in [k]}\gamma_j = \sum_{j \in [k']}\gamma_j = w^*_{i+1}(e^*) = w_i^*(e^*).$$
    Here, the second equality is by the inductive hypothesis and the final equality is by \Cref{lem:WeakToStrong:properties-of-alpha^M-and-alpha^R}. Next, suppose that $e^* = e_i \cup f_i$. Then, we have the following:
    $$\sum_{j \in [k]}\gamma_j = \sum_{j \in [k']}\gamma_j' = w^*_{i+1}(e_i\cup f_i) + \alpha_i^R= w_i^*(e_i\cup f_i),$$
    where the second equality is by the inductive hypothesis.
\end{proof}
\fi

\subsection{Number of Recursive Calls and Hypergraph Support Size}\label{sec:WeakToStrongCover:NumberOfHyperedges}

We first set up some notation that will be used in the remainder of this section. 
By \Cref{lem:weak-to-strong:correctness:main:strong-cover}, the number of recursive calls made by \Cref{alg:WeakToStrongCover} is finite.
We will use $\ell$ to denote the depth of recursion of the algorithm. We will refer to the recursive call at depth $i$ as \emph{recursive call $i$} or the \emph{$i^{th}$ recursive call}. Let $i \in [\ell]$ be a recursive call. We let the hypergraph $(H_i = (V_i, E_i), w_i:E_i \rightarrow \Z_+)$ and the function $p_i : 2^{V_i} \rightarrow \Z$ denote the inputs to the $i^{th}$ recursive call. Furthermore, we let $e_i, f_i, \beta_i^M, \alpha_i^M, (H_i^M, w_i^M), \beta_i^R, \alpha_i^R, (H_i^R, w_i^R), p_i^R, (H_0^{i}, w_0^{i}), E_0^i, \zeros_i, V_i', p'_i$ denote the quantities $e, f, \beta^M, \alpha^M, (H^M, w^M), \beta^R, \alpha^R, (H^R, w^R), p^R, (H_0, w_0), E_0, \zeros, V', p'$ as defined by \Cref{alg:WeakToStrongCover} during the $i^{th}$ recursive call.
%Furthermore, we let $e_i, f_i \in E_i$ denote the hyperedges chosen by the algorithm, $\alpha_i^M, \alpha_i^R$ denote the values $\alpha^M, \alpha^R$, the set $\zeros_i$ denotes the set $\zeros$, $p_i^R \coloneqq  p_i - \alpha_i^R d_{(V, \{e_i\cup f_i\})}$, and $w_i^M, w_i^R$ denote the weight functions $w^M, w^R$ computed by the algorithm at the $i^{th}$ recursive call. 

%We recall that  a set $X\subseteq V$ is said to be \emph{tight} w.r.t. a function $p:2^V\rightarrow\Z$ and a hypergraph $(H = (V, E), w)$ if $p(X) = b_{(H,w)}(X)$. We also recall that $\tightSetFamily{p,H,w}$ denotes the family of tight sets and $\maximalTightSetFamily{p, H, w}$ denotes the family of \emph{inclusion-wise maximal} tight sets w.r.t. the function $p$ and hypergraph $(H,w)$. 
We now define certain set families of that will be crucial in our analysis. We let $\tightSetFamily{i}, \tightSetFamilyAfterMerge{i}, \tightSetFamilyAfterReduce{i}$ denote the families $\tightSetFamily{p_i, H_i, w_i}, \tightSetFamily{p_i, H_i^M, w_i^M}$ and $ \tightSetFamily{p_i^R, H_i^R, w_i^R}$ respectively, and $\maximalTightSetFamily{i}, \maximalTightSetFamilyAfterMerge{i}, \maximalTightSetFamilyAfterReduce{i}$ denote the families $\maximalTightSetFamily{p_i, H_i, w_i}, \maximalTightSetFamily{p_i, H_i^M, w_i^M}$ and $ \maximalTightSetFamily{p_i^R, H_i^R, w_i^R}$ respectively. We define the family of \emph{cumulative projected maximal $(p_i, H_i, w_i)$-tight sets} as:
$$\cumulativeMaximalTightSetFamily{i} \coloneqq  \begin{cases}
    \maximalTightSetFamily{1}& \text{ if $i = 1$,}\\
    
   \cumulativeMaximalTightSetFamily{i-1}|_{V_i} \cup \maximalTightSetFamily{i}& \text{ if $i \geq 2$}.
\end{cases}$$
\iffalse
$$\cumulativeMaximalTightSetFamily{i} \coloneqq  \begin{cases}
    \maximalTightSetFamily{1}& \text{ if $i = 1$,}\\
    
   \cumulativeMaximalTightSetFamily{i-1}|_{V_i} \cup \maximalTightSetFamily{i}& \text{ if $i \geq 2$}.
\end{cases}$$
\fi
The next lemma states two important properties of these set families.
\begin{lemma}\label{lem:WeakToStrong:properties-of-cumulative-tight-set-families}
Suppose that the input to \Cref{alg:WeakToStrongCover} is a non-empty hypergraph $(H = (V, E), w)$ and a symmetric skew-supermodular function $p:2^V\rightarrow \Z$ 
such that every hyperedge $\Tilde{e} \in E$ is contained in some $(p, H, w)$-tight set, $b_{(H,w)}(u) \not = 0$ for every $u \in V$, and $b_{(H,w)}(X) \geq p(X)$ for every $X \subseteq V$. Furthermore, suppose that the execution of \Cref{alg:WeakToStrongCover} on the input instance terminates using $\ell \in \Z_+$ recursive calls. Then, for every $i \in [\ell - 1]$, we have the following:
\iffalse
    Consider an execution of \Cref{alg:WeakToStrongCover} with $\ell \in \Z_+$ recursive calls. Suppose that the input to the $i^{th}$ recursive call is a hypergraph $(H_i = (V_i, E_i), w_i)$ and a function $p_i:2^{V_i}\rightarrow \Z$.
    Furthermore, suppose that  $(H_1, w_1)$ is a non-empty hypergraph and $p_1$ is a symmetric skew-supermodular function %satisfying conditions \ref{cond:WeakToStrong:(a)} and \ref{cond:WeakToStrong:(b)} 
such that every hyperedge $\Tilde{e} \in E_1$ is contained in some $(p_1, H_1, w_1)$-tight set, $b_{H_1, w_1}(u) \not = 0$ for every $u \in V_1$, and $b_{H_1,w_1}(X) \geq p_1(X)$ for every $X \subseteq V_1$.  Then, for every $i \in [\ell - 1]$, we have the following:
\fi
     \begin{enumerate}
        % \item $\tightSetFamily{i} \subseteq \tightSetFamilyAfterMerge{i} \subseteq \tightSetFamilyAfterReduce{i}$
        % \item $\projTightSetFamilyAfterReduce{i}{i+1} \subseteq \tightSetFamily{i+1}$
         \item $\projCumulativeMaximalTightSetFamily{i}{i+1} \subseteq \tightSetFamily{i+1}$, and
         \item the family $\cumulativeMaximalTightSetFamily{i}$ is laminar.
     \end{enumerate}
\end{lemma}
\begin{proof} 
Let $i\in [\ell - 1]$ be a recursive call. We recall that $p_i:2^{V_i}\rightarrow\Z$ and $(H_i, w_i)$ are inputs to the $i^{th}$ recursive call, where $p_1 := p$ and $(H_1, w_1) := (H, w)$.
By Lemmas \ref{lem:weak-to-strong-cover:next-recursive-call-hypothesis-satification:symm-skew-supermod-strong-cover} and \ref{lem:weak-to-strong-cover:next-recursive-call-hypothesis-satification:hyperedge-in-tight-set-positive-coverage} and induction on $i$, we have that $p_i$ is a symmetric skew-supermodular function and $(H_i = (V_i, E_i), w_i)$ is a non-empty hypergraph such that every hyperedge $\Tilde{e} \in E_i$ is contained in some $(p_i, H_i, w_i)$-tight set, $b_{(H_i, w_i)}(u) \not = 0$ for every $u \in V_i$, and $b_{(H_i, w_i)}(X) \geq p_i(X)$ for every $X \subseteq V_i$.
% We recall that the function $p_{i} = \functionContract{(p_{i-1} - b_{(H_0^{i-1}, w_0^{i-1})})}{\zeros_{i-1}}$ for every $i \geq 2$. Since $p_1$ is a symmetric skew-supermodular function, by \Cref{lem:weak-to-strong-cover:next-recursive-call-hypothesis-satification:symm-skew-supermod-strong-cover} and induction on $i$, we have that $p_{i}:2^{V_{i}}\rightarrow\Z$ is also a symmetric skew-supermodular function for every $i\ge 2$. 
% By \Cref{lem:WeakToStrongCover:TightSetFamiliesDuringAlgorithm} and induction on $i$, we have that $T_i\subseteq T_i^M \subseteq T_i^R$ and $T_i^R|_{V_{i+1}}\subseteq T_{i+1}$ for every $i\ge 2$. 
We now prove both claims separately below.
    \begin{enumerate}
        % \item We will first show that $\tightSetFamily{i}\subseteq \tightSetFamilyAfterReduce{i}$. Let $X \in \tightSetFamily{i}$. Here, we observe that $e_i\cap X, f_i\cap X \not = \emptyset$ as otherwise we have that 
        % $b_i^M(X) = b_i(X) - \alpha_i^M = p_i(X) - \alpha_i^M < p_i(X)$, where the first equality is because of the contradiction assumption and the final inequality is by \Cref{lem:WeakToStrong:properties-of-alpha^M-and-alpha^R}. This gives a contradiction. Thus, we have that $p_{i}(X) = b_{i}(X) = b_{i}^M(X)$, i.e. $X \in \tightSetFamilyAfterMerge{i}$.

        % We will now show the second inclusion in the claim, i.e. $\tightSetFamilyAfterMerge{i} \subseteq\tightSetFamilyAfterReduce{i}$. Let $X \in \tightSetFamilyAfterMerge{i}$. First, suppose that $\alpha_i^R = 0$. Then, we have that $b_i^R(X) = b_i^M(X) = p_i(X) = p_i^R(X)$, where the first and final equalities are because $\alpha_i^R = 0$. Thus, $X \in \tightSetFamilyAfterReduce{i}$. Next, suppose that $\alpha_{i}^R > 0$. Here, we consider the case where $e_i\cup f_i \not \in \delta_H(X)$ as otherwise, we have that $b_i^R(X) = b_i^M(X) - \alpha_i^R = p_i(X) - \alpha_i^R = p_i^R(X)$, and the claim holds. Then, we have the following: $b_{i}^R(X) = b_i^M(X) - \alpha_i^R = p_i^M(X) - \alpha_i^R < p_i^M(X) = p_i^R(X)$ which is a contradiction. Here, the final equality is because $d_{(V, \{e_i\cup f_i\})}(X) = 0$ since $e_i\cup f_i \not \in \delta_H(X)$.

        % \item Let $X \in \tightSetFamilyAfterReduce{i}$.  Then, we have the following:
        % $$p_{i+1}(X - \zeros_i) = p_i^R|_{\zeros_i}(X - \zeros_i) \geq p_i^R(X) = b_i^R(X) = b_i^R(X - \zeros_i),$$
        % where the final equality is because $\zeros_i = \{u \in V_i : b_i^R(u) = 0\}$. Thus, the set $X - \zeros_i \in \projTightSetFamilyAfterReduce{i}{i+1}$.
        \item By way of contradiction, let $i \in [\ell - 1]$ be the least index such that $\cumulativeMaximalTightSetFamily{i}|_{V_{i+1}} - T_{i+1} \not = \emptyset$. We note that $i \geq 2$ since $\cumulativeMaximalTightSetFamily{1} = \maximalTightSetFamily{1}\subseteq T_1$ and $T_1|_{V_2}\subseteq T_2$ by Lemma  \ref{lem:WeakToStrongCover:TightSetFamiliesDuringAlgorithm}(1) and \ref{lem:WeakToStrongCover:TightSetFamiliesDuringAlgorithm}(2), and hence, $\cumulativeMaximalTightSetFamily{1}|_{V_2}\subseteq T_2$. Thus, by our choice of index $i$, we have that $\cumulativeMaximalTightSetFamily{i-1}|_{V_i} \subseteq \tightSetFamily{i}$, and so it follows that $\cumulativeMaximalTightSetFamily{i} = \cumulativeMaximalTightSetFamily{i-1}|_{V_i} \cup \maximalTightSetFamily{i} \subseteq \tightSetFamily{i}$. Moreover, $T_i|_{V_{i+1}}\subseteq T_{i+1}$ by Lemma  \ref{lem:WeakToStrongCover:TightSetFamiliesDuringAlgorithm}(1) and \ref{lem:WeakToStrongCover:TightSetFamiliesDuringAlgorithm}(2). Consequently we have that $\cumulativeMaximalTightSetFamily{i}|_{V_{i+1}} \subseteq \tightSetFamily{i}|_{V_{i+1}} \subseteq \tightSetFamily{i+1}$, a contradiction to the choice of the index $i$. 
        %We also have that $\maximalTightSetFamily{i+1}\subseteq \tightSetFamily{i+1}$. Consequently, $\cumulativeMaximalTightSetFamily{i}|_{V_{i+1}}\cup \maximalTightSetFamily{i+1}\subseteq \tightSetFamily{i+1}$, a contradiction to the choice of the index $i$. 

        \iffalse
        We prove the claim by induction on $i$. 
        %We recall that $\cumulativeMaximalTightSetFamily{1} = \maximalTightSetFamily{1}$, and $\cumulativeMaximalTightSetFamily{i} = \left(\cumulativeMaximalTightSetFamily{i-1} \cup \maximalTightSetFamily{i-1}^M \cup \maximalTightSetFamily{i-1}^R\right)|_{V_i}\cup \maximalTightSetFamily{i}$ for every $i \geq 2$. 
        For the base case of $i=1$, we note that $\maximalTightSetFamily{1}\subseteq T_1$ and $T_1|_{V_2}\subseteq T_2$ by Lemma  \ref{lem:WeakToStrongCover:TightSetFamiliesDuringAlgorithm}(1) and \ref{lem:WeakToStrongCover:TightSetFamiliesDuringAlgorithm}(2), and hence, $\maximalTightSetFamily{1}|_{V_2}\subseteq T_2$. Let $i\ge 2$. 
        By the induction hypothesis, we have that $\cumulativeMaximalTightSetFamily{i-1}|_{V_i}\subseteq T_i$. Thus, $\cumulativeMaximalTightSetFamily{i} = \cumulativeMaximalTightSetFamily{i-1}|_{V_i} \cup \maximalTightSetFamily{i} \subseteq \tightSetFamily{i}$. Moreover, $T_i|_{V_{i+1}}\subseteq T_{i+1}$ by Lemma  \ref{lem:WeakToStrongCover:TightSetFamiliesDuringAlgorithm}(1) and \ref{lem:WeakToStrongCover:TightSetFamiliesDuringAlgorithm}(2). Consequently we have that $\cumulativeMaximalTightSetFamily{i}|_{V_{i+1}} \subseteq \tightSetFamily{i}|_{v_{i+1}} \subseteq \tightSetFamily{i+1}$.
        %Hence, $\cumulativeMaximalTightSetFamily{i-1}|_{V_i}|_{V_{i+1}}\subseteq T_{i+1}$. We also have that $\maximalTightSetFamily{i-1}^M\subseteq T_{i-1}^M\subseteq T_{i-1}^R$, $T_{i-1}^R|_{V_i}\subseteq T_i$, and $T_{i}|_{V_{i+1}}\subseteq T_{i+1}$, and hence $\maximalTightSetFamily{i-1}^M|_{V_i}|_{V_{i+1}}\subseteq T_{i+1}$. Similarly, we have that $\maximalTightSetFamily{i-1}^R\subseteq T_{i-1}^R$, $T_{i-1}^R|_{V_i}\subseteq T_i$, and $T_{i}|_{V_{i+1}}\subseteq T_{i+1}$, and hence $\maximalTightSetFamily{i-1}^R|_{V_i}|_{V_{i+1}}\subseteq T_{i+1}$. We have that $\maximalTightSetFamily{i}\subseteq T_i\subseteq T_i^M\subseteq T_i^R$ and $T_i^R|_{V_{i+1}}\subseteq T_{i+1}$, and hence, $\maximalTightSetFamily{i}|_{V_{i+1}}\subseteq T_{i+1}$. 
        \fi
        \item By way of contradiction, let $i \in [\ell - 1]$ be the least index such that the family $\cumulativeMaximalTightSetFamily{i}$  is not laminar. 
        %First, suppose that \cumulativeMaximalTightSetFamily{i} is not laminar.
        We note that $i \geq 2$ since the family $\cumulativeMaximalTightSetFamily{1} = \maximalTightSetFamily{1}$ is disjoint by \Cref{lem:weak-to-strong:uncrossing:maximal-tight-set-family-disjoint}. Thus, we have that $\cumulativeMaximalTightSetFamily{i} = \projCumulativeMaximalTightSetFamily{i-1}{i} \cup \maximalTightSetFamily{i}$.
        We note that the family $\maximalTightSetFamily{i}$ is also disjoint by \Cref{lem:weak-to-strong:uncrossing:maximal-tight-set-family-disjoint}. Furthermore, the family $\cumulativeMaximalTightSetFamily{i-1}$ (and consequently, the family  $\cumulativeMaximalTightSetFamily{i-1}|_{V_i}$) is laminar by our choice of $i\in[\ell - 1]$. Since the family $\cumulativeMaximalTightSetFamily{i}$ is not laminar, there exist distinct sets $X \in \projCumulativeMaximalTightSetFamily{i-1}{i}$ and $Y \in \maximalTightSetFamily{i}$ such that $X - Y, Y - X, X \cap Y \not = \emptyset$. By part (1) of the current lemma (shown above), we have that $X \in \tightSetFamily{i}$. Consequently, there exists $Z\in \maximalTightSetFamily{i}$ such that $X \subseteq Z$. We note that $Y \not = Z$ because $X -Z=\emptyset$ and $X - Y \not = \emptyset$. Thus, by \Cref{lem:weak-to-strong:uncrossing:maximal-tight-set-family-disjoint} we have that $Z\cap Y = \emptyset$. Therefore, $X\cap Y = \emptyset$, contradicting our choice of sets $X$ and $Y$.

        \iffalse
        \shubhang{Old proof with bug below.} By way of contradiction, let $i \in [\ell]$ be the least index such that $\cumulativeMaximalTightSetFamily{i}$ is not laminar. We note that $i \geq 2$ by \Cref{lem:weak-to-strong:uncrossing:maximal-tight-set-family-disjoint}. Then, there exist two sets $X, Y \in \cumulativeMaximalTightSetFamily{i}$ such that $X- Y, Y - X, X\cap Y \not = \emptyset$. We recall that $\cumulativeMaximalTightSetFamily{i} = \left(\cumulativeMaximalTightSetFamily{i-1} \cup \maximalTightSetFamily{i-1}^M \cup \maximalTightSetFamily{i-1}^R\right)|_{V_i}\cup \maximalTightSetFamily{i}$. By \Cref{lem:weak-to-strong:uncrossing:maximal-tight-set-family-disjoint}, both sets $X$ and $Y$ cannot be contained in the family $\maximalTightSetFamily{i}$. 

        First, suppose that $X, Y \in \left(\cumulativeMaximalTightSetFamily{i-1} \cup \maximalTightSetFamily{i-1}^M \cup \maximalTightSetFamily{i-1}^R\right)|_{V_i}$. Then, there exist sets $A_X, A_Y \in \cumulativeMaximalTightSetFamily{i-1} \cup \maximalTightSetFamily{i-1}^M \cup \maximalTightSetFamily{i-1}^R$ such that the sets $A_X-A_X, A_Y-A_X, A_X\cap A_Y\neq \emptyset$. If $A_X, A_Y \in \cumulativeMaximalTightSetFamily{{i-1}}$, then the index $i-1$ contradicts our choice of index $i$. Hence, at most one of the sets among $A_X$ and $A_Y$ can be in $\cumulativeMaximalTightSetFamily{{i-1}}$. 
        %Furthermore, if both sets $X, Y \in \projCumulativeMaximalTightSetFamily{i-1}{i}$, then there exist $A_X, A_Y \in \cumulativeMaximalTightSetFamily{i-1}$ such that $X \subseteq A_X$, $Y \subseteq A_Y$. Consequently the sets $A_X$ and $A_Y$ cross and the index $i-1$ contradicts our choice of index $i$. 
        We first consider the case where exactly one of the sets among $A_X$ and $A_Y$ is in $\cumulativeMaximalTightSetFamily{{i-1}}$. Without loss of generality, suppose that $A_X \in \cumulativeMaximalTightSetFamily{i-1}$. Suppose that $A_Y \in \maximalTightSetFamilyAfterMerge{i-1}$ (or $A_Y\in \maximalTightSetFamilyAfterReduce{i-1}$). Then, by \Cref{lem:WeakToStrongCover:TightSetFamiliesDuringAlgorithm}(1), we have that $\cumulativeMaximalTightSetFamily{i-1} \subseteq \tightSetFamily{i-1} \subseteq \tightSetFamilyAfterMerge{i-1}$ ($\subseteq \tightSetFamilyAfterReduce{i-1}$). In particular, there exists a set $B_X \in \maximalTightSetFamilyAfterMerge{i-1}$ (or $B_X \in \maximalTightSetFamilyAfterReduce{i-1}$) such that $A_X\subseteq B_X$. Since $A_X \cap A_Y \not = \emptyset$, we have that $B_X \cap A_Y \not = \emptyset$. Thus, by \Cref{lem:weak-to-strong:uncrossing:maximal-tight-set-family-disjoint}, we have that $A_X \subseteq B_X = A_Y$, contradicting $A_X-A_Y\neq \emptyset$. 
        Next, we consider the case where both $A_X$ and $A_Y$ are not in $\cumulativeMaximalTightSetFamily{{i-1}}$. By \Cref{lem:weak-to-strong:uncrossing:maximal-tight-set-family-disjoint}, both the sets $A_X$ and $A_Y$ cannot belong to the same family $\maximalTightSetFamilyAfterMerge{i-1}$ or $\maximalTightSetFamilyAfterReduce{i-1}$\shubhang{this is a bug. to use \Cref{lem:weak-to-strong:uncrossing:maximal-tight-set-family-disjoint}, every hyperedge needs to be contained in a tight set and every vertex needs to be contained in a hyperedge...}. Without loss of generality, let $A_X\in \maximalTightSetFamily{i-1}^M $ and $A_Y\in \maximalTightSetFamily{i-1}^R$. Then, by \Cref{lem:WeakToStrongCover:TightSetFamiliesDuringAlgorithm} part (1), we have that $\tightSetFamilyAfterMerge{i-1}\subseteq \tightSetFamilyAfterReduce{i-1}$. In particular, there exists a set $B_X \in \maximalTightSetFamilyAfterReduce{i-1}$ such that $A_X\subseteq B_X$. Since $A_X \cap A_Y \not = \emptyset$, we have that $B_X \cap A_Y \not = \emptyset$. Thus, by \Cref{lem:weak-to-strong:uncrossing:maximal-tight-set-family-disjoint}, we have that $A_X \subseteq B_X = A_Y$, contradicting $A_X-A_Y\neq \emptyset$. 
        %Using similar contradiction arguments, we can show that $A_X, A_Y \not \in \cumulativeMaximalTightSetFamily{i-1} \cup \maximalTightSetFamily{i-1}^M \cup \maximalTightSetFamily{i-1}^R$, a contradiction.
        
        
        Next, without loss of generality, suppose that $X \in \left(\cumulativeMaximalTightSetFamily{i-1} \cup \maximalTightSetFamily{i-1}^M \cup \maximalTightSetFamily{i-1}^R\right)|_{V_i}$ and $Y \in \maximalTightSetFamily{i}$. By \Cref{lem:WeakToStrongCover:TightSetFamiliesDuringAlgorithm} part (2), we have that $X \in \tightSetFamily{i}$. Thus, there exists a set $Z \in \maximalTightSetFamily{i}$ such that $X \subseteq Z$. Since $X-Y\neq \emptyset$, we have that $Z \not = Y$. Then, we have that $X\cap Y \subseteq Z \cap Y = \emptyset$, where the final equality is by \Cref{lem:weak-to-strong:uncrossing:maximal-tight-set-family-disjoint}. This contradicts $X\cap Y\neq \emptyset$.
        \fi
    \end{enumerate} 
\end{proof}

% \begin{enumerate}
%     \item $\tightSetFamily{i} \subseteq \tightSetFamilyAfterCreate{i}$,
%     \item  $\projTightSetFamilyAfterCreate{i}{i+1} \subseteq \tightSetFamily{i+1}$,  
%      \item $\cumulativeMaximalTightSetFamily{i} \subseteq \tightSetFamily{i}$,
%      \item $\cumulativeMaximalTightSetFamilyAfterCreate{i} \subseteq \tightSetFamilyAfterCreate{i}$, and
%      \item the families \cumulativeMaximalTightSetFamily{i} and \cumulativeMaximalTightSetFamilyAfterCreate{i} are laminar.
%  \end{enumerate}
% \begin{proof} 
% We recall that the function $p_{i} = (p_{i-1} - \alpha_{i-1} b_{(V_{i-1}, \{A_{i-1}\})})|_{\zeros_{i-1}}$ for $i \geq 2$. Since $p_1$ is a skew-supermodular function, we have that $p_{i}:2^{V_{i}}\rightarrow\Z$ is a skew-supermodular function for every $i\in [\ell]$. Furthermore, by \Cref{lem:Covering-Algorithm:hypothesis-after-contractionstep}, the functions $p_i$ and 
% $m_i:V_i\rightarrow\Z_+$ 
%     satisfy conditions (a) and (b) of \Cref{thm:WeakCoverViaUniformHypergraph:main} for every $i\in [\ell]$ since the functions $p_1$ and $m_1$ satisfy the two conditions. Let $i \in [\ell]$ be a recursive call. We now prove each claim separately below.
%     \begin{enumerate}
%         \item Let $X\in \tightSetFamily{i}$. We first show that $|A_i \cap X| \leq 1$. By way of contradiction, let $|A_i\cap X| \geq 2$. Then, we have the following:
%         \begin{align*}
%             m_{i}'(X) &= m_i(X) - \alpha_i\chi_{A_i}(X)&\\
%             &= p_i(X) - \alpha_i\chi_{A_i}(X)&\\
%             &\leq p_i(X) - 2\alpha_i&\\
%             &< p_i(X) - \alpha_i b_{(V_i, \{A_i\})}(X)&\\
%             & = p_{i}'(X),&
%         \end{align*}
%     which is a contradiction to \Cref{lem:Covering-Algorithm:hypothesis-after-createstep}. Here, the second equality is because $X \in \tightSetFamily{i}$. The first inequality is by our assumption that $|A_i\cap X|\ge 2$. The second inequality is because $b_{(V_i, \{A_i\})}(X) = 1$. We now show that $X \in \tightSetFamilyAfterCreate{i}$. We have the following: 
%     $$m_i'(X) = m_i(X) - \alpha_i|A_i\cap X| = p_i(X) - \alpha_i b_{(V_i, \{A_i\})}(X) = p_i'(X).$$
%     Here, the second equality holds because $X \in \tightSetFamily{i}$, and $b_{(V_i, \{A_i\})}(X) = |A_i\cap X|$ since $|A_i\cap X|\leq 1$. 

%     \item  We note that each set in the family $\projTightSetFamilyAfterCreate{i}{i+1}$ is of the form $Y - \zeros_i$ for some set $Y \in \tightSetFamilyAfterCreate{i}$. Thus, it suffices to consider an arbitrary set $X \in \tightSetFamilyAfterCreate{i}$ and show that $X - \zeros_i \in \tightSetFamily{i+1}$. We have the following:
%     $$m_{i+1}(X - \zeros_i) = m_i'(X) = p_i'(X) \leq p_{i+1}(X - \zeros_i) \leq m_{i+1}(X - \zeros_i).$$
%         Thus, all inequalities are equations and we have that $X - \zeros \in \tightSetFamily{i+1}$. Here, the first equality is because $m'(\zeros_i) = 0$. The second equality is because $X \in \tightSetFamilyAfterCreate{i}$. The first inequality is because $p_{i+1} = p_i'|_{\zeros_i}$. The final inequality is because $m_{i+1}(Z) \geq p_{i+1}(Z)$ for every $Z\subseteq V_{i+1}$ by \Cref{lem:Covering-Algorithm:hypothesis-after-contractionstep}. We remark that the statement (and proof) of this part can be strengthened to say that $\projTightSetFamilyAfterCreate{i}{i+1} = \tightSetFamily{i+1}$, but the weaker statement suffices for our purposes and so we omit the details. 
        
%         \item By way of contradiction, let $i \in [\ell]$ be the smallest recursive call such that the claim is false. We note that $i \geq 2$ because $\cumulativeMaximalTightSetFamily{1} = \maximalTightSetFamily{1} \subseteq \tightSetFamily{1}$. Thus, we have that $\cumulativeMaximalTightSetFamily{i} = \projCumulativeMaximalTightSetFamily{i-1}{i} \cup \projMaximalTightSetFamilyAfterCreate{i-1}{i} \cup \maximalTightSetFamily{i}$. We observe that  $\cumulativeMaximalTightSetFamily{i-1}\subseteq \tightSetFamily{i-1} \subseteq \tightSetFamilyAfterCreate{i-1}$, where the first inclusion is by our choice of $i \in [\ell]$, and the second inclusion is by part (a). Furthermore, we have that  $\projCumulativeMaximalTightSetFamily{i-1}{i} \cup \projMaximalTightSetFamilyAfterCreate{i-1}{i} \subseteq \projTightSetFamilyAfterCreate{i-1}{i} \subseteq \tightSetFamily{i}$, where the second inclusion is by parts (a) and (b). Thus, $\cumulativeMaximalTightSetFamily{i} \subseteq \tightSetFamily{i}$.

%         \item We recall that $\cumulativeMaximalTightSetFamilyAfterCreate{i} = \cumulativeMaximalTightSetFamily{i} \cup \maximalTightSetFamilyAfterCreate{i}$. Then, we have that $\cumulativeMaximalTightSetFamily{i} \subseteq \tightSetFamilyAfterCreate{i}$ by parts (c) and (a), and $\maximalTightSetFamilyAfterCreate{i} \subseteq \tightSetFamilyAfterCreate{i}$ by definition.
        
%         \item By way of contradiction, let $i \in [\ell]$ be the smallest recursive call such that the family $\cumulativeMaximalTightSetFamily{i}$ or \cumulativeMaximalTightSetFamilyAfterCreate{i} is not laminar. First, suppose that \cumulativeMaximalTightSetFamily{i} is not laminar.
%         We note that $i \geq 2$ as the family $\cumulativeMaximalTightSetFamily{1} = \maximalTightSetFamily{1}$ is disjoint by \Cref{lem:PropertiesOfTightSets:calTpmDisjoint}. Thus, we have that $\cumulativeMaximalTightSetFamily{i} = \projCumulativeMaximalTightSetFamilyAfterCreate{i-1}{i} \cup \maximalTightSetFamily{i}$.
%         We note that the family $\maximalTightSetFamily{i}$ is disjoint by \Cref{lem:PropertiesOfTightSets:calTpmDisjoint}. Furthermore, the family $\cumulativeMaximalTightSetFamilyAfterCreate{i-1}$ is laminar by our choice of $i\in[\ell]$. Since the family $\cumulativeMaximalTightSetFamily{i}$ is not laminar, there exist distinct sets $X \in \cumulativeMaximalTightSetFamily{i-1}$ and $Y \in \maximalTightSetFamily{i}$ such that the sets $X$ and $Y$ cross. By part (c), we have that $X \in \tightSetFamily{i}$. Consequently, there exists $Z\in \maximalTightSetFamily{i}$ such that $X \subseteq Z$. By \Cref{lem:PropertiesOfTightSets:calTpmDisjoint} we have that $Z\cap Y = \emptyset$. Thus, $X\cap Y = \emptyset$, contradicting our choice of sets $X$ and $Y$.

%         Next, suppose that \cumulativeMaximalTightSetFamily{i} is laminar, but \cumulativeMaximalTightSetFamilyAfterCreate{i} is not laminar. We recall that $\cumulativeMaximalTightSetFamilyAfterCreate{i} = \cumulativeMaximalTightSetFamily{i} \cup \maximalTightSetFamilyAfterCreate{i}$. Similar to the previous case, there exist distinct sets  $X \in \cumulativeMaximalTightSetFamily{i}$ and $Y \in \maximalTightSetFamilyAfterCreate{i}$ such that the sets $X$ and $Y$ cross. By parts (c) and (a), we have that $X \in \tightSetFamilyAfterCreate{i}$. Thus, there exists a set $Z \in \maximalTightSetFamilyAfterCreate{i}$ such that $X \subseteq Z$. By \Cref{lem:PropertiesOfTightSets:calTpmDisjoint}, either $Y = Z$ or $Y\cap Z = \emptyset$. In particular, we have that either $X \subseteq Y$ or $X\cap Y = \emptyset$, contradicting our choice of the sets $X,Y$.
%     \end{enumerate}
% \end{proof}

\begin{lemma}\label{lem:WeakToStrongCover:NumberOfHyperedges}
%Consider an execution of \Cref{alg:WeakToStrongCover}.
Suppose that the input to \Cref{alg:WeakToStrongCover} is a non-empty hypergraph $(H = (V, E), w)$ and a symmetric skew-supermodular function $p:2^V\rightarrow \Z$ 
%satisfying conditions \ref{cond:WeakToStrong:(a)} and \ref{cond:WeakToStrong:(b)} 
such that every hyperedge $\Tilde{e} \in E$ is contained in some $(p, H, w)$-tight set, $b_{(H,w)}(u) \not = 0$ for every $u \in V$, and $b_{(H,w)}(X) \geq p(X)$ for every $X \subseteq V$. 
Then, we have the following:
    \begin{enumerate}
        \item the recursion depth $\ell \in \Z_+$ of the algorithm is at most $|E| + 10|V| - 1$, and
        \item $|E^*| \leq |E| + 10|V| - 2$, where $(H^* = (V, E^*), w^*)$ denotes the hypergraph returned by the algorithm.
        %\knote{This lemma needs to be strengthened to say that the number of hyperedges in the hypergraph input to every intermediate recursive call is at most $|E|+10|V|$. This fact is the fact used in the run-time analysis. }
    \end{enumerate}
\end{lemma}

\begin{proof}
Let $i \in [\ell - 1]$ be a recursive call. We recall that $(H_i = (V_i, E_i), w_i:E_i \rightarrow \Z_+)$ and $p_i : 2^{V_i} \rightarrow \Z$ denote the inputs to the $i^{th}$ recursive call, and $(H_i^* = (V, E_i^*), w_i^*)$ denotes the hypergraph returned by the $i^{th}$ recursive call. Furthermore, $e_i, f_i, \beta_i^M, \alpha_i^M, (H_i^M, w_i^M), \beta_i^R, \alpha_i^R, (H_i^R, w_i^R)$, $p_i^R$, $(H_0^{i}, w_0^{i})$, $E_0^i, \zeros_i$ denote the quantities  as defined by \Cref{alg:WeakToStrongCover} during the $i^{th}$ recursive call.
    We prove the claims separately below.
    \begin{enumerate}
        \item We have the following:
        $$\ell \leq \left|\left\{i \in [\ell - 1] : \alpha_i^R = \beta_i^R\right\}\right|  +  \left|\left\{i \in [\ell- 1] : \alpha_i^R = w_i^M(e_i \cup f_i)\right\}\right| + 1 \leq |E| + 10|V| - 1,$$
        where the final inequality is by \Cref{claim:WeakToStrong:hyperedge-support-size:alpha^R=weight-of-added-hyperedge} and \Cref{claim:WeakToStrong:hyperedge-support-size:alpha^R<weight-of-added-hyperedge} below.
        \item 
        %\shubhang{todo: the proof is faulty/incomplete as it does not account for hyperedges that were reduced and added back.} Consider a recursive call $i \in [\ell]$. We recall that by part (4) of \Cref{lem:WeakToStrong:properties-of-alpha^M-and-alpha^R}, we have that $E_{i+1} = E_i - \{e_i\}\cdot \indicator_{\alpha_i^M = w_i(e_i)} - \{f_i\}\cdot \indicator_{\alpha_i^M = w_i(f_i)} + \{e_i\cup f_i\}\cdot \indicator_{\alpha_i^R < w_i^M(e_i\cup f_i)}$. Consequently, regardless of the value of $\alpha_i^M$, we have that $|E_{i+1}| \leq |E_i| + 1$. Furthermore, the inequality is an equality if and only if $\alpha_i^M = w_i(e_i) = w_i (f_i)$ and $\alpha_i^R = \beta_i^R < w_i^M(e_i\cup f_i)$. The number of such recursive calls is at most $5|V| - 1$ by \Cref{claim:WeakToStrong:hyperedge-support-size:alpha^R=weight-of-added-hyperedge}.
        %Let $i \in [\ell]$ be a recursive call. 
        %Let $\left(H_i = (V_i, E_i), w_i\right)$ denote the hypergraph input to the $i^{th}$ recursive call, and $\left(H_i^* = (V_i, E_i^*), w_i^*\right)$ denote the hypergraph returned by the $i^{th}$ recursive call. We recall that by the algorithm, we have that $(H_i^*, w_i^*) = (H_{i+1}^*, w_{i+1}^*) + \left((V_i,\{e_i\cup f_i\}), \alpha^R_i\chi_{e_i\cup f_i}\right)$. Consequently, 
        We have the following:
        $$|E_1^*| \leq |E_2^*| + |E_0^2| =   |E_{2}^*| + \indicator_{\alpha_1^R > 0} = |E_{\ell}^*| + \sum_{i \in [\ell - 1]}\indicator_{\alpha_i^R > 0} \leq \ell - 1\leq |E_1| + 10|V_1| - 2.$$
        Here, the second inequality is because $E_{\ell}^* = \emptyset$ by definition of the \Cref{alg:WeakToStrongCover} and the final inequality is by part (1) of the claim (shown above).
    \end{enumerate}
\end{proof}

\begin{claim}\label{claim:WeakToStrong:hyperedge-support-size:alpha^R=weight-of-added-hyperedge}
$\left|\left\{i \in [\ell-1] : \alpha_i^R = \beta_i^R\right\}\right| \leq 5|V| - 1$.
\end{claim}
\begin{proof}
We define a potential function $\phi:[\ell - 1] \rightarrow\Z_{\geq 0}$ as follows: for every $i \in [\ell]$,
    $$\phi(i) \coloneqq  \left|\cumulativeMaximalTightSetFamily{i}\right| + 3|\cumulativeZeros{i-1}|,$$
    where we define the set $\cumulativeZeros{0} \coloneqq  \emptyset$.
We first show that the function $\phi$ is monotone non-decreasing.
Let $i \in [\ell -1]$ be a recursive call of \Cref{alg:WeakToStrongCover}. If $\zeros_i = \emptyset$, then $|\cumulativeMaximalTightSetFamily{i}| \leq |\cumulativeMaximalTightSetFamily{i+1}|$ since $\cumulativeMaximalTightSetFamily{i+1} = \cumulativeMaximalTightSetFamily{i} |_{V_{i+1}}\cup \maximalTightSetFamily{i+1} = \cumulativeMaximalTightSetFamily{i}  \cup \maximalTightSetFamily{i+1},$   where the second equality is because $V_{i+1} = V_i - \zeros_i = V_i$. Then, we have that $\phi(i+1) \geq \phi(i)$ since $\cumulativeZeros{i-1} \subseteq \cumulativeZeros{i}$ by definition. Thus, we suppose that $\zeros_i \not = \emptyset$. We have the following: 
\begin{align*}
    \phi(i+1) - \phi(i)& = |\cumulativeMaximalTightSetFamily{i+1}| + 3 |\cumulativeZeros{i}| - |\cumulativeMaximalTightSetFamily{i}| - 3|\cumulativeZeros{i-1}|&\\
    & = \left|\cumulativeMaximalTightSetFamily{i}|_{V_{i+1}}\cup \maximalTightSetFamily{i+1}  \right| - |\cumulativeMaximalTightSetFamily{i}| + 3|\zeros_{i}|&\\
    &\geq \left|\cumulativeMaximalTightSetFamily{i}|_{V_{i+1}}\right| - |\cumulativeMaximalTightSetFamily{i}| + 3|\zeros_{i}|&\\
    & \geq  0.&
\end{align*}
Here, the final inequality is because of the following: we recall that the family $\cumulativeMaximalTightSetFamily{i}$ is laminar by \Cref{lem:WeakToStrong:properties-of-cumulative-tight-set-families}. Furthermore, the family $\projCumulativeMaximalTightSetFamily{i}{i+1}$ is the projection of the laminar family $\cumulativeMaximalTightSetFamily{i}$ to the ground set $V_{i+1} = V_i - \zeros_i$. The inequality then follows by \Cref{lem:CoveringAlgorithm:projection-laminar-family}  --  here, we observe that  \Cref{lem:CoveringAlgorithm:projection-laminar-family} says that this inequality is strict because $\zeros_i \not = \emptyset$. We will refer back to this observation in the next part of the proof.

Consider a recursive call $i \in \left\{i \in [\ell - 1] : \alpha_i^R = \beta_i^R\right\}$. We now show that $\phi(i+1) > \phi(i)$, i.e. the potential value strictly increases. 
We suppose that $\zeros_i = \emptyset$ as otherwise $\phi(i+1) > \phi(i)$ by the observation in the final line of the previous paragraph.
Thus, $p_{i+1} = p_i^R$ and $(H', w') = (H^R, w^R)$. 
We note that the hypergraph $(H_i, w_i)$ is non-empty since $i < \ell$.
We consider a set $X\subseteq V_i$ such that $e_i\cup f_i\subseteq X\subseteq V_i$ and $b_i^M(X)-p_i(X)=\beta_i^R$ (such a set exists by the definition of $\beta_i^R$). 
Then, we have the following: $b_{i+1}(X) = b_i^R(X)=b_i^M(X)-\alpha_i^R=p_i(X)=p_i^R(X) = p_{i+1}(X)$, where the first equality is because $(H', w') = (H^R, w^R)$, the second equality is because $\zeros_i = \emptyset$, the third equality is because $\alpha_i = \beta_i$ by our choice of $i \in [\ell - 1]$, and the final equality is because $p_{i+1} = p_{i}^R$. Thus, $X \in \tightSetFamily{i+1}$. Consequently, 
%Since $\alpha_i^R=\beta_i^R$, we have that $b_i^R(X)=b_i^M(X)-\alpha_i^R=p_i(X)=p_i^R(X)$ and hence, $X\in \tightSetFamily{i}^R$. In particular, 
there exists a set $Y \in \maximalTightSetFamily{i+1}$ such that $e_i\cup f_i \subseteq X \subseteq Y$. We note that $Y \not \in \tightSetFamily{i}$ due to the following: suppose (for sake of contradiction) that $Y \in \tightSetFamily{i}$, then 
%In particular, we have that the set $\maximalTightSetContainingHyperedge{i}^R(e_i\cup f_i)\not = \emptyset$ because $e_i\cup f_i \subseteq X$. We note that the set $\maximalTightSetContainingHyperedge{i}^R(e_i\cup f_i) \not \in \tightSetFamily{i}$ due to the following: suppose that the set $Y\coloneqq \maximalTightSetContainingHyperedge{i}^R(e_i\cup f_i) \in \tightSetFamily{i}$, then
$b_i(Y)=p_i(Y)$ and $b_i^M(Y)=b_i(Y)-\alpha_i^M=p_i(Y)-\alpha_i^M<p_i(Y)$, a contradiction to \Cref{lem:WeakToStrong:properties-of-alpha^M-and-alpha^R}(3) -- the first equality is by definition of $b_i^M$ and $e_i\cup f_i\subseteq Y$ and the inequality is because $\alpha_i^M\ge 1$ by \Cref{lem:WeakToStrong:properties-of-alpha^M-and-alpha^R}(1). Thus, $Y \in \maximalTightSetFamily{i+1} - \tightSetFamily{i} \subseteq \maximalTightSetFamily{i+1} - \cumulativeMaximalTightSetFamily{i}$ by \Cref{lem:WeakToStrong:properties-of-cumulative-tight-set-families}(1).
Consequently, we have that
$\cumulativeMaximalTightSetFamily{i} \subsetneq \cumulativeMaximalTightSetFamily{i} \cup \maximalTightSetFamily{i+1} = \cumulativeMaximalTightSetFamily{i+1}.$
%Here, the strict inclusion is because $ \maximalTightSetFamily{i}^R - \cumulativeMaximalTightSetFamily{i} \not = \emptyset$ and the equality is because $V_{i+1} = V_i - \zeros = V_i$.
%since $\cumulativeMaximalTightSetFamily{i} \subseteq \tightSetFamily{i}$ by \Cref{lem:WeakToStrong:properties-of-cumulative-tight-set-families}. The first equality is because $\alpha_i^R = 0$ by \Cref{lem:WeakToStrong:properties-of-alpha^M-and-alpha^R} and $V_{i+1} = V_i - \zeros_i = V_i$. Thus, $\phi(i+1) > \phi(i)$.

We now show the claim. By the previous two properties of our potential function $\phi$, we have the following:
$$\left|\left\{i \in [\ell - 1] : \alpha_i^R = \beta_i^R\right\}\right| \leq \phi(\ell - 1) - \phi(1) \leq |\cumulativeMaximalTightSetFamily{\ell - 1}| + 3|\cumulativeZeros{\ell - 1}| \leq 5|V| - 1,$$
where the final inequality is because the family $\cumulativeMaximalTightSetFamily{\ell - 1}$ is laminar by \Cref{lem:WeakToStrong:properties-of-cumulative-tight-set-families}(2).
\end{proof}

% \begin{proof}
%     Let $\ell \in \Z_+$ denote the number of recursive calls of the execution of \Cref{alg:WeakToStrongCover}. Let $i \in [\ell]$ be a recursive call and let the function $p_i:2^{V_i}\rightarrow \Z$ and the hypergraph $(H_i = (V_i, E_i), w_i)$ denote the inputs to the $i^{th}$ recursive call. 
%     We define a potential function $\phi:[\ell] \rightarrow\Z_{\geq 0}$ as follows: for every $i \in [\ell]$,
%     $$\phi(i) \coloneqq  \left|\cumulativeMaximalTightSetFamily{i}\right| + 3|\cumulativeZeros{i-1}|,$$
%     where we define the set $\cumulativeZeros{0} \coloneqq  \emptyset$.
%      Let $(H^* = (V_i, E_i^*), w_i^*)$ denote the hypergraph returned by the $i^{th}$ recursive call. Then, \Cref{claim:WeakToStrong:hyperedge-support-szie:main} below shows that the function $\phi$ is monotone non-decreasing and that $\phi(i+1) > \phi(i)$ if $|E^*_{i}| > |E^*_{i+1}|$. We note that by the definition of \Cref{alg:WeakToStrongCover}, $|E^*_{i}| - |E^*_{i+1}| \leq 1$.
%     Thus, we have the following:
%     $$|E^*| - |E| \leq \sum_{i \in [\ell - 1]} \left(\phi(i+1) - \phi(i)\right)
%          \leq |\cumulativeMaximalTightSetFamily{\ell}| + 3|\cumulativeZeros{\ell - 1}|\leq 5|V| - 1,$$
% where the second inequality is because the function $\phi$ is non-decreasing and the final inequality is because $|\cumulativeMaximalTightSetFamily{\ell}| \leq 2|V| - 1$ since the family $\cumulativeMaximalTightSetFamily{\ell}$ is laminar by \Cref{lem:WeakToStrong:properties-of-cumulative-tight-set-families}.
% \end{proof}

\begin{claim}\label{claim:WeakToStrong:hyperedge-support-size:alpha^R<weight-of-added-hyperedge}
    $\left|\left\{i \in [\ell - 1] : \alpha_i^R = w_i^M(e_i \cup f_i)\right\}\right| \leq |E| + 5|V| - 1$. %\kristof{Maybe dismiss the index of $E$ and $V$?}
\end{claim}
\begin{proof}
    We define a potential function $\phi:[\ell - 1] \rightarrow \Z_{\geq 0}$ as follows: for every $i \in [\ell]$,
    $$\phi(i) \coloneqq  -|E_i| - \left|\left\{j \in [i,\ell - 1] : \alpha_j^R < w_j^M(e_j\cup f_j)\right\}\right|.$$
    
    First, consider a recursive call $i \in [\ell - 1]$ such that $\alpha_i^R = \beta_i^R < w_i^M(e_i \cup f_i)$. We will show that the potential function $\phi$ does not decrease. By Lemma \ref{lem:WeakToStrong:properties-of-alpha^M-and-alpha^R}(2) and \ref{lem:WeakToStrong:properties-of-alpha^M-and-alpha^R}(5), we have that 
    \begin{align*}
       E_{i+1} & = E_i - \{e_i\}\cdot \indicator_{\alpha_i^M = w_i(e_i)} - \{f_i\}\cdot \indicator_{\alpha_i^M = w_i(f_i)} + \{e_i\cup f_i\}\cdot \indicator_{\alpha_i^R = \beta_i^R < w_i^M(e_i \cup f_i)}&\\
        &= E_i - \{e_i\}\cdot \indicator_{\alpha_i^M = w_i(e_i)} - \{f_i\}\cdot \indicator_{\alpha_i^M = w_i(f_i)} + \{e_i\cup f_i\}.&
    \end{align*}
Consequently, 
    %regardless of the value of $\alpha_i^M$, 
    we have that $|E_{i+1}| \leq |E_i| + 1$. Furthermore, by the choice of recursive call $i$, we have that $\left|\{j \in [i+1, \ell - 1] : \alpha_j^R < w_j^M(e_j\cup f_j)\}\right| = \left|\{j \in [i, \ell - 1] : \alpha_j^R < w_j^M(e_j\cup f_j)\}\right| - 1$. Thus, $\phi(i+1) \geq \phi(i)$.
    
     Next, consider a recursive call $i \in [\ell - 1]$ such that $\alpha_i^R = w_i^M(e_i \cup f_i)$. We will show that the potential function $\phi$ strictly increases. We note that the hypergraph $(H_i, w_i)$ is non-empty since $i < \ell$.  
     By \Cref{lem:WeakToStrong:properties-of-alpha^M-and-alpha^R}(1), we have that $\alpha_i^R = w_i^M(e_i \cup f_i) = \alpha_i^M \geq 1$. 
     Consequently, by  \Cref{lem:WeakToStrong:alpha^R=0-iff-alpha^M=beta^M}, we have that $\alpha_i^M  < \beta_i^M$ and hence, $\alpha_i^M \in \{w_i(e_i), w_i(f_i)\}$ by definition. 
     Suppose that $\alpha_i^M = w_i(e_i)$. 
     Then, we have that $w^R_i(e_i) = w^M_i(e_i) = w_i(e_i) - \alpha_{i}^M = 0$. 
     Thus, we have that $e_i \not \in E_{i+1}$. 
     Then, by our choice of $i$ and by Lemma \ref{lem:WeakToStrong:properties-of-alpha^M-and-alpha^R}(2) and \ref{lem:WeakToStrong:properties-of-alpha^M-and-alpha^R}(5), we have that $E_{i+1} \subsetneq E_i$. 
     We note that this strict inclusion can also be obtained in the case where $\alpha_i^M = w_i(f_i)$ using the same argument. Furthermore, by the choice of recursive call $i$, we have that $\left|\{j \in [i+1, \ell - 1] : \alpha_j^R < w_j^M(e_j\cup f_j)\}\right| = \left|\{j \in [i, \ell - 1] : \alpha_j^R < w_j^M(e_j\cup f_j)\}\right|$.
     Thus, $\phi(i+1) > \phi(i)$.

      The previous two properties of our potential function give us the following:
      $$\left|\left\{i \in [\ell - 1] : \alpha_i^R = w_i^M(e_i \cup f_i)\right\}\right| \leq \phi(\ell - 1) - \phi(1) \leq |E_1| + \left|\left\{j \in [\ell - 1] : \alpha_j^R < w_j^M(e_j\cup f_j)\right\}\right| \leq |E| + 5|V| - 1,$$ 
      %\kristof{Maybe dismiss the index of $E$ and $V$?}
      where the final inequality is by \Cref{claim:WeakToStrong:hyperedge-support-size:alpha^R=weight-of-added-hyperedge} -- we note that if $\alpha_j^R < w_j^M(e_j\cup f_j)$, then $\alpha_j^R = \beta_j^R$ by definition.
\end{proof}

\iffalse
\subsection{Strongly Polynomial Runtime of \Cref{alg:WeakToStrongCover}}\label{sec:WeakToStrongCover:strongly-poly-runtime:main}
\shubhang{Todo.}
In this section, we show that \Cref{alg:WeakToStrongCover} can be implemented to run in strongly polynomial time, given access to \functionMaximizationOracleStrongCover{p}. As a necessary step to show this, the next lemma says that maximization oracle for the input function $p$ during a recursive call of \Cref{alg:WeakToStrongCover} can be used to implement the maximization oracle for the input function $p'$ to the subsequent recursive call in strongly polynomial time. We note that the proof of the lemma is very similar to the proof of \Cref{lem:p''-max-oracle-from-p-max-oracle} in \Cref{sec:CoveringAlgorithm:p''-max-oracle-using-p-max-oracle} and so we omit its proof for brevity.

\begin{lemma}\label{lem:WeakToStrong:p'-oracle-from-p-oracle}
Consider a recursive call of \Cref{alg:WeakToStrongCover}. Suppose that the input to the recursive call is a non-empty hypergraph $(H = (V, E), w)$ and a symmetric skew-supermodular function $p:2^V\rightarrow \Z$ satisfying conditions \ref{cond:WeakToStrong:(a)} and \ref{cond:WeakToStrong:(b)} such that $b_{(H,w)}(X) \geq p(X)$ for every $X \subseteq V$. Let $p':2^{V'}\rightarrow\Z$ be the function constructed by the algorithm as input to the subsequent recursive call. 
Then, a  \functionMaximizationOracleStrongCover{p'} can be implemented using a  \functionMaximizationOracleStrongCover{p} in strongly polynomial time.
\end{lemma}

The next lemma shows that the family $\maximalTightSetFamily{p, H, w}$ can be computed in strongly polynomial number of calls to \functionMaximizationOracleStrongCover{p}. We will subsequently use this in \Cref{lem:WeakToStrongCover:single-recursive-call-strongly-polytime} to show that a recursive call of \Cref{alg:WeakToStrongCover} runs in strongly polynomial time given the maximization oracle of the input function.

 \begin{lemma}\label{lem:WeakToStrong::maximal-tight-set-family-computed-in-strongly-polytime}
 Let $p:2^V\rightarrow\Z$ be a skew-supermodular function and let $\left(H = (V, E), w:E \rightarrow\Z_+\right)$ be a hypergraph satisfying conditions \ref{cond:WeakToStrong:(a)} and \ref{cond:WeakToStrong:(b)} such that $b_{(H,w)}(X) \geq p(X)$ for every $X \subseteq V$. Then, the family $\maximalTightSetFamily{p, H, w}$ can be computed in $\poly(|V|)$ calls to \functionMaximizationOracleStrongCover{p}.
\end{lemma}
\begin{proof}
    We note that we can find a $(p, H, w)$-tight set $X\subseteq V$ 
    %containing the set $Z \subseteq V$ 
    (if one exists) by a single call to \functionMaximizationOracleStrongCover{p} by letting  $(G_0, w_0) \coloneqq  \left((V, \emptyset), \emptyset\right)$ to be the empty hypergraph, the sets $S_0, T_0 \coloneqq  \emptyset$, and the vector $y_0$ to be the all-zeros vector as inputs to $\functionMaximizationOracleStrongCover{{p}}$. 
    Let $X \subseteq V$ be a $(p, H, w)$-tight set. Then, we can find a $(p, H, w)$-tight set $Y \subseteq V$ such that $X \subseteq Y$ (if such a set $Y$ exists) by guessing a vertex $u \in Y - X$ (by enumerating over all vertices $u \in V - X$) and querying $\functionMaximizationOracleStrongCover{p}$ on inputs  $(G_0, w_0) \coloneqq  \left((V, \emptyset), \emptyset\right)$, $S_0 \coloneqq  X \cup \{u\}$, $T_0 \coloneqq  \emptyset$, and the vector $y_0$ as the all-zeros vector. Iteratively applying this procedure will result in a maximal $(p, H, w)$-tight set $Z \in \maximalTightSetFamily{p, H, w}$. By \Cref{lem:weak-to-strong:uncrossing:maximal-tight-set-family-disjoint}, The family $\maximalTightSetFamily{p, H, w}$ is disjoint. Thus, recursively applying the procedure on the set $V - Z$ (i.e. including $Z$ in the parameter $T_0$ to all subsequent queries to \functionMaximizationOracleStrongCover{p}) will give us the family $\maximalTightSetFamily{p, H, w}$. This procedure requires  $\poly(|V|)$ many calls to \functionMaximizationOracleStrongCover{p}.
    %Let $X = \functionMaximizationOracleStrongCover{p}\left((G_0,w_0), S_0, T_0, y_0\right)$ be a $(p, H, w)$-tight set for some choices of inputs $(G_0, w_0), S_0, T_0,$ and $y_0$. Then, we can find a $(p, H, w)$-tight set $Y \subseteq V$ such that $X \subseteq Y$ (if such a set $Y$ exists) by guessing a vertex $u \in Y - X$ (by enumerating over all vertices $u \in V - X$) and computing $\functionMaximizationOracleStrongCover{p}\left((H,w), S\cup \{u\}, T, y\right)$. Iteratively applying this procedure will result in a maximal $(p, H, w)$-tight set. This requires at most $O(|V|)$ calls to \functionMaximizationOracleStrongCover{p}.
    \end{proof}

The next lemma says that \Cref{alg:WeakToStrongCover} can be implemented so that a single recursive call (excluding the time of recursion) runs in strongly polynomial time and makes $\poly(|V|)$ calls to \functionMaximizationOracleStrongCover{p}.
  \begin{lemma}\label{lem:WeakToStrongCover:single-recursive-call-strongly-polytime}
      Let $p:2^V\rightarrow\Z$ be a symmetric skew-supermodular function and let $\left(H = (V, E), w:E \rightarrow\Z_+\right)$ be a hypergraph satisfying conditions \ref{cond:WeakToStrong:(a)} and \ref{cond:WeakToStrong:(b)} such that $b_{(H,w)}(X) \geq p(X)$ for every $X \subseteq V$. Then, 
 \Cref{alg:WeakToStrongCover} can be implemented such that a recursive call of the algorithm on the input $((H,w), p)$ runs in strongly polynomial time and makes $\poly(|V|)$ calls to \functionMaximizationOracleStrongCover{p}.
 %, where $p:2^V\rightarrow \Z$ is a symmetric skew-supermodular function and $(H, w)$ is a hypergraph such that $b_{(H,w)}(Z) \geq p(Z)$ for every $Z \subseteq V$, $b_{(H,w)}(u)>0$ for every $u \in V$, and $\maximalTightSetContainingHyperedge{p}(\Tilde{e}) \not = \emptyset$ for every $\Tilde{e} \in E$.
  \end{lemma}  
  \begin{proof}
  It suffices to show that  \Cref{alg:WeakToStrongCover:(4)}, \Cref{alg:WeakToStrongCover:(5)} and \Cref{alg:WeakToStrongCover:(7)} of \Cref{alg:WeakToStrongCover} can be implemented to run in strongly polynomial time.
  We recall that step \Cref{alg:WeakToStrongCover:(4)} is to compute two hyperedges $e, f \in E$ such that $\maximalTightSetContainingHyperedge{p}(e) \not = \maximalTightSetContainingHyperedge{p}(f)$. We now show how to implement this computation in strongly polynomial time.  By \Cref{lem:weak-to-strong:uncrossing:maximal-tight-set-family-disjoint}, the family $\maximalTightSetFamily{p, H, w}$ is a disjoint family. Thus, $|$\maximalTightSetFamily{p, H, w}$| = O(|V|)$. Furthermore,  by \Cref{lem:WeakToStrong::maximal-tight-set-family-computed-in-strongly-polytime}, we can explicitly compute the family $\maximalTightSetFamily{p, H, w}$ in strongly polynomial time. We note that $\maximalTightSetContainingHyperedge{p}(e) \in \maximalTightSetFamily{p, H, w}$ for every $e \in E$.
      Moreover, by \Cref{cor:weak-to-strong:uncrossing:exist-two-hyperedges-contained-in-disjoint-maximal-tight-sets}, there exist hyperedges $e, f \in E$ such that $\maximalTightSetContainingHyperedge{p}(e) \cap \maximalTightSetContainingHyperedge{p}(f) = \emptyset$. Thus, we can find the required hyperedges for step (4) of \Cref{alg:WeakToStrongCover} by finding a pair of sets $X, Y \in \maximalTightSetFamily{p, H, w}$ that each contain a hyperedge, and define these hyperedges to be $e$ and $f$ respectively.

      We now show how to implement \Cref{alg:WeakToStrongCover:(5)} of \Cref{alg:WeakToStrongCover} in strongly polynomial time. In particular, we show how to compute $\beta^M$ in strongly polynomial time.
We have the following:
\begin{align*}
    \beta^M& = \min\left\{b_{(H,w)}(X) - p(X) : X \subseteq V \text{ and } e \cap X, f \cap X \not = \emptyset\right\} &\\
    & = \min_{u \in e, v \in f}\left\{\min\left\{b_{(H,w)}(X) - p(X) : \{u,v\} \subseteq X \subseteq V\right\}\right\}&\\
    &=: \min_{u \in e, v \in f}\left\{\beta^M_{uv}\right\},&
\end{align*}
where the quantity $\beta^M_{uv}$ is defined by the third equation. Then, $\beta^M$ can be computed in strongly polynomial time because  for every $u \in e$ and $v \in f$, the value $\beta^M_{uv}$ can be computed using a single oracle call to \functionMaximizationOracleStrongCover{p}. Finally, we observe that step \Cref{alg:WeakToStrongCover:(7)} of the algorithm can be computed in strongly polynomial time using a single call to $\functionMaximizationOracleStrongCover{p}$.
  \end{proof}

    We now show the main result of this section which says that \Cref{alg:WeakToStrongCover} can be implemented to run in strongly polynomial time.

\begin{lemma}\label{lem:WeakToStrongCover:strongly-polynomial-runtime:main}
    Let $p:2^V\rightarrow\Z$ be a symmetric skew-supermodular function and let $\left(H = (V, E), w:E \rightarrow\Z_+\right)$ be a hypergraph satisfying conditions \ref{cond:WeakToStrong:(a)} and \ref{cond:WeakToStrong:(b)} such that $b_{(H,w)}(X) \geq p(X)$ for every $X \subseteq V$. Then, \Cref{alg:WeakToStrongCover} can be implemented such that an execution of the algorithm on the input $((H,w), p)$ terminates in strongly polynomial time and makes only $\poly(|V|)$ calls to \functionMaximizationOracleStrongCover{p}.
    %, where $p:2^V\rightarrow \Z$ is a symmetric skew-supermodular function and $(H, w)$ is a hypergraph such that $b_{(H,w)}(Z) \geq p(Z)$ for every $Z \subseteq V$, $b_{(H,w)}(u)>0$ for every $u \in V$, and $\maximalTightSetContainingHyperedge{p}(\Tilde{e}) \not = \emptyset$ for every $\Tilde{e} \in E$.
\end{lemma}
\begin{proof}
 By \Cref{lem:WeakToStrongCover:NumberOfHyperedges} the execution of \Cref{alg:WeakToStrongCover} witnesses at most $\ell = O(|E| + |V|)$ recursive calls. Let $i \in [\ell]$ be a recursive call and let the hypergraph $(H_{i}, w_i)$ and function $p_i$ denote the inputs to the $i^{th}$ recursive call. Then, by \Cref{lem:WeakToStrong:p'-oracle-from-p-oracle}, \functionMaximizationOracleStrongCover{p_i} can be used to implement \functionMaximizationOracleStrongCover{p_{i+1}} in strongly polynomial time. By induction on $i$, we have that \functionMaximizationOracleStrongCover{p_i} can be implemented using \functionMaximizationOracleStrongCover{p} in strongly polynomial time. The claim then follows by \Cref{lem:WeakToStrongCover:single-recursive-call-strongly-polytime}.
\end{proof}
\fi


\subsection{Run-time of the Algorithm}
%\Cref{alg:WeakToStrongCover}}
\label{sec:WeakToStrongCover:strongly-poly-runtime:main}
In this section, we show that \Cref{alg:WeakToStrongCover} can be implemented to run in strongly polynomial time. 
\Cref{lem:WeakToStrong:strongly-polytime:main} is the main lemma of this section. 
In order to show that each recursive call of \Cref{alg:WeakToStrongCover} can be implemented in polynomial time, we will require the ability to compute the family \maximalTightSetFamily{p, H, w}.  
In the following lemma, we show that this family can indeed be computed in polynomial time. We defer the proof of the lemma to \Cref{appendix:sec:Function-Maximization-Oracles:computing-maximal-pHw-tight-sets-family}.

\begin{restatable}{lemma}{lemComputingMaximalTightSetFamily}\label{lem:FunctionMaximizetionOracles:computing-maximal-pHw-tight-sets-with-p-max-oracle}
 Let $p:2^V\rightarrow\Z$ be a skew-supermodular function and let $\left(H = (V, E), w:E \rightarrow\Z_+\right)$ be a hypergraph such that every hyperedge $e \in E$ is contained in some $(p, H, w)$-tight set, $b_{(H,w)}(u) \not = 0$ for every $u \in V$, and $b_{(H,w)}(X) \geq p(X)$ for every $X \subseteq V$.
Then, the family $\maximalTightSetFamily{p, H, w}$ can be computed in $O(|V|^2(|V|+|E|))$ time using $O(|V|^2)$ queries to \functionMaximizationOracle{p}. The run-time includes the time to construct the hypergraphs used as input to the queries to \functionMaximizationOracle{p}. Moreover, for each query to \functionMaximizationOracle{p}, the hypergraph $(G_0, c_0)$ used as input to the query has at most $|V|$ vertices and at most $|E| + 1$ hyperedges.  \end{restatable}

We now bound the run-time of \Cref{alg:WeakToStrongCover}. 

\begin{lemma}\label{lem:WeakToStrong:strongly-polytime:main}
Suppose that the input to \Cref{alg:WeakToStrongCover} is a non-empty hypergraph $(H = (V, E), w)$ and a symmetric skew-supermodular function $p:2^V\rightarrow \Z$ such that every hyperedge $\Tilde{e} \in E$ is contained in some $(p, H, w)$-tight set, $b_{(H,w)}(u) \not = 0$ for every $u \in V$, and $b_{(H,w)}(X) \geq p(X)$ for every $X \subseteq V$.
 Then, \Cref{alg:WeakToStrongCover}  can be implemented to run in $O(|V|^3(|V| + |E|)^2)$ time using $O(|V|^3(|V|+|E|))$ queries to \functionMaximizationOracleStrongCover{p}. The run-time includes the time to construct the hypergraphs used as input to the queries to \functionMaximizationOracleStrongCover{p}.
Moreover, for each query to \functionMaximizationOracleStrongCover{p}, the hypergraph $(G_0, c_0)$ used as input to the query has $O(|V|)$ vertices and $O(|V|+|E|)$ hyperedges.
\end{lemma}
\begin{proof}
Let $\ell$ denote the recursion depth of \Cref{alg:WeakToStrongCover}. By \Cref{lem:WeakToStrongCover:NumberOfHyperedges}, we have that $\ell=O(|E|+|V|)$. 
Let  $i \in [\ell]$ be a recursive call.
 We recall that $p_i:2^{V_i}\rightarrow\Z$ and $(H_i, w_i)$ are inputs to the $i^{th}$ recursive call, where $p_1 := p$ and $(H_1, w_1) := (H, w)$.
By Lemmas \ref{lem:weak-to-strong-cover:next-recursive-call-hypothesis-satification:symm-skew-supermod-strong-cover} and \ref{lem:weak-to-strong-cover:next-recursive-call-hypothesis-satification:hyperedge-in-tight-set-positive-coverage} and induction on $i$, we have that $p_i$ is a symmetric skew-supermodular function and $(H_i = (V_i, E_i), w_i)$ is a (non-empty if $i \leq \ell - 1$) hypergraph such that every hyperedge $\Tilde{e} \in E_i$ is contained in some $(p_i, H_i, w_i)$-tight set, $b_{(H_i, w_i)}(u) \not = 0$ for every $u \in V_i$, and $b_{(H_i, w_i)}(X) \geq p_i(X)$ for every $X \subseteq V_i$. Let $e_i, f_i, \beta_i^M, \alpha_i^M, (H_i^M, w_i^M), \beta_i^R, \alpha_i^R, (H_i^R, w_i^R)$, $p_i^R$, $(H_0^{i}, w_0^{i})$, $E_0^i, \zeros_i$ denote the quantities  as defined by \Cref{alg:WeakToStrongCover} during the $i^{th}$ recursive call. We recall that $V_i \subseteq V$. By Lemma \ref{lem:WeakToStrong:properties-of-alpha^M-and-alpha^R}(2), (5) and induction on $i$, we have that $|E_i| \leq |E| + i \leq |E| + \ell = O(|E| + |V|)$, where the final equality is by \Cref{lem:WeakToStrongCover:NumberOfHyperedges}(1). 
%Consequently, $\poly(|V_i|, |E_i|) = \poly(|V|, |E|)$ and $\poly(|V_i|) = \poly(|V|)$. 
Furthermore, the hypergraph returned in Step \ref{alg:WeakToStrongCover:(14)} contains at most $|V|$ vertices and at most $|E_i| + O(|V_i|)= O(|V| + |E|)$ hyperedges by \Cref{lem:WeakToStrongCover:NumberOfHyperedges}(2) and induction on $i$. 
The following claim shows that all steps except for Steps \ref{alg:WeakToStrongCover:(10)},  \ref{alg:WeakToStrongCover:(13)}, and  \ref{alg:WeakToStrongCover:(14)} of \Cref{alg:WeakToStrongCover} in the $i^{th}$ recursive call can be implemented to run in $O(|V_i|^2(|V_i| + |E_i|))$ time
using $O(|V_i|^2)$ queries \functionMaximizationOracle{p_i}. This includes the run-time to construct the hypergraphs used in the queries to \functionMaximizationOracle{p_i}. We will address the run-time for Steps \ref{alg:WeakToStrongCover:(10)},  \ref{alg:WeakToStrongCover:(13)}, and  \ref{alg:WeakToStrongCover:(14)} later. 

\begin{claim}\label{claim:weak-to-strong-runtime-of-one-recursive-call}
For every $i\in [\ell]$, all steps except for Steps \ref{alg:WeakToStrongCover:(10)},  \ref{alg:WeakToStrongCover:(13)}, and  \ref{alg:WeakToStrongCover:(14)} of \Cref{alg:WeakToStrongCover} in the $i^{th}$ recursive call can be implemented to run in 
$O(|V_i|^2(|V_i| + |E_i|))$ time
using 
$O(|V_i|^2)$ queries to \functionMaximizationOracle{p_i}. The run-time includes the time to construct the hypergraphs used as input to the queries to \functionMaximizationOracle{p_i}. Moreover, for each query to \functionMaximizationOracle{p_i}, the hypergraph $(G_0, c_0)$ used as input to the query has  $|V_i|$ vertices and at most $|E_i| + 1$ hyperedges.
\end{claim}
\begin{proof}
We first show how to compute hyperedges $e_i$ and $f_i$ (Step \ref{alg:WeakToStrongCover:(4)}). We recall that by \Cref{cor:weak-to-strong:uncrossing:exist-two-hyperedges-contained-in-disjoint-maximal-tight-sets}, there exist two hyperedges $e_i, f_i \in E_i$ such that $\maximalTightSetContainingHyperedge{p_i, H_i, w_i}(e_i) \not = \maximalTightSetContainingHyperedge{p_i, H_i, w_i}(f_i)$ and consequently Step \ref{alg:WeakToStrongCover:(4)} is well defined. By \Cref{lem:weak-to-strong:uncrossing:maximal-tight-set-family-disjoint}, the family $\maximalTightSetFamily{p_i, H_i, w_i}$ is a disjoint family. 
%Thus, $|\maximalTightSetFamily{p_i, H_i, w_i}| = O(|V_i|)$. 
Furthermore,  by \Cref{lem:FunctionMaximizetionOracles:computing-maximal-pHw-tight-sets-with-p-max-oracle}, we can compute the family $\maximalTightSetFamily{p_i, H_i, w_i}$ in $O(|V_i|^2(|V_i|+|E_i|))$ time using $O(|V_i|^2)$ queries to \functionMaximizationOracle{p_i}. 
Moreover, for each query to \functionMaximizationOracle{p_i}, the hypergraph $(G_0, c_0)$ used as input to the query has at most $|V_i|$ vertices and at most $|E_i|+1$ hyperedges. 
We note that $\maximalTightSetContainingHyperedge{p_i, H_i, w_i}(\Tilde{e}) \in \maximalTightSetFamily{p_i, H_i, w_i}$ for every $\Tilde{e} \in E$.
Consequently, we can find the required hyperedges for Step \ref{alg:WeakToStrongCover:(4)} by finding a pair of distinct sets $X, Y \in \maximalTightSetFamily{p_i, H_i, w_i}$ that each contains a hyperedge, and define these hyperedges to be $e_i$ and $f_i$ respectively. Thus, Step  \ref{alg:WeakToStrongCover:(4)} can be implemented to run in $O(|V_i|^2(|V_i|+|E_i|))$ time using $O(|V_i|^2)$ queries to \functionMaximizationOracle{p_i}. Moreover, for each query to \functionMaximizationOracle{p_i}, the hypergraph $(G_0, c_0)$ used as input to the query has at most $|V_i|$ vertices and at most $|E_i|+1$ hyperedges. 

Next, we show how to compute the values $\alpha_i^M$ (Step \ref{alg:WeakToStrongCover:(5)}) and $\alpha^R_i$ (Step \ref{alg:WeakToStrongCover:(7)}).  We first show how to compute the value $\beta_i^M$.
We have the following:
\begin{align*}
    \beta_i^M& = \min\left\{b_{(H_i,w_i)}(X) - p_i(X) : X \subseteq V_i \text{ and } e_i \cap X, f_i \cap X \not = \emptyset\right\} &\\
    & = \min_{u \in e_i, v \in f_i}\left\{\min\left\{b_{(H_i,w_i)}(X) - p_i(X) : \{u,v\} \subseteq X \subseteq V_i\right\}\right\}.&
    %&=: \min_{u \in e_i, v \in f_i}\left\{\beta_i^{uv}\right\},&
\end{align*}
For $u \in e_i$ and $v \in f_i$, let $\beta_i^{uv} := \min\{b_{(H_i,w_i)}(X) - p_i(X) : \{u,v\} \subseteq X \subseteq V_i\}.$ Thus, $\beta_i^M = \min_{u \in e_i, v\in f_i}\beta_i^{uv}$.
Then, $\beta_i^M$ (and consequently the value $\alpha_i^M$) can be computed in $O(|V_i|^2(|V_i| + |E_i|))$ time using 
$|V_i|^2$ queries to \functionMaximizationOracle{p_i}, where each query is on a input hypergraph $(G_0, c_0)$ that has at most $|V_i|$ vertices and at most $|E_i|+1$ hyperedges. This is because for every $u \in e_i$ and $v \in f_i$ (there are $O(|V_i|^2)$ choices for $u$ and $v$), the value $\beta_i^{uv}$ can be computed using a single oracle call to \functionMaximizationOracle{p_i} on a hypergraph $(G_0, c_0)$ that has at most $|V_i|$ vertices and at most $|E_i|+1$ hyperedges.
Next, the value $\beta^R_i$, (and consequently the value $\alpha^R_i$)  can be computed 
in $O(|V_i| + |E_i|)$ time using a single query to \functionMaximizationOracle{p_i} on the input hypergraph $(G_0=H_i^M, c_0=w_i^M)$ that has  $|V_i|$ vertices and at most $|E_i|+1$ hyperedges. Here, the hypergraph $(H_i^M, w_i^M)$ has  $|V_i|$ vertices and $|E_i|+1$ hyperedges because of  \Cref{lem:WeakToStrong:properties-of-alpha^M-and-alpha^R}(2).
Next, we show that the hypergraphs $(H_i^M, w_i^M)$ (Step \ref{alg:WeakToStrongCover:(6)}) and $(H_i^R, w_i^R)$ (Step \ref{alg:WeakToStrongCover:(8)}) can be computed in $O(|V_i| + |E_i|)$ time. We recall that the \merge operation decreases the weights of hyperedges $e_i$ and $f_i$ by $\alpha_i^M$,  increases the weight of the hyperedge $e_i\cup f_i$ by $\alpha_i^M$ and discards the resulting zero-weighted edges (if any). This can be implemented to run in $O(|E_i|)$ time. Furthermore, the \reduce operation decreases the weight of the hyperedge $e_i\cup f_i$ by $\alpha^R_i$, and so can also be implemented to run in $O(|E^R_i|) = O(|E_i|)$ time by \Cref{lem:WeakToStrong:properties-of-alpha^M-and-alpha^R}(2) and (5). Finally, the computations of the hypergraph $(H^{i}_0, w^{i}_0)$ (Step \ref{alg:WeakToStrongCover:(9)}), the set $\zeros_i$ (\Cref{alg:WeakToStrongCover:(11)}) and the set $V'_i$  (\Cref{alg:WeakToStrongCover:(12)}) can each be done in $O(|V_i|)$ time. Furthermore, the hypergraph returned in Step \ref{alg:WeakToStrongCover:(14)} contains at most $|V|$ vertices and at most $|E_i| + O(|V_i|)$ hyperedges. Consequently, Steps \ref{alg:WeakToStrongCover:(15)} and \ref{alg:WeakToStrongCover:(16)}  can be implemented to run in $O(|V_i| + |E_i|)$ time. 
\end{proof}

\Cref{claim:weak-to-strong-runtime-of-one-recursive-call} implies that each query to \functionMaximizationOracle{p_i} takes as input a hypergraph 
with at most $|V_i|$ vertices and $O(|E_i|)$ hyperedges.  
Next, we focus on the time to implement a query to \functionMaximizationOracle{p_i} on such hypergraphs. 

\begin{claim}\label{claim:WeakToStrong:p_i-oracle-using-p-oracle}
For every $i\in [\ell]$ and for every given hypergraph $(G = (V_i, E_G), c:E_G\rightarrow\Z_+)$ such that $|E_G| = O(|E_i|)$ AND disjoint subsets $S, T\subseteq V_i$, $\functionMaximizationOracle{p_i}((G, c), S, T, y)$ can be computed in $O(|V|(|V| + |E|))$ time using at most $|V_i|+1$ queries to \functionMaximizationOracleStrongCover{p}. The run-time includes the time to construct the hypergraphs used as input to the queries to \functionMaximizationOracleStrongCover{p}. Moreover, for each query to \functionMaximizationOracleStrongCover{p}, the hypergraph $(G_0, c_0)$ used as input to the query has $O(|V|)$ vertices and $O(|V|+|E|)$ hyperedges.
\end{claim}
\begin{proof}
Let the input to \functionMaximizationOracle{p_i} be a hypergraph $(G = (V_i, E_G), c:E_G\rightarrow\Z_+)$ such that $|E_G| = O(|E_i|)$ and disjoint subsets $S, T\subseteq V_i$. By \Cref{lem:Preliminaries:wc-oracle-from-sc-oracle}, we can implement $\functionMaximizationOracle{p_i}((G, c), S, T, y_0)$ in $O(|V_i|(|V_i| + |E_i|))$ time using at most $|V_i|+1$ queries to $\functionMaximizationOracleStrongCover{p_i}$, where each query to $\functionMaximizationOracleStrongCover{p_i}$ is on a input hypergraph with $|V_i|+1$ vertices and $O(|E_i|)$ hyperedges. Next, we show that the query $\functionMaximizationOracleStrongCover{p_i}((G, c), S, T)$ where $(G = (V_i, E_G), c:E_G\rightarrow\Z_+)$ is a hypergraph with $|E_G| = O(|E_i|)$ and $S, T\subseteq V_i$ are disjoint subsets, can be implemented in $O(|V| + |E|)$ time using one query to  $\functionMaximizationOracleStrongCover{p}$ on an input hypergraph $(G_0, c_0)$ that has $O(|V|)$ vertices and $O(|V|+|E|)$ hyperedges. 

We first recall some notation and define some quantities for convenience. For $i\in [\ell]$, we recall that $(H_0^i, w_0^i)$ is the hypergraph $(H_0, w_0)$, $p^R_i, p_i'$ are the functions $p^R, p'$ respectively, and $\zeros_i$ is the set $\zeros$ defined in the $i^{th}$ recursive call of the algorithm. We note that $p_i' =p_{i+1}$ for every $i\in [\ell-1]$. Furthermore, the hypergraph $(H_0^i, w_0^i)$ contains exactly one hyperedge for every $i\in [\ell-1]$. 
For $i\in [\ell-1]$, let $(H_1^i, w_1^i)$ be the hypergraph obtained from $(H_0^i, w_0^i)$ by adding the vertices $\cup_{j=1}^{i-1}\zeros_{j}$.

Let $i\in [\ell]$. Let the input to \functionMaximizationOracleStrongCover{p_i} be a hypergraph $(G = (V_i, E_G), c:E_G\rightarrow\Z_+)$ such that $|E_G| = O(|E_i|)$ and disjoint subsets $S_0, T_0\subseteq V_i$. 
We note that $p_i=p_{i-1}'$ and hence, it suffices to answer $\functionMaximizationOracleStrongCover{p_{i-1}'}((G, c), S_0, T_0)$. 
Construct the hypergraph $(G_1, c_1)$ to be the hypergraph obtained from $(G,c)$ by adding vertices $\cup_{j=1}^{i-1}\zeros_{j}$. 
Construct the hypergraph $(\Tilde{H}, \Tilde{w})=(\sum_{j=1}^{i-1}H_1^j, \sum_{j=1}^i w_1^j)$. The hypergraph $(\Tilde{H}, \Tilde{w})$ contains at most $i - 1 \leq \ell = O(|E| + |V|)$ hyperedges, and can be constructed in time $O(|V|+|E|)$. 
Next, construct the hypergraph $(G, c)=(G_1+\Tilde{H}, c_1+\Tilde{w})$. The hypergraph $(G, c)$ contains $O(|E|+|V|)$ hyperedges and can be constructed in time $O(|V|+|E|)$. 
%Let $y_1\in \R^{V}$ be the vector obtained by setting $y_1(u):=y_0(u)$ for every $u\in V_{i}$ and $y_1(u):=0$ for every $u\in V\setminus V_{i}$. We note that the vector $y_1$ can be constructed in $O(|V|)$ time.  
Obtain $(Z, p(Z))$ by querying $\functionMaximizationOracleStrongCover{p}((G, c), S_0, T_0)$ and return $(Z-\cup_{j=1}^{i-1}\zeros_{j}, p(Z))$. 
\end{proof}

By \Cref{claim:weak-to-strong-runtime-of-one-recursive-call}, each query to \functionMaximizationOracle{p_i} takes as input a hypergraph with $|V_i|\le |V|$ vertices and $O(|E_i|) = O(|E|+|V|)$ hyperedges. 
By \Cref{claim:WeakToStrong:p_i-oracle-using-p-oracle}, a query to \functionMaximizationOracle{p_i} on a input hypergraph with at most $|V_i|$ vertices and at most $O(|E_i|)$ hyperedges 
can be implemented in $O(|V|(|V| + |E|))$  time using $O(|V|)$ queries to \functionMaximizationOracleStrongCover{p}, where each query to \functionMaximizationOracleStrongCover{p} is on a input hypergraph $(G_0, c_0)$ that has $O(|V|)$ vertices and $O(|E|+|V|)$ hyperedges. 
Thus, all steps except for Steps \ref{alg:WeakToStrongCover:(10)}, \ref{alg:WeakToStrongCover:(13)}, and \ref{alg:WeakToStrongCover:(14)} of \Cref{alg:WeakToStrongCover} in the $i^{th}$ recursive call can be implemented to run in $O(|V|^3(|V| + |E|))$ time using $O(|V|^3)$ queries to \functionMaximizationOracleStrongCover{p}, where each query to \functionMaximizationOracleStrongCover{p} is on a input hypergraph $(G_0, c_0)$ that has $O(|V|)$ vertices and $O(|E|+|V|)$ hyperedges. 

We note that Steps \ref{alg:WeakToStrongCover:(10)} and \ref{alg:WeakToStrongCover:(13)} need not be implemented explicitly for the purposes of the algorithm. Instead, \functionMaximizationOracle{p_{i+1}} on input hypergraphs with $|V_{i+1}|\le |V|$ vertices and at most $|E_{i+1}|+1$ hyperedges can be used to execute all steps except for Steps \ref{alg:WeakToStrongCover:(10)}, \ref{alg:WeakToStrongCover:(13)}, and \ref{alg:WeakToStrongCover:(14)} in the $(i+1)^{th}$ recursive call. 
We have already seen in \Cref{claim:WeakToStrong:p_i-oracle-using-p-oracle} that a query to \functionMaximizationOracle{p_{i+1}} on input hypergraphs with at most $|V|$ vertices and $O(|E_{i+1}|) = O(|V| + |E|)$ hyperedges can be implemented in  $O(|V|(|V|+|E|))$ time using $O(|V|)$ queries to \functionMaximizationOracleStrongCover{p}. 
Step \ref{alg:WeakToStrongCover:(14)} is the recursive call. We note that the hypergraph returned in Step \ref{alg:WeakToStrongCover:(14)} contains at most $|V|$ vertices and at most $|E_i| + O(|V_i|)= O(|V| + |E|)$ hyperedges. 

Thus, we have shown that each recursive call can be implemented to run in time $O(|V|^3(|V| + |E|))$ using $O(|V|^3)$ queries to \functionMaximizationOracleStrongCover{p}, where each query to \functionMaximizationOracleStrongCover{p} is on an input hypergraph $(G_0, c_0)$ that has $O(|V|)$ vertices and $O(|E|+|V|)$ hyperedges. The number of recursive calls is $O(|V|+|E|)$. 
Hence, the algorithm can be implemented to run in $O(|V|^3(|V| + |E|)^2)$ time using $O(|V|^3(|V|+|E|))$ queries to \functionMaximizationOracleStrongCover{p}, where each query to \functionMaximizationOracleStrongCover{p} is on a input hypergraph $(G_0, c_0)$ that has $O(|V|)$ vertices and $O(|E|+|V|)$ hyperedges. 
\end{proof}
\section{Conclusion}\label{sec:conclusion}
We introduced a splitting-off operation in (weighted) hypergraphs. Our contribution on this front is conceptualizing an appropriate notion of splitting-off in hypergraphs. 
Next, we proved that for every hypergraph there exists a local-connectivity preserving complete h-splitting-off at an arbitrary vertex from the hypergraph.  
%with polynomial number of additional hyperedges 
%and gave a strongly polynomial-time algorithm to find such a complete splitting-off at a vertex from a given \emph{weighted} hypergraph. 
Although our proof of existence follows from previously known existential results of Bern\'{a}th and Kir\'{a}ly \cite{Bernath-Kiraly} for an abstract function cover problem, our main contribution is identifying that our newly introduced notion of hypergraph splitting-off falls within the framework of their abstract function cover problem. 
Next, we designed a strongly polynomial-time algorithm to find a local-connectivity preserving complete splitting-off at a vertex. This involved substantial technical challenges to overcome. In particular, our main contribution towards the strongly polynomial-time algorithm is a strengthening of the above-mentioned existential result: we showed that there exists a local connectivity preserving complete h-splitting-off at an arbitrary vertex from the hypergraph 
\emph{in which the number of additional hyperedges is polynomial in the number of vertices}. 
%\emph{with polynomial number of additional hyperedges}. 
The strongly polynomial-time algorithm follows from our techniques to achieve this stronger existential result via standard algorithmic tools in submodularity. 
Finally, we illustrated the usefulness of the existence of local connectivity preserving complete h-splitting-off at an arbitrary vertex from a hypergraph by presenting two applications. 
Our first application is to give a constructive characterization of $k$-hyperedge-connected hypergraphs. Our second application is to give an alternative proof of an approximate min-max relation for 
the max Steiner rooted-connected orientation problem in graphs and hypergraphs. 
Two notable features of our proof of the approximate min-max relation for \emph{graphs} are that (1) it avoids the strong orientation theorem of Nash-Williams and (2) it proves a result for graphs using tools developed for hypergraphs. %We believe that these features suggest the strong 
%Our proof of existence and our algorithm to compute a local connectivity preserving complete splitting-off at an arbitrary vertex from the hypergraph are based on an algorithmic proof for the more general problem of converting a weak cover of a symmetric skew-supermodular function to a strong cover of the same function by merging disjoint hyperedges. 

Local and global connectivity preserving complete splitting-off at a vertex from a graph is a powerful operation for graphs. It finds applications in a variety of graph problems including graph orientation \cite{Mad78, Frank93:Applications-of-submod-functions, FK02, Kiraly-Lau}, connectivity augmentation \cite{Fra92, Fra94, JFJ95, Frank-book, nagamochi_ibaraki_2008_book}, 
minimum cuts enumeration \cite{NNI97, HW96, Goe01}, network design \cite{Lov76, Jor03, CS08, CFLY11}, tree packing \cite{BHTP03, Kri03, JMS03, Lau07}, congruency-constrained cuts \cite{NZ20}, and approximation algorithms for TSP \cite{GB93, BN23}. 
%These far-reaching applications have led to sustained interests in designing fast algorithms for local connectivity preserving splitting-off in graphs \cite{LY13, CLP22-soda, CLP22-stoc, CHLP23}. 
%\cite{Mad78, Frank93:Applications-of-submod-functions, FK02, Kiraly-Lau, Fra92, Fra94, JFJ95, Frank-book, nagamochi_ibaraki_2008_book, NNI97, HW96, Goe01, Jor03, CFLY11, BHTP03, Lau07, NZ20, GB93, BN23}
We believe that our local connectivity preserving complete splitting-off results for hypergraphs is likely to find future applications akin to its counterpart in graphs. 
We mention some of the open questions raised by our work:
\begin{enumerate}
    \item Our work focused on local connectivity preserving complete \emph{h-splitting-off} at a vertex from a hypergraph. We gave an example showing that local/global connectivity preserving complete \emph{g-splitting-off} at a vertex from a hypergraph may not exist (Figure \ref{figure:hypergraph-splitting-off-example-2}). 
    Are there sufficient conditions to guarantee local/global connectivity preserving complete \emph{g-splitting-off} at a vertex from a hypergraph? We recall that Lov\'{a}sz's \cite{Lov74, Lovasz-problems-book} and Mader's \cite{Mad78} results give sufficient conditions to guarantee local and global connectivity preserving complete g-splitting-off at a vertex from a \emph{graph}. 
    \iffalse
    \knote{Include these two in journal version}
    \item The \emph{rank} of a hypergraph is the size of its largest hyperedge. Hypergraphs of rank $2$ are graphs. 
    We note that complete g-splitting-off at a vertex from a graph is necessarily a graph. However, for $r\ge 3$, complete g-splitting-off/h-splitting-off at a vertex from an $r$-rank hypergraph is not necessarily an $r$-rank hypergraph---our splitting-off disjoint hyperedges operations could increase the rank. Is there a notion of splitting-off that preserves rank? Are there sufficient conditions to guarantee splitting-off operations that preserve local/global connectivity as well as rank? Such a result may help in giving a constructive characterization of $r$-rank $k$-hyperedge-connected hypergraphs. 
    \item Mader \cite{Mad82, Mad83} introduced a notion of splitting-off in directed graphs and showed that there exists global/Steiner connectivity preserving splitting-off at a vertex from a directed graph under mild conditions. This result has far reaching implications---e.g., it can be used to prove Edmonds' arborescence packing theorem \cite{Frank-book} and also finds applications in approximation algorithms for certain directed orienteering problems \cite{NR11}. Is there a generalization/counterpart of Mader's splitting-off result for directed graphs to directed hypergraphs?
    \fi
    \item As one of the applications of our splitting-off result, we presented an alternative proof of an approximate min-max relation for the max Steiner \emph{rooted}-connected orientation problem in hypergraphs. The computational complexity of a closely related hypergraph orientation problem is open: In max Steiner connected orientation problem in hypergraphs, the input is a hypergraph $G=(V, E)$ and a subset $T$ of terminals. The goal is to find the maximum $k$ and an orientation $\overrightarrow{G}$ of $G$ such that $\overrightarrow{G}$ contains $k$ hyperarc-disjoint paths from $u$ to $v$ for every pair of distinct terminals $u, v\in T$. Max Steiner connected orientation problem in \emph{graphs} is solvable in polynomial time via the Nash-Williams' strong orientation theorem. 
    Is max Steiner connected orientation problem in \emph{hypergraphs} solvable in polynomial time? 
    %Is this problem solvable in polynomial time in hypergraphs?
    %Does this problem admit a constant factor approximation in hypergraphs? 
    %\item \shubhang{Splitting off in directed hypergraphs? Such a result would help in constructive characterizations of $k$-arc-connected directed hypergraphs and $(k, \ell)$-partition-connected hypergraphs}
\end{enumerate} 

\iffalse
\knote{More questions:
\begin{enumerate}
\item Faster run-time for hypergraph splitting-off?
\item Can we use hypergraph splitting off to design an alternate algorithm to enumerate all min-cut-sets in a hypergraph? Theorem \ref{thm:splitting-off-or-trim-to-preserve-connectivity} is perhaps useful for this purpose? 
\item The \textsc{Max Steiner Rooted-Connected Orientation} Problem is $2$-approximable. What is the inapproximability? Conjecture is that it should be $(2-\epsilon)$-inapproximable. Intuition: the $2$-factor gap between connectivity and weak-partition-connectivity.
\end{enumerate}
}

\shubhang{Assuming polynomially bounded hyperedge weights, the implementation of \functionMaximizationOracleStrongCover{p} can be improved to $\Tilde{O}(|V|)$ max flow computations by computing $|V|$ Gomory-Hu trees, each using $\Tilde{O}(1)$ max-flow computations via Abboud et. al.. Perhaps a single Gomory-Hu tree of a modified hypergraph suffices? If so, then we can compute the required Gomory-Hu tree (without needing the polynomially bounded hyperedge-weights assumption) using the classical $O(|V|)$ max-flows algorithm.}
\fi


\medskip
\paragraph{Acknowledgements.} 
Karthekeyan thanks Eklavya Sharma for engaging in preliminary discussions on hypergraph splitting-off. 
Karthekeyan and Shubhang were supported in part by NSF grants CCF-1814613 and CCF-1907937. 
Karthekeyan was supported in part by the Distinguished Guest Scientist Fellowship of the Hungarian Academy of Sciences -- grant number VK-6/1/2022. 
%This research has been implemented with the support provided 
Krist\'{o}f and Tam\'{a}s were supported in part 
by the Lend\"ulet Programme of the Hungarian Academy of Sciences -- grant number LP2021-1/2021, by the Ministry of Innovation and Technology of Hungary from the National Research, Development and Innovation Fund -- grant number ELTE TKP 2021-NKTA-62 funding scheme, and by the Dynasnet European Research Council Synergy project -- grant number ERC-2018-SYG 810115.

\bibliographystyle{abbrv}
\begin{thebibliography}{10}

\bibitem{egres-open-problem-on-head-disjoint-sco}
{Head-disjoint strongly connected orientations, {\url{http://lemon.cs.elte.hu/egres/open/Head-disjoint_strongly_connected_orientations}}}.

\bibitem{AGL23}
O.~Alrabiah, V.~Guruswami, and R.~Li.
\newblock {Randomly punctured Reed–Solomon codes achieve list-decoding capacity over linear-sized fields}.
\newblock Preprint on arXiv: 2304.09445, 2023.

\bibitem{JFJ95}
J.~Bang-Jensen, A.~Frank, and B.~Jackson.
\newblock Preserving and increasing local edge-connectivity in mixed graphs.
\newblock {\em SIAM Journal on Discrete Mathematics}, 8(2):155--178, 1995.

\bibitem{BJ99}
J.~Bang-Jenson and B.~Jackson.
\newblock Augmenting hypergraphs by edges of size two.
\newblock {\em Mathematical Programming}, 84:467--481, 1999.

\bibitem{BST19}
N.~Bansal, O.~Svensson, and L.~Trevisan.
\newblock New notions and constructions of sparsification for graphs and hypergraphs.
\newblock In {\em IEEE 60th Annual Symposium on Foundations of Computer Science}, FOCS, pages 910--928, 2019.

\bibitem{BCCKM-private}
K.~B{\'e}rczi, K.~Chandrasekaran, C.~Chekuri, T.~Kir{\'a}ly, and S.~Kulkarni.
\newblock Private communication.

\bibitem{Bernath-Kiraly}
A.~Bern\'{a}th and T.~Kir\'{a}ly.
\newblock Covering skew-supermodular functions by hypergraphs of minimum total size.
\newblock {\em Operations Research Letters}, 37(5):345--350, 2009.

\bibitem{BK12}
A.~Bern\'{a}th and T.~Kir\'{a}ly.
\newblock A unifying approach to splitting-off.
\newblock {\em Combinatorica}, 32:373--–401, 2012.

\bibitem{BHTP03}
A.~Bhalgat, R.~Hariharan, T.~Kavitha, and D.~Panigrahi.
\newblock {Fast Edge Splitting and Edmonds' Arborescence Construction for Unweighted Graphs}.
\newblock In {\em Proceedings of the 19th Annual ACM-SIAM Symposium on Discrete Algorithms}, SODA, pages 455--464, 2008.

\bibitem{BN23}
J.~Blauth and M.~N\"{a}gele.
\newblock {An improved approximation guarantee for Prize-Collecting TSP}.
\newblock Preprint on arXiv: 2212.03776, 2023.

\bibitem{CHLP23}
R.~Cen, W.~He, J.~Li, and D.~Panigrahi.
\newblock Steiner connectivity augmentation and splitting-off in poly-logarithmic maximum flows.
\newblock In {\em Proceedings of the 34th Annual ACM-SIAM Symposium on Discrete Algorithms}, SODA, pages 2449--2488, 2023.

\bibitem{CLP22-soda}
R.~Cen, J.~Li, and D.~Panigrahi.
\newblock Augmenting edge connectivity via isolating cuts.
\newblock In {\em Proceedings of the 33rd Annual ACM-SIAM Symposium on Discrete Algorithms}, SODA, pages 3237--3252, 2022.

\bibitem{CLP22-stoc}
R.~Cen, J.~Li, and D.~Panigrahi.
\newblock Edge connectivity augmentation in near-linear time.
\newblock In {\em Proceedings of the 54th Annual ACM SIGACT Symposium on Theory of Computing}, STOC, pages 137--150, 2022.

\bibitem{CFLY11}
Y.~H. Chan, W.~S. Fung, L.~C. Lau, and C.~K. Yung.
\newblock {Degree Bounded Network Design with Metric Costs}.
\newblock {\em SIAM Journal on Computing}, 40(4):953--980, 2011.

\bibitem{CC22}
K.~Chandrasekaran and C.~Chekuri.
\newblock Hypergraph $k$-cut for fixed $k$ in deterministic polynomial time.
\newblock {\em Mathematics of Operations Research}, 47(4), 2022.
\newblock Prelim. version in FOCS 2020.

\bibitem{CC23}
K.~Chandrasekaran and C.~Chekuri.
\newblock Min-max partitioning of hypergraphs and symmetric submodular functions.
\newblock {\em Combinatorica}, 43:455--477, 2023.
\newblock Prelim. version in SODA 2021.

\bibitem{CXY19-j}
K.~Chandrasekaran, C.~Xu, and X.~Yu.
\newblock {Hypergraph $k$-Cut in randomized polynomial time}.
\newblock {\em Mathematical Programming}, 186:85--113, 2021.
\newblock Prelim. version in SODA 2018.

\bibitem{CK14}
C.~Chekuri and N.~Korula.
\newblock {A Graph Reduction Step Preserving Element-Connectivity and Packing Steiner Trees and Forests}.
\newblock {\em SIAM Journal on Discrete Mathematics}, 28(2):577--597, 2014.

\bibitem{CS08}
C.~Chekuri and B.~Shepherd.
\newblock {Approximate Integer Decompositions for Undirected Network Design Problems}.
\newblock {\em SIAM J. Discrete Math.}, 23:163--177, 2008.

\bibitem{CX18}
C.~Chekuri and C.~Xu.
\newblock Minimum cuts and sparsification in hypergraphs.
\newblock {\em SIAM Journal on Computing}, 47(6):2118--2156, 2018.

\bibitem{CKN20}
Y.~Chen, S.~Khanna, and A.~Nagda.
\newblock Near-linear size hypergraph cut sparsifiers.
\newblock In {\em IEEE 61st Annual Symposium on Foundations of Computer Science}, FOCS, pages 61--72, 2020.

\bibitem{Cosh-thesis}
B.~Cosh.
\newblock {\em Vertex splitting and connectivity augmentation in hypergraphs}.
\newblock PhD thesis, {University of London}, 2000.

\bibitem{FPZ23}
K.~Fox, D.~Panigrahi, and F.~Zhang.
\newblock Minimum cut and minimum k-cut in hypergraphs via branching contractions.
\newblock {\em ACM Trans. Algorithms}, 19(2), 2023.
\newblock Prelim. version in SODA 2019.

\bibitem{Fra92}
A.~Frank.
\newblock Augmenting graphs to meet edge-connectivity requirements.
\newblock {\em SIAM Journal on Discrete Mathematics}, 5(1):25--53, 1992.

\bibitem{Frank93:Applications-of-submod-functions}
A.~Frank.
\newblock Applications of submodular functions.
\newblock {\em Surveys in Combinatorics, 1993 (Keele)}, pages 85--136, 1993.

\bibitem{Fra94}
A.~Frank.
\newblock Connectivity augmentation problems in network design.
\newblock {\em {Mathematical Programming: State of the Art 1994}}, pages 34--63, 1994.

\bibitem{Frank-book}
A.~Frank.
\newblock {\em Connections in combinatorial optimization}.
\newblock Oxford University Press, Oxford, 2011.

\bibitem{FK02}
A.~Frank and Z.~Kir\'{a}ly.
\newblock Graph orientations with edge-connection and parity constraints.
\newblock {\em Combinatorica}, 22:47--70, 2002.

\bibitem{Gab94}
H.~N. Gabow.
\newblock Efficient splitting off algorithms for graphs.
\newblock In {\em Proceedings of the 26th Annual ACM Symposium on Theory of Computing}, STOC, pages 696--–705, 1994.

\bibitem{GKP17}
M.~Ghaffari, D.~Karger, and D.~Panigrahi.
\newblock {Random Contractions and Sampling for Hypergraph and Hedge Connectivity}.
\newblock In {\em Proceedings of the 28th annual ACM-SIAM Symposium on Discrete Algorithms}, SODA, pages 1101--1114, 2017.

\bibitem{GB93}
M.~Goemans and D.~Bertsimas.
\newblock Survivable networks, linear programming relaxations and the parsimonious property.
\newblock {\em Mathematical Programming}, 60:145--166, 1993.

\bibitem{Goe01}
M.~X. Goemans.
\newblock {Approximate Edge Splitting}.
\newblock {\em SIAM Journal on Discrete Mathematics}, 14(1):138--141, 2001.

\bibitem{GLSTW22}
Z.~Guo, R.~Li, C.~Shangguan, I.~Tamo, and M.~Wootters.
\newblock {Improved List-Decodability and List-Recoverability of Reed-Solomon Codes via Tree Packings}.
\newblock In {\em IEEE 62nd Annual Symposium on Foundations of Computer Science}, FOCS, pages 708--719, 2022.

\bibitem{HW96}
M.~Henzinger and D.~Williamson.
\newblock On the number of small cuts in a graph.
\newblock {\em Information Processing Letters}, 59:41--44, 1996.

\bibitem{JMS03}
K.~Jain, M.~Mahdian, and M.~R. Salavatipour.
\newblock {Packing Steiner Trees}.
\newblock In {\em Proceedings of the 14th Annual ACM-SIAM Symposium on Discrete Algorithms}, SODA, pages 266--274, 2003.

\bibitem{JLS22}
A.~Jambulapati, Y.~P. Liu, and A.~Sidford.
\newblock Chaining, group leverage score overestimates, and fast spectral hypergraph sparsification.
\newblock In {\em Proceedings of the 55th Annual ACM Symposium on Theory of Computing}, STOC, pages 196--206, 2023.

\bibitem{Jor03}
T.~Jord\'{a}n.
\newblock {On minimally $k$-edge-connected graphs and shortest $k$-edge-connected Steiner networks}.
\newblock {\em Discrete Applied Mathematics}, 131(2):421--432, 2003.

\bibitem{KKTY21-STOC}
M.~Kapralov, R.~Krauthgamer, J.~Tardos, and Y.~Yoshida.
\newblock Towards tight bounds for spectral sparsification of hypergraphs.
\newblock In {\em Proceedings of the 53rd Annual ACM SIGACT Symposium on Theory of Computing}, STOC, pages 598--611, 2021.

\bibitem{KKTY21-FOCS}
M.~Kapralov, R.~Krauthgamer, J.~Tardos, and Y.~Yoshida.
\newblock {Spectral Hypergraph Sparsifiers of Nearly Linear Size}.
\newblock In {\em IEEE 62nd Annual Symposium on Foundations of Computer Science}, FOCS, pages 1159--1170, 2022.

\bibitem{Kiraly-Lau}
T.~Király and L.~C. Lau.
\newblock {Approximate min–max theorems for Steiner rooted-orientations of graphs and hypergraphs}.
\newblock {\em Journal of Combinatorial Theory, Series B}, 98:1233--1252, 11 2008.
\newblock Preliminary version in FOCS 2006.

\bibitem{KK15}
D.~Kogan and R.~Krauthgamer.
\newblock Sketching cuts in graphs and hypergraphs.
\newblock In {\em Proceedings of the 2015 Conference on Innovations in Theoretical Computer Science}, ITCS, pages 367--376, 2015.

\bibitem{Kri03}
M.~Kriesell.
\newblock {Edge-disjoint trees containing some given vertices in a graph}.
\newblock {\em J. Comb. Theory Ser. B}, 88:53--63, 2003.

\bibitem{Lau07}
L.~C. Lau.
\newblock {An Approximate Max-Steiner-Tree-Packing Min-Steiner-Cut Theorem}.
\newblock {\em Combinatorica}, 27:71--90, 2007.

\bibitem{LY13}
L.~C. Lau and C.~K. Yung.
\newblock {Efficient Edge Splitting-Off Algorithms Maintaining All-Pairs Edge-Connectivities}.
\newblock {\em SIAM Journal on Computing}, 42(3):1185--1200, 2013.

\bibitem{Lee22}
J.~R. Lee.
\newblock Spectral hypergraph sparsification via chaining.
\newblock In {\em Proceedings of the 55th Annual ACM Symposium on Theory of Computing}, STOC, pages 207--218, 2023.

\bibitem{LiM17}
P.~Li and O.~Milenkovic.
\newblock Inhomogeneous hypergraph clustering with applications.
\newblock {\em Advances in neural information processing systems}, 30, 2017.

\bibitem{Lov74}
L.~Lov{\'a}sz.
\newblock Lecture.
\newblock {\em Presented in a Conference on Graph Theory, Prague}, 1974.

\bibitem{Lov76}
L.~Lov\'{a}sz.
\newblock {On some connectivity properties of Eulerian graphs}.
\newblock {\em Acta Math. Hungarica}, 28:129--138, 1976.

\bibitem{Lovasz-problems-book}
L.~Lov{\'a}sz.
\newblock {\em {Combinatorial Problems and Exercises, Second Edition}}.
\newblock American Mathematical Soc., 1993.
\newblock First Edition: North-Holland, Amsterdam, 1979.

\bibitem{Mad78}
W.~Mader.
\newblock A reduction method for edge-connectivity in graphs.
\newblock In {\em Annals of Discrete Mathematics}, volume~3, pages 145--164. Elsevier, 1978.

\bibitem{nagamochi_ibaraki_2008_book}
H.~Nagamochi and T.~Ibaraki.
\newblock {\em Algorithmic Aspects of Graph Connectivity}.
\newblock Encyclopedia of Mathematics and its Applications. Cambridge University Press, 2008.

\bibitem{NNI97}
H.~Nagamochi, K.~Nishimura, and T.~Ibaraki.
\newblock Computing all small cuts in an undirected network.
\newblock {\em SIAM Journal on Discrete Mathematics}, 10(3):469--481, 1997.

\bibitem{NZ20}
M.~N\"{a}gele and R.~Zenklusen.
\newblock A new contraction technique with applications to congruency-constrained cuts.
\newblock {\em Mathematical Programming}, 183(2):455--481, 2020.
\newblock Prelim. version in IPCO 2019.

\bibitem{NW60}
C.~S. J.~A. Nash-Williams.
\newblock On orientations, connectivity and odd vertex pairings in finite graphs.
\newblock {\em Canad. J. Math}, 12:555--567, 1960.

\bibitem{NW61}
C.~S. J.~A. Nash-Williams.
\newblock {Edge disjoint spanning trees of finite graphs}.
\newblock {\em J. London Math. Soc.}, 36:445--450, 1961.

\bibitem{Ornes21}
S.~Ornes.
\newblock How big data carried graph theory into new dimensions.
\newblock {\em Quanta Magazine}, 2021.
\newblock \url{https://www.quantamagazine.org/how-big-data-carried-graph-theory-into-new-dimensions-20210819/}.

\bibitem{Quanrud22}
K.~Quanrud.
\newblock Quotient sparsification for submodular functions.
\newblock Manuscript available at \url{kentquanrud.com}, November 2022.

\bibitem{Rob39}
H.~E. Robbins.
\newblock {A Theorem on Graphs with an Application to a Problem of Traffic Control}.
\newblock {\em Amer. Math. Monthly}, 46:281--283, 1939.

\bibitem{Schlagetal23}
S.~Schlag, T.~Heuer, L.~Gottesb{\"u}ren, Y.~Akhremtsev, C.~Schulz, and P.~Sanders.
\newblock High-quality hypergraph partitioning.
\newblock {\em ACM Journal of Experimental Algorithmics}, 27:1--39, 2023.

\bibitem{SY19}
T.~Soma and Y.~Yoshida.
\newblock Spectral sparsification of hypergraphs.
\newblock In {\em Proceedings of the 30th Annual ACM-SIAM Symposium on Discrete Algorithms}, SODA, pages 2570--2581, 2019.

\bibitem{Tu61}
W.~T. Tutte.
\newblock {On the problem of decomposing a graph into $n$ connected factors}.
\newblock {\em J. London Math. Soc.}, 36:221--230, 1961.

\bibitem{VeldtBK22}
N.~Veldt, A.~R. Benson, and J.~Kleinberg.
\newblock Hypergraph cuts with general splitting functions.
\newblock {\em SIAM Review}, 64(3):650--685, 2022.

\end{thebibliography}

\appendix
\section{An Example for Weak To Strong Cover with Quadratic New Hyperedges}\label{sec:appendix:tight-example-weak-to-strong}
%\shubhang{ Question: Is there a ``tight'' example for property (3) \Cref{thm:WeakToStrongCover:main}? By tight, I mean an instance $(H, w)$ and $p$ such that $(H, w)$ weakly covers $p$, but every hypergraph $(H^* = (V, E^*), w^*)$ that strongly covers $p$ has $|E^*| = |E| + \Omega(|V|)$?}

\begin{lemma}
There exists a symmetric skew-supermodular function  $p:2^V\rightarrow\Z$ and a hypergraph\linebreak $\left(H = (V, E), w:E \rightarrow\Z_+\right)$ satisfying $b_{(H, w)}(X) \geq p(X)$ for every $X \subseteq V$ such that for every hypergraph $\left(H^* = (V, E^*), w^*:E^* \rightarrow\Z_+\right)$ obtained by merging hyperedges of the hypergraph $(H, w)$ such that $d_{(H^*, w^*)}(X) \geq p(X)$ for every $X\subseteq V$, we have that 
\[
|E^*-E|=\Omega(|V|^2). 
\]
\end{lemma}
\begin{proof}
Let $u\in V$ be a fixed vertex. Let $p:2^V\rightarrow\Z\cup\{-\infty\}$ be a symmetric skew-supermodular function defined as follows,
$$p(X) \coloneqq  \begin{cases}
    {\binom{|V| - 1}{2}}& \text{ if } X \in \left\{\{u\}, V - u\right\},\\
    - \infty &\text{otherwise.}
\end{cases}$$
Furthermore, let $(H = (V, E), w:E\rightarrow\Z_+)$ be a hypergraph such that $E \coloneqq  {\binom{V - u}{2}}\cup \{\{u\}\}$, and the weight function $w$ is defined as
$$w(e) \coloneqq  \begin{cases}
    1& \text{ if } e \in {\binom{V - u}{2}}\\
    {\binom{|V| - 1}{2}}& \text{otherwise (i.e., if $e = \{u\}$)}.
\end{cases}$$

We note that $b_{H, w}(X) \geq p(X)$ for every $X \subseteq V$. Next, consider the hypergraph $(H^* = (V, E^*), w^*)$, where the hyperedge set $E^* \coloneqq  \left\{\{u, a, b\} : ab \in {\binom{V-u}{2}}\right\}$, and the weight function $w^*(e) \coloneqq  1$ for every $e \in E^*$. Then, the hypergraph $(H^*, w^*)$ strongly covers the function $p$ because $d_{H^*, w^*}(X) \geq p(X)$ for every $X \subseteq V$. Furthermore, the hypergraph $(H^*, w^*)$ is the \emph{unique} hypergraph that can be obtained by merging hyperedges of $(H, w)$ to strongly cover the function $p$. We observe that $|E^*-E| = \Omega(|V|^2)$. 
\end{proof}
\section{Projection of Laminar Family}\label{appendix:sec:uncrossing-properties:projection-of-laminar-families}
We recall that 
%two sets $X, Y \subseteq V$ are said to cross if the sets $X - Y, Y - X, X \cap Y$ are non-empty. A
a family $\calL \subseteq 2^V$ is \emph{laminar} if for every pair of sets $X, Y\calL$, we have that either $X\subseteq Y$ or $Y\subseteq X$ or $X\cap Y=\emptyset$  --  we use \emph{family} to refer to a set (and not a multiset) of subsets of the ground set $V$. 
%\kristof{I would leave out this sentence}\shubhang{would it be clear from context? I do not want multiplicities in the family, but I have never seen ``laminar set'' being used.} \knote{Explicitly say that a family is a set and not a multiset.} 
For a set $\calZ\subseteq V$, and a family $\calL \subseteq 2^V$,
the \emph{projection} of the family $\calL$ to the ground set $V - \calZ$ is the family formed by restricting each set of $\calL$ to the ground set $V - \calZ$, i.e. $\{X - \calZ : X \in \calL\}$.
The next lemma says that if the family $\calL$ is laminar, then the size of the projection %$\calL'$ 
is comparable to that of the original family $\calL$.


\projectionLaminarFamily*
\begin{proof}
    We first show that the family $\calL'$ is a laminar family. By way of contradiction, let $X,Y \in \calL'$ be two distinct sets such that $X-Y, Y-X, X\cap Y \not = \emptyset$. Then, there exist distinct sets $A, B \in \calL$ such that $A - \calZ = X$ and $B - \calZ = Y$. Since the family $\calL$ is laminar, we have that either $A\cap B = \emptyset$, $A \subseteq B$ or $B\subseteq A$. We note that if $A\subseteq B$, then we have that $X = A - \calZ \subseteq B - \calZ = Y$, which contradicts our choice of sets $X$ and $Y$. Similarly, if $B\subseteq A$, then we obtain that $Y \subseteq X$, contradicting our choice of sets $X$ and $Y$. Thus, we have that $A\cap B = \emptyset$, and so $X\cap Y = (A - \calZ)\cap(B - \calZ) = \emptyset$, once again contradicting our choice of sets $X$ and $Y$. 

    Next, we show the second part of the claim, i.e. $|\calL| \leq |\calL'| + 3|\calZ|$. For this, we first  partition the family $\calL$ into two subfamilies $\calL_1 : = \{A \in \calL : A \subseteq \calZ\}$ and $\calL_2 \coloneqq  \calL - \calL_1$. Furthermore, for every $X \in \calL'$, we let $S_X \coloneqq  \{A \in \calL_2 : A - \calZ = X\}$ denote the family of sets in $\calL_2$ that map to the (non-empty) set $X$ under projection with the set $\calZ$. 
    We note that for every $A \in \calL_2$, we have that $\emptyset \not = A - \calZ \in \calL'$ and hence $A \in S_{A - \calZ}$. Consequently, the subfamilies $S_X$ for $X \in \calL'$ partition the family $\calL_2$ 
    %\knote{Why is $S_X\cap S_Y=\emptyset$ for $X, Y\in \calL'$? Also, explain why $S_X\subseteq \calL_2$ for every $X\in \calL'$.}
    . In particular, we have that $\calL = \calL_1 \uplus \left(\biguplus_{X \in \calL'} S_X \right)$. \Cref{claim:Laminar-family-projection:S_X-is-a-chain} below shows that the family $S_X$ is a chain family for every $X \in \calL'$. Next, for every $X \in \calL'$, we define three sets: we let $A_X \coloneqq  \arg\max\{|A| : A \in S_X\}$, $B_{X} \coloneqq  \arg\min\{|B| : B \in S_X\}$, and $\calZ_X\coloneqq A_X-B_X$. 
    %denote the largest and smallest sets in the family $S_X$ respectively. 
    We note that $X \subseteq B_X \subseteq A_X$, where the second containment is because the family $S_X$ is a chain family. 
    %We also let $\calZ_X \coloneqq  A_X - B_X$. 
    Since $A_X - B_X \subseteq A_X - X \subseteq \calZ$, we have that $\calZ_X \subseteq \calZ$. \Cref{claim:Laminar-family-projection:|S_X|-atmost-|Z_X|+1} below shows that $|S_X| \leq |\calZ_X| + 1$. Furthermore, \Cref{claim:Laminar-family-projection:Z_X-Z_Y-disjoint} below shows that $\calZ_X$ and $\calZ_Y$ are disjoint for distinct $X,Y \in \calL'$. Then, we have the following:
\begin{align*}
        |\calL| & = |\calL_1| + \sum_{X \in \calL'} |S_X|&\\
        & \leq 2|\calZ| - 1 + \sum_{X \in \calL'}\left(|\calZ_X| + 1\right)&\\
        &\le 3|\calZ| + |\calL'|-1.&
    \end{align*}
    Here, the first equality is because $\calL = \calL_1 \uplus \left(\biguplus_{X \in \calL'} S_X \right)$. The first inequality is because the family $\calL_1$ is a laminar family on the ground set $\calZ$, and by \Cref{claim:Laminar-family-projection:|S_X|-atmost-|Z_X|+1}. The final inequality is because the sets $\calZ_X$ for $X \in\calL'$ form a subpartition of the set $\calZ$ by \Cref{claim:Laminar-family-projection:Z_X-Z_Y-disjoint}.
\end{proof}

\begin{claim}\label{claim:Laminar-family-projection:S_X-is-a-chain}
        For every $X \in\calL'$, the family $S_X$ is a chain family.
    \end{claim}
    \begin{proof}
        By way of contradiction, let $X \in \calL'$ be such that the family $S_X$ is not a chain. Then, there exist sets $A, B \in S_X$ such that $A - B, B - A \not = \emptyset$. Since $A$ and $B$ are sets in the laminar family $\calL$, we have that the sets $A$ and $B$ are disjoint. However, both the sets $A$ and $B$ contain the set $X$, and so $A\cap B \not = \emptyset$, a contradiction to the fact that $A$ and $B$ are sets in the laminar family $\calL$.
    \end{proof}

\begin{claim}\label{claim:Laminar-family-projection:|S_X|-atmost-|Z_X|+1}
    For every $X \in \calL'$, we have that $|S_X| \leq |\calZ_{X}| + 1$.
\end{claim}
\begin{proof}
    We may assume that $|S_X| \geq 2$, i.e.,  $A_X \not = B_X$, as otherwise the claim holds. Then, we have the following:
    $|S_X| \leq |A_X| - |B_X| + 1 \leq |\calZ_X| +1,$ where the first inequality is by \Cref{claim:Laminar-family-projection:S_X-is-a-chain}.
\end{proof}

\begin{claim}\label{claim:Laminar-family-projection:Z_X-Z_Y-disjoint}
    For distinct non-empty sets $X, Y \in \calL'$, we have that $\calZ_{X}\cap\calZ_{Y} = \emptyset$.
\end{claim}
\begin{proof}
     Let $X, Y \in\calL'$ be distinct sets. Since the family $\calL'$ is laminar, we have that either $X \subsetneq Y$ or $Y\subsetneq X$ or $X\cap Y = \emptyset$.

    First, suppose that $X \subsetneq Y$ (we note that the case of $Y\subsetneq X$ can be show by a similar argument). Here, we note that $X \subseteq A_X \cap B_Y $ since $X \subseteq A_X$ and $X \subseteq Y \subseteq B_Y$. Consequently, $A_X\cap B_Y\neq \emptyset$. 
    Since the family $\calL$ is laminar and $A_X, B_Y \in \calL$, we have that either $A_X \subseteq B_Y$ or $B_Y \subseteq A_X$. We note that if $A_X \subseteq B_Y$, then the claim holds since $\calZ_X \subseteq A_X \subseteq B_{Y} \subseteq V - \calZ_Y$.
    We now show that $B_Y - A_X \not = \emptyset$, i.e. the scenario $B_Y \subseteq A_X$ cannot happen. 
    Consider a vertex $v \in Y - X$. We note that 
    $v \not \in \calZ$
    since $v \in Y$. Furthermore, $v \not \in A_X$ since $v \not \in X$ and $A_X - X \subseteq \calZ$.
    Thus, we have that $v \in B_Y - A_X \not = \emptyset$.

    Next, suppose that the sets $X$ and $Y$ are disjoint. We also suppose that $A_X \cap A_Y \not = \emptyset$ as otherwise, we have that $\calZ_X \cap \calZ_Y \subseteq A_X \cap A_Y = \emptyset$ and the claim holds. Since the family $\calL$ is laminar and $A_X, A_Y \in \calL$, either $A_X \subseteq A_Y$ or $A_Y \subseteq A_X$. Suppose that  $A_X \subseteq A_Y$ (we note that the case of $A_Y \subseteq A_X$ can be show by a similar argument).
    We also suppose that $A_X - B_Y \not = \emptyset$, as otherwise $\calZ_X = A_X - B_X \subseteq B_Y \subseteq V - (A_Y - B_Y) = V - \calZ_Y$, and the claim holds. 
    %By way of contradiction, suppose that $\calZ_X \cap \calZ_Y \not = \emptyset$.
    The rest of the proof is devoted to showing that this case is impossible. In particular, we show that
    the sets $A_X$ and $B_Y$ are such that $A_X-B_Y, B_Y-A_X, A_X\cap B_Y\neq \emptyset$, contradicting the laminarity of the family $\calL$. First, we show that $B_Y - A_X \not = \emptyset$. By way of contradiction, suppose that $B_Y \subseteq A_X$. Then, we have that $Y \subseteq A_X$ since $Y \subseteq B_Y$. Furthermore,  we recall that $Y\cap \calZ = \emptyset$. Thus, we have that $\emptyset = Y \cap X = Y \cap (A_X - \calZ) = Y\cap A_X$. Thus, $Y\subseteq A_X$ and $Y\cap A_X=\emptyset$ together imply that $Y = \emptyset$, contradicting $Y \in \calL'$. Second, we show that $B_Y\cap A_X \not = \emptyset$. By way of contradiction, suppose that $B_Y\cap A_X = \emptyset$. We recall that $A_X\subseteq A_Y$. Consequently, we have that $X \subseteq A_X \subseteq A_Y - B_Y \subseteq \calZ$, contradicting $X \in \calL'$. Thus the sets $A_X$ and $B_Y$ are such that $A_X-B_Y, B_Y-A_X, A_X\cap B_Y\neq \emptyset$, giving us the required contradiction.
\end{proof}
\section{Optimization Problems using Function Maximization Oracle}\label{appendix:sec:Function-Maximization-Oracles}
In this section, we prove Lemmas 
\ref{lem:Preliminaries:wc-oracle-from-sc-oracle}, \ref{lemma:helper-for-p-max-oracle}, and \ref{lem:FunctionMaximizetionOracles:computing-maximal-pHw-tight-sets-with-p-max-oracle}. The proofs of these lemmas follow from known tools in the area of submodular functions. We present their proofs for the sake of completeness. 



\subsection{\functionMaximizationOracle{p} using \functionMaximizationOracleStrongCover{p}}\label{sec:wc-oracle-from-sc-oracle}
In this section, we prove \Cref{lem:Preliminaries:wc-oracle-from-sc-oracle}. 
\lemmaWCoraclefromSCoracle*
\begin{proof}
Let $u \in V$ be a vertex. We define the hypergraph $G^u := (V, E^u := \{e \cup \{u\} : e \in E\})$ and the weight function $w^u : E^u \rightarrow\Z_+$ as $w^u(e \cup \{u\}) := w(e)$ for every $e \in E$. We observe that $b_{(G, w)}(X) = d_{(G^u, w^u)}(X)$ for every set $X \subseteq V - \{u\}$. We also let $(Z^u, p(Z^u))$ denote the tuple obtained by querying \functionMaximizationOracleStrongCover{p}$((G_0, c_0) := (G^u, w^u), S_0 := S, T_0 := T\cup \{u\})$.

We now describe a procedure to compute an answer $(Z, p(Z))$ to the query $\functionMaximizationOracle{p}((G,w), S, T)$. We consider two cases to define our set $Z$. If $T \not= \emptyset$, then we choose an arbitrary vertex $t \in T$ and define $Z := Z^t$. Otherwise, $T = \emptyset$ and we define the set $Z$ to be an arbitrary optimum solution to $\max\{p(Z) - b_{(G, w)}(Z) : Z \in \calZ \}$, where $\calZ := \{Z^u: u \in V\} \cup \{ V\}$.  We return the tuple $(Z, p(Z))$ as the answer to the query $\functionMaximizationOracle{p}((G, w), S, T)$.

We now prove correctness of the above procedure. Let $Z$ be the set returned by the above procedure. Then, by definition $S\subseteq V\subseteq V-T$. Our goal is to show that $Z$ is an optimum solution to the following problem:
\begin{align}
    \max\left\{p(X)-b_{(G, w)}(X): S\subseteq X\subseteq V-T\right\} \label{eq:wc-oracle-from-sc-oracle}
\end{align}
Let $Y$ be an optimum solution to \eqref{eq:wc-oracle-from-sc-oracle}. It suffices to show that $p(Z)-b_{(G, w)}(Z) \ge p(Y)-b_{(G, w)}(Y) $.   
We consider two cases: suppose that $T\neq \emptyset$. Then, we have that 
\begin{align*}
p(Z)-b_{(G, w)}(Z)
&=p(Z^t)-d_{(G^t, w^t)}(Z^t)\\
&\ge p(Y)-d_{(G^t, w^t)}(Y) \\
&= p(Y) - b_{(G, w)}(Y) .
\end{align*} 
In the above, the first equality is because $Z=Z^t$, the inequality is because $Y\subseteq V-\{v\}$ and by definition of $Z^t$, and the last equality is because $b_{(G, w)}(Y) = d_{(G^v, w^v)}(Y)$ since $Y\subseteq V-\{t\}$. 
Next, suppose that $T=\emptyset$. We consider two subcases. Firstly, suppose that $V$ is the unique optimum solution to \eqref{eq:wc-oracle-from-sc-oracle}. Then, $Z=Y=V$ by definition and we are done. 
Secondly, suppose that $Y\neq V$. Let $v\in V-Y$.  
Then, we have the following:
\begin{align*}
p(Z)-b_{(G, w)}(Z)
%&=p(Z^u)-b_{(G, w)}(Z^u)\\
&\ge p(Z^v) - b_{(G, w)}(Z^v) \\
&=p(Z^v)-d_{(G^v, w^v)}(Z^v) \\
&\ge p(Y)-d_{(G^v, w^v)}(Y) \\
&=p(Y)-b_{(G, w)}(Y).
\end{align*} 
In the above sequence, the first inequality is because of the choice of the set $Z$ (i.e., $Z$ is an optimum solution to $\max\{p(Z)-b_{(G, w)}(Z): Z\in \calZ\}$ and $Z^v \in \calZ$), the first equality is because  $b_{(G, w)}(Z^v) = d_{(G^v, w^v)}(Z^v)$ since $Z^v\subseteq V-\{v\}$, the second inequality is by definition of $Z^v$, and the second equality is because $b_{(G, w)}(Y) = d_{(G^v, w^v)}(Y)$ since $Y\subseteq V-\{v\}$. Thus, all inequalities in the above sequence are equalities, and we have that $p(Z)-b_{(G, w)}(Z) = p(Y)-b_{(G, w)}(Y)$. Thus, in both cases, the set $Z$ returned by our procedure is an optimum solution to \eqref{eq:wc-oracle-from-sc-oracle}.

We now analyze the runtime of the procedure and the number of \functionMaximizationOracleStrongCover{p} queries. Let $u \in V$ be a vertex. Constructing the hypergraph $(G^{u}, w^{u})$ takes $O(|V| + |E|)$. Thus, a tuple $(Z^{u}, p(Z^u))$ takes $O(|V|+|E|)$ time and one query to \functionMaximizationOracleStrongCover{p}. The runtime includes the time to construct the inputs to the query. Moreover, the input hypergraph to the query has $|V|$ vertices and at most $|E|$ hyperedges. Thus, if $T \not = \emptyset$, the procedure makes a single \functionMaximizationOracleStrongCover{p} query and runs in $O(|V| + |E|)$ time. Alternatively, if  $T  = \emptyset$, then the procedure makes $|V| + 1$ \functionMaximizationOracleStrongCover{p} queries and runs in $O(|V|(|V| + |E|))$ time. In both cases, the runtime includes the time to construct the inputs for the \functionMaximizationOracleStrongCover{p} queries. Moreover, the input hypergraph to each query has $|V|$ vertices and at most $|E|$ hyperedges.
\end{proof}

\subsection{Skew-supermodularity of a function \(p\)}\label{sec:helper-for-p-max-oracle}
In this section, we prove \Cref{lemma:helper-for-p-max-oracle}. 
\lempmaxoracle*
\begin{proof}
    For ease of notation, let $p=p_{(G, c, r)}$. 
    The function $R:2^V\rightarrow \Z_{\ge 0}$ is symmetric skew-supermodular (see Lemma 8.1.9 of \cite{Frank-book}) and the function $d_{(G,c)}:2^V\rightarrow \Z_{\ge 0}$ is symmetric submodular. The function $p$ is obtained by subtracting a symmetric submodular function from a symmetric skew-supermodular and is hence, symmetric skew-supermodular. 

    Next, we show that \functionMaximizationOracleStrongCover{p} can be implemented in strongly polynomial-time. 
    Let the input to the oracle be the hypergraph $(G_0=(V, E_0), c_0: E_0\rightarrow \Z_+)$ and disjoint subsets $S_0, T_0\subseteq V$. 
    We construct the hypergraph $(G+G_0, c+c_0)$. Moreover, for distinct vertices $u, v\in V$ with $u\in V-T_0$ and $v\in V-S_0$, we compute  
    \begin{align*}
    \lambda_{(G+G_0, c+c_0)}(u, v)&\coloneqq \min\left\{d_{(G+G_0, c+c_0)}(Z): S_0\cup \{u\}\subseteq Z\subseteq V-(T_0\cup\{v\})\right\} 
    \end{align*}
    and the set $Z(u, v)$ that achieves the minimum for $\lambda_{(G+G_0, c+c_0)}(u, v)$. We note that $\lambda_{(G+G_0, c+c_0)}(u, v)$ and $Z(u,v)$ can be computed using a single call to a min $(s,t)$-cut oracle on the hypergraph $(G+G_0, c+c_0)$. This hypergraph has $|V|$ vertices and at most $|E_0|+|E|$ hyperedges.  

    To implement the \functionMaximizationOracleStrongCover{p}, we observe that 
    \begin{align*}
    \max&\left\{p(Z)-d_{(G_0, c_0)}(Z): S_0\subseteq Z\subseteq V-T_0\right\}\\
    &=\max\left\{R(Z)-d_{(G, c)}(Z)-d_{(G_0, c_0)}(Z): S_0\subseteq Z\subseteq V-T_0\right\}\\
    &=\max\left\{r(u, v)-\lambda_{((G, c), (G_0, c_0))}(u, v): u\in V-T_0, v\in V-S_0, u\neq v\right\}. 
    \end{align*}
    The RHS problem can be solved by iterating over all distinct vertices $u\in V-T_0$ and $v\in V-S_0$ and returning a pair for which the objective function $r(u, v)-\lambda_{((G, c), (G_0, c_0))}(u, v)$ is maximum. Let $u\in V-T_0$ and $v\in V-S_0$ be a pair of vertices that achieves the maximum in the RHS problem. Then, we observe that the set $Z(u, v)$ is a maximizer for the LHS problem. Thus, it suffices to return $Z(u, v)$ and the value $p(Z(u, v))$. We observe that the value $p(Z(u, v))$ can be computed by subtracting $d_{(G_0, c_0)}(Z(u,v))$ from the objective value of the RHS problem.

    The above procedure can be implemented in $O(|V|^2(|V| + |E| + |E_0|))$ time and $O(|V|^2)$ calls to a min $(s,t)$-cut oracle on hypergraphs with $|V|$ vertices and $|E| + |E_0|$ hyperedges. Then, the claimed runtime is obtained by using a crude upper bound on the run-time of standard deterministic min $(s,t)$-cut algorithms for hypergraphs---we recall that min $(s,t)$-cut in a $n$-vertex $m$-hyperedge hypergraph can be solved in $O((m+n)mn)$.
\end{proof}


\subsection{Computing the Family of Maximal \((p, H, w)\)-Tight Sets}\label{appendix:sec:Function-Maximization-Oracles:computing-maximal-pHw-tight-sets-family}
In this section, we prove \Cref{lem:FunctionMaximizetionOracles:computing-maximal-pHw-tight-sets-with-p-max-oracle}. 
We recall that for a set function $p:2^V\rightarrow\Z$ and hypergraph $(H = (V, E), w)$, a set $X\subseteq V$ is $(p, H, w)$-tight if $p(X) = b_{(H,w)}(X)$. The family $\maximalTightSetFamily{p, H, w}$ denotes the family of all inclusion-wise maximal $(p, H, w)$-tight sets. In this section, we show that the family $\maximalTightSetFamily{p, H, w}$ can be computed in $O(|V|^2(|V| + |E|))$ time using $O(|V|^2)$ queries to \functionMaximizationEmptyOracle{p} if the function $p$ and hypergraph $(H, w)$ satisfy certain properties.

\lemComputingMaximalTightSetFamily*
\begin{proof}
We describe a procedure that recursively constructs \maximalTightSetFamily{p, H, w}. The procedure finds a $(p, H, w)$-tight set $X\subseteq V$, and then searches for a maximal $(p, H, w)$-tight set $Y$ containing the set $X$ --- we describe this computation in the paragraph below. 
The procedure is recursive and it reaches the base case if 
the computed set $X$ is not a $(p, H, w)$-tight set or $X,Y = \emptyset$. If $X$ is not a $(p, H, w)$-tight set, then the procedure returns the empty set, and if $X,Y = \emptyset$, then the procedure returns $\{\emptyset\}$.
If both of the base case conditions are not satisfied, the procedure recurses on the function $p':2^{V'}\rightarrow\Z$ defined on the ground set $V':=V - Y$ as $p'(Z) = p(Z)$ for every $Z \subseteq V'$, and the hypergraph $(H' = (V', E'), w')$ where $E' := \{e \cap V' : e \in E\}$ and $w'(e'):= \sum_{e \in E}w(e)\cdot\indicator_{e\cap V' = e'}$ for every $e' \in E'$, to obtain the family $\maximalTightSetFamily{p', H', w'}$.
It then returns the family $\{Y\} \cup (\maximalTightSetFamily{p', H', w'} - \{\emptyset\})$. 
The procedure terminates within $|V|$ recursive calls because the size of the ground set decreases during every recursive call (i.e. $|V'| < |V|$). Furthermore, the correctness of the procedure follows from disjointness of the family $\maximalTightSetFamily{p, H, w}$ by \Cref{lem:weak-to-strong:uncrossing:maximal-tight-set-family-disjoint} and because $b_{(H,w)}(X) = b_{(H', w')}(X)$ for each $X \subseteq V'$, where $(H', w')$ is the hypergraph constructed as input to the subsequent recursive call.
    
We now bound the run-time and the number of queries required by the procedure described above. Consider a recursive call of the procedure.
We note that a $(p, H, w)$-tight set exists if and only if the set returned by the query $\functionMaximizationOracle{p}((G_0, c_0):= (H, w), S_0:=\emptyset, T_0:=\emptyset)$ is a $(p, H, w)$-tight set.
Furthermore, verifying whether the returned set $X$ is $(p, H, w)$-tight can be done in $O(|V| + |E|)$ time by checking whether $p(X) = b_{(H,w)}(X)$.
Thus, a $(p, H, w)$-tight set $X\subseteq V$ can be found in  $O(|V| + |E|)$ time using one query to \functionMaximizationOracle{p}. Moreover, the input hypergraph to \functionMaximizationOracle{p} query has $|V|$ vertices and $|E|$ hyperedges.
Next, \Cref{claim:FunctionMaximizetionOracles:maximal-tight-set-containing-given-tight-set} below shows that we can find a maximal $(p, H, w)$-tight set $Y$ containing the set $X$ in $O(|V|(|V|+|E|))$ time using $O(|V|)$ queries to \functionMaximizationOracle{p} --- here, the runtime includes the time required to construct the inputs to \functionMaximizationOracle{p}; moreover, each query to \functionMaximizationOracle{p} takes the hypergraph $(H, w)$ as input.
Let $p':2^{V'}\rightarrow\Z$ be the input function to the subsequent recursive call. Let $(G = (V', E_G), c)$ be a hypergraph such that $|E_G| \leq |E|$ and $S, T \subseteq V'$ be disjoints sets. 
Then, \Cref{claim:FunctionMaximizetionOracles:computing-maximal-tight-sets:maximization-oracle-for-restriction} below shows that a query to $\functionMaximizationOracle{p'}((G, c), S, T)$ can be implemented in $O(|V| + |E|)$ time using one query to \functionMaximizationOracle{p} --- here, the runtime includes the time required to construct the input to \functionMaximizationOracle{p} query; moreover the hypergraph input to \functionMaximizationOracle{p} query has $|V|$ vertices and $|E_G|$ hyperedges. We note that all hypergraphs $(G = (V', E_G), c)$ that are used as inputs to \functionMaximizationOracle{p'} during subsequent recursive calls satisfy $|E_G| \leq |E|$.
Thus, every recursive call of the procedure can be implemented to run in $O(|V|(|V|+|E|))$ time and make $O(|V|)$ calls to \functionMaximizationOracle{p}, where the runtime includes the time to construct the inputs to each query to \functionMaximizationOracle{p}. Consequently, the procedure runs in time $O(|V|^2(|V|+|E|))$ and makes $O(|V|^2)$ queries to $\functionMaximizationOracle{p}$, where the runtime includes the time to construct the inputs to the \functionMaximizationOracle{p} queries.
\end{proof}

\begin{claim}\label{claim:FunctionMaximizetionOracles:maximal-tight-set-containing-given-tight-set}
         Let $X \subseteq V$ be a $(p, H, w)$-tight set. Then, a maximal $(p, H, w)$-tight set containing the set $X$ can be computed in $O(|V|(|V|+|E|))$ time using $O(|V|)$ queries to \functionMaximizationOracle{p}. The runtime includes the time required to construct the inputs to \functionMaximizationOracle{p}. Moreover, each query to \functionMaximizationOracle{p} takes the hypergraph $(H, w)$ as input.
\end{claim}
\begin{proof}
We describe a recursive procedure to compute a maximal $(p, H, w)$-tight set containing the set $X$. For every $u \in V - X$, let $Z^u$ be the set returned by querying $\functionMaximizationOracle{p}((H_0, w_0):= (H, w), S_0 := X\cup\{u\}, T_0 := \emptyset)$. If there exists a vertex $u \in X$ such that the set $Z^u$ is a $(p, H, w)$-tight set, then recursively compute a $(p, H, w)$-tight set containing the set $X':= Z^u$. Otherwise, return the set $X$. 

The correctness of the above procedure is because of the following. The above procedure terminates (in at most $|V|$ recursive calls) because the cardinality of the input set $X$ increases during each recursive call (i.e., $|X'| > |X|$). Let $Y \subseteq V$ be a maximal $(p, H, w)$-tight set such that $X \subseteq Y$. If $Y = X$, then there are no $(p, H, w)$-tight sets that strictly contain the set $X$ and the procedure returns the set $X$. Suppose that $X \subsetneq Y$ and let $u \in Y - X$. Then, there exists a $(p, H, w)$-tight set $Z$ such that $X\cup\{u\} \subsetneq Z \subseteq Y$. Thus, the set $Z^u$ computed by the procedure will be a $(p, H, w)$-tight set that strictly contains the set $X$. Using this fact, the procedure can be shown to return a maximal $(p, H, w)$-tight set containing the set $X$ by induction on $|X|$.   
    
We now describe a slight modification of the above procedure that leads to an improved runtime and a smaller number of queries to \functionMaximizationOracle{p}. Let $u \in V - X$ be a vertex. If the set $Z^{u}$ is $(p, H, w)$-tight, then we include $u$ in the input set $S_0$ to all $\functionMaximizationOracle{p}$ queries in subsequent recursive calls of the procedure. Alternatively, if the set $Z^{u}$ is not $(p, H, w)$-tight, then we include $u$ in the input set $T_0$ to all $\functionMaximizationOracle{p}$ queries in subsequent recursive calls procedure. Consequently, the number of queries to $\functionMaximizationOracle{p}$ (and number of recursive calls of the procedure) is at most $|V|$. We note that constructing the input to each query and verifying whether the set returned by the query is $(p, H, w)$-tight can be done in $O(|V| + |E|)$ time. Thus, the procedure terminates in $O(|V|(|V|+|E|))$ time. Here, the runtime includes the time to construct the inputs to the \functionMaximizationOracle{p} queries. Moreover, each query takes the hypergraph $(H,w)$ as input. 
\end{proof}
    
\begin{claim}\label{claim:FunctionMaximizetionOracles:computing-maximal-tight-sets:maximization-oracle-for-restriction}
    Let $p:2^V\rightarrow\Z$ be a function. Let $Y \subseteq V$ and let $p'$ denote the function $p$ restricted to the ground set $V':=V - Y$. Then, for a given hypergraph $(G = (V', E_G), c:E_G\rightarrow\Z_+)$ such that $|E_G| \leq |E|$ and disjoint subsets $S, T \subseteq V'$, $\functionMaximizationOracle{p'}((G,c), S, T, y)$ can be computed in $O(|V| + |E|)$ time using one query to $\functionMaximizationOracle{p}$.
    The runtime includes the time required to construct the input to \functionMaximizationOracle{p} query. Moreover, the hypergraph input to \functionMaximizationOracle{p} query has $|V|$ vertices and $|E_G|$ hyperedges.
\end{claim}
\begin{proof}
    Let $(G_1, c_1)$ be the hypergraph obtained from $(G,c)$ by adding the vertices $Y$. Then, the answer to the query $\functionMaximizationOracle{p'}((G, c), S, T)$ is the same as the answer to the query $\functionMaximizationOracle{p}((G_1, c_1), S, T\cup Y)$.
\end{proof}






\end{document}
